\documentclass[11pt,a4paper,oneside]{book}

\usepackage[T1]{fontenc}
\usepackage[utf8]{inputenc}
\usepackage{libertinus}
\usepackage{amsmath,amssymb,amsthm,mathtools,mathrsfs,bm}
\usepackage[margin=2.4cm]{geometry}
\usepackage{microtype}
\usepackage{array,booktabs,longtable,tabularx,multirow}
\usepackage{enumitem}
\usepackage[dvipsnames,svgnames]{xcolor}
\usepackage{listings}
\usepackage{tcolorbox}
\usepackage{fancyhdr}
\usepackage{needspace}
\usepackage[hidelinks]{hyperref}
\usepackage{xurl}
\usepackage{cleveref}

\tcbuselibrary{skins,breakable}

\hypersetup{
  pdftitle={Asymptotic Transseries and Inversion of Four Combinatorial Sequences},
  pdfsubject={Transseries, Lambert-W reversion, and discrete inversion},
  bookmarksnumbered=true,
}

\pagestyle{fancy}
\fancyhf{}
\fancyhead[L]{\small\nouppercase\leftmark}
\fancyhead[R]{\small\thepage}
\renewcommand{\headrulewidth}{0.4pt}

\lstset{
  basicstyle=\ttfamily\footnotesize,
  breaklines=true,
  columns=fullflexible,
  keepspaces=true,
  showstringspaces=false,
  frame=single,
  framesep=3pt,
  rulecolor=\color{black!25},
}

\theoremstyle{plain}
\newtheorem{theorem}{Theorem}[chapter]
\newtheorem{proposition}[theorem]{Proposition}
\newtheorem{lemma}[theorem]{Lemma}
\newtheorem{corollary}[theorem]{Corollary}
\newtheorem{conjecture}[theorem]{Conjecture}
\newtheorem{identity}[theorem]{Identity}
\theoremstyle{definition}
\newtheorem{definition}[theorem]{Definition}
\newtheorem{algorithm}[theorem]{Algorithm}
\newtheorem{example}[theorem]{Example}
\newtheorem{problem}[theorem]{Problem}
\newtheorem{convention}[theorem]{Convention}
\theoremstyle{remark}
\newtheorem{remark}[theorem]{Remark}
\newtheorem{warning}[theorem]{Warning}
\newtheorem{historical}[theorem]{Historical note}

\crefname{chapter}{Part}{Parts}

% `book` provides no abstract environment.
\providecommand{\abstractname}{Abstract}
\newenvironment{abstract}%
  {\begin{center}\bfseries\abstractname\end{center}%
   \begin{list}{}{\setlength{\leftmargin}{2em}\setlength{\rightmargin}{2em}}%
   \item[]\small}%
  {\end{list}}

% ---------------------------------------------------------------- house macros
% Number systems.
\newcommand{\Nat}{\mathbb{N}}
\newcommand{\Int}{\mathbb{Z}}
\newcommand{\Rat}{\mathbb{Q}}
\newcommand{\Real}{\mathbb{R}}
\newcommand{\Cplx}{\mathbb{C}}

% Upright constants and differentials.
\newcommand{\e}{\mathrm{e}}
\newcommand{\ii}{\mathrm{i}}
\newcommand{\dd}{\mathrm{d}}

% Coefficient extraction, asymptotic symbols, operators.
\newcommand{\Coef}[1]{[#1]\,}
\newcommand{\bigO}{\mathcal{O}}
\newcommand{\littleo}{o}
\DeclareMathOperator{\Res}{Res}
\DeclareMathOperator{\sgn}{sgn}
\DeclareMathOperator{\Li}{Li}
\DeclareMathOperator*{\argmin}{arg\,min}

% Bell polynomials: ordinary partial (\OB) and exponential partial (\EB).
\newcommand{\OB}{\widehat{B}}
\newcommand{\EB}{B}

% The Lambert function, its branches, and the Lambert carrier of Part 0.
\newcommand{\Wlam}{W}
\newcommand{\Xcore}{X}

% Stirling numbers.
\newcommand{\stirlingfirst}[2]{\genfrac{[}{]}{0pt}{}{#1}{#2}}
\newcommand{\stirlingsecond}[2]{\genfrac{\{}{\}}{0pt}{}{#1}{#2}}

% Editorial boxes.
\newtcolorbox{sourcebox}{colback=black!3,colframe=black!35,boxrule=0.4pt,
  left=4pt,right=4pt,top=3pt,bottom=3pt,breakable}
\newtcolorbox{repairbox}{colback=Goldenrod!12,colframe=Goldenrod!60!black,
  boxrule=0.4pt,left=4pt,right=4pt,top=3pt,bottom=3pt,breakable}

\allowdisplaybreaks[1]
\setlength{\parskip}{2pt plus 1pt}

\begin{document}
\frontmatter

\title{\textbf{Asymptotic Transseries and Inversion\\ of Four Combinatorial Sequences}\\[0.35em]
  \Large Rooted Trees, the Double Factorial, the Partition Numbers,
  and the Swing Factorial:\\ one Lambert-cored reversion calculus and its four specializations}
\author{Vladimir Reshetnikov\\
  \small ProveIt formal-development synthesis}
\date{3 September 2026}

\maketitle

\begin{abstract}
Four combinatorial sequences are studied here to all orders and then
inverted: the numbers \(a_n\) of unlabelled rooted trees on \(n\) vertices
(OEIS A000081, the trees underlying Butcher series), the double factorial
\(n!!\), the partition numbers \(p(n)\) (A000041), and the swing factorial
\(\operatorname{sw}(n)=n!/\lfloor n/2\rfloor!^{2}\) (A056040).  For each we
give the complete forward asymptotic transseries and then the asymptotic
inverse: given a value, recover the index, to all orders and with an explicit
error theory.

The four subjects look unrelated and their forward analyses genuinely are ---
a Pólya-tree singularity analysis, a pair of gamma-ratio branches, Rademacher's
convergent root-of-unity series, and a second pair of gamma-ratio branches.
Their \emph{inversions} are not.  Every one of them, after a monomial change
of variable, is the inversion of an exponential--power model
\[
  F(x)=C\exp(ax^{\alpha})\,x^{b}\exp Q(x^{-\delta}),
\]
whose dominant phase is inverted exactly by the Lambert \(W\)-function and
whose remaining corrections are generated by one triangular Bell-polynomial
recurrence.  Chapter~0 develops that apparatus once, in the generality the
four cases require, with complete proofs; Chapters~1--4 supply only the
parameter dictionary for each sequence and the mathematics genuinely specific
to it.  Chapter~5 compares the four and records what the common frame does and
does not explain.

This volume is a consolidation of twelve independently written articles,
three on each subject, filed in this repository on 3 September 2026.  Where
the three treatments of a subject agreed, one presentation was kept; where
they disagreed, the disagreement is stated in the text and resolved; where a
result was asserted without proof, a proof is supplied or the statement is
demoted to a remark that says what is actually established.  Corrections to
the sources are boxed in the text and listed in
Appendix~\ref{app:repairs}; provenance is in Appendix~\ref{app:provenance}.
\end{abstract}

\tableofcontents

\mainmatter

%% ==================== fragment p0 ====================
\chapter{The common apparatus}\label{p0:sec:top}

\begin{sourcebox}
All twelve filed articles carry a private copy of the material of this
chapter.  The presentation below follows, in the main, the two sections
\emph{A universal coefficient calculus for exponential--power inversion} and
\emph{Lambert-normalized inverse transseries} of
\textsf{Swing\_Factorial\_Transseries\_Article}, which state the reversion in
the greatest generality of the twelve.  Four things have been changed.
(i)~The Lagrange fixed-point lemma is reproved: the swing article's proof
appeals to Lagrange--B\"urmann inversion in a setting where the required
invertibility hypothesis fails, and the gap is closed here by a universality
argument (\cref{p0:lem:LB-universal}).
(ii)~The exponential--power model is separated into a \emph{formal} and an
\emph{analytic} reading, and every asymptotic claim is given an explicit
remainder hypothesis; the sources routinely pass from a coefficient recurrence
to a Poincar\'e statement without one.
(iii)~The dominant core is axiomatized (\cref{p0:def:core}) so that the
factorial-type phase of \cref{p2:sec:top} --- which does \emph{not} belong to
the exponential--power model --- is covered by the same reversion theorem.
(iv)~The discrete-inverse material of
\textsf{Butcher\_Tree\_Transseries-2},
\textsf{double\_factorial\_transseries-2} and the three partition articles is
merged into one statement with a residue-class lattice.
\end{sourcebox}

\section{Standing conventions}\label{p0:sec:conventions}

Throughout the volume $R$ denotes a commutative $\Rat$-algebra, and
$R[[t]]$, $R[[t,u]]$ formal power series rings over it.  For a formal series
$f$ we write $\Coef{t^n}f$ for the coefficient of $t^n$.  The symbols
$\e,\ii,\dd$ are the upright constants and differential; $\Nat$, $\Int$,
$\Rat$, $\Real$, $\Cplx$ are the usual number systems, and
$\Nat=\{0,1,2,\dots\}$.  Asymptotic statements use $\bigO$ and
$\littleo$ in the Landau sense as the stated variable tends to
$+\infty$; an implied constant is always allowed to depend on every quantity
held fixed in the statement (in particular on the truncation index $N$) and on
nothing else.

The sources disagree on three points of bookkeeping.  We fix them once.

\begin{remark}[Conventions adopted, and where the sources differ]
\label{p0:rem:conventions}
\leavevmode
\begin{enumerate}
\item \emph{Ordinary versus exponential Bell polynomials.}  The partition
      articles and \textsf{Butcher\_Tree\_Transseries-2} take the
      \emph{exponential} partial Bell polynomials $\EB_{n,k}$ as primitive and
      convert; \textsf{Swing\_Factorial\_Transseries\_Article} and
      \textsf{double\_factorial\_transseries-2} take the \emph{ordinary} ones
      $\OB_{n,k}$ as primitive.  The volume defines both
      (\cref{p0:def:obell,p0:def:ebell}), records the conversion
      (\cref{p0:lem:bell-conversion}), and uses $\OB$ in every reversion
      formula and $\EB$ in every exponentiation or logarithm of a
      \emph{unit} series.  The reason is not taste: reversion multiplies a
      correction series whose coefficients carry no factorials, and $\OB$ is
      the transform that leaves them alone.
\item \emph{Sign of the logarithmic phase.}  We write the dominant phase as
      \[
        L=aX+b\log X ,
      \]
      with $b$ a \emph{signed} constant.  \textsf{Butcher\_Tree\_Transseries-2}
      writes $Z=X-p\log X$ and the partition articles write $\rho-2\log\rho=L$;
      both are the case $b=-p<0$.  Every formula below is stated in the signed
      convention, and the dictionary of
      \cref{p0:tab:dictionary} gives $b$ for each chapter.  One dividend of the
      signed convention is \cref{p0:cor:branch-rule}: the Lambert branch is
      $\Wlam_0$ when $b>0$ and $\Wlam_{-1}$ when $b<0$, with no case analysis
      beyond the sign.
\item \emph{Truncation index.}  A truncation ``of order $N$'' retains the terms
      with index $n=1,\dots,N$ \emph{inclusive}, together with the leading
      term.  \textsf{Swing\_Factorial\_Transseries\_Article} uses the exclusive
      convention for its forward table (``the truncation labelled $N$ retains
      $A_0,\dots,A_{N-1}$'') and the inclusive one for its inverse table
      (``the first $N$ correction terms are retained''); the two tables of that
      article are therefore not on a common scale.  The volume is inclusive
      throughout.
\end{enumerate}
\end{remark}

\section{The exponential--power model}\label{p0:sec:model}

\subsection{Formal and analytic readings}

\begin{definition}[The exponential--power model]\label{p0:def:model}
Fix constants $C>0$, $a>0$, $\alpha\in(0,1]$, $\delta>0$, $b\in\Real$, and a
formal series without constant term,
\[
 Q(t)=\sum_{j\ge1}q_jt^j\in t\,\Real[[t]] .
\]
The associated \emph{exponential--power model} is the expression
\begin{equation}\label{p0:eq:model}
 F(x)=C\exp\bigl(ax^\alpha\bigr)\,x^{b}\,\exp Q\bigl(x^{-\delta}\bigr),
\end{equation}
understood through its logarithm
\begin{equation}\label{p0:eq:model-log}
 \log F(x)=\log C+ax^{\alpha}+b\log x+Q\bigl(x^{-\delta}\bigr).
\end{equation}
We call $(C,a,\alpha,b,\delta,Q)$ the \emph{model data}.

A positive function $F\in C^1([x_0,\infty))$ is said to be \emph{of
exponential--power type with data $(C,a,\alpha,b,\delta,Q)$} when, for every
$M\in\Nat$, there is a constant $K_M$ with
\begin{equation}\label{p0:eq:model-analytic}
 \Biggl|\log F(x)-\log C-ax^{\alpha}-b\log x
        -\sum_{j=1}^{M}q_jx^{-j\delta}\Biggr|
 \le K_Mx^{-(M+1)\delta}
 \qquad(x\ge x_0),
\end{equation}
and \emph{differentiably} of that type when in addition
\begin{equation}\label{p0:eq:model-derivative}
 (\log F)'(x)=a\alpha x^{\alpha-1}+\frac bx+\bigO\bigl(x^{-1-\delta}\bigr)
 \qquad(x\to\infty).
\end{equation}
\end{definition}

\begin{remark}[Why the two readings are kept apart]\label{p0:rem:formal-analytic}
The series $Q$ is \emph{not} assumed convergent.  In
\cref{p1:sec:top,p2:sec:top,p4:sec:top} it is divergent of Gevrey type, which
is what makes \cref{p0:thm:optimal-truncation} bite there.  It is \emph{not}
divergent in \cref{p3:sec:top}: for the partition numbers
$Q(t)=\log(1-t)$ has radius of convergence $1$, every algebraic expansion of
that chapter converges, and \cref{p0:thm:optimal-truncation} correspondingly
has no content there --- the error floor of \cref{p3:sec:top} is set by the
exponentially small Rademacher remainder and not by a least term.  This is
precisely why the two readings are kept apart: convergence of $Q$ is a
property of the subject, not of the frame.  Every purely algebraic assertion
below --- coefficient recurrences, closed coefficient formulae, degree bounds
--- is a statement about \eqref{p0:eq:model-log} as a formal object and needs
no analytic hypothesis at all.  Every assertion with an error term uses
\eqref{p0:eq:model-analytic}, and every assertion that inverts uses
\eqref{p0:eq:model-derivative} as well: an asymptotic expansion may not be
differentiated without a hypothesis, and the inversion of $F$ is controlled by
$(\log F)'$.  In each chapter \eqref{p0:eq:model-derivative} is verified
directly from a digamma, Laplace or Rademacher representation, never inferred
from \eqref{p0:eq:model-analytic}.
\end{remark}

\subsection{Reduction to a linear phase}

\begin{lemma}[Monomial $\alpha$-reduction]\label{p0:lem:alpha-reduction}
Let $F$ be of exponential--power type with data
$(C,a,\alpha,b,\delta,Q)$ on $[x_0,\infty)$, and put
\begin{equation}\label{p0:eq:alpha-substitution}
 \xi=x^{\alpha},\qquad x=\xi^{1/\alpha},\qquad
 G(\xi):=F\bigl(\xi^{1/\alpha}\bigr).
\end{equation}
Then $G$ is of exponential--power type on $[x_0^\alpha,\infty)$ with data
\begin{equation}\label{p0:eq:alpha-reduced-data}
 \Bigl(C,\;a,\;1,\;\frac b\alpha,\;\frac\delta\alpha,\;Q\Bigr),
\end{equation}
that is,
\begin{equation}\label{p0:eq:alpha-reduced-log}
 \log G(\xi)=\log C+a\xi+\frac b\alpha\log\xi
             +Q\bigl(\xi^{-\delta/\alpha}\bigr).
\end{equation}
The coefficients $q_j$ are unchanged.  If moreover $F$ is differentiably of
its type then $G$ is differentiably of its type.  Finally $x\mapsto x^\alpha$
is an increasing bijection of $[x_0,\infty)$ onto $[x_0^\alpha,\infty)$, so
$F$ is strictly increasing if and only if $G$ is, and then
\begin{equation}\label{p0:eq:alpha-inverse-transport}
 F^{-1}(y)=\bigl(G^{-1}(y)\bigr)^{1/\alpha}.
\end{equation}
\end{lemma}

\begin{proof}
Substituting $x=\xi^{1/\alpha}$ into \eqref{p0:eq:model-log} and using
$x^{\alpha}=\xi$, $\log x=\alpha^{-1}\log\xi$ and
$x^{-j\delta}=\xi^{-j\delta/\alpha}$ gives
\eqref{p0:eq:alpha-reduced-log} term by term; the truncated inequality
\eqref{p0:eq:model-analytic} transforms into the corresponding inequality for
$G$ with the same constants $K_M$, because
$x^{-(M+1)\delta}=\xi^{-(M+1)\delta/\alpha}$.  For the derivative,
$\dd x/\dd\xi=\alpha^{-1}\xi^{1/\alpha-1}$, so
\[
 (\log G)'(\xi)=(\log F)'\bigl(\xi^{1/\alpha}\bigr)
 \cdot\frac{\xi^{1/\alpha-1}}{\alpha}
 =\Bigl(a\alpha\xi^{1-1/\alpha}+\frac{b}{\xi^{1/\alpha}}
        +\bigO\bigl(\xi^{-(1+\delta)/\alpha}\bigr)\Bigr)
   \frac{\xi^{1/\alpha-1}}{\alpha}
 =a+\frac{b/\alpha}{\xi}+\bigO\bigl(\xi^{-1-\delta/\alpha}\bigr),
\]
which is \eqref{p0:eq:model-derivative} for $G$ with $\alpha$ replaced by $1$.
The last two assertions are immediate from strict monotonicity of
$x\mapsto x^\alpha$ on $(0,\infty)$.
\end{proof}

\begin{remark}[Normalization used from here on]\label{p0:rem:alpha-one}
By \cref{p0:lem:alpha-reduction} every statement about the model
\eqref{p0:eq:model} reduces to the case $\alpha=1$, and we assume $\alpha=1$
from \cref{p0:sec:core} onwards.  We assume in addition $\delta=1$; this is the
value of $\delta/\alpha$ in every chapter of the volume
(\cref{p0:tab:dictionary}).  \Cref{p0:prop:two-grid} states the general-$\delta$
reversion, on the two-generator monomial grid it requires.
\end{remark}

\begin{remark}[The affine shift is taken \emph{before} the reduction]
\label{p0:rem:affine-shift}
Chapters~\ref{p2:sec:top} and \ref{p3:sec:top} apply the model not to the
combinatorial index $n$ but to a shifted variable $x=n-x_\ast$ ($x_\ast=1/24$
for the partition numbers, $x_\ast=-1$ for the double factorial).  This is not
a loss of generality but a genuine choice of normal form, and the order of
operations matters.  If $F$ has data $(C,a,\alpha,b,\delta,Q)$ in $x$ and one
re-expands in $n=x+x_\ast$, then for $\alpha<1$
\[
 a(n-x_\ast)^{\alpha}
 =an^{\alpha}+a\sum_{k\ge1}\binom{\alpha}{k}(-x_\ast)^kn^{\alpha-k},
\]
so the phase acquires an entire new tower of powers $n^{\alpha-k}$: the
re-expanded object no longer has model data with a single exponent grid.  The
partition articles say this explicitly --- their infinite algebraic
coefficient arrays in the variable $n^{-1/2}$ ``arise entirely from
re-expanding the shift $N=n-1/24$''.  The volume therefore normalizes the
shift first, inverts in $x$, and returns to the index by
$n=x+x_\ast$, an exact identity that costs nothing.
\end{remark}

\begin{remark}[Which chapters the model covers, and which it does not]
\label{p0:rem:model-scope}
Chapters~\ref{p1:sec:top}, \ref{p3:sec:top} and \ref{p4:sec:top} are governed by
\cref{p0:def:model}: the rooted-tree counts and the swing factorial with
$\alpha=1$, the partition numbers with $\alpha=1/2$
(\cref{p0:tab:dictionary}).  \Cref{p2:sec:top} is not.
\end{remark}

\begin{repairbox}
\textbf{Claim repaired.}  It is natural, and it is asserted in the merge brief
for this volume, that the case $\alpha=1$ of \cref{p0:def:model} covers the
double factorial alongside the swing factorial and the rooted-tree counts.
It does not.  For the swing factorial
$n!/\lfloor n/2\rfloor!^{2}$ the growth is $2^{n}n^{\pm1/2}$ and the model
applies with $a=\log 2$.  For the double factorial,
\[
 \log\bigl((x+1)!!\bigr)=\frac{x+1}{2}\bigl(\log(x+1)-1\bigr)+\bigO(\log x),
\]
so the dominant phase is $\tfrac12s\log s$ with $s=x+1$, not $ax$: no choice of
the data $(C,a,\alpha,b,\delta,Q)$ with $\alpha\le1$ reproduces it, since
$s\log s/s^{\alpha}\to\infty$ for every $\alpha\le1$.  The three double
factorial articles are consistent with this; they never claim membership of the
exponential--power family, and they invert the phase
$\tfrac{s}{2}(\log s-1)$ directly.
\textbf{Correction.}  The chapter states the reversion for an axiomatized
\emph{dominant core} (\cref{p0:def:core,p0:thm:core-reversion}) of which both
the linear-plus-logarithmic phase $aX+b\log X$ and the factorial phase
$X(\kappa\log X+d)$ are instances, both exactly invertible by Lambert $\Wlam$
(\cref{p0:thm:lambert-core,p0:prop:factorial-core}).  The theorem named in the
Part~0 interface, \cref{p0:thm:lambert-centered}, is stated for the
exponential--power model as specified, and \cref{p0:thm:core-reversion} is the
version \cref{p2:sec:top} cites.
\end{repairbox}

\begin{definition}[Admissible dominant core]\label{p0:def:core}
Let $X_\ast\ge0$.  A function $\Phi_0\in C^{2}((X_\ast,\infty))$ is an
\emph{admissible core} when
\begin{enumerate}
\item $\Phi_0'(X)>0$ for $X>X_\ast$ and $\Phi_0(X)\to+\infty$ as
      $X\to+\infty$;
\item $\Phi_0'$ is eventually monotone and $\Phi_0'(X)\to\Lambda_\infty\in(0,+\infty]$;
\item $X\Phi_0'(X)\to+\infty$.
\end{enumerate}
Its \emph{core slope} is $\Phi_0'$.  Given a target $L$ in the range of
$\Phi_0$, the \emph{core solution} $X=X(L)$ is the unique large solution of
$\Phi_0(X)=L$.
\end{definition}

The two cores used in this volume are
\begin{equation}\label{p0:eq:two-cores}
 \Phi_0(X)=aX+b\log X
 \quad(\text{\cref{p1:sec:top},\ \ref{p3:sec:top},\ \ref{p4:sec:top}}),
 \qquad
 \Phi_0(X)=X\,(\kappa\log X+d)
 \quad(\text{\cref{p2:sec:top}}),
\end{equation}
with $a,\kappa>0$.  Both are admissible on a suitable half-line, and both are
inverted exactly by $\Wlam$ in \cref{p0:sec:core}.

\section{The coefficient calculus}\label{p0:sec:calculus}

\subsection{Two families of partial Bell polynomials}

\begin{definition}[Ordinary partial Bell polynomials]\label{p0:def:obell}
For a sequence $\gamma=(\gamma_1,\gamma_2,\dots)$ in $R$ and $n,k\in\Nat$ put
\begin{equation}\label{p0:eq:obell-def}
 \OB_{n,k}(\gamma):=\Coef{t^n}\Bigl(\sum_{r\ge1}\gamma_rt^r\Bigr)^{k}.
\end{equation}
Explicitly,
\begin{equation}\label{p0:eq:obell-multinomial}
 \OB_{n,k}(\gamma)
 =\sum_{\substack{m_1,m_2,\dots\ge0\\ \sum_rm_r=k,\ \sum_rrm_r=n}}
   \frac{k!}{\prod_{r\ge1}m_r!}\prod_{r\ge1}\gamma_r^{m_r}.
\end{equation}
In particular $\OB_{0,0}=1$, $\OB_{n,0}=0$ for $n\ge1$, $\OB_{n,k}=0$ for
$k>n$, $\OB_{n,n}=\gamma_1^{\,n}$, and $\OB_{n,k}$ depends only on
$\gamma_1,\dots,\gamma_{n-k+1}$.
\end{definition}

\begin{definition}[Exponential partial Bell polynomials]\label{p0:def:ebell}
For $x=(x_1,x_2,\dots)$ in $R$ and $n\ge k\ge0$ put
\begin{equation}\label{p0:eq:ebell-def}
 \EB_{n,k}(x_1,x_2,\dots)
 :=n!\sum_{\substack{m_1,m_2,\dots\ge0\\ \sum_jm_j=k,\ \sum_jjm_j=n}}
   \prod_{j\ge1}\frac1{m_j!}\Bigl(\frac{x_j}{j!}\Bigr)^{m_j},
\end{equation}
and let $\EB_n:=\sum_{k=0}^{n}\EB_{n,k}$ be the complete polynomial.  These are
characterized by the generating identities
\begin{align}
 \frac1{k!}\Bigl(\sum_{j\ge1}x_j\frac{t^j}{j!}\Bigr)^{k}
 &=\sum_{n\ge k}\EB_{n,k}(x_1,x_2,\dots)\frac{t^n}{n!},
 \label{p0:eq:ebell-egf-partial}\\
 \exp\Bigl(\sum_{j\ge1}x_j\frac{t^j}{j!}\Bigr)
 &=\sum_{n\ge0}\EB_n(x_1,\dots,x_n)\frac{t^n}{n!}.
 \label{p0:eq:ebell-egf-complete}
\end{align}
\end{definition}

\begin{lemma}[Conversion]\label{p0:lem:bell-conversion}
For all $n\ge k\ge0$,
\begin{equation}\label{p0:eq:bell-conversion}
 \OB_{n,k}(\gamma_1,\gamma_2,\dots)
 =\frac{k!}{n!}\,\EB_{n,k}(1!\,\gamma_1,\,2!\,\gamma_2,\dots).
\end{equation}
\end{lemma}

\begin{proof}
Set $x_j=j!\,\gamma_j$ in \eqref{p0:eq:ebell-egf-partial}.  The left-hand side
becomes $(k!)^{-1}\bigl(\sum_j\gamma_jt^j\bigr)^k$, whose coefficient of $t^n$
is $\OB_{n,k}(\gamma)/k!$; the right-hand side has coefficient
$\EB_{n,k}(1!\gamma_1,\dots)/n!$.
\end{proof}

\subsection{Powers and logarithms of a unit series}

\begin{lemma}[Powers and logarithms]\label{p0:lem:power-log}
Let $w=\sum_{r\ge1}\gamma_rt^r\in tR[[t]]$ and $n\ge1$.
\begin{enumerate}
\item For every $\beta$ in $R$ (or in any commutative $\Rat$-algebra containing
      $R$),
      \begin{equation}\label{p0:eq:power-general}
       \Coef{t^n}(1+w)^{\beta}
       =\sum_{k=1}^{n}\binom{\beta}{k}\OB_{n,k}(\gamma),
       \qquad
       \binom\beta k=\frac{\beta(\beta-1)\cdots(\beta-k+1)}{k!}.
      \end{equation}
\item For an integer $j\ge1$,
      \begin{equation}\label{p0:eq:power-negative}
       \Coef{t^n}(1+w)^{-j}
       =\sum_{k=0}^{n}(-1)^{k}\binom{j+k-1}{k}\OB_{n,k}(\gamma).
      \end{equation}
\item
      \begin{equation}\label{p0:eq:log-obell}
       \Coef{t^n}\log(1+w)
       =\sum_{k=1}^{n}\frac{(-1)^{k+1}}{k}\OB_{n,k}(\gamma).
      \end{equation}
\end{enumerate}
Equivalently, in exponential form, for $U(t)=1+\sum_{j\ge1}u_jt^j$,
\begin{align}
 \Coef{t^N}U(t)^{\beta}
 &=\frac1{N!}\sum_{k=0}^{N}\beta^{\underline k}\,
   \EB_{N,k}(1!u_1,2!u_2,\dots),
 \qquad \beta^{\underline k}=\beta(\beta-1)\cdots(\beta-k+1),
 \label{p0:eq:unit-power-ebell}\\
 \Coef{t^N}\log U(t)
 &=\frac1{N!}\sum_{k=1}^{N}(-1)^{k-1}(k-1)!\,
   \EB_{N,k}(1!u_1,2!u_2,\dots).
 \label{p0:eq:log-ebell}
\end{align}
\end{lemma}

\begin{proof}
All four identities are the substitution of $w$ into a formal series in one
variable.  Since $w\in tR[[t]]$, the family $(w^k)_{k\ge0}$ is summable in the
$t$-adic topology and $\Coef{t^n}w^k=0$ for $k>n$, so every sum displayed is
finite and the substitution is legitimate.  Now
$(1+w)^\beta=\sum_{k\ge0}\binom\beta kw^k$ gives
\eqref{p0:eq:power-general} on applying \cref{p0:def:obell};
$\binom{-j}{k}=(-1)^k\binom{j+k-1}{k}$ gives
\eqref{p0:eq:power-negative}; and
$\log(1+w)=\sum_{k\ge1}(-1)^{k+1}w^k/k$ gives \eqref{p0:eq:log-obell}.
For the exponential forms write $U=1+w$ with $\gamma_j=u_j$ and substitute
\eqref{p0:eq:bell-conversion} into \eqref{p0:eq:power-general} and
\eqref{p0:eq:log-obell}, using
$\binom\beta k=\beta^{\underline k}/k!$ and $(-1)^{k+1}/k=(-1)^{k-1}(k-1)!/k!$.
\end{proof}

\begin{corollary}[Exponential and logarithmic jets]\label{p0:cor:exp-log-jets}
Let $D(t)=\sum_{j\ge1}d_jt^j\in tR[[t]]$ and $A(t)=\exp D(t)=\sum_{n\ge0}A_nt^n$.
Then $A_0=1$,
\begin{equation}\label{p0:eq:exp-bell}
 A_n=\frac1{n!}\EB_n(1!d_1,2!d_2,\dots,n!d_n)
    =\sum_{\substack{k_1,k_2,\dots\ge0\\ \sum_jjk_j=n}}\prod_{j\ge1}\frac{d_j^{k_j}}{k_j!},
\end{equation}
and the coefficients obey the triangular recurrence
\begin{equation}\label{p0:eq:exp-recurrence}
 nA_n=\sum_{j=1}^{n}j\,d_j\,A_{n-j}\qquad(n\ge1).
\end{equation}
Conversely, if $A=1+\sum_{n\ge1}A_nt^n$ is a unit series then
$d_n=\Coef{t^n}\log A$ is given by \eqref{p0:eq:log-ebell}.
\end{corollary}

\begin{proof}
The first equality is \eqref{p0:eq:ebell-egf-complete} with $x_j=j!d_j$; the
second is the expansion of $\exp$ as a product over $j$.  For
\eqref{p0:eq:exp-recurrence} differentiate $A=\exp D$ to get $A'=D'A$ and
compare coefficients of $t^{n-1}$.  The converse is
\eqref{p0:eq:log-ebell}.
\end{proof}

\subsection{Lagrange--B\"urmann inversion and the fixed-point form}

\begin{theorem}[Lagrange--B\"urmann]\label{p0:thm:LB}
Let $\varphi\in R[[w]]$ with $\varphi(0)$ invertible in $R$.  There is a unique
$W\in tR[[t]]$ with
\begin{equation}\label{p0:eq:LB-equation}
 W=t\,\varphi(W),
\end{equation}
and for every $G\in R[[w]]$ and every $n\ge1$,
\begin{equation}\label{p0:eq:LB-general}
 \Coef{t^n}G\bigl(W(t)\bigr)
 =\frac1n\Coef{w^{n-1}}\Bigl(G'(w)\,\varphi(w)^n\Bigr).
\end{equation}
In particular
\begin{equation}\label{p0:eq:LB-basic}
 \Coef{t^n}W(t)=\frac1n\Coef{w^{n-1}}\varphi(w)^n .
\end{equation}
\end{theorem}

\begin{proof}
Existence and uniqueness: \eqref{p0:eq:LB-equation} determines
$\Coef{t^n}W$ from $\Coef{t^{<n}}W$, because $t\varphi(W)$ modulo $t^{n+1}$
depends only on $W$ modulo $t^{n}$.

For \eqref{p0:eq:LB-general} work in the field of formal Laurent series
$R((w))$ and let $\Res$ denote the coefficient of $w^{-1}$.  Two facts are
needed.

\emph{(a) $\Res f'=0$ for every $f\in R((w))$}, since $w^{-1}$ has no
antiderivative among the monomials: $(w^{k})'=kw^{k-1}$ is a multiple of
$w^{-1}$ only for $k=0$, when it vanishes.

\emph{(b) Change of variables.}  Put $\psi(w):=w/\varphi(w)$, a series in
$wR[[w]]$ with $\psi'(0)=\varphi(0)^{-1}$ invertible; then $W$ is the
compositional inverse of $\psi$, so $\psi(W(t))=t$.  For $f\in R((w))$,
\begin{equation}\label{p0:eq:res-change}
 \Res_w\bigl(f(w)\bigr)=\Res_t\bigl(f(W(t))\,W'(t)\bigr).
\end{equation}
Indeed both sides are $R$-linear and continuous, so it suffices to check
$f=w^k$.  For $k\ne-1$, $W^kW'=\bigl(W^{k+1}/(k+1)\bigr)'$ and (a) gives
$0=\Res_w w^k$.  For $k=-1$, write $W=t\vartheta(t)$ with $\vartheta$ a unit;
then $W'/W=t^{-1}+\vartheta'/\vartheta$ and the second summand lies in
$R[[t]]$, so $\Res_t(W'/W)=1=\Res_ww^{-1}$.

Now apply \eqref{p0:eq:res-change} to
$f(w):=G(w)\psi(w)^{-n-1}\psi'(w)$, noting that
$f(W(t))W'(t)=G(W(t))\,t^{-n-1}\,\tfrac{\dd}{\dd t}\psi(W(t))
 =G(W(t))\,t^{-n-1}$ because $\psi(W(t))=t$:
\[
 \Coef{t^n}G(W(t))=\Res_t\bigl(G(W(t))\,t^{-n-1}\bigr)
 =\Res_w\bigl(G(w)\,\psi(w)^{-n-1}\psi'(w)\bigr).
\]
Since $\psi^{-n-1}\psi'=-\tfrac1n\bigl(\psi^{-n}\bigr)'$, fact (a) and the
product rule give
\[
 \Res_w\bigl(G\,\psi^{-n-1}\psi'\bigr)
 =-\frac1n\Res_w\bigl(G\cdot(\psi^{-n})'\bigr)
 =\frac1n\Res_w\bigl(G'\,\psi^{-n}\bigr)
 =\frac1n\Res_w\bigl(G'(w)\,\varphi(w)^{n}w^{-n}\bigr),
\]
which is \eqref{p0:eq:LB-general}.  Taking $G(w)=w$ gives
\eqref{p0:eq:LB-basic}.
\end{proof}

The reversions of this volume all take the form $U=\Psi(t,U)$ in which the
r\^ole of $\varphi(0)$ is played by $\Psi(t,0)$, an element of $tR[[t]]$, which
is \emph{never} invertible.  \Cref{p0:thm:LB} therefore does not apply
directly, and the following universality lemma is what makes the passage
legitimate.

\begin{lemma}[Lagrange's formula without an invertibility hypothesis]
\label{p0:lem:LB-universal}
Let $S$ be any commutative $\Rat$-algebra and $\varphi\in S[[w]]$
\emph{arbitrary} (in particular $\varphi(0)$ need not be invertible).  Let
$W\in zS[[z]]$ be the unique solution of $W=z\varphi(W)$.  Then for every
$m\ge1$,
\begin{equation}\label{p0:eq:LB-universal}
 \Coef{z^m}W=\frac1m\Coef{w^{m-1}}\varphi(w)^m .
\end{equation}
\end{lemma}

\begin{proof}
Existence and uniqueness of $W$ hold as in \cref{p0:thm:LB} (the recursion is
$z$-adic and uses no invertibility).  Fix $m$.  Iterating $W=z\varphi(W)$
shows that $\Coef{z^m}W$ is a universal polynomial with integer coefficients in
$\varphi_0,\varphi_1,\dots,\varphi_{m-1}$, where $\varphi=\sum_i\varphi_iw^i$;
the right-hand side of \eqref{p0:eq:LB-universal} is a universal polynomial
with rational coefficients in the same variables.  Both sides are therefore
elements $P_m,\;\widetilde P_m$ of the polynomial ring
$\mathcal P:=\Rat[\varphi_0,\dots,\varphi_{m-1}]$, and they are obtained from
$\varphi$ by the unique $\Rat$-algebra homomorphism
$\mathcal P\to S$ sending $\varphi_i\mapsto\varphi_i$.  It thus suffices to
prove $P_m=\widetilde P_m$ in $\mathcal P$.  The localization map
$\mathcal P\to\mathcal P[\varphi_0^{-1}]$ is injective ($\mathcal P$ is an
integral domain), and in $\mathcal P[\varphi_0^{-1}]$ the element $\varphi_0$
is invertible, so \cref{p0:thm:LB} applies and gives
$P_m=\widetilde P_m$ there.  Injectivity transports the equality back to
$\mathcal P$, and the homomorphism transports it to $S$.
\end{proof}

\begin{lemma}[Lagrange fixed-point formula]\label{p0:lem:lagrange-fp}
Let $\Psi\in tR[[t,u]]$, that is, a formal series in $t$ and $u$ each of whose
monomials has $t$-degree at least one.  Then the equation
\begin{equation}\label{p0:eq:fp-equation}
 U(t)=\Psi\bigl(t,U(t)\bigr)
\end{equation}
has a unique solution $U\in tR[[t]]$, and for every $n\ge1$
\begin{equation}\label{p0:eq:fp-coefficient}
 \Coef{t^n}U=\sum_{m=1}^{n}\frac1m\Coef{t^nu^{m-1}}\Psi(t,u)^m .
\end{equation}
\end{lemma}

\begin{proof}
Uniqueness and existence: $\Psi(t,U)$ modulo $t^{n+1}$ depends only on $U$
modulo $t^{n}$ because $\Psi$ is divisible by $t$; so
\eqref{p0:eq:fp-equation} determines the coefficients of $U$ recursively, and
the resulting series solves the equation by $t$-adic continuity.

For \eqref{p0:eq:fp-coefficient} apply \cref{p0:lem:LB-universal} over the
ring $S:=R[[t]]$ with $\varphi(w):=\Psi(t,w)\in S[[w]]$ and marker $z$: the
unique $\widetilde U\in zS[[z]]$ with $\widetilde U=z\varphi(\widetilde U)$
satisfies
\[
 \Coef{z^m}\widetilde U=\frac1m\Coef{w^{m-1}}\Psi(t,w)^m\in S .
\]
Since $\Psi\in tR[[t,u]]$, we have $\Psi^m\in t^mR[[t,u]]$, whence
$\Coef{t^n}\Coef{w^{m-1}}\Psi(t,w)^m=0$ for $m>n$.  Consequently the family
$\bigl(\Coef{z^m}\widetilde U\bigr)_{m\ge1}$ is summable in the $t$-adic
topology of $S$, the specialization $z=1$ is defined coefficientwise in $t$,
and the resulting series $U:=\sum_{m\ge1}\Coef{z^m}\widetilde U$ satisfies
$U=\Psi(t,U)$ by the same summability.  Uniqueness identifies it with the
solution of \eqref{p0:eq:fp-equation}, and extracting $\Coef{t^n}$ gives
\eqref{p0:eq:fp-coefficient} with the sum truncated at $m=n$.
\end{proof}

\begin{repairbox}
\textbf{Source claim.}  \textsf{Swing\_Factorial\_Transseries\_Article},
Lemma \emph{Lagrange fixed-point formula}, proves
\eqref{p0:eq:fp-coefficient} by writing ``Ordinary Lagrange--B\"urmann
inversion in $z$ gives $U=\sum_{m\ge1}z^mm^{-1}[u^{m-1}]\Psi(t,u)^m$''.  The
same step, with the same silence, appears in
\textsf{Partition\_Number\_Transseries\_and\_Inverse} (``Introduce
$u=z\Phi_\rho(u)$.  Lagrange--B\"urmann gives \dots''), in
\textsf{Partition\_Numbers\_Transseries\_and\_Inverse}, and in
\textsf{double\_factorial\_transseries-2}.
\textbf{Defect.}  The classical Lagrange--B\"urmann theorem in the form
\eqref{p0:eq:LB-basic} requires $\varphi(0)$ to be invertible in the
coefficient ring.  Here $\varphi(0)=\Psi(t,0)\in tR[[t]]$, a non-unit, and the
residue proof of \cref{p0:thm:LB} breaks at the change of variables: $\psi=w/\varphi$
is not a well-defined element of $wR[[t]][[w]]$.
\textbf{Repair.}  \Cref{p0:lem:LB-universal}: the identity is a polynomial
identity in the coefficients of $\varphi$ over $\Rat$, hence holds in the
polynomial ring by injectivity of localization at $\varphi_0$, hence in every
$\Rat$-algebra.  \textsf{double\_factorial\_transseries-2} comes closest to
noticing the issue --- ``Although $A$ is a divergent Gevrey series, it is a
valid formal power series with nonzero linear coefficient'' --- but that
remark addresses convergence, not invertibility, and its own $G_\Lambda$ does
have an invertible linear coefficient, so the gap there is only in the
inherited lemma.
\end{repairbox}


\begin{lemma}[Fixed-point formula with a quadratic self-interaction]
\label{p0:lem:lagrange-fp-quadratic}
Let $\Psi=h+f$ with $h\in u^2R[[u]]$ and $f\in tR[[t,u]]$.  Then
$U=\Psi(t,U)$ has a unique solution $U\in tR[[t]]$, and for every $n\ge1$
\begin{equation}\label{p0:eq:fp-quadratic}
 \Coef{t^n}U=\sum_{m=1}^{2n-1}\frac1m\Coef{t^nu^{m-1}}\Psi(t,u)^m .
\end{equation}
\end{lemma}

\begin{proof}
\emph{Existence and uniqueness.}  Write $U=U_{<n}+e_nt^n+\cdots$.  Then
$h(U)=h(U_{<n})+h'(U_{<n})\,e_nt^n+\cdots$, and $h'(U_{<n})\in tR[[t]]$ because
$h'\in uR[[u]]$ and $U_{<n}\in tR[[t]]$; so $e_n$ does not occur in
$\Coef{t^n}h(U)$.  It does not occur in $\Coef{t^n}f(t,U)$ either, since $f$ is
divisible by $t$.  Hence the equation determines $e_n$, as
$\Coef{t^n}\Psi(t,U_{<n})$, from $U_{<n}$; the resulting series solves the
equation by $t$-adic continuity.

\emph{The closed formula.}  Introduce a marker $z$ and let
$\widetilde U\in zS[[z]]$, $S:=R[[t]]$, solve
$\widetilde U=z\Psi(t,\widetilde U)$.  \Cref{p0:lem:LB-universal} gives
$\Coef{z^m}\widetilde U=m^{-1}\Coef{u^{m-1}}\Psi(t,u)^m$.  We claim
\begin{equation}\label{p0:eq:quadratic-vanishing}
 \Coef{t^nu^{m-1}}\Psi(t,u)^m=0\qquad\text{whenever } m>2n-1 .
\end{equation}
Expand $\Psi^m=(h+f)^m$ into $2^m$ products; a product using $h$ in $\sigma$ of
the $m$ slots and $f$ in the remaining $m-\sigma$ has $u$-degree at least
$2\sigma$ and $t$-degree at least $m-\sigma$.  For the $u$-degree to equal
$m-1$ we need $2\sigma\le m-1$, hence the $t$-degree is at least
$m-\tfrac{m-1}2=\tfrac{m+1}2$; so $\Coef{t^n}$ of that product vanishes unless
$n\ge(m+1)/2$, which is \eqref{p0:eq:quadratic-vanishing}.  Consequently the
family $(\Coef{z^m}\widetilde U)_{m\ge1}$ is $t$-adically summable, the
specialization $z=1$ is defined coefficientwise, and, again by summability, the
resulting series solves $U=\Psi(t,U)$.  Uniqueness identifies it with $U$, and
extracting $\Coef{t^n}$ gives \eqref{p0:eq:fp-quadratic}.
\end{proof}

\Cref{p0:lem:lagrange-fp} is the case $h=0$ of
\cref{p0:lem:lagrange-fp-quadratic}, in which the argument above sharpens the
range of $m$ from $2n-1$ to $n$.

\section{The Lambert core}\label{p0:sec:core}

\subsection{Exact inversion of the linear--logarithmic phase}

\begin{theorem}[Exact solution of the dominant block]\label{p0:thm:lambert-core}
Let $a>0$, $b\in\Real$, and consider, for $X>0$,
\begin{equation}\label{p0:eq:core-equation}
 \Phi_0(X):=aX+b\log X=L .
\end{equation}
\begin{enumerate}
\item If $b=0$ the unique solution is $X=L/a$.
\item If $b\ne0$, put $W:=(a/b)X$.  Then \eqref{p0:eq:core-equation} holds if
      and only if
      \begin{equation}\label{p0:eq:core-W-equation}
       W\e^{W}=z(L),\qquad z(L):=\frac ab\,\e^{L/b},
      \end{equation}
      so that the solutions of \eqref{p0:eq:core-equation} are exactly
      \begin{equation}\label{p0:eq:core-lambert}
       X=\frac ba\,\Wlam_k\!\left(\frac ab\,\e^{L/b}\right)
      \end{equation}
      as $k$ ranges over the branches of the Lambert function for which
      $\Wlam_k(z(L))$ is defined.
\item \emph{Branch rule.}  Suppose $b\ne0$.
      \begin{itemize}
      \item If $b>0$ then $z(L)>0$, equation \eqref{p0:eq:core-equation} has
            exactly one positive solution for every real $L$, $\Phi_0$ is
            strictly increasing on $(0,\infty)$, and
            $X=(b/a)\Wlam_0(z(L))$.
      \item If $b<0$ then $z(L)=-(a/|b|)\e^{-L/|b|}$ lies in $(-\e^{-1},0)$
            precisely when
            \begin{equation}\label{p0:eq:core-threshold}
             L>L_c:=|b|\Bigl(1-\log\frac{|b|}{a}\Bigr)
                  =\Phi_0\!\left(\frac{|b|}{a}\right),
            \end{equation}
            and then \eqref{p0:eq:core-equation} has exactly two positive
            solutions, namely
            $(b/a)\Wlam_0(z(L))<|b|/a<(b/a)\Wlam_{-1}(z(L))$.
            $\Phi_0$ is strictly increasing on $(|b|/a,\infty)$ and the large
            solution is
            \begin{equation}\label{p0:eq:core-lambert-minus}
             X=\frac ba\,\Wlam_{-1}\!\left(\frac ab\,\e^{L/b}\right)
              =\frac{|b|}{a}\Bigl(-\Wlam_{-1}\bigl(z(L)\bigr)\Bigr).
            \end{equation}
      \end{itemize}
      In both cases the large solution satisfies $X\to+\infty$ and
      $X=L/a+\bigO(\log L)$ as $L\to+\infty$.
\item \emph{Slope identities.}  With $W=(a/b)X$ and $b\ne0$,
      \begin{equation}\label{p0:eq:core-slope}
       \Phi_0'(X)=a+\frac bX=a\,\frac{1+W}{W},
       \qquad
       \frac{\dd X}{\dd L}=\frac1{\Phi_0'(X)}=\frac1a\,\frac{W}{1+W} .
      \end{equation}
\end{enumerate}
\end{theorem}

\begin{proof}
(1) is immediate.  For (2), multiply \eqref{p0:eq:core-equation} by $1/b$ and
exponentiate:
\[
 \e^{L/b}=\e^{(a/b)X}\,X=\e^{W}\cdot\frac ba\,W ,
\]
because $X=(b/a)W$.  Multiplying by $a/b$ gives
\eqref{p0:eq:core-W-equation}; the step is reversible, and no logarithm of a
negative number occurs, which is what makes the argument uniform in the sign of
$b$.  By definition, the solutions of $W\e^{W}=z$ are the values $\Wlam_k(z)$,
whence \eqref{p0:eq:core-lambert}.

(3) We have $\Phi_0'(X)=a+b/X$.  For $b>0$ this is positive on $(0,\infty)$
and $\Phi_0$ increases from $-\infty$ to $+\infty$, so there is exactly one
solution; correspondingly $W=(a/b)X>0$ and $\Wlam_0$ is the unique real branch
on $(0,\infty)$.  For $b<0$, $\Phi_0'$ vanishes only at $X_\ast=|b|/a$, is
negative on $(0,X_\ast)$ and positive on $(X_\ast,\infty)$; $\Phi_0$ decreases
from $+\infty$ to $\Phi_0(X_\ast)=|b|-|b|\log(|b|/a)=L_c$ and then increases
back to $+\infty$.  Hence there are two positive solutions exactly when
$L>L_c$, one in each monotonicity interval.  On the $W$ side $W=(a/b)X<0$, and
the two real branches of $\Wlam$ on $(-\e^{-1},0)$ are $\Wlam_0$, with values
in $(-1,0)$, and $\Wlam_{-1}$, with values in $(-\infty,-1)$; the map
$W\mapsto(b/a)W$ reverses order, so $\Wlam_{-1}$ produces the larger $X$,
namely the one exceeding $|b|/a$.  Also $z(L)>-\e^{-1}$ is equivalent to
$(a/|b|)\e^{-L/|b|}<\e^{-1}$, that is to \eqref{p0:eq:core-threshold}.
Finally $aX=L-b\log X$ with $X\to\infty$ gives $aX=L+\bigO(\log L)$.

(4) $a+b/X=a+b\cdot\dfrac{a}{bW}=a\bigl(1+W^{-1}\bigr)$, and $\dd X/\dd L$ is
the reciprocal of $\Phi_0'$ by the inverse function theorem, applicable
because $\Phi_0'\ne0$ on the relevant interval.
\end{proof}

\begin{corollary}[Branch rule]\label{p0:cor:branch-rule}
For $b\ne0$ the branch of the Lambert function carrying the large solution of
$L=aX+b\log X$ is determined by the sign of $b$ alone: $\Wlam_0$ if $b>0$,
$\Wlam_{-1}$ if $b<0$.
\end{corollary}

\begin{remark}[Why the branch is not a detail]\label{p0:rem:branch}
Three of the four chapters have $b<0$ and therefore use $\Wlam_{-1}$; the odd
phase of the swing factorial has $b=+1/2$ and uses $\Wlam_0$.
\textsf{Swing\_Factorial\_Transseries\_Article} puts the point sharply: the
principal branch ``would produce the small root of the simplified equation and
is irrelevant to the large-index inverse''.  For $b<0$ the two roots are the
two sides of a genuine fold at $X_\ast=|b|/a$, and choosing wrongly yields an
object all of whose asymptotic corrections are then computed correctly about
the wrong point.
\end{remark}

\subsection{The factorial core}

\begin{proposition}[Exact inversion of a log-linear phase]
\label{p0:prop:factorial-core}
Let $\kappa>0$ and $d\in\Real$, and consider
\begin{equation}\label{p0:eq:factorial-core-equation}
 \Phi_0(X):=X\bigl(\kappa\log X+d\bigr)=L,\qquad X>0 .
\end{equation}
For $L>0$ put
\begin{equation}\label{p0:eq:factorial-core-solution}
 v:=\Wlam_0\!\left(\frac L\kappa\,\e^{d/\kappa}\right),
 \qquad
 X:=\exp\Bigl(v-\frac d\kappa\Bigr) .
\end{equation}
Then $X$ is the unique solution of \eqref{p0:eq:factorial-core-equation} with
$\kappa\log X+d>0$, and
\begin{equation}\label{p0:eq:factorial-core-slope}
 \Phi_0'(X)=\kappa\log X+d+\kappa=\kappa\,(v+1),
 \qquad
 \frac{\dd X}{\dd L}=\frac1{\kappa(v+1)} .
\end{equation}
Moreover $\Phi_0$ is an admissible core in the sense of \cref{p0:def:core} on
$\bigl(\e^{-(d+\kappa)/\kappa},\infty\bigr)$, with $\Phi_0'(X)\to+\infty$.
\end{proposition}

\begin{proof}
Put $s=\log X$.  Then \eqref{p0:eq:factorial-core-equation} reads
$\e^{s}(\kappa s+d)=L$, that is
$\bigl(s+\tfrac d\kappa\bigr)\e^{s+d/\kappa}=(L/\kappa)\e^{d/\kappa}$.  The
right-hand side is positive, so $\Wlam_0$ is the unique real branch and
$s+d/\kappa=v$, which is \eqref{p0:eq:factorial-core-solution}.  Uniqueness
under the stated side condition holds because $\kappa s+d>0$ forces
$s+d/\kappa>0$, and $\sigma\mapsto\sigma\e^{\sigma}$ is injective on
$(0,\infty)$.  Differentiating,
$\Phi_0'(X)=\kappa\log X+d+\kappa=\kappa(s+d/\kappa)+\kappa=\kappa(v+1)$.
Admissibility: $\Phi_0'$ is increasing, tends to $+\infty$, is positive exactly
for $\log X>-(d+\kappa)/\kappa$, and $X\Phi_0'(X)\to\infty$.
\end{proof}

\begin{remark}[The double factorial normalization]\label{p0:rem:factorial-core}
\Cref{p2:sec:top} uses $\kappa=\tfrac12$, $d=-\tfrac12$ and the centred
variable $s=x+1$, so that $\Phi_0(s)=\tfrac s2(\log s-1)$,
$(L/\kappa)\e^{d/\kappa}=2L/\e$, $X=\e^{v+1}$ with $v=\Wlam_0(2L/\e)$, and
$\Phi_0'(X)=\tfrac12(v+1)=\tfrac12\log X$.  This reproduces the core of all
three double factorial articles; the quantity they call $\Lambda$ is
$v+1=\log X$.
\end{remark}

\section{All-orders reversion around the core}\label{p0:sec:reversion}

\subsection{The Lambert-centred inverse}

\begin{theorem}[All-orders reversion around the Lambert core]
\label{p0:thm:lambert-centered}
Let $R$ be a commutative $\Rat$-algebra, let $a\in R$ be invertible, let
$b\in R$, and let $Q=\sum_{j\ge1}q_jt^j\in tR[[t]]$.  Define
\begin{equation}\label{p0:eq:Psi-general}
 \Psi(t,u):=-\frac1a\left[\,b\log(1+tu)
   +Q\!\left(\frac{t}{1+tu}\right)\right].
\end{equation}
\begin{enumerate}
\item $\Psi\in tR[[t,u]]$.
\item There is a unique $U=\sum_{n\ge1}c_nt^n\in tR[[t]]$ with
      $U=\Psi(t,U)$, and
      \begin{equation}\label{p0:eq:c-lagrange}
       c_n=\sum_{m=1}^{n}\frac1m\Coef{t^nu^{m-1}}\Psi(t,u)^m
       \qquad(n\ge1).
      \end{equation}
\item \emph{Triangular Bell recurrence.}  Put $\gamma_1:=0$ and
      $\gamma_r:=c_{r-1}$ for $r\ge2$.  Then, for $n\ge1$,
      \begin{equation}\label{p0:eq:c-bell}
       c_n=-\frac1a\left[
         b\sum_{k=1}^{n}\frac{(-1)^{k+1}}{k}\OB_{n,k}(\gamma)
         +\sum_{j=1}^{n}q_j\sum_{k=0}^{n-j}(-1)^{k}
            \binom{j+k-1}{k}\OB_{n-j,k}(\gamma)
       \right],
      \end{equation}
      and only $c_1,\dots,c_{n-1}$ occur on the right-hand side.  In
      particular $c_n$ lies in the $\Rat$-subalgebra of $R$ generated by
      $a^{-1}$, $b$ and $q_1,\dots,q_n$; no logarithm appears.
\item \emph{Formal exactness.}  In $t^{-1}R[[t]]$ put
      \begin{equation}\label{p0:eq:centered-ansatz}
       \mathsf x:=t^{-1}+U(t)=t^{-1}+\sum_{n\ge1}c_nt^n .
      \end{equation}
      Then $t\mathsf x=1+tU\in1+tR[[t]]$ and $1/\mathsf x=t/(1+tU)\in tR[[t]]$,
      and the following identity holds in $R[[t]]$:
      \begin{equation}\label{p0:eq:formal-exactness}
       a\bigl(\mathsf x-t^{-1}\bigr)+b\log(t\mathsf x)
       +Q\bigl(1/\mathsf x\bigr)=0 .
      \end{equation}
      Equivalently: writing $t=X^{-1}$ and
      $x=X+\sum_{n\ge1}c_nX^{-n}$, the quantity $ax+b\log x+Q(1/x)$ equals
      $aX+b\log X$ to all orders.
\item \emph{Analytic version.}  Let $F$ be differentiably of
      exponential--power type with data $(C,a,1,b,1,Q)$ on $[x_0,\infty)$
      (\cref{p0:def:model}), strictly increasing there, with
      $F(x)\to\infty$.  For $y$ in the range of $F$ set $L=\log(y/C)$ and let
      $X=X(y)$ be the large core solution of $L=aX+b\log X$ furnished by
      \cref{p0:thm:lambert-core}.  Then $X(y)\to\infty$ and, for every fixed
      $N\ge0$,
      \begin{equation}\label{p0:eq:centered-asymptotic}
       F^{-1}(y)=X+\sum_{n=1}^{N}\frac{c_n}{X^{n}}
                 +\bigO\bigl(X^{-N-1}\bigr)
       \qquad(y\to\infty).
      \end{equation}
\end{enumerate}
\end{theorem}

\begin{proof}
(1) $\log(1+tu)=\sum_{k\ge1}(-1)^{k+1}(tu)^k/k\in tR[[t,u]]$, and
$t/(1+tu)=t\sum_{k\ge0}(-tu)^k\in tR[[t,u]]$, so
$(t/(1+tu))^j\in t^jR[[t,u]]$ and the sum over $j$ is summable; hence
$Q(t/(1+tu))\in tR[[t,u]]$.

(2) is \cref{p0:lem:lagrange-fp}.

(3) Put $w:=tU=\sum_{r\ge1}\gamma_rt^r$ with $\gamma$ as stated; $\gamma_1=0$
because $U\in tR[[t]]$.  Then $\Coef{t^n}\log(1+w)$ is given by
\eqref{p0:eq:log-obell}, while
\[
 \Coef{t^n}Q\!\left(\frac t{1+w}\right)
 =\sum_{j\ge1}q_j\,\Coef{t^{n-j}}(1+w)^{-j}
 =\sum_{j=1}^{n}q_j\sum_{k=0}^{n-j}(-1)^k\binom{j+k-1}{k}\OB_{n-j,k}(\gamma)
\]
by \eqref{p0:eq:power-negative}.  Since $c_n=\Coef{t^n}\Psi(t,U)$, this is
\eqref{p0:eq:c-bell}.  For triangularity, $\OB_{n,k}(\gamma)$ depends only on
$\gamma_1,\dots,\gamma_{n-k+1}$, that is on $c_1,\dots,c_{n-k}$; in the first
sum $k\ge1$ and in the second $j\ge1$, so $c_n$ never occurs.  (In fact
$\OB_{n,k}(\gamma)=0$ for $2k>n$, because $\gamma_1=0$.)

(4) With $\mathsf x=t^{-1}(1+w)$, $w=tU$, one has
$a(\mathsf x-t^{-1})=aU$, $\log(t\mathsf x)=\log(1+w)$ and
$1/\mathsf x=t/(1+w)$, so the left-hand side of
\eqref{p0:eq:formal-exactness} equals
$aU+b\log(1+tU)+Q(t/(1+tU))=a\bigl(U-\Psi(t,U)\bigr)=0$.

(5) Fix $N\ge0$, set $t:=1/X$, $U_N(t):=\sum_{n=1}^{N}c_nt^n$ and
$x_N:=X+U_N(t)=X\bigl(1+tU_N(t)\bigr)$.  By
\cref{p0:thm:lambert-core}(3), $X\to\infty$ and $X\sim L/a$; since
$tU_N(t)=\bigO(t^2)$ we get $x_N=X\bigl(1+\bigO(X^{-2})\bigr)$, so
$x_N\to\infty$ and $x_N\ge x_0$ for all large $y$.

\emph{Residual.}  Apply \eqref{p0:eq:model-analytic} with $M=N$:
\[
 \log F(x_N)=\log C+ax_N+b\log x_N+\sum_{j=1}^{N}q_jx_N^{-j}+E_N,
 \qquad |E_N|\le K_Nx_N^{-N-1}.
\]
Since $\log y=\log C+aX+b\log X$ and $\log x_N=\log X+\log(1+tU_N(t))$,
\[
 r_N:=\log F(x_N)-\log y=g(t)+E_N,
 \qquad
 g(t):=aU_N(t)+b\log\bigl(1+tU_N(t)\bigr)
   +\sum_{j=1}^{N}q_j\left(\frac{t}{1+tU_N(t)}\right)^{j}.
\]
Here $g$ is real-analytic on a neighbourhood of $t=0$, being a finite
composition of analytic maps with $1+tU_N(t)\to1$; its Taylor series at $t=0$
is therefore the corresponding formal series.  By (4) the formal series
$aU+b\log(1+tU)+Q(t/(1+tU))$ vanishes identically.  Replacing $U$ by $U_N$
changes it by an element of $t^{N+1}R[[t]]$, because $U-U_N\in t^{N+1}R[[t]]$
and the maps $u\mapsto\log(1+tu)$, $u\mapsto(t/(1+tu))^j$ are $t$-adically
continuous; and discarding the terms of $Q$ with $j>N$ changes it by an element
of $t^{N+1}R[[t]]$ as well.  Hence the Taylor series of $g$ vanishes through
order $t^{N}$, so $g(t)=\bigO(t^{N+1})$ and
\[
 |r_N|\le\bigl|g(1/X)\bigr|+K_Nx_N^{-N-1}=\bigO\bigl(X^{-N-1}\bigr).
\]

\emph{Slope.}  By \eqref{p0:eq:model-derivative} with $\alpha=1$,
$(\log F)'(x)=a+\bigO(x^{-1})$, so there is $x_1\ge x_0$ with
$(\log F)'\ge a/2$ on $[x_1,\infty)$.  For all large $y$ both $x_N$ and
$F^{-1}(y)$ lie in $[x_1,\infty)$.  \Cref{p0:thm:backward-error}, applied to
$\log F$ on that half-line, gives
\[
 \bigl|x_N-F^{-1}(y)\bigr|\le\frac{|r_N|}{a/2}
 =\bigO\bigl(X^{-N-1}\bigr),
\]
which is \eqref{p0:eq:centered-asymptotic}.
\end{proof}

\begin{remark}[No logarithms in the coefficients]\label{p0:rem:no-logs}
Part~(3) is the structural reason for centring at the Lambert core: the whole
logarithmic tower has been resummed into $X$, so the $c_n$ are constants of the
coefficient ring.  The alternative expansion of \cref{p0:sec:flattening} has
coefficients that are polynomials in $\log L$ whose degree grows with the
order.  This is an algebraic statement and by itself says nothing about
numerical accuracy; the numerical comparison is a separate, chapterwise matter
(see \cref{p0:rem:accuracy-claim}).
\end{remark}

\begin{proposition}[General exponent grid]\label{p0:prop:two-grid}
Let $\delta>0$ and keep the remaining hypotheses of
\cref{p0:thm:lambert-centered}.  Let $s$ be a second formal variable and put
\begin{equation}\label{p0:eq:Psi-two-grid}
 \Psi_\delta(t,s,u):=-\frac1a\left[\,b\log(1+tu)
   +Q\!\left(\frac{s}{(1+tu)^{\delta}}\right)\right],
\end{equation}
where $(1+tu)^{-\delta}$ is the formal binomial series
\eqref{p0:eq:power-general}.  Then $\Psi_\delta$ lies in the ideal $(t,s)$ of
$R[[t,s,u]]$, the equation $U=\Psi_\delta(t,s,U)$ has a unique solution
$U=\sum_{i+j\ge1}c_{i,j}t^is^j$ in the ideal $(t,s)$ of $R[[t,s]]$, and the
substitution $t=X^{-1}$, $s=X^{-\delta}$ reverts \eqref{p0:eq:model-log}:
formally
\begin{equation}\label{p0:eq:two-grid-exactness}
 ax+b\log x+Q\bigl(x^{-\delta}\bigr)=aX+b\log X,
 \qquad x=X+\sum_{i+j\ge1}c_{i,j}X^{-i-j\delta}.
\end{equation}
For $\delta=1$ the two markers coincide and
$\Psi_\delta(t,t,u)=\Psi(t,u)$ of \eqref{p0:eq:Psi-general}.
\end{proposition}

\begin{proof}
The membership $\Psi_\delta\in(t,s)$ and the existence and uniqueness of $U$
follow exactly as in parts (1) and (2) of \cref{p0:thm:lambert-centered}, with
the $t$-adic filtration replaced by the $(t,s)$-adic one: $\Psi_\delta$ lies in
$(t,s)$, so $\Psi_\delta(t,s,U)$ modulo $(t,s)^{d+1}$ depends only on $U$
modulo $(t,s)^{d}$.  The identity \eqref{p0:eq:two-grid-exactness} is the
computation of part (4) with $1/\mathsf x=t/(1+tU)$ replaced by
$\mathsf x^{-\delta}=s(1+tU)^{-\delta}$.  Finally the substitution is
legitimate because, for each real exponent $\theta$, only finitely many pairs
$(i,j)$ with $i,j\ge0$ satisfy $i+j\delta=\theta$ when $\delta>0$.
\end{proof}

\begin{remark}[When the grid degenerates]\label{p0:rem:grid}
If $\delta$ is irrational the exponents $i+j\delta$ are pairwise distinct and
the monomial scale is a free monoid on two generators.  If $\delta=p/q$ in
lowest terms, different pairs $(i,j)$ can give the same exponent and the
correct statement is that the reversion lives in $R[[X^{-1/q}]]$.  Every
chapter of this volume has $\delta/\alpha=1$ after
\cref{p0:lem:alpha-reduction}, so the degeneracy is never met here.
\end{remark}

\subsection{Reversion about an arbitrary admissible core}

\begin{theorem}[Reversion about an arbitrary core]\label{p0:thm:core-reversion}
Let $R$ be a commutative $\Rat$-algebra and $\Lambda\in R$ invertible.  Consider
\begin{equation}\label{p0:eq:core-master}
 \Lambda E+h(E)+f(t,E)=0,
 \qquad h\in u^2R[[u]],\quad f\in tR[[t,u]] .
\end{equation}
Then \eqref{p0:eq:core-master} has a unique solution
$E=\sum_{n\ge1}e_nt^n\in tR[[t]]$; it satisfies the triangular recurrence
\begin{equation}\label{p0:eq:core-triangular}
 e_n=-\frac1\Lambda\Coef{t^n}
   \Bigl[h\bigl(E_{<n}\bigr)+f\bigl(t,E_{<n}\bigr)\Bigr],
 \qquad E_{<n}:=\sum_{j=1}^{n-1}e_jt^j ,
\end{equation}
and the closed formula
\begin{equation}\label{p0:eq:core-lagrange}
 e_n=\sum_{m=1}^{2n-1}\frac1m
   \Coef{t^nu^{m-1}}\Bigl(-\Lambda^{-1}\bigl[h(u)+f(t,u)\bigr]\Bigr)^{m}.
\end{equation}
If moreover $h(u)=\sum_{m\ge2}h_mu^m$ and
$f(t,u)=\sum_{k\ge1}\phi_k(t)\,(1+u)^{-\mu_k}$ with $\phi_k\in tR[[t]]$ and
$\mu_k\in\Nat$, then
\begin{equation}\label{p0:eq:core-bell}
 e_n=-\frac1\Lambda\left[
  \sum_{m=2}^{n}h_m\,\OB_{n,m}(e)
  +\sum_{k\ge1}\sum_{i=1}^{n}\bigl(\Coef{t^i}\phi_k\bigr)
    \sum_{m=0}^{n-i}(-1)^m\binom{\mu_k+m-1}{m}\OB_{n-i,m}(e)
 \right],
\end{equation}
in which only $e_1,\dots,e_{n-1}$ occur.
\end{theorem}

\begin{proof}
Existence, uniqueness and \eqref{p0:eq:core-lagrange} are
\cref{p0:lem:lagrange-fp-quadratic} applied to $\Psi:=-\Lambda^{-1}[h+f]$.
For \eqref{p0:eq:core-triangular}, note as in the proof of that lemma that
$e_n$ occurs neither in $\Coef{t^n}h(E)$ nor in $\Coef{t^n}f(t,E)$, so both may
be computed from $E_{<n}$.  For \eqref{p0:eq:core-bell} apply
\cref{p0:def:obell} to $h$ and \eqref{p0:eq:power-negative} to each
$(1+E)^{-\mu_k}$, then convolve with $\phi_k$; the displayed ranges are those
outside which the ordinary Bell polynomials vanish.
\end{proof}

\begin{remark}[How the chapters instantiate it]\label{p0:rem:core-instances}
For the linear--logarithmic core the shorter route of
\cref{p0:thm:lambert-centered} is available, because measuring the displacement
in units of $X$ makes the entire perturbation divisible by $t$.
\Cref{p2:sec:top} needs the general form: there $\Lambda$ is the normalized
core slope $\log X$ of \cref{p0:prop:factorial-core},
$h(E)=(1+E)\log(1+E)-E$ and
$f(t,E)=\sum_{k\ge1}\alpha_kt^{2k}(1+E)^{-(2k-1)}$, and
\eqref{p0:eq:core-bell} is exactly the Bell recurrence of
\textsf{double\_factorial\_transseries-2}.  The coefficient ring there is
$\Rat[\Lambda^{-1}]$: the $e_n$ are \emph{not} constants but polynomials in
$\Lambda^{-1}$, of degree at most $2n-1$ and with no constant term.
\end{remark}

\subsection{The perturbation-operator form}

\begin{proposition}[Inverse perturbation operator]\label{p0:prop:operator-form}
Let $\Phi_0$ be a phase with $\Phi_0'$ nonvanishing near the point considered,
let $X$ solve $\Phi_0(X)=L$, and consider
\begin{equation}\label{p0:eq:perturbed-problem}
 L=\Phi_0(x)+\varepsilon\,\varrho(x).
\end{equation}
Put $\mathscr D_0:=\Phi_0'(X)^{-1}\,\dd/\dd X$.  Then, as a formal power series
in $\varepsilon$,
\begin{equation}\label{p0:eq:operator-series}
 x=X+\sum_{m\ge1}\frac{(-\varepsilon)^m}{m!}\,
   \mathscr D_0^{\,m-1}\!\left[\frac{\varrho(X)^m}{\Phi_0'(X)}\right].
\end{equation}
\end{proposition}

\begin{proof}
Change coordinate to $\lambda:=\Phi_0(x)$, so that $x=\Phi_0^{-1}(\lambda)$ and
\eqref{p0:eq:perturbed-problem} becomes
$L=\lambda+\varepsilon\varrho_0(\lambda)$ with
$\varrho_0:=\varrho\circ\Phi_0^{-1}$.  This is a near-identity reversion in
$\lambda$ with small parameter $\varepsilon$; Lagrange--B\"urmann in the form
\eqref{p0:eq:LB-general}, with the observable $G=\Phi_0^{-1}$, gives
\[
 \Phi_0^{-1}(\lambda)
 =\Phi_0^{-1}(L)+\sum_{m\ge1}\frac{(-\varepsilon)^m}{m!}
  \left(\frac{\dd}{\dd L}\right)^{m-1}
  \Bigl[(\Phi_0^{-1})'(L)\,\varrho_0(L)^m\Bigr].
\]
Finally $(\Phi_0^{-1})'(L)=1/\Phi_0'(X)$, $\varrho_0(L)=\varrho(X)$ and
$\dd/\dd L=\mathscr D_0$ by the chain rule.
\end{proof}

\begin{remark}[What the operator form is for]\label{p0:rem:operator}
This is the form in which an \emph{exponentially small} forward correction is
transported through the inverse map: with $\varrho$ a Stokes discontinuity or a
terminant remainder rather than a power series, $-\varepsilon\varrho(X)/\Phi_0'(X)$
is the leading displacement and the higher terms generate the nonlinear sectors.
As a statement about a formal series in $\varepsilon$ it is exact; as an
asymptotic statement about a particular $\varrho$ it needs the hypotheses of
\cref{p0:thm:remainder-transport}.  Expanding $\varrho$ in inverse powers of $X$
reproduces \eqref{p0:eq:c-lagrange}.
\end{remark}

\section{Flattening the core}\label{p0:sec:flattening}

\Cref{p0:thm:lambert-centered} hides the explicit scale of the answer inside
$X$.  Unfolding $X$ produces the polynomial--logarithmic transseries in which
the four chapters state their headline formulae.  Throughout this section $H$
is an indeterminate; in the analytic application it is specialized to
$\log(L/a)$.

\begin{theorem}[Direct polynomial--logarithmic reversion]\label{p0:thm:flattening}
Let $R$ be a commutative $\Rat$-algebra, $a\in R$ invertible, $b\in R$,
$Q=\sum_{j\ge1}q_jt^j\in tR[[t]]$, and let $H$ be an indeterminate.  Set
\begin{equation}\label{p0:eq:Phi-direct}
 \Phi(t,u;H):=-b\log\bigl(1+t(u-bH)\bigr)
   -Q\!\left(\frac{at}{1+t(u-bH)}\right)
 \ \in\ t\,R[H][[t,u]].
\end{equation}
\begin{enumerate}
\item There is a unique $V=\sum_{n\ge1}P_n(H)\,t^n\in t\,R[H][[t]]$ with
      $V=\Phi(t,V;H)$, and
      \begin{equation}\label{p0:eq:P-lagrange}
       P_n(H)=\sum_{m=1}^{n}\frac1m\Coef{t^nu^{m-1}}\Phi(t,u;H)^m .
      \end{equation}
\item \emph{Triangular Bell recurrence.}  Put $\gamma_1:=-bH$ and
      $\gamma_r:=P_{r-1}(H)$ for $r\ge2$.  Then
      \begin{equation}\label{p0:eq:P-bell}
       P_n(H)=-b\sum_{k=1}^{n}\frac{(-1)^{k+1}}{k}\OB_{n,k}(\gamma)
        -\sum_{j=1}^{n}q_ja^{\,j}\sum_{k=0}^{n-j}(-1)^k\binom{j+k-1}{k}
          \OB_{n-j,k}(\gamma),
      \end{equation}
      in which only $P_1,\dots,P_{n-1}$ occur.
\item \emph{Shape.}  $P_n\in R[H]$ with
      \begin{equation}\label{p0:eq:P-shape}
       \deg_HP_n\le n,\qquad
       \Coef{H^n}P_n=\frac{b^{\,n+1}}{n},\qquad
       P_n(0)=a^{\,n+1}c_n ,
      \end{equation}
      where $c_n$ are the Lambert-centred coefficients of
      \cref{p0:thm:lambert-centered}.  In particular $\deg_HP_n=n$ exactly
      when $b$ is not a zero divisor, and the coefficients of $P_n$ lie in the
      $\Rat$-subalgebra generated by $a$, $a^{-1}$, $b$, $q_1,\dots,q_n$.
\item \emph{Formal exactness and uniqueness.}  Put $t=L^{-1}$ and
      $H=\log(L/a)$, and
      \begin{equation}\label{p0:eq:direct-ansatz}
       x:=\frac1a\left[L-bH+\sum_{n\ge1}\frac{P_n(H)}{L^{n}}\right].
      \end{equation}
      Then $ax+b\log x+Q(1/x)=L$ as an identity of formal series in $L^{-1}$
      with polynomial coefficients in $H$.  Conversely, \eqref{p0:eq:direct-ansatz}
      is the \emph{only} expansion of the form
      $x=a^{-1}\bigl[L-bH+\sum_{n\ge1}\Pi_n(H)L^{-n}\bigr]$,
      with $\Pi_n\in R[H]$, that solves the equation coefficientwise.
\item \emph{Analytic version.}  Under the hypotheses of
      \cref{p0:thm:lambert-centered}(5), with $L=\log(y/C)$ and
      $H=\log(L/a)$, for every fixed $N\ge0$,
      \begin{equation}\label{p0:eq:direct-asymptotic}
       F^{-1}(y)=\frac1a\left[L-bH+\sum_{n=1}^{N}\frac{P_n(H)}{L^{n}}\right]
        +\bigO\!\left(\frac{(1+H)^{N+1}}{L^{N+1}}\right)
       \qquad(y\to\infty).
      \end{equation}
\end{enumerate}
\end{theorem}

\begin{proof}
(1) Membership of $\Phi$ in $tR[H][[t,u]]$ is checked as in
\cref{p0:thm:lambert-centered}(1), the extra additive term $-bH$ inside
$1+t(u-bH)$ being harmless because it is multiplied by $t$.
\Cref{p0:lem:lagrange-fp} over the ring $R[H]$ then gives existence,
uniqueness and \eqref{p0:eq:P-lagrange}.

(2) Put $w:=t(V-bH)=\sum_{r\ge1}\gamma_rt^r$ with $\gamma$ as stated; apply
\eqref{p0:eq:log-obell} to $\log(1+w)$ and \eqref{p0:eq:power-negative} to
$(1+w)^{-j}$, and convolve the latter with the factor $(at)^j$.  Only
$\gamma_1,\dots,\gamma_{n-k+1}$, that is $H$ and $P_1,\dots,P_{n-k}$, occur.

(4) Set $t=L^{-1}$.  With $x$ as in \eqref{p0:eq:direct-ansatz},
\[
 x=\frac1{at}\bigl(1+t(V-bH)\bigr)=\frac{1+w}{at},
 \qquad ax=\frac1t+(V-bH),\qquad
 \log x=H+\log(1+w),\qquad \frac1x=\frac{at}{1+w},
\]
where $\log x$ uses $\log\bigl(1/(at)\bigr)=\log(L/a)=H$.  Hence
\[
 ax+b\log x+Q(1/x)-L
 =(V-bH)+b\bigl(H+\log(1+w)\bigr)+Q\!\left(\frac{at}{1+w}\right)
 =V-\Phi(t,V;H)=0 .
\]
For uniqueness, substituting a general $\Pi=\sum\Pi_n(H)t^n$ in place of $V$
turns the equation into $\Pi=\Phi(t,\Pi;H)$, and
\cref{p0:lem:lagrange-fp} has a unique solution.

(3) We show by strong induction that $\deg_HP_n\le n$ and, in addition, that
$\deg_H\gamma_r\le r-1$ for $r\ge2$.  The base is
$\gamma_1=-bH$, of degree $1$.  Assume $\deg_HP_j\le j$ for all $j<n$, so
$\deg_H\gamma_r=\deg_HP_{r-1}\le r-1$ for $2\le r\le n$.  A monomial of
$\OB_{n,k}(\gamma)$ is a product $\gamma_{r_1}\cdots\gamma_{r_k}$ with
$r_1+\dots+r_k=n$ and $r_\nu\ge1$, hence of $H$-degree at most
$\sum_\nu r_\nu=n$; so the first sum of \eqref{p0:eq:P-bell} has $H$-degree at
most $n$.  In the second sum the ordinary Bell polynomial has first index
$n-j$ with $j\ge1$, so its $H$-degree is at most $n-1$.  Hence
$\deg_HP_n\le n$, completing the induction.

A monomial of $H$-degree exactly $n$ requires $\deg_H\gamma_{r_\nu}=r_\nu$ for
every $\nu$, which forces $r_\nu=1$ for all $\nu$ because
$\deg_H\gamma_r\le r-1$ for $r\ge2$.  Thus $k=n$, the monomial is
$\gamma_1^{\,n}=(-bH)^n$, and $\Coef{H^n}\OB_{n,k}(\gamma)=0$ for $k<n$ while
$\OB_{n,n}(\gamma)=(-bH)^n$.  Only the term $k=n$ of the first sum of
\eqref{p0:eq:P-bell} therefore survives, and
\[
 \Coef{H^n}P_n=-b\cdot\frac{(-1)^{n+1}}{n}\cdot(-b)^n
 =\frac{(-1)^{2n+2}\,b^{\,n+1}}{n}=\frac{b^{\,n+1}}{n}.
\]
For the constant term put $H=0$ in \eqref{p0:eq:Phi-direct}:
\[
 \Phi(t,u;0)=-b\log(1+tu)-Q\!\left(\frac{at}{1+tu}\right).
\]
If $U(t)=\sum c_nt^n$ solves $U=\Psi(t,U)$ with $\Psi$ of
\eqref{p0:eq:Psi-general}, put $\widetilde V(t):=a\,U(at)$.  Then
$t\widetilde V(t)=(at)U(at)$, so
\[
 \Phi\bigl(t,\widetilde V(t);0\bigr)
 =-b\log\bigl(1+(at)U(at)\bigr)-Q\!\left(\frac{at}{1+(at)U(at)}\right)
 =a\,\Psi\bigl(at,U(at)\bigr)=a\,U(at)=\widetilde V(t),
\]
so $\widetilde V$ solves the $H=0$ fixed-point equation; by uniqueness
$P_n(0)=\Coef{t^n}\widetilde V=a\cdot a^{n}c_n=a^{n+1}c_n$.

(5) Fix $N$.  Put $x_N$ for the truncation displayed in
\eqref{p0:eq:direct-asymptotic}.  Exactly as in the proof of
\cref{p0:thm:lambert-centered}(5), the residual
$r_N:=\log F(x_N)-\log y$ splits as $g(t)+E_N$ with
$|E_N|\le K_Nx_N^{-N-1}$ and $g$ analytic in $t=1/L$ \emph{with coefficients
polynomial in $H$}; by (4) the Taylor coefficients of $g$ vanish through order
$N$, and the coefficient of $t^{N+1}$ is a polynomial in $H$ of degree at most
$N+1$.  Hence $|r_N|=\bigO\bigl((1+H)^{N+1}L^{-N-1}\bigr)$.  The slope
bound $(\log F)'\ge a/2$ and \cref{p0:thm:backward-error} finish the proof,
after noting $L\asymp aX$ and $H=\log(L/a)$ so that $x_N\to\infty$.
\end{proof}

\begin{theorem}[The pure Lambert block in closed form]\label{p0:thm:pure-stirling}
Let $\stirlingfirst{n}{m}$ be the \emph{unsigned} Stirling numbers of the first
kind, so that the signed ones are
$(-1)^{n-m}\stirlingfirst{n}{m}$ and
$\log(1+z)^m=m!\sum_{n\ge m}(-1)^{n-m}\stirlingfirst{n}{m}z^n/n!$.  When $Q=0$
the polynomials of \cref{p0:thm:flattening} are
\begin{equation}\label{p0:eq:pure-stirling}
 \boxed{\;
 P_n^{(0)}(H)=b^{\,n+1}\sum_{m=1}^{n}(-1)^{m+1}
   \frac{\stirlingfirst{n}{m}}{(n-m+1)!}\,H^{\,n-m+1}\;}
\end{equation}
for $n\ge1$.  Equivalently, clearing the binomial,
\begin{equation}\label{p0:eq:pure-stirling-unsigned}
 P_n^{(0)}(H)=\frac{b^{\,n+1}}{n!}\sum_{m=1}^{n}(-1)^{m+1}(m-1)!\,
  \stirlingfirst{n}{m}\binom{n}{m-1}H^{\,n-m+1}.
\end{equation}
In particular
\begin{equation}\label{p0:eq:pure-lead}
 \Coef{H^n}P_n^{(0)}=\frac{b^{\,n+1}}{n},
 \qquad P_n^{(0)}(0)=0,
\end{equation}
and the first four are
\begin{equation}\label{p0:eq:pure-first}
 P_1^{(0)}=b^2H,\quad
 P_2^{(0)}=\frac{b^3}{2}\bigl(H^2-2H\bigr),\quad
 P_3^{(0)}=\frac{b^4}{6}\bigl(2H^3-9H^2+6H\bigr),\quad
 P_4^{(0)}=\frac{b^5}{12}\bigl(3H^4-22H^3+36H^2-12H\bigr).
\end{equation}
Consequently the Lambert core itself unfolds as
\begin{equation}\label{p0:eq:core-unfolded}
 X=\frac1a\left[L-bH+\sum_{n\ge1}\frac{P_n^{(0)}(H)}{L^n}\right],
 \qquad H=\log\frac La ,
\end{equation}
formally, and with remainder $\bigO\bigl((1+H)^{N+1}L^{-N-1}\bigr)$ after
truncation at $n=N$.
\end{theorem}

\begin{proof}
With $Q=0$, \eqref{p0:eq:P-lagrange} reads
\[
 P_n^{(0)}(H)=\sum_{m=1}^{n}\frac{(-b)^m}{m}
  \Coef{t^nu^{m-1}}\log\bigl(1+t(u-bH)\bigr)^m .
\]
Insert $\log(1+z)^m=m!\sum_{r\ge m}(-1)^{r-m}\stirlingfirst{r}{m}z^r/r!$ with
$z=t(u-bH)$; the coefficient of $t^n$ is
$m!\,(-1)^{n-m}\stirlingfirst{n}{m}(u-bH)^n/n!$, and
$\Coef{u^{m-1}}(u-bH)^n=\binom{n}{m-1}(-bH)^{n-m+1}$.  Hence
\[
 P_n^{(0)}(H)
 =\sum_{m=1}^{n}\frac{(-b)^m}{m}\cdot
   \frac{m!\,(-1)^{n-m}\stirlingfirst{n}{m}}{n!}
   \binom{n}{m-1}(-b)^{n-m+1}H^{n-m+1},
\]
and $(-b)^{m}(-b)^{n-m+1}(-1)^{n-m}=(-1)^{n+1}(-1)^{n-m}b^{n+1}
=(-1)^{m+1}b^{n+1}$, which gives
\eqref{p0:eq:pure-stirling-unsigned}.  Since
$(m-1)!\binom{n}{m-1}/n!=1/(n-m+1)!$, the two displayed forms agree and
\eqref{p0:eq:pure-stirling} follows.  Every term carries $H^{n-m+1}$ with
$m\le n$, so no constant term occurs and $P_n^{(0)}(0)=0$; the term $m=1$ has
$\stirlingfirst{n}{1}=(n-1)!$ and gives
$\Coef{H^n}P_n^{(0)}=b^{n+1}(n-1)!/n!=b^{n+1}/n$.
Finally \eqref{p0:eq:core-unfolded} is \cref{p0:thm:flattening}(4) and (5)
with $Q=0$, in which case the unique solution of $ax+b\log x=L$ is the core
$X$ of \cref{p0:thm:lambert-core}.
\end{proof}

\begin{remark}[Three source formulae, one identity]\label{p0:rem:stirling-crosswalk}
The closed form \eqref{p0:eq:pure-stirling} appears in the sources in three
different dresses, all with the signed Stirling numbers
$s(n,m)=(-1)^{n-m}\stirlingfirst{n}{m}$.
\textsf{Swing\_Factorial\_Transseries\_Article} states
\eqref{p0:eq:pure-stirling-unsigned};
\textsf{Partition\_Number\_Transseries\_and\_Asymptotic\_Inverse} states
$P_j(L)=2^{j+1}\sum_{q\le j}s(j,q)L^{j+1-q}/(j+1-q)!$;
\textsf{Partition\_Number\_Transseries\_and\_Inverse} states
$P_m(\ell)=\frac1{m!}\sum_j2^j(j-1)!\,s(m,j)\binom m{j-1}(2\ell)^{m-j+1}$.
The last two are the specialization $a=1$, $b=-2$ of
\eqref{p0:eq:pure-stirling}, and all three agree with it as exact rational
polynomials for $1\le n\le8$, which was checked in exact rational arithmetic
(\cref{p0:rem:verification}).  The volume uses
\eqref{p0:eq:pure-stirling}: it is the shortest of the four, it avoids the
letter $s$, which \cref{p3:sec:top} needs for Dedekind sums, and it makes the
leading coefficient \eqref{p0:eq:pure-lead} a one-line consequence.
\end{remark}

\begin{remark}[Both normalizations, and what is and is not claimed]
\label{p0:rem:accuracy-claim}
The Lambert-centred expansion \eqref{p0:eq:centered-asymptotic} and the
flattened expansion \eqref{p0:eq:direct-asymptotic} are two truncations of the
same object: by \eqref{p0:eq:core-unfolded}, expanding $X$ and then each power
$X^{-n}$ in the flattened scale returns \eqref{p0:eq:direct-asymptotic}, and
conversely collecting in the flattened series all terms generated by
$b\log x$ alone rebuilds $X$.  What differs is the coefficient ring: the
$c_n$ are constants, the $P_n$ are polynomials of degree $n$ in
$H=\log(L/a)$.  Several of the sources add that Lambert centring is
numerically more accurate at equal truncation order and support this with
tables.  Part~0 makes no such claim: the comparison depends on the size of $H$
relative to $L$ and on the particular data, it is an empirical statement about
specific truncations, and the supporting tables belong to the chapters, where
they are reported with their provenance.  What Part~0 does assert is the
algebraic statement of \cref{p0:rem:no-logs}, and the remainder bounds
\eqref{p0:eq:centered-asymptotic} and \eqref{p0:eq:direct-asymptotic}, whose
right-hand sides differ by the factor $(1+H)^{N+1}$ --- and that factor is the
only rigorous difference between the two at fixed order.
\end{remark}

\section{The three inverses of a discrete sequence}\label{p0:sec:discrete}

\begin{definition}[The three inverse objects]\label{p0:def:three-inverses}
Let $(A_n)_{n\ge n_0}$ be real numbers, strictly increasing for $n\ge n_1\ge n_0$,
with $A_n\to+\infty$.
\begin{enumerate}
\item The \emph{integer staircase}, or generalized inverse, is
      \begin{equation}\label{p0:eq:staircase-def}
       N_\ast(y):=\min\{n\ge n_1:A_n\ge y\},\qquad y\le\sup_nA_n .
      \end{equation}
      It is an integer-valued, nondecreasing step function, constant on each
      interval $(A_{n-1},A_n]$.
\item Given a continuous, strictly increasing
      $\mathcal A:[n_1,\infty)\to\Real$ with $\mathcal A(n)=A_n$ for every
      integer $n\ge n_1$ --- an \emph{admissible interpolation} --- the
      \emph{interpolated inverse} is $\nu_{\mathcal A}:=\mathcal A^{-1}$,
      defined on $[A_{n_1},\infty)$.
\item The \emph{asymptotic inverse} is the formal transseries produced by
      \cref{p0:thm:lambert-centered} or \cref{p0:thm:flattening} from a model
      $F$ for which the chapter in question proves an asymptotic identity
      along the sequence.  It is attached to no particular interpolation.
\end{enumerate}
\end{definition}

\begin{remark}[The three are genuinely different]\label{p0:rem:three-differ}
Object (1) is canonical but not smooth; object (2) is smooth but not canonical
--- distinct admissible interpolations agree at the integers and differ in
between, hence have different inverses; object (3) is canonical as a formal
series but determines a function only after a summation convention is fixed.
The three coincide only in the following precise senses:
$N_\ast=\lceil\nu_{\mathcal A}\rceil$ for \emph{every} admissible
$\mathcal A$ (\cref{p0:thm:staircase}(1)), and
$\nu_{\mathcal A}(A_n)=n$ for every admissible $\mathcal A$ and every integer
$n\ge n_1$.  The asymptotic series (3) is independent of the choice in
$\mathcal A$ at every algebraic order, but the exponentially small layers are
not: the partition articles make this explicit for the Rademacher amplitudes,
which are periodic in the index and therefore genuinely depend on how the
sequence is continued off the lattice.
\end{remark}

\begin{theorem}[Staircase recovery and the separation condition]
\label{p0:thm:staircase}
Notation as in \cref{p0:def:three-inverses}, with $\mathcal A$ admissible and
$\nu:=\nu_{\mathcal A}$.
\begin{enumerate}
\item \emph{Exact rounding.}  For every $y$ in the range of $\mathcal A$,
      that is for $A_{n_1}\le y<\sup_nA_n$,
      \begin{equation}\label{p0:eq:ceil-identity}
       N_\ast(y)=\bigl\lceil\nu(y)\bigr\rceil,
      \end{equation}
      with $\lceil m\rceil=m$ for integers $m$.  The identity holds for every
      admissible interpolation.
\item \emph{Separation condition.}  Let $x_\ast\in\Real$ and $\eta\ge0$ satisfy
      $|x_\ast-\nu(y)|\le\eta$.  If
      \begin{equation}\label{p0:eq:separation}
       \bigl\lceil x_\ast-\eta\bigr\rceil=\bigl\lceil x_\ast+\eta\bigr\rceil,
      \end{equation}
      then $N_\ast(y)$ equals that common value.  Condition
      \eqref{p0:eq:separation} holds if and only if the interval
      $[x_\ast-\eta,x_\ast+\eta]$ is contained in a single half-open interval
      $(m-1,m]$; it can fail for arbitrarily small $\eta>0$, namely when
      $\nu(y)$ is within $\eta$ of an integer from above.
\item \emph{Inverse along the range.}  If $y=A_n$ for an integer $n\ge n_1$ and
      $|x_\ast-\nu(y)|\le\eta<\tfrac12$, then $n=\lfloor x_\ast+\tfrac12\rfloor$.
      This case is unconditional: no separation hypothesis is needed, because
      $\nu(A_n)=n$ exactly.
\item \emph{Residue classes.}  Suppose $r\ge1$ and that for each
      $\rho\in\{0,1,\dots,r-1\}$ the subsequence $(A_n)_{n\equiv\rho\ (r)}$ is
      strictly increasing for $n\ge n_1$ and admits an admissible
      interpolation $\mathcal A_\rho$ on the class, with inverse $\nu_\rho$.
      Then, for $y$ large enough,
      \begin{equation}\label{p0:eq:residue-staircase}
       N_\ast(y)=\min_{0\le\rho<r}
        \left(\rho+r\left\lceil\frac{\nu_\rho(y)-\rho}{r}\right\rceil\right).
      \end{equation}
\item \emph{No uniform smooth representation.}  For every $Y$ in the range and
      every continuous $g:[Y,\infty)\to\Real$,
      \begin{equation}\label{p0:eq:no-smooth}
       \sup_{y\ge Y}\bigl|N_\ast(y)-g(y)\bigr|\ \ge\ \tfrac12 .
      \end{equation}
      Equivalently, the $\bigO(1)$ periodic layer
      $\lceil x\rceil-x-\tfrac12=\pi^{-1}\sum_{m\ge1}m^{-1}\sin(2\pi mx)$
      (valid for $x\notin\Int$) is not a refinement that can be dropped: it is
      the leading discrete transmonomial of the generalized inverse.
\end{enumerate}
\end{theorem}

\begin{proof}
(1) For an integer $n\ge n_1$, $A_n\ge y$ is equivalent to
$\mathcal A(n)\ge y$, hence to $n\ge\nu(y)$ by strict monotonicity of
$\mathcal A$.  The least such integer is $\lceil\nu(y)\rceil$.

(2) The function $\lceil\cdot\rceil$ is constant on each $(m-1,m]$ and
nondecreasing; since $\nu(y)\in[x_\ast-\eta,x_\ast+\eta]$,
\eqref{p0:eq:separation} forces $\lceil\nu(y)\rceil$ to equal the common
value, and (1) applies.  The characterization of
\eqref{p0:eq:separation} and the failure mode are immediate from the same
description of $\lceil\cdot\rceil$.

(3) $\nu(A_n)=n$ because $\mathcal A(n)=A_n$; hence $|x_\ast-n|\le\eta<1/2$ and
rounding to the nearest integer returns $n$.

(4) Fix $\rho$.  Among the integers $n\equiv\rho\pmod r$ with $n\ge n_1$, the
condition $A_n\ge y$ is, by (1) applied to the class, equivalent to
$n\ge\nu_\rho(y)$.  Writing $n=\rho+rk$ with $k\in\Int$, the least admissible
$k$ is $\lceil(\nu_\rho(y)-\rho)/r\rceil$, so the least admissible $n$ in the
class is $\rho+r\lceil(\nu_\rho(y)-\rho)/r\rceil$.  Taking the minimum over
$\rho$ gives the least admissible $n$ without a congruence restriction, which
is $N_\ast(y)$.

(5) Fix $n$ with $A_n>Y$.  On $(A_{n-1},A_n]$ we have $N_\ast=n$ and on
$(A_n,A_{n+1}]$ we have $N_\ast=n+1$.  For $\epsilon>0$ small,
\[
 1=\bigl|N_\ast(A_n+\epsilon)-N_\ast(A_n)\bigr|
 \le\bigl|N_\ast(A_n+\epsilon)-g(A_n+\epsilon)\bigr|
  +\bigl|g(A_n+\epsilon)-g(A_n)\bigr|
  +\bigl|g(A_n)-N_\ast(A_n)\bigr| .
\]
Letting $\epsilon\downarrow0$ and using continuity of $g$ gives
$1\le2\sup_{y\ge Y}|N_\ast-g|$.  For the Fourier identity, the function
$x\mapsto\lceil x\rceil-x-\tfrac12$ is $1$-periodic, equal to
$-(x-\lfloor x\rfloor)+\tfrac12$ shifted, i.e.\ minus the standard sawtooth,
whose Fourier series is the displayed one and converges pointwise off
$\Int$.
\end{proof}

\begin{proposition}[Certified discrete inversion]\label{p0:prop:certified-rounding}
In the setting of \cref{p0:thm:staircase}(4), suppose that for each $\rho$ an
approximation $x_\rho$ with $|x_\rho-\nu_\rho(y)|\le\eta_\rho$ is available.
Then $N_\ast(y)$ is determined by at most
$\sum_\rho\bigl(\lfloor2\eta_\rho/r\rfloor+2\bigr)$ exact evaluations of the
sequence: for each $\rho$ list the integers $n\equiv\rho\pmod r$ in
$[x_\rho-\eta_\rho,\,x_\rho+\eta_\rho+r]$, evaluate $A_n$, and take the least
$n$ over all classes with $A_n\ge y$.
\end{proposition}

\begin{proof}
By \cref{p0:thm:staircase}(4) the class-$\rho$ candidate is
$\rho+r\lceil(\nu_\rho(y)-\rho)/r\rceil$, and $\nu_\rho(y)$ lies in
$[x_\rho-\eta_\rho,x_\rho+\eta_\rho]$; the least integer of the class that is
$\ge\nu_\rho(y)$ therefore lies in $[x_\rho-\eta_\rho,x_\rho+\eta_\rho+r]$.
Since the class subsequence is increasing, the test $A_n\ge y$ identifies it,
and the outer minimum over $\rho$ is \eqref{p0:eq:residue-staircase}.
\end{proof}

\begin{remark}[The parity instances]\label{p0:rem:parity-instances}
\Cref{p4:sec:top} and \cref{p2:sec:top} are the case $r=2$.  The formula
$\rho+r\lceil(\nu_\rho-\rho)/r\rceil$ specializes to
$2\lceil\nu_0/2\rceil$ and $2\lceil(\nu_1-1)/2\rceil+1$, which is exactly the
pair of parity candidates of \textsf{double\_factorial\_transseries-2}, and
$N_\ast=\min$ of the two is its ``exact generalized inverse'' theorem.
\textsf{Swing\_Factorial\_Transseries\_Article} states the corresponding
\emph{floor} formulae for the threshold query ``the greatest index in the class
$\rho$ whose sequence value is at most $y$''; those are
$\rho+r\lfloor(\nu_\rho-\rho)/r\rfloor$, obtained from
\eqref{p0:eq:residue-staircase} by replacing $\lceil\cdot\rceil$ with
$\lfloor\cdot\rfloor$ and $\min$ with $\max$, with the same proof.  The
partition chapter is the case $r=1$, where
\eqref{p0:eq:residue-staircase} reduces to \eqref{p0:eq:ceil-identity}.
\end{remark}

\begin{repairbox}
\textbf{Source claim.}  Several sources write, or invite the reader to write,
$N_\ast(y)=\nu(y)+\littleo(1)$.
\textbf{Defect.}  This is false uniformly in $y$, by
\cref{p0:thm:staircase}(5).  \textsf{Butcher\_Tree\_Transseries-2} already
flags it (``An expansion such as $N_\ast(y)=\nu(y)+\littleo(1)$ is
false uniformly in $y$''), and
\textsf{Partition\_Number\_Transseries\_and\_Asymptotic\_Inverse} likewise
separates the three objects; but the qualification is absent from the summary
statements of the other articles.
\textbf{Correction.}  The volume uses only the three statements
\cref{p0:thm:staircase}(1), (3), (4): an exact ceiling identity valid for every
admissible interpolation, an unconditional nearest-integer recovery
\emph{along the range} $y=A_n$, and a residue-class refinement.  Nothing else
about $N_\ast$ is asserted.
\end{repairbox}

\section{Error control, transport, and optimal truncation}
\label{p0:sec:errors}

\subsection{Backward error certifies forward error}

\begin{theorem}[Residual-to-root conversion]\label{p0:thm:backward-error}
Let $I\subseteq\Real$ be an interval and $f:I\to\Real$ differentiable with
\begin{equation}\label{p0:eq:slope-bounds}
 0<m\le f'(\xi)\le M\qquad(\xi\in I).
\end{equation}
Let $x_\dagger\in I$ satisfy $f(x_\dagger)=\lambda$, let $x_\ast\in I$, and set
the \emph{residual} $r_\ast:=f(x_\ast)-\lambda$.  Then
\begin{equation}\label{p0:eq:backward-two-sided}
 \frac{|r_\ast|}{M}\ \le\ |x_\ast-x_\dagger|\ \le\ \frac{|r_\ast|}{m} .
\end{equation}
Moreover, if $f$ is continuous on $I$ and
$[x_\ast-\varrho,\,x_\ast+\varrho]\subseteq I$ with $\varrho:=|r_\ast|/m$,
then $f$ takes the value $\lambda$ somewhere in that interval; the residual is
therefore a \emph{certificate}, not merely a bound conditional on the
existence of a nearby root.
\end{theorem}

\begin{proof}
By the mean value theorem, $r_\ast=f(x_\ast)-f(x_\dagger)=f'(\xi)(x_\ast-x_\dagger)$
for some $\xi$ between the two points, and \eqref{p0:eq:slope-bounds} gives
\eqref{p0:eq:backward-two-sided}.  For the certificate, suppose $r_\ast>0$
(the case $r_\ast<0$ is symmetric and $r_\ast=0$ is trivial).  Then
$f(x_\ast)=\lambda+r_\ast>\lambda$, while
\[
 f(x_\ast-\varrho)\le f(x_\ast)-m\varrho=\lambda+r_\ast-|r_\ast|=\lambda,
\]
using $f'\ge m$.  The intermediate value theorem supplies a point of
$[x_\ast-\varrho,x_\ast]$ where $f=\lambda$.
\end{proof}

\begin{remark}[Use in the logarithmic variable]\label{p0:rem:log-residual}
In every application $f=\log F$ and $\lambda=\log y$: $F$ itself overflows any
fixed-precision arithmetic long before the asymptotic regime begins, whereas
$\log F$ is computable from a digamma or Laplace representation throughout.
The lower bound in \eqref{p0:eq:backward-two-sided} makes the residual test
sharp --- a small residual is necessary as well as sufficient.
\end{remark}

\subsection{Transport of a controlled forward remainder}

\begin{theorem}[Remainder transport]\label{p0:thm:remainder-transport}
Let $I=[x_1,\infty)$, let $f_0:I\to\Real$ be differentiable with
$f_0'\ge m>0$, and let $\varepsilon:I\to\Real$ be differentiable with
$\sup_I|\varepsilon'|\le\theta<m$.  Put $f:=f_0+\varepsilon$.  For $\lambda$ in
the common range let $x_0=x_0(\lambda)\in I$ solve $f_0(x_0)=\lambda$ and
$x=x(\lambda)\in I$ solve $f(x)=\lambda$.  Then
\begin{enumerate}
\item \emph{Rigorous bound.}
      \begin{equation}\label{p0:eq:transport-bound}
       \bigl|x(\lambda)-x_0(\lambda)\bigr|
       \le\frac{\bigl|\varepsilon\bigl(x_0(\lambda)\bigr)\bigr|}{m-\theta}.
      \end{equation}
\item \emph{First-order law with explicit error.}  If moreover $f_0\in C^2(I)$
      with $\sup_I|f_0''|\le M_2$, then
      \begin{equation}\label{p0:eq:transport-first-order}
       x(\lambda)-x_0(\lambda)
       =-\frac{\varepsilon(x_0)}{f_0'(x_0)}+E,
       \qquad
       |E|\le\frac{|\varepsilon(x_0)|}{m\,(m-\theta)}
        \left(\frac{M_2\,|\varepsilon(x_0)|}{m-\theta}+\theta\right).
      \end{equation}
\end{enumerate}
In particular, if $\varepsilon(x)=\bigO(x^{-M-1})$ and
$\varepsilon'(x)=\littleo(1)$ as $x\to\infty$, then
$x(\lambda)-x_0(\lambda)=-\varepsilon(x_0)/f_0'(x_0)\,(1+\littleo(1))
=\bigO(x_0^{-M-1})$.
\end{theorem}

\begin{proof}
Since $f'=f_0'+\varepsilon'\ge m-\theta>0$, $f$ is strictly increasing and
$x(\lambda)$ is unique.  The residual of $x_0$ in $f$ is
$f(x_0)-\lambda=f_0(x_0)+\varepsilon(x_0)-\lambda=\varepsilon(x_0)$, so
\cref{p0:thm:backward-error} applied to $f$ with $m$ replaced by $m-\theta$
gives \eqref{p0:eq:transport-bound}.

For (2), write $d:=x-x_0$.  By the mean value theorem there is $\xi$ between
$x_0$ and $x$ with $0=f(x)-\lambda=\varepsilon(x_0)+f'(\xi)d$, so
$d=-\varepsilon(x_0)/f'(\xi)$ and
\[
 d+\frac{\varepsilon(x_0)}{f_0'(x_0)}
 =\varepsilon(x_0)\left(\frac1{f_0'(x_0)}-\frac1{f'(\xi)}\right)
 =\varepsilon(x_0)\,\frac{f'(\xi)-f_0'(x_0)}{f_0'(x_0)\,f'(\xi)} .
\]
Now $|f'(\xi)-f_0'(x_0)|\le|f_0'(\xi)-f_0'(x_0)|+|\varepsilon'(\xi)|
\le M_2|d|+\theta\le M_2|\varepsilon(x_0)|/(m-\theta)+\theta$ by
\eqref{p0:eq:transport-bound}, while $f_0'(x_0)\ge m$ and
$f'(\xi)\ge m-\theta$.  Substituting gives
\eqref{p0:eq:transport-first-order}.  The final assertion follows since then
$\theta\to0$ and $|\varepsilon(x_0)|\to0$.
\end{proof}

\begin{corollary}[Forward truncation to inverse accuracy]
\label{p0:cor:forward-to-inverse}
Let $F$ be differentiably of exponential--power type with data
$(C,a,1,b,1,Q)$, strictly increasing on $[x_1,\infty)$ with
$(\log F)'\ge m>0$ there.  Let $x_\ast$ solve exactly the truncated equation
\[
 \log C+ax_\ast+b\log x_\ast+\sum_{j=1}^{M}q_jx_\ast^{-j}=\log y .
\]
Then $|x_\ast-F^{-1}(y)|\le K_Mx_\ast^{-M-1}/m$ with $K_M$ the constant of
\eqref{p0:eq:model-analytic}.
\end{corollary}

\begin{proof}
The residual of $x_\ast$ in $\log F$ is exactly the forward remainder, of
modulus at most $K_Mx_\ast^{-M-1}$; apply
\cref{p0:thm:backward-error}.
\end{proof}

\subsection{Least-term truncation and the beyond-all-orders floor}

\begin{theorem}[Optimal truncation and the exponential floor]
\label{p0:thm:optimal-truncation}
Let $A>0$, $\beta\in\Real$, $K>0$, $x_0\ge1$, and let
$\varphi:[x_0,\infty)\to\Real$ satisfy the \emph{Gevrey-$1$ remainder bound}
\begin{equation}\label{p0:eq:gevrey-remainder}
 \Bigl|\varphi(x)-\sum_{j=1}^{M}q_jx^{-j}\Bigr|
 \ \le\ K\,\Gamma(M+\beta)\,A^{-M}x^{-M}
 \qquad(M\ge1,\ x\ge x_0).
\end{equation}
\begin{enumerate}
\item \emph{Least-term envelope.}  There are $x_1\ge x_0$ and
      $K'=K'(K,A,\beta)$ such that, choosing $M=M(x):=\lfloor Ax\rfloor$,
      \begin{equation}\label{p0:eq:least-term}
       \Bigl|\varphi(x)-\sum_{j=1}^{M(x)}q_jx^{-j}\Bigr|
       \ \le\ K'\,x^{\beta-1/2}\,\e^{-Ax}
       \qquad(x\ge x_1).
      \end{equation}
\item \emph{Transported inverse error.}  Let $F$ be as in
      \cref{p0:cor:forward-to-inverse}, with
      $\log F(x)=\log C+ax+b\log x+\varphi(x)$ and $(\log F)'\ge m>0$ on
      $[x_1,\infty)$.  If $x_\ast$ solves the optimally truncated equation
      exactly, then
      \begin{equation}\label{p0:eq:optimal-inverse}
       \bigl|x_\ast-F^{-1}(y)\bigr|
       \ \le\ \frac{K'}{m}\,x_\ast^{\beta-1/2}\,\e^{-Ax_\ast}.
      \end{equation}
\item \emph{The floor is real.}  Suppose in addition that there are
      $M_0\ge1$ and $0<k\le K'' $ with
      \begin{equation}\label{p0:eq:two-sided-coefficients}
       k\,\Gamma(M+\beta)A^{-M}\ \le\ |q_M|\ \le\ K''\,\Gamma(M+\beta)A^{-M}
       \qquad(M\ge M_0).
      \end{equation}
      Then
      \begin{equation}\label{p0:eq:floor}
       \log\ \min_{M\ge M_0}\ \bigl|q_M\bigr|x^{-M}
       \ =\ -Ax+\bigO(\log x)
       \qquad(x\to\infty),
      \end{equation}
      so that \emph{no} truncation of the algebraic series achieves an error
      of smaller exponential order than $\e^{-Ax}$.
\end{enumerate}
\end{theorem}

\begin{proof}
(1) Stirling's formula gives constants $C_\beta$ and $M_\beta\ge1$ with
$\Gamma(M+\beta)\le C_\beta M^{M+\beta-1/2}\e^{-M}$ for all $M\ge M_\beta$.
With $M=\lfloor Ax\rfloor$ and $x$ so large that $M\ge M_\beta$,
\[
 \Gamma(M+\beta)A^{-M}x^{-M}
 \le C_\beta M^{\beta-1/2}\e^{-M}\left(\frac{M}{Ax}\right)^{M}
 \le C_\beta M^{\beta-1/2}\e^{-M},
\]
because $M\le Ax$.  Since $M>Ax-1$, we have $\e^{-M}<\e\cdot\e^{-Ax}$, and
$M^{\beta-1/2}\le C_\beta'x^{\beta-1/2}$ for $x\ge x_1$ (with
$C_\beta'$ depending only on $A$ and $\beta$, since
$Ax-1\le M\le Ax$).  Combining with \eqref{p0:eq:gevrey-remainder} gives
\eqref{p0:eq:least-term}.

(2) The residual of $x_\ast$ in $\log F$ is exactly the left-hand side of
\eqref{p0:eq:least-term} evaluated at $x_\ast$; apply
\cref{p0:thm:backward-error}.

(3) Enlarging $M_0$ we may assume $M_0\ge M_\beta$.  Write $u:=Ax$ and
$g(M):=(M/u)^{M}\e^{-M}$, so that by Stirling there are $0<c_\beta\le C_\beta$
with
\begin{equation}\label{p0:eq:stirling-sandwich}
 c_\beta M^{\beta-1/2}g(M)\ \le\ \Gamma(M+\beta)(Ax)^{-M}
 \ \le\ C_\beta M^{\beta-1/2}g(M)
 \qquad(M\ge M_\beta).
\end{equation}

\emph{Upper bound.}  Take $M=\lfloor Ax\rfloor$ in the right-hand inequality of
\eqref{p0:eq:two-sided-coefficients} and repeat the estimate of part~(1); this
gives $\min_{M\ge M_0}|q_M|x^{-M}\le K'''x^{\beta-1/2}\e^{-Ax}$ once
$\lfloor Ax\rfloor\ge M_0$.

\emph{Lower bound.}  Note first that $\log g(M)=M\log(M/u)-M$ has derivative
$\log(M/u)$, so $g$ decreases on $(0,u)$, increases on $(u,\infty)$, and
\begin{equation}\label{p0:eq:g-min}
 g(M)\ \ge\ g(u)=\e^{-u}=\e^{-Ax}\qquad(M>0).
\end{equation}
Set $h(M):=M^{\beta-1/2}g(M)$; then
$(\log h)'(M)=(\beta-\tfrac12)/M+\log(M/u)$, which is positive as soon as
$M\ge2u$ and $2u\ge2|\beta-\tfrac12|/\log2$, so $h$ is increasing on
$[2u,\infty)$ for all large $x$.  Consequently
\[
 \min_{M\ge M_0}h(M)
 =\min_{M_0\le M\le2u}h(M)
 \ \ge\ \Bigl(\min_{M_0\le M\le2u}M^{\beta-1/2}\Bigr)\e^{-Ax}
 \ \ge\ c\,x^{-|\beta-1/2|}\,\e^{-Ax}
\]
for $x\ge x_1$, using \eqref{p0:eq:g-min} and
$\min_{M_0\le M\le2u}M^{\beta-1/2}\ge\min\bigl(M_0^{\beta-1/2},(2Ax)^{\beta-1/2}\bigr)$.
Combining with the left-hand inequalities of
\eqref{p0:eq:two-sided-coefficients} and \eqref{p0:eq:stirling-sandwich},
$|q_M|x^{-M}\ge k\,c_\beta\,c\,x^{-|\beta-1/2|}\e^{-Ax}$ for every
$M\ge M_0$.  Taking logarithms of the two bounds gives
\eqref{p0:eq:floor}.
\end{proof}

\begin{remark}[What may and may not be appended]\label{p0:rem:no-bare-exponential}
\Cref{p0:thm:optimal-truncation} bounds the size of the optimal remainder; it
does not produce it.  In particular, adding a term $c\,\e^{-Ax}$ with an
arbitrarily chosen constant $c$ to the divergent series is \emph{not} a valid
transseries completion.  The double factorial and swing sources make the point
explicitly: on the positive real axis the Borel singularities of these phases
lie at $\pm\ii\pi m$, so the exponentially small sectors are
terminant-weighted combinations of $\e^{\mp\ii\pi mx}$, whose optimally
truncated magnitude is $\e^{-\pi x}$ but which are not real exponentials.  The
unambiguous completion on the positive axis is the exact Borel--Laplace
representation, when the chapter provides one.  Where a chapter does supply a
rigorously normalized exponential sector, its transport through the inverse map
is \cref{p0:prop:operator-form} at first order and
\cref{p0:thm:remainder-transport} with an explicit error.
\end{remark}

\section{The parameter dictionary}\label{p0:sec:dictionary}

\begin{table}[htbp]
\centering
\caption{Model data of the four chapters, in the normal form of
\cref{p0:def:model} and after the $\alpha$-reduction of
\cref{p0:lem:alpha-reduction}.  ``Core'' names the admissible core of
\cref{p0:def:core} and the Lambert branch selected by
\cref{p0:cor:branch-rule}.  The row for \cref{p2:sec:top} is
\emph{not} an instance of \cref{p0:def:model}; see the repair box in
\cref{p0:sec:model}.}
\label{p0:tab:dictionary}
\begin{tabular}{llcccl}
\toprule
Chapter & sequence & variable $x$ & $\alpha$ & $b$ & core, branch\\
\midrule
\ref{p1:sec:top} & rooted trees, A000081 & $n$
 & $1$ & $-3/2$ & $ax+b\log x$, $\Wlam_{-1}$\\
\ref{p2:sec:top} & double factorial, A006882 & $n+1$
 & --- & --- & $\tfrac x2(\log x-1)$, $\Wlam_{0}$\\
\ref{p3:sec:top} & partitions, A000041 & $n-\tfrac1{24}$
 & $1/2$ & $-1$ & $ax^{1/2}+b\log x$, $\Wlam_{-1}$\\
\ref{p4:sec:top} & swing factorial, A056040 & $n$
 & $1$ & $\sigma/2$ & $ax+b\log x$, $\Wlam_{0}$ or $\Wlam_{-1}$\\
\bottomrule
\end{tabular}
\end{table}

The constants $a$ are: $a=\log\alpha_{\mathrm{Otter}}$ for
\cref{p1:sec:top}; $a=\pi\sqrt{2/3}$ for \cref{p3:sec:top}, in which
$\alpha=1/2$ and \cref{p0:lem:alpha-reduction} substitutes
$\xi=(n-1/24)^{1/2}$, after which $b/\alpha=-2$ and the core equation is
$\xi\mapsto a\xi-2\log\xi$ --- the equation the three partition articles write
as $\rho-2\log\rho=L$ after rescaling $a$ to $1$; and $a=\log2$ for
\cref{p4:sec:top}, with $\sigma=(-1)^{n+1}$ the parity transmonomial.  Each
chapter verifies \eqref{p0:eq:model-derivative} directly and states its own
$q_j$.

\begin{remark}[Reproducibility of the Part~0 identities]\label{p0:rem:verification}
Every coefficient identity of this chapter was checked in exact rational
arithmetic, with no floating-point step, over six independent choices of the
data $(a,b,q_1,\dots,q_N)$ with $a,b,q_j\in\Rat$, $N\in\{10,12\}$, using a
truncated-series engine over $\Rat[H]$ built on Python's
\texttt{fractions.Fraction}.  The identities checked were:
the Lagrange closed formula \eqref{p0:eq:c-lagrange} against the defining
fixed point; the ordinary-Bell recurrence \eqref{p0:eq:c-bell}; the vanishing
of the residual $ax+b\log x+Q(1/x)-L$ under the flattened ansatz
\eqref{p0:eq:direct-ansatz} through order $L^{-N}$; the shape statements
\eqref{p0:eq:P-shape}; and the closed form
\eqref{p0:eq:pure-stirling} together with its agreement with the three source
formulae listed in \cref{p0:rem:stirling-crosswalk}.  All checks passed.  No
numerical constant is asserted anywhere in this chapter.
\end{remark}

%% ==================== fragment p1 ====================
\chapter{Butcher--P\'olya rooted trees (A000081)}
\label{p1:sec:top}

\begin{sourcebox}
This chapter merges the three independently written treatments
\emph{Butcher\_Tree\_Transseries}, \emph{Butcher\_Tree\_Transseries-2} and
\emph{Butcher\_Tree\_Transseries-3} of the special-function-inversion draft set. Below they
are called \textsf{S1}, \textsf{S2} and \textsf{S3}. \textsf{S2} is a book-length treatment
and is the most complete on the inverse and on the inherited singularity at $-\sqrt\rho$;
\textsf{S3} is the most careful about what is and is not proved beyond the dominant sector;
\textsf{S1} has the shortest route to the universal Cayley coefficients and the only treatment
of the compositional inverse of the generating function. Where they disagree the disagreement
is stated in the text.  Every numerical constant retained here has been recomputed from the
exact divisor recurrence at $60$--$110$ working decimal digits; see \cref{p1:sub:audit}.
\end{sourcebox}

The subject of this chapter is the sequence
\begin{equation}
 0,\;1,\;1,\;2,\;4,\;9,\;20,\;48,\;115,\;286,\;719,\;1842,\;4766,
 \;12486,\;32973,\;87811,\;\ldots
 \label{p1:eq:sequence}
\end{equation}
(OEIS A000081), where for $n\ge1$ the term $a_n$ counts isomorphism classes of rooted,
unlabelled, non-plane trees on $n$ vertices.  These are the \emph{P\'olya trees}; in numerical
analysis the same shapes index Butcher series, elementary differentials and Runge--Kutta order
conditions, which is the origin of the name \emph{Butcher tree numbers}.  We set $a_0=0$
throughout.

The sequence grows like $C\alpha^nn^{-3/2}$ with $\alpha=\rho^{-1}$ and $\rho$ the Otter
singularity.  Thus, in the language of \cref{p0:def:model}, its asymptotic interpolation is an
exponential--power model with
\begin{equation}
 \boxed{\;\alpha_{\mathrm{model}}=1,\qquad
 a_{\mathrm{model}}=\lambda:=\log(1/\rho),\qquad
 b_{\mathrm{model}}=-\tfrac32,\qquad
 \delta=1.\;}
 \label{p1:eq:dictionary-preview}
\end{equation}
This chapter is therefore the \emph{linear-phase} instance of the Part~0 apparatus: no
$\alpha$-reduction (\cref{p0:lem:alpha-reduction}) is needed, the Lambert core
(\cref{p0:thm:lambert-core}) applies verbatim, and the entire reversion is a substitution into
\cref{p0:thm:lambert-centered}.  Everything specific to rooted trees is the computation of the
model's data $C$, $\lambda$ and $Q$; that computation is the substance of
\crefrange{p1:sec:polya}{p1:sec:forward}, and the inversion in \cref{p1:sec:inverse} is short
precisely because Part~0 has already been proved.

\section{The functional equation and the exact recurrence}
\label{p1:sec:polya}

\subsection{P\'olya's multiset equation}

Let $\mathcal T$ be the class of unlabelled non-plane rooted trees and let $\mathcal Z$ be one
atom (one vertex).  Deleting the root of a tree leaves an unordered multiset of rooted trees,
and grafting a multiset of rooted trees onto a fresh root inverts this.  Hence
\begin{equation}
 \mathcal T=\mathcal Z\times\operatorname{MSET}(\mathcal T).
 \label{p1:eq:species}
\end{equation}

\begin{proposition}[P\'olya's equation and the Euler product]
\label{p1:prop:polya}
Let $T(z)=\sum_{n\ge1}a_nz^n\in\Int[[z]]$.  Then, as identities of formal power series,
\begin{align}
 T(z)&=z\prod_{r\ge1}(1-z^r)^{-a_r},
 \label{p1:eq:euler-product}\\
 T(z)&=z\exp\!\left(\sum_{k\ge1}\frac{T(z^k)}{k}\right).
 \label{p1:eq:polya}
\end{align}
\end{proposition}

\begin{proof}
A multiset of rooted trees that contains $m_j$ copies of the $j$th shape has size
$\sum_jm_j|j|$ and contributes $z^{\sum_jm_j|j|}$.  For one fixed shape of size $r$, unbounded
multiplicity contributes $\sum_{m\ge0}z^{rm}=(1-z^r)^{-1}$; there are $a_r$ shapes of size
$r$, and distinct shapes are chosen independently, so the forest hanging from the root has
generating function $\prod_{r\ge1}(1-z^r)^{-a_r}$.  Every factor is a unit of $\Int[[z]]$ and
the product is $z$-adically convergent because the $r$th factor is $1+\bigO(z^r)$. Multiplying
by the root gives \cref{p1:eq:euler-product}.  Taking the formal logarithm and using
$-\log(1-w)=\sum_{k\ge1}w^k/k$,
\[
\log\frac{T(z)}{z} =\sum_{r\ge1}a_r\sum_{k\ge1}\frac{z^{rk}}{k}
=\sum_{k\ge1}\frac1k\sum_{r\ge1}a_r(z^k)^r =\sum_{k\ge1}\frac{T(z^k)}{k},
\]
the rearrangement being legitimate because for each fixed power of $z$ only finitely many
pairs $(r,k)$ contribute.  Exponentiating gives \cref{p1:eq:polya}.
\end{proof}

\begin{remark}
\Cref{p1:eq:polya} is purely formal; no analytic hypothesis is used, and none is available at
this stage.  All three siblings state it correctly, but only \textsf{S2} says explicitly that
the identity holds in $\Int[[z]]$ before any radius of convergence has been produced.  We keep
that discipline: \cref{p1:sec:dominant} is the first place where analysis enters.
\end{remark}

\subsection{The divisor recurrence}

\begin{definition}[Divisor transform]
\label{p1:def:divisor}
For $m\ge1$ put
\begin{equation}
 b_m:=\sum_{d\mid m}d\,a_d .
 \label{p1:eq:divisor}
\end{equation}
\end{definition}

\begin{theorem}[Otter--P\'olya recurrence]
\label{p1:thm:recurrence}
With $a_1=1$ one has, for every $n\ge2$,
\begin{equation}
 (n-1)\,a_n=\sum_{j=1}^{n-1}b_j\,a_{n-j}.
 \label{p1:eq:recurrence}
\end{equation}
\end{theorem}

\begin{proof}
Write $G(z)=T(z)/z=\sum_{n\ge0}a_{n+1}z^n$, a unit of $\Rat[[z]]$ with $G(0)=1$. From
\cref{p1:eq:euler-product} and the computation in the proof of \cref{p1:prop:polya},
\[
\log G(z)=\sum_{r\ge1}a_r\sum_{k\ge1}\frac{z^{rk}}{k} =\sum_{m\ge1}\frac{z^m}{m}\sum_{d\mid
m}d\,a_d =\sum_{m\ge1}\frac{b_m}{m}z^m,
\]
where the inner rearrangement groups the pairs $(r,k)$ with $rk=m$ and uses $1/k=d/m$ for
$d=r$.  Formal logarithmic differentiation gives $G'(z)=G(z)\sum_{m\ge1}b_mz^{m-1}$.
Extracting $\Coef{z^{n-1}}$,
\[
n\,a_{n+1}=\sum_{j=1}^{n}b_j\,a_{n-j+1},
\]
and replacing $n$ by $n-1$ yields \cref{p1:eq:recurrence}.
\end{proof}

\begin{remark}[Why the recurrence is triangular, and its cost]
\label{p1:rem:triangular}
Although $b_n$ contains the summand $n\,a_n$, the right-hand side of \cref{p1:eq:recurrence}
uses only $b_j$ with $j<n$, so $a_n$ is determined by $a_1,\ldots,a_{n-1}$.  Implemented as a
sieve --- after computing $a_n$, add $n\,a_n$ to $b_n,b_{2n},b_{3n},\ldots$ --- the divisor
transform costs $\bigO(N\log N)$ additions and the convolutions cost $\bigO(N^2)$ integer
operations for $a_1,\ldots,a_N$.  The integrality assertion $(n-1)\mid\sum_jb_ja_{n-j}$ is a
free exact self-test on every step, and was used as such in the computations of
\cref{p1:sub:audit}.
\end{remark}

\subsection{The self-interaction and its coefficients}

Separating the $k=1$ term of \cref{p1:eq:polya} is the structural move on which everything
else rests.

\begin{definition}[Disturbance and Cayley argument]
\label{p1:def:R-zeta}
Put
\begin{equation}
 R(z):=\sum_{k\ge2}\frac{T(z^k)}{k},
 \qquad
 \zeta(z):=z\,\e^{R(z)} .
 \label{p1:eq:R-zeta}
\end{equation}
\end{definition}

With this notation \cref{p1:eq:polya} reads
\begin{equation}
 T(z)=\zeta(z)\,\e^{T(z)} .
 \label{p1:eq:T-zeta}
\end{equation}

\begin{lemma}[Coefficients of the disturbance]
\label{p1:lem:R-coefficients}
$R(z)=\sum_{r\ge2}\varrho_rz^r$ with
\begin{equation}
 \varrho_r=\frac1r\sum_{\substack{d\mid r\\ d<r}}d\,a_d
 =\frac{b_r-r\,a_r}{r}\in\Rat_{\ge0}.
 \label{p1:eq:varrho}
\end{equation}
In particular $\varrho_r$ depends only on $a_d$ with $d\le r/2$, and $\varrho_2=\tfrac12$,
$\varrho_3=\tfrac13$, $\varrho_4=\tfrac34$, $\varrho_5=\tfrac15$, $\varrho_6=\tfrac32$,
$\varrho_7=\tfrac17$, $\varrho_8=\tfrac{19}{8}$.
\end{lemma}

\begin{proof}
$R(z)=\sum_{k\ge2}\sum_{d\ge1}\frac{a_d}{k}z^{kd}$; collecting the terms with $kd=r$ gives
$\varrho_r=\sum_{d\mid r,\,d<r}a_d/(r/d)=\frac1r\sum_{d\mid r,\,d<r}d\,a_d$, which is
$\left(b_r-ra_r\right)/r$ by \cref{p1:eq:divisor}. Every divisor $d<r$ of $r$ satisfies $d\le
r/2$.  The listed values follow from $a_1=a_2=1$, $a_3=2$, $a_4=4$: e.g.\
$\varrho_4=(1\cdot1+2\cdot1)/4=\tfrac34$ and $\varrho_6=(1+2\cdot1+3\cdot2)/6=\tfrac32$.
\end{proof}

\Cref{p1:lem:R-coefficients} is due to \textsf{S2}; \textsf{S1} and \textsf{S3} carry $R$ only
through the double sum $\sum_{k\ge2}\sum_na_n\rho^{kn}(\cdots)$. The single-sum form is used
throughout this chapter because it turns every quantity below into one geometrically
convergent series in the exact integers $a_r$ (\cref{p1:lem:Rhat}).

\section{The Cayley--Lambert normal form and the Otter singularity}
\label{p1:sec:dominant}

\subsection{Factorisation through the Cayley tree function}

\begin{definition}[Cayley tree function]
\label{p1:def:cayley}
Let $\mathcal C(w)$ be the unique formal solution of $\mathcal C=w\e^{\mathcal C}$ with
$\mathcal C(0)=0$.  It is analytic on $|w|<\e^{-1}$ and
\begin{equation}
 \mathcal C(w)=-\Wlam_0(-w),
 \label{p1:eq:cayley-lambert}
\end{equation}
where $\Wlam_0$ is the principal Lambert branch.
\end{definition}

Comparing \cref{p1:eq:T-zeta} with \cref{p1:def:cayley},
\begin{equation}
 \boxed{\;T(z)=\mathcal C\bigl(\zeta(z)\bigr)
 =-\Wlam_0\bigl(-\zeta(z)\bigr).\;}
 \label{p1:eq:factorisation}
\end{equation}
This holds as formal series, and as an identity of analytic functions on the common domain of
the branches selected at the origin.  The point of the factorisation is a clean separation of
labour: \emph{all} of the singular mechanism is the universal square-root branch point of
$\mathcal C$ at $w=\e^{-1}$, and \emph{all} of the tree-specific information is packed into
the map $\zeta$, which --- as we now show --- is analytic on a strictly larger disc than $T$
itself.

\begin{lemma}[Elementary bounds on the radius]
\label{p1:lem:radius-bounds}
The radius of convergence $\rho$ of $T$ satisfies $\tfrac14\le\rho\le\e^{-1}$.
\end{lemma}

\begin{proof}
Since $a_n\ge1$ for $n\ge1$ we have $\rho\le1$, so for $0<x<\rho$ every $x^{k}$ with $k\ge2$
lies in $(0,\rho)$ as well.  Upper bound: for such $x$ the series $T(x)$ and $R(x)$ converge,
all terms are nonnegative, and $\zeta(x)=x\e^{R(x)}\ge x$; so \cref{p1:eq:T-zeta} gives
$T(x)\ge x\e^{T(x)}$, i.e.\ $x\le T(x)\e^{-T(x)}\le\sup_{u\ge0}u\e^{-u}=\e^{-1}$.  Letting
$x\uparrow\rho$ gives $\rho\le\e^{-1}$.  Lower bound: forgetting the identification of
isomorphic children embeds the shapes counted by $a_n$ into the rooted \emph{plane} trees on
$n$ vertices, of which there are $\frac1n\binom{2n-2}{n-1}\le4^{n}$; hence $\limsup
a_n^{1/n}\le4$ and $\rho\ge\tfrac14$.
\end{proof}

\begin{proposition}[Analyticity of the disturbance]
\label{p1:prop:R-analytic}
$R$, and hence $\zeta$, is analytic on $|z|<\sqrt\rho$, and the series $\sum_{k\ge2}T(z^k)/k$
converges normally on compact subsets together with all its derivatives.  In particular $R$
and $\zeta$ are analytic at $z=\rho$, since $\rho<\sqrt\rho$ by \cref{p1:lem:radius-bounds}.
\end{proposition}

\begin{proof}
Fix $r<\sqrt\rho$, so $r^2<\rho$.  Since $T$ has radius $\rho$ and $T(w)=w+\bigO(w^2)$, there
are $M>0$ and $r_0\in(0,1)$ with $|T(w)|\le M|w|$ for $|w|\le r_0$.  For $|z|\le r$, only
finitely many $k$ have $r^k>r_0$, and for each such $k$ we have $r^k\le r^2<\rho$, so $T(z^k)$
is defined and bounded there.  For the remaining $k$, $|T(z^k)/k|\le Mr^k/k$, which is
summable. Differentiating $j$ times multiplies the $k$th tail term by a factor $\bigO(k^j
r^{-j})$, still summable.  Weierstrass' theorem gives analyticity and termwise
differentiation.
\end{proof}

\begin{proposition}[Critical characterisation of $\rho$]
\label{p1:prop:critical}
With $\rho\in[\tfrac14,\e^{-1}]$ the radius of convergence of $T$,
\begin{equation}
 T(\rho)=1,\qquad \zeta(\rho)=\e^{-1},\qquad
 \log\rho+1+R(\rho)=0.
 \label{p1:eq:critical}
\end{equation}
Moreover $\zeta$ is strictly increasing on $(0,\sqrt\rho)$ and
\begin{equation}
 \theta_1:=\e\rho\,\zeta'(\rho)=1+\rho R'(\rho)>0 .
 \label{p1:eq:theta1}
\end{equation}
\end{proposition}

\begin{proof}
All coefficients of $T$, hence of $R$ and of $\zeta=z\e^{R}$, are nonnegative, and
$\zeta(z)=z+\tfrac12z^{3}+\tfrac13z^{4}+\cdots$ has positive linear coefficient; therefore
$\zeta'(x)>0$ on $(0,\sqrt\rho)$, and $\zeta$ is strictly increasing there.

\emph{Step 1: $T(x)<1$ on $[0,\rho)$, and $T=\mathcal C\circ\zeta$ there.}  On $[0,\rho)$ the
function $T$ is continuous, increasing, with $T(0)=0$, and \cref{p1:eq:T-zeta} gives
$\zeta(x)=T(x)\e^{-T(x)}\le\e^{-1}$.  Let $\rho_1:=\sup\{x\in[0,\rho):T(x)<1\}$.  On
$[0,\rho_1)$ we have $\zeta(x)<\e^{-1}$, so $\mathcal C$ is analytic at $\zeta(x)$ and
$\mathcal C\circ\zeta$ is an analytic solution of \cref{p1:eq:T-zeta} agreeing with $T$ at
$0$; by uniqueness of that germ and analytic continuation, $T=\mathcal C\circ\zeta$ on
$[0,\rho_1)$.  If $\rho_1<\rho$ then $T(\rho_1)=1$ and $\zeta(\rho_1)=\e^{-1}$; but $\zeta$ is
strictly increasing, so $\zeta(x)>\e^{-1}$ for $x\in(\rho_1,\rho)$, contradicting
$\zeta(x)\le\e^{-1}$.  Hence $\rho_1=\rho$.

\emph{Step 2: $\zeta(\rho)=\e^{-1}$.}  By \cref{p1:prop:R-analytic} the function $\zeta$ is
analytic at $\rho$, and by Step~1 and monotonicity $\zeta(\rho)=\lim_{x\uparrow\rho}
\zeta(x)\le\e^{-1}$.  If the inequality were strict, $\mathcal C$ would be analytic at
$\zeta(\rho)$ and $\mathcal C\circ\zeta$ would be analytic in a full neighbourhood of $\rho$
and would agree with $T$ on $[0,\rho)$; then $T$ would continue analytically through $\rho$,
contradicting Pringsheim's theorem, which makes the positive point $z=\rho$ a singularity of a
series with nonnegative coefficients.  Hence $\zeta(\rho)=\e^{-1}$ and, by
\cref{p1:def:cayley}, $T(\rho)=\mathcal C(\e^{-1})=1$.

\emph{Step 3.}  Taking logarithms in $\zeta(\rho)=\rho\e^{R(\rho)}=\e^{-1}$ gives the scalar
equation; equivalently $R(\rho)=\lambda-1$ with $\lambda=-\log\rho$, and $R(\rho)>0$
independently re-proves $\rho<\e^{-1}$.  Finally $\zeta'=(1/z+R')\zeta$, so
$\e\rho\zeta'(\rho)=\e\zeta(\rho)\bigl(1+\rho R'(\rho)\bigr)=1+\rho R'(\rho)$, positive
because $R$ has nonnegative coefficients.
\end{proof}

\begin{proposition}[Uniqueness of the dominant singularity]
\label{p1:prop:unique-dominant}
$z=\rho$ is the only singularity of $T$ on the circle $|z|=\rho$.
\end{proposition}

\begin{proof}
Take $z=\rho\e^{\ii\vartheta}$ with $\vartheta\not\equiv0$.  The series
$\zeta(z)=z+\tfrac12z^3+\tfrac13z^4+\cdots$ has nonnegative coefficients, so
$|\zeta(z)|\le\zeta(\rho)=\e^{-1}$, with equality only if the three displayed nonzero terms
$z$, $\tfrac12z^3$, $\tfrac13z^4$ share one argument, i.e.\
$\e^{2\ii\vartheta}=\e^{3\ii\vartheta}=1$, hence $\e^{\ii\vartheta}=1$ and $\vartheta\equiv0$.
So $|\zeta(z)|<\e^{-1}$ strictly, $\mathcal C$ is analytic at $\zeta(z)$, and $\zeta$ is
analytic at $z$ because $\rho<\sqrt\rho$ (\cref{p1:prop:R-analytic}).  Hence $T=\mathcal
C\circ\zeta$ is analytic at $z$.
\end{proof}

\begin{remark}
This is \textsf{S3}'s argument, and it is the only one of the three that is self-contained:
\textsf{S1} and \textsf{S2} appeal to ``positivity and aperiodicity'' without exhibiting the
aperiodicity witness.  The three coefficients $z,z^3,z^4$ of $\zeta$ furnish that witness
explicitly, because $\gcd(3-1,4-1)=1$.
\end{remark}

\subsection{The constants}
\label{p1:sub:constants}

\begin{convention}[Growth constants]
\label{p1:def:constants}
Throughout the chapter
\begin{equation}
 \alpha:=\rho^{-1},\qquad \lambda:=\log\alpha=-\log\rho,\qquad p:=\tfrac32 .
 \label{p1:eq:alpha-lambda}
\end{equation}
\end{convention}

\begin{repairbox}
\textbf{Notation collision between the siblings.}  \textsf{S2} and \textsf{S3} write
$\alpha=\rho^{-1}$ and $\lambda=\log\alpha$; \textsf{S1} writes exactly the opposite,
$\lambda=\rho^{-1}$ and $\alpha=-\log\rho$.  The symbols therefore cross-denote between the
drafts, and \textsf{S1}'s $\Omega_-/\alpha$, $\e^{\alpha n}$ and $d_m$ formulas must be read
with $\alpha=\lambda_{\text{here}}$.  This volume adopts the two-to-one majority reading
\cref{p1:eq:alpha-lambda}; every formula imported from \textsf{S1} has been translated.
\end{repairbox}

Solving $\log\rho+1+R(\rho)=0$ from \cref{p1:eq:varrho} with $700$ exact tree numbers at $110$
working digits (see \cref{p1:sub:audit}) gives
\begin{align}
 \rho&=0.338321856899207695196112625717017053183774607532967795572303776
 \ldots,
 \label{p1:eq:rho-numeric}\\
 \alpha&=2.955765285651994974714817524123194588375492304663596595350472478
 \ldots,
 \label{p1:eq:alpha-numeric}\\
 \lambda&=1.083757597244624660482711980896174279049662124202672234914427065
 \ldots,
 \label{p1:eq:lambda-numeric}\\
 R(\rho)&=0.083757597244624660482711980896174279049662124202672234914427065
 \ldots,
 \label{p1:eq:Rrho-numeric}
\end{align}
and, as an independent consistency check on the whole pipeline,
\begin{equation}
 \lambda=1+R(\rho),
 \label{p1:eq:lambda-check}
\end{equation}
which is nothing but \cref{p1:eq:critical} rearranged; the two computed values agree to all
$70$ digits carried.  Two further constants are recorded here because they are quoted
repeatedly:
\begin{equation}
 \theta_1=1.216004561824048922157212619228187444\ldots,
 \qquad
 \zeta'(\rho)=1.322241142696930264451284889\ldots .
 \label{p1:eq:theta1-numeric}
\end{equation}

\section{The universal Cayley branch germ}
\label{p1:sec:cayley}

The next result is universal: it depends on the Cayley function alone, not on rooted trees,
and applies verbatim to \emph{any} class satisfying $T=\zeta\e^{T}$ with analytic $\zeta$.

\begin{definition}[Cayley kernel]
\label{p1:def:cayley-kernel}
Put
\begin{equation}
 \Phi(v):=\frac{2\bigl(1-(1-v)\e^{v}\bigr)}{v^{2}}
 =\sum_{m\ge0}\frac{2(m+1)}{(m+2)!}\,v^{m}
 =\sum_{m\ge0}\frac{2}{(m+2)\,m!}\,v^{m}
 =1+\frac23v+\frac14v^{2}+\frac1{15}v^{3}+\cdots .
 \label{p1:eq:Phi}
\end{equation}
\end{definition}

The two displayed coefficient forms agree because $2(m+1)/(m+2)!=2/((m+2)\,m!)$; \textsf{S1}
and \textsf{S3} use the first, \textsf{S2} the second.

\begin{theorem}[Universal Cayley coefficients]
\label{p1:thm:omega}
Set $w=(1-\upsilon)/\e$ and $\mathcal C(w)=1-v$ on the branch reached from $0<w<\e^{-1}$. Then
$\upsilon=\tfrac12v^{2}\Phi(v)$, the inverse germ is
$v=\sum_{m\ge1}\omega_m(2\upsilon)^{m/2}$, and for every $m\ge1$
\begin{equation}
 \boxed{\;\omega_m=\frac1m\,\Coef{v^{m-1}}\Phi(v)^{-m/2}
 =\frac1m\sum_{k=0}^{m-1}\binom{-m/2}{k}
 \OB_{m-1,k}\!\left(\tfrac23,\tfrac14,\tfrac1{15},\ldots\right)
 \in\Rat .\;}
 \label{p1:eq:omega-master}
\end{equation}
Consequently, on that branch,
\begin{equation}
 \mathcal C\!\left(\frac{1-\upsilon}{\e}\right)
 =1+\sum_{r\ge1}c_r\,\upsilon^{r/2},
 \qquad
 c_r=-2^{r/2}\,\omega_r .
 \label{p1:eq:cayley-puiseux}
\end{equation}
\end{theorem}

\begin{proof}
From $\mathcal C=w\e^{\mathcal C}$ we get $\e w=\mathcal C\e^{1-\mathcal C}$, so with
$\mathcal C=1-v$,
\[
1-\upsilon=\e w=(1-v)\e^{v}, \qquad\text{hence}\qquad \upsilon=1-(1-v)\e^{v}
=\sum_{m\ge2}\frac{m-1}{m!}v^{m} =\frac{v^{2}}{2}\Phi(v),
\]
the middle identity by expanding $\e^v$ and collecting.  Put $\varsigma=\sqrt{2\upsilon}$ with
the determination $\varsigma/v\to1$; then $\varsigma=v\,\Phi(v)^{1/2}$ is an ordinary power
series in $v$ with nonzero linear coefficient, so \cref{p0:thm:LB} applies with
$\phi(v)=\Phi(v)^{1/2}$ and yields
\[
\Coef{\varsigma^{m}}v(\varsigma) =\frac1m\Coef{v^{m-1}}\phi(v)^{-m}
=\frac1m\Coef{v^{m-1}}\Phi(v)^{-m/2}=\omega_m .
\]
The Bell form is \cref{p0:lem:power-log} applied to the unit series $\Phi$, whose coefficients
are $2(m+1)/(m+2)!$ for $m\ge1$.  Rationality follows because $\binom{-m/2}{k}\in\Rat$ and
$\OB_{m-1,k}$ is a polynomial with rational coefficients in rational arguments.  Finally
$\mathcal C=1-v=1-\sum_m\omega_m\varsigma^m =1-\sum_m\omega_m2^{m/2}\upsilon^{m/2}$, which is
\cref{p1:eq:cayley-puiseux}.
\end{proof}

\begin{remark}[Two normalisations, one object]
\label{p1:rem:omega-vs-c}
\textsf{S2} works with $\omega_m$ (expanding in $\varsigma=\sqrt{2\upsilon}$) while
\textsf{S1} and \textsf{S3} work with $c_r$ (expanding in $\upsilon^{1/2}$).  The dictionary
is $c_r=-2^{r/2}\omega_r$, verified here for $1\le r\le18$ by exact rational arithmetic.  The
$\omega_m$ are rational, the $c_r$ alternate between rationals and rational multiples of
$\sqrt2$; this is the only reason to prefer $\omega$ for tables.
\end{remark}

\begin{table}[htbp]
\centering
\caption{Exact universal Cayley coefficients $\omega_m$ of
\cref{p1:eq:omega-master}, recomputed by rational arithmetic from
\cref{p1:eq:Phi}.}
\label{p1:tab:omega}
\footnotesize
\setlength{\tabcolsep}{5pt}
\begin{tabular}{r@{\quad}l@{\qquad}r@{\quad}l}
\toprule
$m$ & $\omega_m$ & $m$ & $\omega_m$\\
\midrule
$1$ & $1$ & $10$ & $-5776369/1515591000$\\
$2$ & $-1/3$ & $11$ & $169709463197/69528040243200$\\
$3$ & $11/72$ & $12$ & $-1118511313/709296588000$\\
$4$ & $-43/540$ & $13$ & $667874164916771/650782456676352000$\\
$5$ & $769/17280$ & $14$ & $-500525573/744761417400$\\
$6$ & $-221/8505$ & $15$ & $103663334225097487/234281684403486720000$\\
$7$ & $680863/43545600$ & $16$ & $-466901817532379/1595278956070800000$\\
$8$ & $-1963/204120$ & $17$ & $21235294185086305043/109242202556140093440000$\\
$9$ & $226287557/37623398400$ & $18$ & $-106040742894306601/818378104464320400000$\\
\bottomrule
\end{tabular}
\end{table}

The corresponding $c_r$ are
\begin{equation}
\begin{aligned}
 c_1&=-\sqrt2, & c_2&=\tfrac23, &
 c_3&=-\tfrac{11\sqrt2}{36}, & c_4&=\tfrac{43}{135},\\
 c_5&=-\tfrac{769\sqrt2}{4320}, & c_6&=\tfrac{1768}{8505}, &
 c_7&=-\tfrac{680863\sqrt2}{5443200}, & c_8&=\tfrac{3926}{25515}.
\end{aligned}
\label{p1:eq:c-values}
\end{equation}

\begin{remark}[On the observed sign pattern]
\label{p1:rem:omega-signs}
All three siblings display $\omega_m$ with strictly alternating signs and none proves it.  The
data in \cref{p1:tab:omega} are consistent with $\sgn\omega_m=(-1)^{m+1}$ for $1\le m\le18$,
and this is what one expects from the fact that the inverse function $v(\varsigma)$ of
$\varsigma=v\Phi(v)^{1/2}$ has its nearest singularity on the negative $\varsigma$-axis; but
no proof is offered here and the statement is recorded as an observation, not as a result.
\end{remark}

\section{The Puiseux expansion at the Otter singularity}
\label{p1:sec:puiseux}

\subsection{The local disturbance series}

\begin{definition}[Local coordinate and disturbance]
\label{p1:def:local}
Write
\begin{equation}
 s:=1-\frac{z}{\rho},
 \qquad
 \theta(s):=1-\e\,\zeta\bigl(\rho(1-s)\bigr)=\sum_{m\ge1}\theta_ms^{m}.
 \label{p1:eq:theta-def}
\end{equation}
The constant term vanishes by \cref{p1:eq:critical}.  By \cref{p1:prop:R-analytic}, $\theta$
is analytic on the image of $|z|<\sqrt\rho$ under $z\mapsto s$, namely the disc
$|s-1|<\rho^{-1/2}$; in particular its Taylor series at $s=0$ converges for
$|s|<\rho^{-1/2}-1=0.71923\ldots$.
\end{definition}

Taylor's formula gives the derivative form
\begin{equation}
 \theta_m=\frac{(-1)^{m+1}\e\rho^{m}}{m!}\,\zeta^{(m)}(\rho),
 \qquad m\ge1,
 \label{p1:eq:theta-derivative}
\end{equation}
consistent with \cref{p1:eq:theta1}: $\theta_1=\e\rho\zeta'(\rho)$.  For computation the
following coefficient form is preferable, since it needs nothing but the integers $a_r$ and
the single scalar $\rho$.

\begin{lemma}[Local coefficients of $R$ from the tree numbers]
\label{p1:lem:Rhat}
Let $\widehat R_m:=\Coef{s^{m}}R\bigl(\rho(1-s)\bigr)$.  Then
\begin{equation}
 \boxed{\;
 \widehat R_m=(-1)^{m}\sum_{r\ge2}\varrho_r\binom{r}{m}\rho^{r}
 =(-1)^{m}\,\frac{\rho^{m}R^{(m)}(\rho)}{m!},
 \;}
 \label{p1:eq:Rhat}
\end{equation}
with $\varrho_r$ as in \cref{p1:eq:varrho}.  The series converges geometrically: for fixed $m$
the general term is $\bigO\bigl(r^{m}\rho^{r/2}\bigr)$.
\end{lemma}

\begin{proof}
By \cref{p1:lem:R-coefficients}, $R(\rho(1-s))=\sum_{r\ge2}\varrho_r\rho^{r}(1-s)^{r}$, and
$\Coef{s^m}(1-s)^r=(-1)^m\binom rm$.  The middle expression is a Taylor coefficient of $R$ at
$\rho$ in the variable $\rho s$, giving the right-hand form.  For the convergence rate, every
divisor $d<r$ of $r$ satisfies $d\le r/2$, so $\varrho_r\le a_{\lfloor r/2\rfloor}\,d(r)$ with
$d(r)$ the number of divisors of $r$; since $\limsup_na_n^{1/n}=\rho^{-1}$, for every
$\epsilon>0$ there is $K_\epsilon$ with $\varrho_r\rho^{r}\le
K_\epsilon(\sqrt\rho+\epsilon)^{r}$.  With $\binom rm\le r^{m}$ the claim follows.
\end{proof}

\begin{remark}[A dedup]
\label{p1:rem:single-sum}
\textsf{S1} (its $\ell_m$) and \textsf{S3} (its $\Coef{s^m}R(\rho(1-s))$) both give the
\emph{double} sum $(-1)^m\sum_{k\ge2}\frac1k\sum_{n\ge1}a_n\binom{kn}{m}\rho^{kn}$.
\Cref{p1:eq:Rhat} is the same number after the inner divisor collection of
\cref{p1:lem:R-coefficients}, and is one sum rather than two.  The two forms were checked
against each other numerically to $70$ digits for $0\le m\le30$.
\end{remark}

\begin{lemma}[The disturbance coefficients]
\label{p1:lem:theta}
With $\widehat R_m$ as in \cref{p1:lem:Rhat}, put $\widehat R_{+}(s):=\sum_{m\ge1}\widehat
R_ms^{m}$ and $E(s):=\exp\widehat R_{+}(s)=\sum_{m\ge0}E_ms^{m}$.  Then $E_0=1$,
\begin{equation}
 E_m=\sum_{k=0}^{m}\frac1{k!}\,\OB_{m,k}
 \bigl(\widehat R_1,\widehat R_2,\ldots\bigr)
 =\frac1m\sum_{k=1}^{m}k\,\widehat R_k\,E_{m-k},
 \label{p1:eq:E-bell}
\end{equation}
and
\begin{equation}
 \boxed{\;\theta_m=E_{m-1}-E_m,\qquad m\ge1.\;}
 \label{p1:eq:theta-bell}
\end{equation}
\end{lemma}

\begin{proof}
By \cref{p1:eq:critical}, $\e\zeta(\rho)=1$, so
\[
\e\,\zeta\bigl(\rho(1-s)\bigr) =\frac{\rho(1-s)\e^{R(\rho(1-s))}}{\rho\,\e^{R(\rho)}}
=(1-s)\,\e^{\widehat R_{+}(s)}=(1-s)E(s),
\]
because $\widehat R_0=R(\rho)$.  Hence $\theta(s)=1-(1-s)E(s)$ and $\theta_m=-(E_m-E_{m-1})$
for $m\ge1$, while the constant term is $1-E_0=0$. The two expressions for $E_m$ are the
ordinary Bell expansion of an exponential (\cref{p0:def:obell}) and the standard triangular
exponential recurrence obtained from $E'=\widehat R_{+}'E$.
\end{proof}

\begin{remark}[The jet route]
\label{p1:rem:jets}
\textsf{S2} computes $\theta_m$ instead from the \emph{scaled logarithmic jets}
\begin{equation}
 \ell_j:=\rho^{j}\left.\frac{\dd^{j}}{\dd z^{j}}\log\zeta(z)\right|_{z=\rho}
 =(-1)^{j-1}(j-1)!+\sum_{r\ge j}\varrho_r\,r^{\underline{j}}\,\rho^{r},
 \label{p1:eq:jets}
\end{equation}
via $\theta_m=\frac{(-1)^{m+1}}{m!}\EB_m(\ell_1,\ldots,\ell_m)$, where $\EB_m$ is the complete
exponential Bell polynomial and $r^{\underline{j}}=r(r-1)\cdots(r-j+1)$.  This is equivalent
to \cref{p1:lem:theta} through $\widehat R_m=(-1)^m\rho^mR^{(m)}(\rho)/m!$ and
$\ell_j=(-1)^{j}j!\,\widehat R_j$ for $j\ge1$ up to the $\log z$ term; both routes were
implemented and agree.  \Cref{p1:lem:theta} is preferred because it avoids the factorial
rescaling and its intermediate quantities stay $\bigO(1)$. \textsf{S2}'s normalisation
$s_j=2\theta_j$ (it expands $2(1-\e\zeta)$ rather than $1-\e\zeta$) accounts for every factor
of $2$ that differs between the siblings' formulas for $t_r$.
\end{remark}

The first disturbance coefficients, recomputed, are
\begin{equation}
\begin{aligned}
 \theta_1&=\phantom{-}1.216004561824048922157212619\ldots,&
 \theta_2&=-0.458415729856391719180847717\ldots,\\
 \theta_3&=\phantom{-}0.448104774389535043665734122\ldots,&
 \theta_4&=-0.399561679956209677406231980\ldots,\\
 \theta_5&=\phantom{-}0.385251341946613000719546396\ldots,&
 \theta_6&=-0.389806561959203341949164640\ldots .
\end{aligned}
\label{p1:eq:theta-numeric}
\end{equation}

\subsection{All local coefficients}

\begin{theorem}[Every Puiseux coefficient at $\rho$]
\label{p1:thm:t-coefficients}
Factor $\theta(s)=\theta_1s\bigl(1+U(s)\bigr)$ with $U(s)=\sum_{j\ge1}u_js^{j}$,
$u_j=\theta_{j+1}/\theta_1$.  Then in a slit neighbourhood of $z=\rho$ there is a convergent
Puiseux expansion
\begin{equation}
 T(z)=1+\sum_{n\ge1}t_n\,s^{n/2},
 \qquad s=1-\frac z\rho,
 \label{p1:eq:T-puiseux}
\end{equation}
and for every $n\ge1$
\begin{equation}
 \boxed{\;
 t_n=\sum_{\substack{1\le r\le n\\ r\equiv n\ (2)}}
 c_r\,\theta_1^{r/2}\,
 \Coef{s^{(n-r)/2}}\bigl(1+U(s)\bigr)^{r/2}
 =\sum_{\substack{1\le r\le n\\ r\equiv n\ (2)}}
 c_r\,\theta_1^{r/2}
 \sum_{k=0}^{(n-r)/2}\binom{r/2}{k}\OB_{\frac{n-r}2,k}(u_1,u_2,\ldots).\;}
 \label{p1:eq:t-master}
\end{equation}
Thus $t_n$ is a finite expression in $\omega_1,\ldots,\omega_n$, $\theta_1^{\pm1/2}$ and
$\theta_2,\ldots,\theta_{\lceil (n+1)/2\rceil}$.
\end{theorem}

\begin{proof}
Near $z=\rho$ we have $\e\zeta=1-\theta(s)$ with $\theta(0)=0$, so by
\cref{p1:eq:factorisation} and \cref{p1:eq:cayley-puiseux},
$T=1+\sum_{r\ge1}c_r\theta(s)^{r/2}$, the branch being the one with $T\to1^{-}$ as
$s\to0^{+}$.  Substituting $\theta(s)^{r/2}=\theta_1^{r/2}s^{r/2}(1+U(s))^{r/2}$, a term of
index $r$ can contribute to $s^{n/2}$ only when $n-r$ is a nonnegative even integer.
Extracting $\Coef{s^{(n-r)/2}}$ and applying \cref{p0:lem:power-log} to the unit series $1+U$
gives \cref{p1:eq:t-master}.  Convergence: $\theta$ is analytic and has a simple zero at $s=0$
with $\theta_1>0$ (\cref{p1:eq:theta1}), so $U$ is analytic with $U(0)=0$; choosing a slit
neighbourhood on which $|U|<1$ and a consistent square root, the convergent series
\cref{p1:eq:cayley-puiseux} may be composed with the analytic map $s\mapsto\theta(s)$.
\end{proof}

\begin{corollary}
\label{p1:cor:t-low}
$t_1=-\sqrt{2\theta_1}$, $t_2=\tfrac23\theta_1$, and
$t_3=-\tfrac{11\sqrt2}{36}\theta_1^{3/2}-\tfrac{\sqrt2}{2}\theta_2 \theta_1^{-1/2}$.
\end{corollary}

\begin{proof}
Take $n=1,2,3$ in \cref{p1:eq:t-master} with $c_1=-\sqrt2$, $c_2=\tfrac23$,
$c_3=-\tfrac{11\sqrt2}{36}$ and $\OB_{0,0}=1$, $\OB_{1,1}(u)=u_1$, $\binom{1/2}{1}=\tfrac12$.
\end{proof}

\begin{remark}[Sibling cross-check]
\label{p1:rem:t3-cross}
In \textsf{S2}'s normalisation $s_1=2\theta_1$, $s_2=2\theta_2$, the same coefficient reads
$t_3=-(11s_1^{2}+36s_2)/(72\sqrt{s_1})$; substituting $s_j=2\theta_j$ reproduces
\cref{p1:cor:t-low} exactly.  Likewise $t_1=-\sqrt{s_1}$ and $t_2=s_1/3$.  Both were checked
numerically.
\end{remark}

\begin{table}[htbp]
\centering
\caption{Puiseux coefficients in $T(z)=1+\sum_{n\ge1}t_n(1-z/\rho)^{n/2}$,
recomputed from \cref{p1:eq:t-master} at $90$ working digits.}
\label{p1:tab:t}
\small
\setlength{\tabcolsep}{6pt}
\begin{tabular}{r@{\quad}r@{\qquad}r@{\quad}r}
\toprule
$n$ & $t_n$ & $n$ & $t_n$\\
\midrule
$1$ & $-1.559490020374640885542$ & $10$ & $\phantom{-}0.205984831277907418736$\\
$2$ & $\phantom{-}0.810669707882699281438$ & $11$ & $-0.052724708498190562315$\\
$3$ & $-0.285487021612845605304$ & $12$ & $\phantom{-}0.017026568754953669446$\\
$4$ & $\phantom{-}0.165372365712083893477$ & $13$ & $-0.152370624366325396286$\\
$5$ & $-0.342459970402154201042$ & $14$ & $\phantom{-}0.173702883299850462894$\\
$6$ & $\phantom{-}0.317407225946528563568$ & $15$ & $-0.014473703739527044857$\\
$7$ & $-0.107778800291631009182$ & $16$ & $-0.021899517611215562233$\\
$8$ & $\phantom{-}0.061384957055835104689$ & $17$ & $-0.144547193570909705477$\\
$9$ & $-0.195212383597556463613$ & $18$ & $\phantom{-}0.176077108885017779062$\\
\bottomrule
\end{tabular}
\end{table}

\begin{remark}[Analytic versus singular local terms]
\label{p1:rem:even-odd}
The even part $\sum_jt_{2j}s^{j}$ of \cref{p1:eq:T-puiseux} is analytic at $z=\rho$.  A fixed
analytic monomial $(1-z/\rho)^{j}$ is a polynomial in $z$ and contributes nothing to
$\Coef{z^n}$ once $n>j$; hence only the odd coefficients $t_{2j+1}$ enter the algebraic
$\rho^{-n}n^{-\bullet}$ asymptotics of \cref{p1:sec:forward}.  This is not a discarding of
information: the analytic germ continues through $\rho$ and reappears in the exponentially
smaller sectors of \cref{p1:sec:sectors}.  The sign pattern of \cref{p1:tab:t} is not monotone
and should not be over-read; these are local analytic data, not the final coefficient
constants.
\end{remark}

In Otter's coordinate $\rho-z=\rho s$ the leading behaviour reads $T(z)=1-\mathsf
b\,(\rho-z)^{1/2}+\bigO(\rho-z)$ with
\begin{equation}
 \mathsf b=-\frac{t_1}{\sqrt\rho}=2.68112814726711223857732878\ldots .
 \label{p1:eq:otter-b}
\end{equation}

\begin{theorem}[Local expansion with an explicit remainder]
\label{p1:thm:local-remainder}
For every $N\ge1$ there are a dented (``$\Delta$-domain'') neighbourhood $\Delta$ of $\rho$
and a constant $K_N$ such that
\begin{equation}
 \left|T(z)-1-\sum_{n=1}^{N}t_n\left(1-\frac z\rho\right)^{n/2}\right|
 \le K_N\left|1-\frac z\rho\right|^{(N+1)/2},
 \qquad z\in\Delta .
 \label{p1:eq:local-remainder}
\end{equation}
\end{theorem}

\begin{proof}
By \cref{p1:prop:R-analytic} the map $s\mapsto\theta(s)$ is analytic on a disc
$|s|<\varepsilon_0$; by \cref{p1:eq:theta1} it has a simple zero at $s=0$ with $\theta_1>0$.
Shrinking $\varepsilon_0$, $|U(s)|<\tfrac12$ there. The series \cref{p1:eq:cayley-puiseux}
converges for $|\upsilon|$ small; a Cauchy estimate on the analytic function $\upsilon\mapsto
\mathcal C((1-\upsilon)/\e)$ restricted to a slit disc gives $|c_r|\le K\varepsilon_1^{-r/2}$.
Composing convergent expansions and using Taylor's theorem with remainder on $(1+U)^{r/2}$
gives \cref{p1:eq:local-remainder} on the slit disc $\{|s|<\varepsilon,\ |\arg s|<\pi-\eta\}$,
which is a $\Delta$-domain at $\rho$ in the $z$-variable.
\end{proof}

\section{Transfer to the all-orders coefficient asymptotics}
\label{p1:sec:forward}

\subsection{Exact transfer of a singular monomial}

For $\beta\notin\Nat_0$ and $n\ge0$,
\begin{equation}
 \Coef{z^{n}}\left(1-\frac z\rho\right)^{\beta}
 =\rho^{-n}(-1)^{n}\binom{\beta}{n}
 =\rho^{-n}\frac{\Gamma(n-\beta)}{\Gamma(-\beta)\Gamma(n+1)} .
 \label{p1:eq:exact-transfer}
\end{equation}
For $\beta\in\Nat_0$ the left-hand side vanishes for $n>\beta$ --- the analytic-germ statement
of \cref{p1:rem:even-odd}.  It remains to expand the gamma ratio to all orders.

\begin{lemma}[Gamma-ratio coefficients]
\label{p1:lem:gamma-ratio}
Let $B_m(x)$ be the Bernoulli polynomials, and for $\beta\in\Cplx$ set
\begin{equation}
 \varkappa_m(\beta):=\frac{(-1)^{m+1}}{m(m+1)}
 \bigl(B_{m+1}(-\beta)-B_{m+1}(1)\bigr),\qquad m\ge1,
 \label{p1:eq:varkappa}
\end{equation}
and define $g_r(\beta)$ by
$\exp\bigl(\sum_{m\ge1}\varkappa_m(\beta)v^{m}\bigr)=\sum_{r\ge0}g_r(\beta)v^{r}$,
equivalently
\begin{equation}
 g_0(\beta)=1,\qquad
 g_r(\beta)=\frac1r\sum_{m=1}^{r}m\,\varkappa_m(\beta)\,g_{r-m}(\beta)
 =\sum_{k=0}^{r}\frac1{k!}\OB_{r,k}
 \bigl(\varkappa_1(\beta),\varkappa_2(\beta),\ldots\bigr).
 \label{p1:eq:g-recurrence}
\end{equation}
Then, for fixed $\beta$ and $n\to\infty$ in any sector avoiding the negative real axis,
\begin{equation}
 \frac{\Gamma(n-\beta)}{\Gamma(n+1)}
 \sim n^{-\beta-1}\sum_{r\ge0}\frac{g_r(\beta)}{n^{r}} .
 \label{p1:eq:gamma-ratio}
\end{equation}
Each $g_r$ is a polynomial of degree $2r$ in $\beta$ with rational coefficients; the first
five are
\begin{align}
 g_0&=1,
 &g_1(\beta)&=\frac{\beta(\beta+1)}2,
 &g_2(\beta)&=\frac{\beta(\beta+1)(\beta+2)(3\beta+1)}{24},
 \notag\\
 g_3(\beta)&=\frac{\beta^{2}(\beta+1)^{2}(\beta+2)(\beta+3)}{48},
 \label{p1:eq:g-first}
\end{align}
\begin{equation}
 g_4(\beta)=\frac{\beta(\beta+1)(\beta+2)(\beta+3)(\beta+4)
 \bigl(15\beta^{3}+30\beta^{2}+5\beta-2\bigr)}{5760}.
 \label{p1:eq:g4}
\end{equation}
In particular $\bigl(g_r(\tfrac12)\bigr)_{r\ge0}
=1,\tfrac38,\tfrac{25}{128},\tfrac{105}{1024},
\tfrac{1659}{32768},\tfrac{6237}{262144},\ldots$
\end{lemma}

\begin{proof}
The Euler--Maclaurin/Bernoulli form of Stirling's expansion gives, for fixed $a,b$,
\[
\log\frac{\Gamma(n+a)}{\Gamma(n+b)\,n^{a-b}} \sim\sum_{m\ge1}\frac{(-1)^{m+1}}{m(m+1)}
\frac{B_{m+1}(a)-B_{m+1}(b)}{n^{m}}
\]
as $n\to\infty$ in the stated sectors.  Taking $a=-\beta$, $b=1$ yields
$\log\bigl(\Gamma(n-\beta)/\Gamma(n+1)\bigr)\sim-(\beta+1)\log n
+\sum_m\varkappa_m(\beta)n^{-m}$, and exponentiating an asymptotic series term by term is
legitimate, producing exactly the coefficients of $\exp\bigl(\sum\varkappa_mv^m\bigr)$.  The
recurrence follows from differentiating that generating identity; the Bell form is
\cref{p0:def:obell}. Degree: $\varkappa_m$ has degree $m+1$ in $\beta$ and vanishes at
$\beta=0$, so a monomial $\prod\varkappa_{m_i}$ with $\sum m_i=r$ has degree $\le
r+\#\{i\}\le2r$; the displayed $g_1,\ldots,g_4$ were recomputed here from
\cref{p1:eq:varkappa} by exact rational arithmetic.
\end{proof}

\begin{remark}[Part~0 candidate]
\label{p1:rem:gamma-part0}
\Cref{p1:lem:gamma-ratio} is not specific to rooted trees and is used verbatim wherever a
chapter transfers an algebraic singularity to coefficient asymptotics.  It is listed in the
report as a candidate for promotion into Part~0.  All three siblings state it; \textsf{S2} and
\textsf{S3} give the same recurrence, \textsf{S1} the same object under the name
$\gamma_m(\beta)$ with $\eta_r(\beta)$ for $\varkappa_r(\beta)$.
\end{remark}

\subsection{The master coefficient formula}

\begin{theorem}[All-orders dominant expansion of A000081]
\label{p1:thm:forward}
For every fixed $M\ge0$,
\begin{equation}
 \boxed{\;
 a_n=\frac{\rho^{-n}}{\sqrt{\pi}\,n^{3/2}}
 \left(\sum_{m=0}^{M}\frac{\tau_m}{n^{m}}+\bigO(n^{-M-1})\right),
 \qquad n\to\infty,\;}
 \label{p1:eq:forward}
\end{equation}
where
\begin{equation}
 \boxed{\;
 \tau_m=\sum_{j=0}^{m}\Delta_j\,t_{2j+1}\,
 g_{m-j}\!\left(j+\tfrac12\right),
 \qquad
 \Delta_j:=\frac{(-1)^{j+1}(2j+1)!!}{2^{j+1}}
 =\frac{\sqrt\pi}{\Gamma(-j-\tfrac12)} .\;}
 \label{p1:eq:tau-master}
\end{equation}
Every ingredient is a finite coefficient extraction: $t_{2j+1}$ by \cref{p1:eq:t-master},
$g_r$ by \cref{p1:eq:g-recurrence}, $\theta_m$ by \cref{p1:eq:theta-bell}, $\omega_m$ by
\cref{p1:eq:omega-master}, $\varrho_r$ by \cref{p1:eq:varrho}.
\end{theorem}

\begin{proof}
Fix $M$ and apply \cref{p1:thm:local-remainder} with $N=2M+2$: on a $\Delta$-domain at $\rho$,
\[
T(z)=A(z)+\sum_{j=0}^{M}t_{2j+1}\left(1-\frac z\rho\right)^{j+1/2} +\bigO\!\left(\left|1-\frac
z\rho\right|^{M+3/2}\right),
\]
with $A$ the polynomial collecting the even (analytic) terms.  By
\cref{p1:prop:unique-dominant} the only singularity on $|z|=\rho$ is $\rho$ itself, so the
Cauchy contour may be deformed into the standard Flajolet--Odlyzko $\Delta$-contour.  The
polynomial $A$ contributes nothing to $\Coef{z^n}$ for $n>\deg A$.  Each retained monomial
contributes exactly $\rho^{-n}t_{2j+1}\Gamma(n-j-\tfrac12)/
\bigl(\Gamma(-j-\tfrac12)\Gamma(n+1)\bigr)$ by \cref{p1:eq:exact-transfer}, whose asymptotic
expansion is \cref{p1:eq:gamma-ratio}; the $\bigO$-term transfers to
$\bigO(\rho^{-n}n^{-M-5/2})$ by the standard Hankel-contour estimate. Collecting the
coefficient of $\rho^{-n}n^{-m-3/2}$ from the pairs $j+r=m$ and multiplying through by
$\sqrt\pi$ gives \cref{p1:eq:tau-master}, the identity
$\Gamma(-j-\tfrac12)=(-4)^{j+1}(j+1)!\sqrt\pi/(2j+2)! =(-1)^{j+1}2^{j+1}\sqrt\pi/(2j+1)!!$
supplying $\Delta_j$.
\end{proof}

\begin{corollary}[Otter's constant]
\label{p1:cor:otter}
With $C:=\tau_0/\sqrt\pi$,
\begin{equation}
 \tau_0=-\frac{t_1}{2}=\sqrt{\frac{\theta_1}{2}},
 \qquad
 C=\sqrt{\frac{\theta_1}{2\pi}}=\sqrt{\frac{\e\rho\,\zeta'(\rho)}{2\pi}}
 =\sqrt{\frac{1+\rho R'(\rho)}{2\pi}},
 \qquad
 a_n\sim C\,\alpha^{n}n^{-3/2}.
 \label{p1:eq:otter}
\end{equation}
Numerically
\begin{equation}
 \tau_0=0.7797450101873204427711036979\ldots,\qquad
 C=0.4399240125710253040409033914345447648\ldots .
 \label{p1:eq:C-numeric}
\end{equation}
\end{corollary}

\begin{proof}
Set $m=0$ in \cref{p1:eq:tau-master}: $\Delta_0=-\tfrac12$, $g_0=1$, so
$\tau_0=-t_1/2=\sqrt{2\theta_1}/2=\sqrt{\theta_1/2}$ by \cref{p1:cor:t-low}. Divide by
$\sqrt\pi$ and use \cref{p1:eq:theta1}.
\end{proof}

The identity $C^{2}=\theta_1/(2\pi)$ is a genuinely independent check: the left side comes
through the whole Puiseux/transfer chain, the right side directly from one derivative of
$\zeta$.  The two agree to all $55$ digits carried.

Expanding \cref{p1:eq:tau-master} for small $m$ gives the transfer identities
\begin{equation}
\begin{aligned}
 \tau_0&=-\tfrac12t_1,\\
 \tau_1&=-\tfrac3{16}\bigl(t_1-4t_3\bigr)
 =-\tfrac3{16}t_1+\tfrac34t_3,\\
 \tau_2&=-\tfrac5{256}\bigl(5t_1-72t_3+96t_5\bigr)
 =-\tfrac{25}{256}t_1+\tfrac{45}{32}t_3-\tfrac{15}8t_5,\\
 \tau_3&=-\tfrac{105}{2048}\bigl(t_1-44t_3+160t_5-128t_7\bigr)
 =-\tfrac{105}{2048}t_1+\tfrac{1155}{512}t_3-\tfrac{525}{64}t_5
 +\tfrac{105}{16}t_7,\\
 \tau_4&=-\tfrac{21}{65536}
 \bigl(79t_1-10800t_3+81600t_5-161280t_7+92160t_9\bigr).
\end{aligned}
\label{p1:eq:tau-expanded}
\end{equation}
\textsf{S2} states the right-hand forms, \textsf{S3} the factored left-hand forms; they are
identical after clearing denominators, and both were re-derived here symbolically from
\cref{p1:eq:tau-master}.  For any $m$ beyond $4$ the compact convolution is far preferable to
the expanded form.

\begin{table}[htbp]
\centering
\caption{Coefficients of the dominant expansion
$a_n\sim\rho^{-n}(\pi n^{3})^{-1/2}\sum_{m\ge0}\tau_mn^{-m}$ and their
normalisations $f_m=\tau_m/\tau_0$, recomputed here.}
\label{p1:tab:tau}
\small
\setlength{\tabcolsep}{6pt}
\begin{tabular}{r@{\quad}r@{\qquad}r}
\toprule
$m$ & $\tau_m$ & $f_m=\tau_m/\tau_0$\\
\midrule
$0$ & $0.7797450101873204427711$ & $1$\\
$1$ & $0.07828911261061096206128$ & $0.1004034800964014347637$\\
$2$ & $0.3929402676631860184743$ & $0.5039343151023035331022$\\
$3$ & $1.537879315978838091102$ & $1.972284908382278469206$\\
$4$ & $8.200844090435596159220$ & $10.51734090413157399274$\\
$5$ & $57.29291473494343787740$ & $73.47647498402047106417$\\
$6$ & $503.0445050262735818248$ & $645.1397552456612928153$\\
$7$ & $5359.600933884326033403$ & $6873.530274463408619941$\\
$8$ & $67342.06920114653048341$ & $86364.21948371140797700$\\
$9$ & $975425.4970695924732213$ & $1250954.458605978293018$\\
$10$ & $15996932.49293173312961$ & $20515594.56480361813667$\\
$11$ & $292822531.3353392231146$ & ---\\
$12$ & $5914523441.293932519117$ & ---\\
\bottomrule
\end{tabular}
\end{table}

Higher values continue smoothly; for reference $\tau_{15}=8.078305401468884289\times10^{13}$,
$\tau_{18}=2.073828085209893749\times10^{18}$.

Writing
\begin{equation}
 F(v):=1+\sum_{m\ge1}f_mv^{m},\qquad f_m=\frac{\tau_m}{\tau_0},
 \label{p1:eq:F-def}
\end{equation}
the inversion-ready normal form is
\begin{equation}
 \boxed{\;a_n\sim C\,\e^{\lambda n}\,n^{-p}\,F(1/n),
 \qquad p=\tfrac32 .\;}
 \label{p1:eq:forward-normal}
\end{equation}

\begin{remark}[Divergence and least-term truncation]
\label{p1:rem:divergence}
\Cref{p1:eq:forward} asserts a Poincar\'e expansion at every fixed order; it does \emph{not}
assert convergence, and \cref{p1:tab:tau} shows the coefficients growing rapidly.  The optimal
truncation and the size of the resulting floor are governed by
\cref{p0:thm:optimal-truncation}; the relevant nearest competing exponential scale is
identified in \cref{p1:sec:sectors} as $\rho^{-n/2}$, so the least-term floor at index $n$ is
$\bigO\bigl(\rho^{n/2}\bigr)$ relative to the leading term --- consistent with the error
tables of \cref{p1:sub:numerics}, where the improvement stops at $n=20$ around $M=6$.
\textsf{S3} states this correctly as a numerical strategy rather than a theorem; \textsf{S1}
and \textsf{S2} call the late-term growth ``the perturbative shadow of farther
singularities'', which is the right heuristic but is not proved anywhere in the three drafts
and is not proved here.
\end{remark}

\section{The multiplicative reciprocal}
\label{p1:sec:reciprocal}

If ``inverse'' is read as the reciprocal sequence, the answer is immediate from
\cref{p1:eq:forward-normal} and requires no inversion theory at all.

\begin{proposition}[All-orders reciprocal]
\label{p1:prop:reciprocal}
Write $1/F(v)=1+\sum_{m\ge1}\bar f_mv^{m}$.  Then
\begin{equation}
 \bar f_m=\sum_{k=1}^{m}(-1)^{k}\,\OB_{m,k}(f_1,f_2,\ldots)
 =-\sum_{k=1}^{m}f_k\,\bar f_{m-k}\qquad(\bar f_0=1),
 \label{p1:eq:recip-coeff}
\end{equation}
and for every fixed $M\ge0$
\begin{equation}
 \frac1{a_n}=\frac{\sqrt\pi\,\rho^{n}n^{3/2}}{\tau_0}
 \left(\sum_{m=0}^{M}\frac{\bar f_m}{n^{m}}+\bigO(n^{-M-1})\right).
 \label{p1:eq:recip-asymptotic}
\end{equation}
\end{proposition}

\begin{proof}
The first equality in \cref{p1:eq:recip-coeff} is \cref{p0:lem:power-log} with exponent $-1$;
the second is the Cauchy product $F\cdot F^{-1}=1$.  For \cref{p1:eq:recip-asymptotic}, note
that $a_n>0$ and, by \cref{p1:eq:forward}, $a_n=C\alpha^nn^{-p}(1+\varepsilon_n)$ with
$\varepsilon_n\to0$; hence $1/a_n$ is defined and $1/(1+\varepsilon_n)$ has the asymptotic
expansion obtained by formally inverting the (asymptotic) series for $1+\varepsilon_n$ --- a
legitimate operation because the leading term is $1\ne0$ and each truncation error is
$\bigO(n^{-M-1})$.
\end{proof}

The first reciprocal coefficients, recomputed, are
\begin{equation}
\begin{aligned}
 \bar f_1&=-0.1004034800964014347637,&
 \bar f_2&=-0.4938534562868350540384,\\
 \bar f_3&=-1.872103543737176348103,&
 \bar f_4&=-9.882481221441551170916,\\
 \bar f_5&=-69.51082491336461883211,&
 \bar f_6&=-616.9168645978433816315 .
\end{aligned}
\label{p1:eq:recip-numeric}
\end{equation}

\begin{remark}[Reciprocal of the full transseries]
\label{p1:rem:recip-multisector}
Writing $a_n=\mathcal M_n(1+\mathcal E_n)$ with $\mathcal M_n$ the complete dominant sector
and $\mathcal E_n$ the sum of relative exponentially small sectors of \cref{p1:sec:sectors},
one has $1/a_n=\mathcal M_n^{-1}\sum_{q\ge0}(-\mathcal E_n)^{q}$; the products generate
exactly the additive semigroup of the singular actions, with coefficients given by finite
Cauchy products.  This is a formal statement about the transseries, and inherits whatever
status \cref{p1:thm:finite-level} has on the chosen continuation domain.
\end{remark}

\begin{warning}
\label{p1:rem:reciprocal}
The reciprocal is \emph{not} a functional inverse.  Its exponential scale is $\rho^{n}$, while
the functional inverse of \cref{p1:sec:inverse} grows only logarithmically in its argument.
All three siblings make this point; it is repeated here because the ambiguity of the word
``inverse'' is exactly what \cref{p0:def:three-inverses} exists to remove.
\end{warning}

\section{The asymptotic inverse: the Part~0 dictionary}
\label{p1:sec:inverse}

\subsection{The parameter dictionary}

Everything needed to invert \cref{p1:eq:forward-normal} has been proved in Part~0.  What this
chapter must supply is only the dictionary and the verification of Part~0's hypotheses.

\begin{proposition}[A000081 as an instance of the exponential--power model]
\label{p1:prop:dictionary}
Let
\begin{equation}
 \mathcal A(x):=C\,\e^{\lambda x}x^{-p}F(x^{-1}),
 \qquad p=\tfrac32 ,
 \label{p1:eq:interpolant}
\end{equation}
with $C$ from \cref{p1:cor:otter} and $F$ from \cref{p1:eq:F-def}.  Then $\mathcal A$ is an
instance of \cref{p0:def:model} under
\begin{equation}
 \boxed{\;
 \begin{array}{c|cccccc}
  \text{Part 0} & C & a & \alpha & b & \delta & Q(t)\\\hline
  \text{here} & C=\sqrt{\theta_1/(2\pi)} & \lambda=\log(1/\rho) & 1
  & -\tfrac32 & 1 & H(t)=\log F(t)
 \end{array}\;}
 \label{p1:eq:dictionary}
\end{equation}
with $H\in t\Real[[t]]$.  In particular $a=\lambda>0$ and $\alpha=1$, so $\mathcal A$ is
already in the linear-phase normal form and \cref{p0:lem:alpha-reduction} is the identity
substitution.
\end{proposition}

\begin{proof}
\Cref{p1:eq:forward-normal} is \cref{p0:def:model} with the stated substitutions; $\lambda>0$
because $\rho<1$ (\cref{p1:prop:critical}), $C>0$ because $\theta_1>0$, and $H(0)=\log
F(0)=\log1=0$ so $H\in t\Real[[t]]$.
\end{proof}

\begin{definition}[Additive phase coefficients]
\label{p1:def:phi}
Write $H(v)=\log F(v)=\sum_{m\ge1}\varphi_mv^{m}$.  By \cref{p0:lem:power-log},
\begin{equation}
 \varphi_m=\sum_{k=1}^{m}\frac{(-1)^{k-1}}{k}\OB_{m,k}(f_1,f_2,\ldots)
 =f_m-\frac1m\sum_{k=1}^{m-1}k\,\varphi_k\,f_{m-k}.
 \label{p1:eq:phi}
\end{equation}
\end{definition}

Taking logarithms in \cref{p1:eq:forward} gives the form used for inversion,
\begin{equation}
 \log a_n=\lambda n-\tfrac32\log n+\log C
 +\sum_{m=1}^{M}\frac{\varphi_m}{n^{m}}+\bigO(n^{-M-1}),
 \label{p1:eq:log-a}
\end{equation}
with, recomputed here,
\begin{equation}
\begin{aligned}
 \varphi_1&=0.1004034800964014347637,&
 \varphi_2&=0.4988938856945692935703,\\
 \varphi_3&=1.922025533841821821591,&
 \varphi_4&=10.19739642337376887053,\\
 \varphi_5&=71.47146706207073604412,&
 \varphi_6&=630.8599376073372478899 .
\end{aligned}
\label{p1:eq:phi-numeric}
\end{equation}

\subsection{The exact Lambert carrier}

\begin{corollary}[Carrier; specialisation of \cref{p0:thm:lambert-core}]
\label{p1:cor:carrier}
Let $y>y_\ast:=C\left(\e\lambda/p\right)^{p} =1.2108232025837572603\ldots$. Then the equation
$y=C\e^{\lambda x}x^{-p}$ has exactly one solution $x_0=x_0(y)>p/\lambda$, namely
\begin{equation}
 \boxed{\;
 x_0(y)=-\frac p\lambda\,
 \Wlam_{-1}\!\left(-\frac\lambda p\left(\frac Cy\right)^{1/p}\right)
 =-\frac{3}{2\lambda}\,
 \Wlam_{-1}\!\left(-\frac{2\lambda}{3}\left(\frac Cy\right)^{2/3}\right).\;}
 \label{p1:eq:carrier}
\end{equation}
\end{corollary}

\begin{proof}
Raise $y=C\e^{\lambda x}x^{-p}$ to the power $-1/p$ and multiply by $-\lambda/p$:
\[
-\frac\lambda p\left(\frac Cy\right)^{1/p} =-\frac{\lambda x}{p}\,\e^{-\lambda x/p}
=\xi\e^{\xi},\qquad \xi:=-\frac{\lambda x}{p}.
\]
Write $\epsilon:=-\frac\lambda p(C/y)^{1/p}$.  The map $x\mapsto C\e^{\lambda x}x^{-p}$ is
strictly increasing on $(p/\lambda,\infty)$ (its logarithmic derivative is $\lambda-p/x>0$
there) with infimum $y_\ast$ attained as $x\downarrow p/\lambda$, so $y>y_\ast$ gives a unique
root $x_0>p/\lambda$, i.e.\ $\xi<-1$; and $y>y_\ast$ is exactly $\epsilon\in(-\e^{-1},0)$.  On
that range $\xi\mapsto\xi\e^{\xi}$ has the two real inverse branches $\Wlam_0\in(-1,0)$ and
$\Wlam_{-1}<-1$, and the required root is the second. This is \cref{p0:thm:lambert-core} with
$L=\log(y/C)$, $a=\lambda$, $b=-p$ and branch index $k=-1$; the branch rule there selects
$k=-1$ under exactly the condition $x>-b/a$.
\end{proof}

\begin{remark}
The principal branch $\Wlam_0$ returns the other, small root of the same equation, which is
the spurious solution: all three siblings say so, and the threshold $y_\ast$ makes the
dichotomy quantitative.  For A000081 the threshold lies below $a_3=2$, so \cref{p1:eq:carrier}
may be used at every sequence value from $n=3$ on.
\end{remark}

\subsection{All algebraic corrections}

\begin{theorem}[Corrections; specialisation of
\cref{p0:thm:lambert-centered}]
\label{p1:thm:d}
Let
\begin{equation}
 \Psi(t,u):=\frac p\lambda\log(1+tu)
 -\frac1\lambda H\!\left(\frac{t}{1+tu}\right)
 \in t\Real[[t,u]] .
 \label{p1:eq:Psi}
\end{equation}
Then there is a unique $D(t)=\sum_{m\ge1}d_mt^{m}\in t\Real[[t]]$ with
$D(t)=\Psi\bigl(t,D(t)\bigr)$, its coefficients are given by \cref{p0:lem:lagrange-fp},
\begin{equation}
 \boxed{\;
 d_m=\sum_{k=1}^{m}\frac1k\,\Coef{t^{m}u^{k-1}}\Psi(t,u)^{k},\;}
 \label{p1:eq:d-master}
\end{equation}
equivalently by the triangular extraction
\begin{equation}
 d_m=\Coef{t^{m}}\left\{
 \frac p\lambda\log\bigl(1+tD_{<m}(t)\bigr)
 -\frac1\lambda H\!\left(\frac{t}{1+tD_{<m}(t)}\right)\right\},
 \qquad D_{<m}=\sum_{j<m}d_jt^{j},
 \label{p1:eq:d-triangular}
\end{equation}
or by the completely explicit ordinary-Bell form
\begin{equation}
 d_m=\frac p\lambda\sum_{r=1}^{m}\frac{(-1)^{r-1}}{r}
 \OB_{m-r,r}(d_1,d_2,\ldots)
 -\frac1\lambda\sum_{k=1}^{m}\varphi_k
 \sum_{r=0}^{m-k}(-1)^{r}\binom{k+r-1}{r}
 \OB_{m-k-r,r}(d_1,d_2,\ldots),
 \label{p1:eq:d-bell}
\end{equation}
in which the index constraints force only $d_1,\ldots,d_{m-1}$ to occur on the right.  With
$x_0$ from \cref{p1:eq:carrier} the complete algebraic inverse is
\begin{equation}
 \boxed{\;\nu(y)\sim x_0(y)+\sum_{m\ge1}\frac{d_m}{x_0(y)^{m}} .\;}
 \label{p1:eq:inverse-series}
\end{equation}
\end{theorem}

\begin{proof}
\Cref{p0:thm:lambert-centered} states exactly this for a model
$C\e^{ax}x^{b}\e^{Q(x^{-\delta})}$ with $\Psi(t,u)
=-\frac1a\bigl[b\log(1+tu)+Q\bigl(t/(1+tu)\bigr)\bigr]$; substituting the dictionary
\cref{p1:eq:dictionary} ($a=\lambda$, $b=-p$, $\delta=1$, $Q=H$) gives \cref{p1:eq:Psi}.
Concretely: putting $x=x_0+D$ into $\log(y/C)=\lambda x-p\log x+H(x^{-1})$ and subtracting the
carrier equation $\log(y/C)=\lambda x_0-p\log x_0$ gives, with $t=x_0^{-1}$,
\[
\lambda D-p\log(1+tD)+H\!\left(\frac{t}{1+tD}\right)=0,
\]
which is $D=\Psi(t,D)$.  Since $\Psi\in t\Real[[t,u]]$, the formal implicit-function theorem
gives a unique solution in $t\Real[[t]]$; \cref{p0:lem:lagrange-fp} converts it into
\cref{p1:eq:d-master}, and \cref{p1:eq:d-triangular} is the same statement read as a
recursion. \Cref{p1:eq:d-bell} follows by expanding
$\log(1+tD)=\sum_r\frac{(-1)^{r-1}}{r}t^rD^r$ and
$H\bigl(t/(1+tD)\bigr)=\sum_k\varphi_kt^k\sum_r(-1)^r\binom{k+r-1}{r}t^rD^r$ and applying
\cref{p0:def:obell} to the powers $D^r$.
\end{proof}

\begin{repairbox}
\textbf{A formal-to-analytic slide in \textsf{S1}.}  \textsf{S1}'s ``General coefficient
formula for the inverse'' is stated as $x(y)\sim X(y)+\sum_md_mX(y)^{-m}$, an asymptotic
assertion, but its proof ends with ``substitution \dots\ proves formal inversion order by
order'' and supplies no remainder.  As stated, \cref{p1:thm:d} is a theorem about formal
series only; the asymptotic reading requires \cref{p1:thm:inverse-remainder} below, whose
argument is \textsf{S3}'s --- the only one of the three drafts that separates the two
statements.  The volume keeps them separate.
\end{repairbox}

\begin{theorem}[Remainder; specialisation of
\cref{p0:thm:remainder-transport}]
\label{p1:thm:inverse-remainder}
Fix $M\ge1$ and set $\nu_M(y):=x_0(y)+\sum_{m=1}^{M-1}d_mx_0(y)^{-m}$.  Let $f(x)=\log
C+\lambda x-p\log x+H(x^{-1})$ be the (asymptotic) logarithm of the forward interpolant.  Then
\begin{equation}
 f\bigl(\nu_M(y)\bigr)-\log y=\bigO\bigl(x_0(y)^{-M}\bigr),
 \qquad
 \nu(y)-\nu_M(y)=\bigO\bigl(x_0(y)^{-M}\bigr)=
 \bigO\bigl((\log y)^{-M}\bigr).
 \label{p1:eq:inverse-remainder}
\end{equation}
\end{theorem}

\begin{proof}
By \cref{p1:thm:d} the coefficients of $t^{1},\ldots,t^{M-1}$ in $D-\Psi(t,D)$ vanish, so
truncating leaves a residual $\bigO(t^{M})$; this is the first claim. Since
$f'(x)=\lambda-p/x+\bigO(x^{-2})$ is bounded away from $0$ for large $x$,
\cref{p0:thm:backward-error} converts the residual into a forward error of the same order
divided by $f'$; this is \cref{p0:thm:remainder-transport} with $a=\lambda$.  Finally
$x_0(y)=\lambda^{-1}\log y\,(1+\littleo(1))$ by \cref{p1:eq:carrier}.
\end{proof}

\subsection{Coefficients}

Expanding \cref{p1:eq:d-triangular} through $t^{5}$ --- and verifying the result symbolically
--- gives the closed forms
\begin{align}
 d_1&=-\frac{\varphi_1}{\lambda},
 \label{p1:eq:d1}\\
 d_2&=-\frac{p\varphi_1+\lambda\varphi_2}{\lambda^{2}},
 \label{p1:eq:d2}\\
 d_3&=-\frac{\lambda\varphi_1^{2}+p^{2}\varphi_1+p\lambda\varphi_2
 +\lambda^{2}\varphi_3}{\lambda^{3}},
 \label{p1:eq:d3}\\
 d_4&=-\frac{1}{2\lambda^{4}}
 \Bigl(5p\lambda\varphi_1^{2}+6\lambda^{2}\varphi_1\varphi_2
 +2p^{3}\varphi_1+2p^{2}\lambda\varphi_2+2p\lambda^{2}\varphi_3
 +2\lambda^{3}\varphi_4\Bigr),
 \label{p1:eq:d4}\\
 d_5&=-\frac{1}{2\lambda^{5}}
 \Bigl(4\lambda^{2}\varphi_1^{3}+9p^{2}\lambda\varphi_1^{2}
 +14p\lambda^{2}\varphi_1\varphi_2+8\lambda^{3}\varphi_1\varphi_3
 +4\lambda^{3}\varphi_2^{2}\notag\\
 &\hspace{16mm}
 +2p^{4}\varphi_1+2p^{3}\lambda\varphi_2+2p^{2}\lambda^{2}\varphi_3
 +2p\lambda^{3}\varphi_4+2\lambda^{4}\varphi_5\Bigr).
 \label{p1:eq:d5}
\end{align}
Beyond $m=5$ the expanded expressions grow fast, and \cref{p1:eq:d-master,p1:eq:d-triangular}
are both shorter and less error-prone.

\begin{table}[htbp]
\centering
\caption{Additive phase coefficients $\varphi_m$ and Lambert-carrier
corrections $d_m$, recomputed here.  The last column is
\textsf{S2}'s normalisation $\mathrm D_m=\lambda^{m+1}d_m$
(see \cref{p1:rem:normalisations}).}
\label{p1:tab:inverse}
\small
\setlength{\tabcolsep}{6pt}
\begin{tabular}{r@{\quad}r@{\qquad}r@{\qquad}r}
\toprule
$m$ & $\varphi_m$ & $d_m$ & $\mathrm D_m=\lambda^{m+1}d_m$\\
\midrule
$1$ & $0.1004034800964014348$ & $-0.09264385352561312984$ & $-0.1088130343442745145$\\
$2$ & $0.4988938856945692936$ & $-0.5885630399313472628$ & $-0.7491856512881932495$\\
$3$ & $1.922025533841821822$ & $-2.596680167407531021$ & $-3.582177326832277789$\\
$4$ & $10.19739642337376887$ & $-13.14905339996706569$ & $-19.65872121138776004$\\
$5$ & $71.47146706207073604$ & $-85.49870151462556652$ & $-138.5327478518560861$\\
$6$ & $630.8599376073372479$ & $-711.2629148329465609$ & $-1248.979326371692854$\\
$7$ & $6753.629059387619229$ & $-7306.788270284967257$ & $-13905.40884474814835$\\
$8$ & $85171.04383262353790$ & $-89561.76088263724860$ & $-184719.1911265080905$\\
$9$ & $1236987.135642947528$ & $-1274829.503083708903$ & ---\\
$10$ & $20325777.10708850337$ & $-20641099.00032138544$ & ---\\
\bottomrule
\end{tabular}
\end{table}

\begin{repairbox}
\textbf{Numerical error in \textsf{S2}'s audit appendix.}  \textsf{S2}'s appendix table of
$\eta_j$ and $D_j$ reports $D_3=-3.582177326832123\ldots$, disagreeing with its own main text
($-3.582177326832277789103009844$) from the fourteenth significant digit. Recomputation
confirms the main-text value, $\mathrm D_3=-3.58217732683227778910301\ldots$; the appendix
entry is wrong. Every other entry of that table was checked and is correct to the digits
shown.
\end{repairbox}

\begin{remark}[Three normalisations of the same inverse]
\label{p1:rem:normalisations}
The siblings work in three different variables, which is why their coefficient lists look
different:
\begin{itemize}
\item \textsf{S1} and \textsf{S3} expand in the index variable $x$ itself and
      obtain the $d_m$ of \cref{p1:tab:inverse}; their carriers agree with
      \cref{p1:eq:carrier}.
\item \textsf{S2} rescales to $X=\lambda x$, so that its carrier is
      $X_0=-p\Wlam_{-1}\bigl(-\tfrac1p\e^{-Z/p}\bigr)$ with
      $Z=\log\bigl(y/(C\lambda^{p})\bigr)$, its corrections are
      $\mathrm D_m=\lambda^{m+1}d_m$, and its phase coefficients are
      $\eta_m^{\textsf{S2}}=\lambda^{m}\varphi_m$.  Both dictionaries were
      verified numerically for $1\le m\le8$, and $X_0=\lambda x_0$ holds
      identically because the two carrier equations differ by that
      substitution.
\item \textsf{S3} also uses $x$ but writes the polynomial--logarithmic form
      in terms of $\mathsf q=p/\lambda$ and
      $\eta_m^{\textsf{S3}}=\varphi_m/\lambda$.
\end{itemize}
The constant appearing in \textsf{S2}'s $Z$ is
$C\lambda^{p}=0.4963361416595997474727369464\ldots$.
\end{remark}

\subsection{Polynomial--logarithmic form}

Eliminating the Lambert function gives the expansion in the canonical logarithmic scale.  That
is \cref{p0:thm:flattening}; only the dictionary is recorded here.

\begin{corollary}[Flattened inverse]
\label{p1:cor:flattening}
Put
\begin{equation}
 \mathcal X:=\frac1\lambda\log\frac yC,\qquad
 L:=\log\mathcal X,\qquad
 \mathsf q:=\frac p\lambda,\qquad
 \eta_k:=\frac{\varphi_k}{\lambda} .
 \label{p1:eq:flatten-dict}
\end{equation}
Then the inverse problem is $x-\mathsf q\log x+\sum_{k\ge1}\eta_kx^{-k} =\mathcal X$, and
\cref{p0:thm:flattening} gives
\begin{equation}
 \boxed{\;
 \nu(y)\sim\mathcal X+\mathsf qL
 +\sum_{m\ge1}\frac{P_m(L)}{\mathcal X^{m}},\;}
 \label{p1:eq:flattened}
\end{equation}
where each $P_m$ is a polynomial of degree $m$ over $\Rat[\mathsf q,\eta_1,\ldots,\eta_m]$
with
\begin{equation}
 \Coef{L^{m}}P_m(L)=\frac{(-1)^{m+1}}{m}\mathsf q^{m+1},
 \qquad P_m(0)=d_m .
 \label{p1:eq:P-shape}
\end{equation}
\end{corollary}

The identity $P_m(0)=d_m$ is not an accident: setting $L=0$ in the flattening recursion
reduces it to the fixed-point equation solved by the Lambert carrier, and formal uniqueness in
\cref{p0:thm:lambert-centered} identifies the two.  The first two polynomials are
\begin{align}
 P_1(L)&=\mathsf q^{2}L-\eta_1
 =1.9156590172195452616\,L-0.0926438535256131298,
 \label{p1:eq:P1}\\
 P_2(L)&=-\tfrac12\mathsf q^{3}L^{2}
 +\bigl(\mathsf q^{3}+\mathsf q\eta_1\bigr)L
 -\mathsf q\eta_1-\eta_2\notag\\
 &=-1.3257062894576032164\,L^{2}
 +2.7896427837264004647\,L
 -0.5885630399313472628 .
 \label{p1:eq:P2}
\end{align}

\begin{remark}[The two flattenings are not the same polynomials]
\label{p1:rem:two-flattenings}
\textsf{S2} flattens in $Z=\log\bigl(y/(C\lambda^{p})\bigr)$ and $\log Z$; \textsf{S3}
flattens in $\mathcal X=\lambda^{-1}\log(y/C)$ and $\log\mathcal X$. The two recursions are
formally the \emph{same} recursion with different parameters --- \textsf{S3}'s is
\textsf{S2}'s under $(p,\eta^{\textsf{S2}}_k)\mapsto(\mathsf q,\eta^{\textsf{S3}}_k)$ --- but
they expand the same function in \emph{different} variables, so their coefficient polynomials
genuinely differ.  Neither draft is in error; the caution is only that the two lists must not
be compared term by term.  We adopt \textsf{S3}'s normalisation because it keeps the index
variable $x$ throughout and so matches \cref{p1:eq:inverse-series}.
\end{remark}

Retaining only the first two scales,
\begin{equation}
 \nu(y)=\frac{\log y}{\lambda}+\frac p\lambda\log\log y
 -\frac1\lambda\log\bigl(C\lambda^{p}\bigr)+\littleo(1),
 \label{p1:eq:inverse-elementary}
\end{equation}
that is
\begin{equation}
 \nu(y)=0.9227155616186015248\log y
 +1.3840733424279022872\log\log y
 +0.6463639827002088073+\littleo(1).
 \label{p1:eq:inverse-elementary-numeric}
\end{equation}
\Cref{p1:eq:inverse-series} is preferable numerically because $\Wlam_{-1}$ resums the entire
logarithmic tower exactly before any truncation is made.

\section{The discrete inverse}
\label{p1:sec:discrete}

The objects distinguished in \cref{p0:def:three-inverses} are, for this sequence: the smooth
inverse $\nu$ of \cref{p1:sec:inverse}; the integer staircase
\begin{equation}
 N(y):=\min\{n\ge1: a_n\ge y\};
 \label{p1:eq:staircase}
\end{equation}
and the round-to-nearest generalized inverse.  Applying \cref{p0:thm:staircase} requires its
separation hypothesis, which for A000081 is supplied by the next two lemmas.

\begin{lemma}[Strict monotonicity]
\label{p1:lem:monotone}
$a_1=a_2=1$ and $a_{n+1}>a_n$ for every $n\ge2$.
\end{lemma}

\begin{proof}
Monotonicity: the map sending a rooted tree $\mathfrak t$ on $n$ vertices to the rooted tree
whose root has the single child subtree $\mathfrak t$ is an injection from the $n$-vertex
shapes into the $(n+1)$-vertex shapes, so $a_{n+1}\ge a_n$.  Strictness for $n\ge2$: when
$n+1\ge3$, the star $K_{1,n}$ rooted at its centre has $n\ge2$ children and is therefore not
in the image of that injection, whence $a_{n+1}\ge a_n+1$.  The values $a_1=a_2=1$ are
immediate.
\end{proof}

\begin{lemma}[Separation]
\label{p1:lem:separation}
$a_{n+1}/a_n\to\alpha=2.9557\ldots>1$; more precisely
\begin{equation}
 \log\frac{a_{n+1}}{a_n}
 =\lambda-\frac{3}{2n}+\frac{\tfrac34-\varphi_1}{n^{2}}
 +\frac{-\tfrac12+\varphi_1-2\varphi_2}{n^{3}}+\bigO(n^{-4}).
 \label{p1:eq:ratio}
\end{equation}
\end{lemma}

\begin{proof}
Subtract \cref{p1:eq:log-a} at $n$ from the same relation at $n+1$, and expand $\log(n+1)=\log
n+n^{-1}-\tfrac12n^{-2}+\tfrac13n^{-3}-\cdots$ and $(n+1)^{-m}=n^{-m}(1-mn^{-1}+\cdots)$.
Collecting powers of $n^{-1}$ gives \cref{p1:eq:ratio}; the limit statement is its leading
term.
\end{proof}

\begin{corollary}[Staircase recovery]
\label{p1:cor:staircase}
Let $\nu_M$ be as in \cref{p1:thm:inverse-remainder}.  For every $M\ge2$ and all sufficiently
large $n$, $\bigl|\nu_M(a_n)-n\bigr|<\tfrac12$, so $n$ is recovered from $a_n$ by rounding.
For a general real target $y>1$, $N(y)=\nu(y)+\bigO(1)$, and certifying $N(y)$ exactly
requires the two-sided test $a_{m-1}<y\le a_m$, which \cref{p1:eq:inverse-series} reduces to a
bounded neighbourhood of $m\approx\nu(y)$.
\end{corollary}

\begin{proof}
The first claim is \cref{p1:thm:inverse-remainder} with $M\ge2$, since $\bigO\bigl((\log
a_n)^{-2}\bigr)=\bigO(n^{-2})$.  For the second, apply \cref{p0:thm:staircase}, whose
separation hypothesis holds by \cref{p1:lem:monotone,p1:lem:separation}: the logarithmic
derivative of the interpolation tends to $\lambda>0$, so multiplying $y$ by a bounded factor
moves $\nu(y)$ by a bounded amount, while consecutive sequence values differ by the
asymptotically fixed factor $\alpha$.  The final sentence is the definition of $N$.
\end{proof}

\begin{warning}[The $\bigO(1)$ rounding layer is not a small correction]
\label{p1:rem:sawtooth}
An assertion of the form $N(y)=\nu(y)+\littleo(1)$ \emph{uniformly in real $y$} is false, and
no descending power--logarithmic series can be true uniformly: arbitrarily small changes of
$y$ across a value $a_n$ change $N(y)$ by one. Once a monotone interpolation $\mathcal
A_{\mathrm{ex}}$ with $\mathcal A_{\mathrm{ex}}(n)=a_n$ has been fixed, one has exactly
$N(y)=\lceil\nu_{\mathrm{ex}}(y)\rceil$ with $\nu_{\mathrm{ex}}$ its ordinary inverse, and for
non-integral $\chi$
\begin{equation}
 \lceil\chi\rceil=\chi+\frac12
 +\frac1\pi\sum_{m\ge1}\frac{\sin(2\pi m\chi)}{m},
 \label{p1:eq:sawtooth}
\end{equation}
an $\bigO(1)$ periodic layer that never tends to zero.  The correct statements are: about
$\nu$; or at the special values $y=a_n$, where $\nu(a_n)=n+\bigO(n^{-M})$ for every fixed $M$;
or about a rounded exact interpolation.  \textsf{S2} states this sharply and \textsf{S3}
restates it; \textsf{S1}'s phrase ``can never be resolved more accurately than one unit'' is
the same point stated loosely, and is corrected here.  The layer \cref{p1:eq:sawtooth} must
not be confused with the exponentially small oscillation of \cref{p1:sec:sectors}: the latter
is genuinely $\littleo(1)$ and has an entirely different origin.
\end{warning}

\begin{remark}[A definitional disagreement]
\label{p1:rem:staircase-index}
\textsf{S1} and \textsf{S2} define the staircase as $\min\{n\ge2:a_n\ge y\}$, \textsf{S3} as
$\min\{n\ge1:a_n\ge y\}$.  The two differ only at $y\le1$, where the tie $a_1=a_2=1$ makes the
choice a pure convention.  This volume uses \cref{p1:eq:staircase} and confines all statements
to $y>1$, where the definitions agree.
\end{remark}

\section{Exponentially small sectors}
\label{p1:sec:sectors}

\subsection{Why one exponential scale is not the whole story}

\Cref{p1:eq:forward} is complete in inverse powers of $n$ at the scale $\rho^{-n}$.  It cannot
see a term $\sigma^{-n}$ with $|\sigma|>\rho$, because $(\rho/|\sigma|)^{n}$ decays faster
than every power of $n^{-1}$. Such terms come from singularities of the analytic continuation
of $T$, and the functional equation exposes exactly two mechanisms for producing them:
\begin{enumerate}
\item \textbf{inherited preimages}: if $T$ is singular at $\tau$ and
      $\sigma^{k}=\tau$ for some $k\ge2$, the summand $T(z^{k})/k$ of $R$ can
      make $R$, hence $T$, singular at $\sigma$;
\item \textbf{new Cayley critical points}: on a continued sheet where $R$ is
      analytic, any $\sigma$ with $\zeta(\sigma)=\e^{-1}$ and
      $\zeta'(\sigma)\ne0$ is a square-root branch point of $T$ by
      \cref{p1:thm:omega}.
\end{enumerate}
Iterating the first mechanism from the dominant point produces the \emph{root-lift skeleton}
\begin{equation}
 \mathscr R:=\bigl\{\rho^{1/m}\e^{2\pi\ii j/m}\;:\;m\ge1,\ 0\le j<m\bigr\},
 \label{p1:eq:skeleton}
\end{equation}
whose radii increase to $1$ and whose arguments become dense.  Relative to the dominant sector
the level-$m$ points carry the exponential weight
\begin{equation}
 \sigma_{m,j}^{-n}=\rho^{-n/m}\e^{-2\pi\ii jn/m},
 \qquad\text{i.e.}\qquad
 \exp\!\left[-\lambda\left(1-\frac1m\right)n\right]
 \text{ relative to }\rho^{-n}.
 \label{p1:eq:skeleton-scale}
\end{equation}

\begin{warning}[What the skeleton is and is not]
\label{p1:rem:skeleton}
\Cref{p1:eq:skeleton} is a list of \emph{candidate} inherited germs generated by the
recursion.  It is not a theorem that every listed point is singular on a given sheet, that
every singularity is listed, or that a listed point is reachable along a prescribed path.
Continuation can also meet solutions of $\zeta(z)=\e^{-1}$ that are not in $\mathscr R$, and a
critical point can coalesce with an inherited germ (which, by \cref{p1:eq:critical-outer}
below, halves an exponent and can create quarter powers).
\end{warning}

\begin{remark}[A disagreement about the unit circle]
\label{p1:rem:natural-boundary}
\textsf{S1} and \textsf{S2} assert that $|z|=1$ is a natural boundary for $T$ and that the
continued singularities accumulate at every point of it, citing Drmota--Gittenberger.
\textsf{S3} makes no such assertion and states only the candidate skeleton with the caveat of
\cref{p1:rem:skeleton}.  This volume records the natural-boundary claim as \emph{reported} and
does not use it: every statement in this section is conditioned on the explicit finite-radius
hypothesis $\mathsf H(\mathcal R)$ of \cref{p1:def:HR}, and nothing below depends on what
happens as $\mathcal R\uparrow1$.
\end{remark}

\subsection{Transfer from a general singular germ}

\begin{theorem}[Universal sector coefficients]
\label{p1:thm:general-sector}
Let $\sigma\ne0$ be an accessible singularity on a specified sheet, put $\mathsf
u=1-z/\sigma$, and suppose that in a $\Delta$-domain at $\sigma$
\begin{equation}
 T(z)=A_\sigma(\mathsf u)
 +\sum_{\beta\in E_\sigma}c_{\sigma,\beta}\,\mathsf u^{\beta}
 +\bigO\bigl(|\mathsf u|^{\beta_{\max}}\bigr),
 \label{p1:eq:general-germ}
\end{equation}
with $A_\sigma$ analytic, $E_\sigma\subset\Rat_{>0}\setminus\Nat_0$ finite and
$\beta_{\max}>\max E_\sigma$.  Then the contribution of \cref{p1:eq:general-germ} to
$\Coef{z^{n}}T$ is
\begin{equation}
 \boxed{\;
 \mathfrak S_\sigma(n)\sim\sigma^{-n}
 \sum_{\beta\in E_\sigma}\frac{c_{\sigma,\beta}}{\Gamma(-\beta)}
 \,n^{-\beta-1}\sum_{r\ge0}\frac{g_r(\beta)}{n^{r}},\;}
 \label{p1:eq:general-sector}
\end{equation}
so the coefficient of $\sigma^{-n}n^{-\beta-1-r}$ is exactly
$c_{\sigma,\beta}\,g_r(\beta)/\Gamma(-\beta)$.
\end{theorem}

\begin{proof}
Apply \cref{p1:eq:exact-transfer} with $\rho$ replaced by $\sigma$ to each displayed monomial
and expand the gamma ratio by \cref{p1:lem:gamma-ratio}.  The analytic part $A_\sigma$ has no
Hankel jump and contributes no large-$n$ tail; the $\bigO$-term transfers by the usual Hankel
estimate.
\end{proof}

\Cref{p1:eq:general-sector} is the exact analogue of \cref{p1:eq:tau-master} and allows
arbitrary ramification exponents, not only half-integers.  Relative to the dominant sector one
introduces the \emph{action}
\begin{equation}
 \Omega_\sigma:=\operatorname{Log}\frac{\sigma}{\rho},
 \qquad\text{so that}\qquad
 \sigma^{-n}=\rho^{-n}\e^{-n\Omega_\sigma}.
 \label{p1:eq:action}
\end{equation}

\begin{definition}[Continuation hypothesis $\mathsf H(\mathcal R)$]
\label{p1:def:HR}
Fix $\rho<\mathcal R<1$.  We say $\mathsf H(\mathcal R)$ holds for a chosen continuation of
$T$ if that continuation is analytic on $|z|<\mathcal R$ after removing finitely many pairwise
disjoint cuts issuing from a finite set $\Sigma_{\mathcal R}$; every
$\sigma\in\Sigma_{\mathcal R}$ carries a germ of the form \cref{p1:eq:general-germ} in a
$\Delta$-domain; and the deformed coefficient contour may be taken on a circle of radius at
least $\mathcal R$, up to harmless indentations, on which $T$ has at most polynomial growth.
\end{definition}

\begin{theorem}[Finite-level coefficient transseries]
\label{p1:thm:finite-level}
Assume $\mathsf H(\mathcal R)$ and let $\mathfrak S_\sigma$ be the multiplier-weighted sector
of \cref{p1:eq:general-sector} attached to the local Hankel cycle actually traversed.  Then,
for any finite truncation of the exponent and inverse-power sums,
\begin{equation}
 a_n=\sum_{\sigma\in\Sigma_{\mathcal R}}
 \mathfrak S_\sigma(n)+\bigO\bigl(\mathcal R^{-n}n^{K}\bigr)
 \label{p1:eq:finite-level}
\end{equation}
for a finite $K$ determined by the boundary growth.
\end{theorem}

\begin{proof}
Deform the Cauchy contour from a small circle about $0$ into the union of small Hankel
contours around the cuts issuing from $\Sigma_{\mathcal R}$ plus an outer contour of radius
$\mathcal R$.  On each Hankel contour apply \cref{p1:thm:general-sector}; the outer contour is
bounded by $\mathcal R^{-n}$ times the allowed polynomial factor.  Only finitely many
singularities and finitely many local terms are retained, so the deformation and the
term-by-term summation are legitimate.  Passing to a different sheet changes the homology
coefficients of the local cycles; those coefficients are the Stokes/intersection multipliers
absorbed into $\mathfrak S_\sigma$.
\end{proof}

\begin{repairbox}
\textbf{\textsf{S1}'s global transseries is not a theorem.}  \textsf{S1} displays
$a_n\sim\mathcal A_\rho(n)+\sum_{\sigma\in\mathfrak S\setminus \{\rho\}}S_\sigma\mathcal
A_\sigma(n)$ as a numbered displayed result over ``the singular set encountered by the chosen
contour continuation'', with no hypothesis attached, while conceding in its introduction that
globally the formula ``is best understood as a formal resurgent/Darboux transseries''. The
volume replaces it by \cref{p1:thm:finite-level}, which is \textsf{S3}'s formulation and the
only one of the three carrying an explicit hypothesis; the unrestricted sum over an infinite
$\mathfrak S$ is not asserted anywhere here. The correct reading of a ``global'' transseries
in the presence of accumulating singularities is the compatible family of the finite-radius
expansions \cref{p1:eq:finite-level} as $\mathcal R$ grows, which is exactly what
\textsf{S2}'s ``projective meaning'' remark says.
\end{repairbox}

\subsection{The local recursion produces every germ}

\begin{lemma}[Pull-back coordinates]
\label{p1:lem:pullback}
Let $z=\sigma(1-\mathsf u)$ and $\tau=\sigma^{k}$.  Then
\begin{equation}
 1-\frac{z^{k}}{\tau}=1-(1-\mathsf u)^{k}
 =k\,\mathsf u\left(1+\sum_{r\ge1}p_{k,r}\mathsf u^{r}\right),
 \qquad
 p_{k,r}=\frac{(-1)^{r}}{k}\binom{k}{r+1},
 \label{p1:eq:pullback}
\end{equation}
and consequently, for a parent monomial,
\begin{equation}
 \left(1-\frac{z^{k}}{\tau}\right)^{\beta}
 =(k\mathsf u)^{\beta}\sum_{N\ge0}\mathsf u^{N}
 \sum_{j=0}^{N}\binom{\beta}{j}
 \OB_{N,j}(p_{k,1},p_{k,2},\ldots).
 \label{p1:eq:pullback-coeff}
\end{equation}
\end{lemma}

\begin{proof}
$1-(1-\mathsf u)^{k}=-\sum_{r\ge1}\binom kr(-\mathsf u)^{r} =k\mathsf u-\sum_{r\ge2}\binom
kr(-1)^{r}\mathsf u^{r}$; dividing by $k\mathsf u$ and re-indexing $r\mapsto r+1$ gives
$p_{k,r}$. \Cref{p1:eq:pullback-coeff} is \cref{p0:lem:power-log} applied to the unit series
$1+\sum_rp_{k,r}\mathsf u^{r}$.
\end{proof}

\begin{theorem}[Recursive completeness of the local calculus]
\label{p1:thm:local-completeness}
Fix a sheet and a point $\sigma\ne0$, and choose a local uniformiser $z=\sigma(1-v^{\mathsf
L})$ in which every parent germ needed in $R(z)$ is an ordinary power series in $v$.  Then the
local germ of $T$ at $\sigma$ is computable to any finite order by the following finite
operations, and only these:
\begin{enumerate}
\item pull back each singular summand $T(z^{k})/k$ by
      \cref{p1:eq:pullback-coeff};
\item add the regular Taylor parts to obtain $R$, and exponentiate by
      \cref{p1:eq:E-bell} to obtain $\zeta=z\e^{R}$;
\item if $\zeta_0:=\zeta(\sigma)\ne\e^{-1}$ on the reached branch, so that
      the selected Cayley branch $\mathcal C_\flat$ is analytic at
      $\zeta_0$, compose regularly:
      \begin{equation}
       T=\sum_{r\ge0}\frac{\mathcal C_\flat^{(r)}(\zeta_0)}{r!}
       \bigl(\zeta-\zeta_0\bigr)^{r},
       \qquad
       \mathcal C_\flat'(w)=\frac{\mathcal C_\flat(w)}
       {w\bigl(1-\mathcal C_\flat(w)\bigr)},
       \label{p1:eq:regular-outer}
      \end{equation}
      the higher derivatives following by Fa\`a di Bruno;
\item if $\zeta_0=\e^{-1}$ and the reached branch has $T(\sigma)=1$, compose
      critically with the universal germ of \cref{p1:thm:omega}:
      \begin{equation}
       T=1+\sum_{r\ge1}c_r\,\theta_\sigma^{r/2},
       \qquad \theta_\sigma:=1-\e\zeta .
       \label{p1:eq:critical-outer}
      \end{equation}
\end{enumerate}
If $\theta_\sigma$ vanishes to order $\varsigma$ in $v$, then $\theta_\sigma^{1/2}$ begins at
$v^{\varsigma/2}$ and the ramification index doubles when $\varsigma$ is odd; this is the only
new ramification mechanism.
\end{theorem}

\begin{proof}
Induction on the depth of $\sigma$ in the pull-back graph.  At depth $0$,
\cref{p1:thm:t-coefficients} gives the dominant germ.  At the induction step, $z\mapsto z^{k}$
converts each already-known germ into a Puiseux series in a common uniformiser by
\cref{p1:eq:pullback-coeff}; the ring $\Cplx[[v]]$ is closed under the sum defining $R$, under
exponentiation and under multiplication by $z=\sigma(1-v^{\mathsf L})$; and the outer
composition is either \cref{p1:eq:regular-outer} or \cref{p1:eq:critical-outer}, both of which
determine every coefficient at finite order by \cref{p0:def:obell,p0:def:ebell}. The
derivative formula in \cref{p1:eq:regular-outer} follows by differentiating $\mathcal
C=w\e^{\mathcal C}$.
\end{proof}

\begin{remark}[A dedup]
\label{p1:rem:generalised-bell}
\textsf{S1} introduces a ``generalized Bell algebra'' $\mathfrak B_{\gamma,m}$ on a
semigroup-graded power-series ring for this step. \Cref{p1:thm:local-completeness} avoids it:
after a common uniformiser $v$ has been chosen, everything is an \emph{ordinary} power series
in $v$ and the standard $\OB$, $\EB$ of \cref{p0:def:obell,p0:def:ebell} suffice.  That is
\textsf{S3}'s organisation and is adopted here; nothing is lost, because the choice of
$\mathsf L$ is exactly what the semigroup grading was encoding.
\end{remark}

\subsection{The first inherited sector, at $-\sqrt\rho$}
\label{p1:sub:minus}

The point
\begin{equation}
 \sigma_-:=-\sqrt\rho=-0.5816544136333942538101\ldots
 \label{p1:eq:sigma-minus}
\end{equation}
is reached by continuing the germ at $0$ along the negative real axis, which never meets the
dominant positive singularity.  At $z=\sigma_-$ the summand $T(z^{2})/2$ of $R$ reaches
$T(\rho)/2$ and inherits the dominant square-root branch, while every $T(z^{k})$ with $k\ge3$
is evaluated at $|\sigma_-|^{k}\le\rho^{3/2}<\rho$, strictly inside the original disc.  Take
the uniformiser
\begin{equation}
 z=\sigma_-(1-v^{2}),\qquad v=\Bigl(1-\tfrac{z}{\sigma_-}\Bigr)^{1/2},
 \qquad\text{so}\qquad
 1-\frac{z^{2}}{\rho}=v^{2}(2-v^{2}).
 \label{p1:eq:v-minus}
\end{equation}

\begin{proposition}[Complete local data at $\sigma_-$]
\label{p1:prop:minus}
Write $R\bigl(\sigma_-(1-v^{2})\bigr)=\sum_{m\ge0}r^-_mv^{m}$ and
$\zeta\bigl(\sigma_-(1-v^{2})\bigr)=\sum_{m\ge0}\zeta^-_mv^{m}$.  Then, with $t_0:=1$,
\begin{equation}
 r^-_m=\frac12\!\!
 \sum_{\substack{0\le r\le m\\ r\equiv m\,(2)}}\!\!
 t_r(-1)^{\frac{m-r}{2}}2^{\frac r2-\frac{m-r}{2}}
 \binom{r/2}{\tfrac{m-r}{2}}
 +\mathbf 1_{2\mid m}(-1)^{m/2}
 \sum_{k\ge3}\frac1k\sum_{n\ge1}a_n\sigma_-^{kn}\binom{kn}{m/2},
 \label{p1:eq:rminus}
\end{equation}
both parts converging geometrically.  Defining $\mathcal E^-_m$ by $\sum_{m\ge0}\mathcal
E^-_mv^{m}=\exp\bigl(\sum_{m\ge0}r^-_mv^{m}\bigr)$, so that $\mathcal
E^-_m=\e^{r^-_0}\sum_{k\ge0}\frac1{k!}\OB_{m,k}(r^-_1,r^-_2,\ldots)$, one has
\begin{equation}
 \zeta^-_m=\sigma_-\bigl(\mathcal E^-_m-\mathcal E^-_{m-2}\bigr),
 \qquad \mathcal E^-_{<0}:=0 .
 \label{p1:eq:zminus}
\end{equation}
Setting $y_-:=\mathcal C_0(\zeta^-_0)=-\Wlam_0(-\zeta^-_0)$, the outer composition is regular
because $y_-\ne1$, and writing $T\bigl(\sigma_-(1-v^{2})\bigr)=y_-+\sum_{m\ge1}t^-_mv^{m}$ one
has the closed form
\begin{equation}
 \boxed{\;r^-_1=\frac{t_1}{\sqrt2},
 \qquad
 t^-_1=\frac{y_-}{1-y_-}\cdot\frac{t_1}{\sqrt2}.\;}
 \label{p1:eq:bminus1}
\end{equation}
\end{proposition}

\begin{proof}
The $k=2$ contribution to $R$ is $\tfrac12T\bigl(\rho(1-v^{2}(2-v^{2}))\bigr)$; inserting
\cref{p1:eq:T-puiseux} with $s=v^{2}(2-v^{2})$, so $s^{r/2}=v^{r}2^{r/2}(1-v^{2}/2)^{r/2}$,
and expanding $(1-v^{2}/2)^{r/2}$ binomially gives the first sum.  The $k\ge3$ contributions
are $\sum_{k\ge3}\frac1k\sum_na_n\sigma_-^{kn}(1-v^{2})^{kn}$, whose $v^{2j}$ coefficient is
the second sum; convergence is geometric already at $k=3$ because $|\sigma_-|^{3}<\rho$.
\Cref{p1:eq:zminus} is $\zeta=\sigma_-(1-v^{2})\e^{R}$. For \cref{p1:eq:bminus1}: only $r=1$
contributes to $v^{1}$ in the first sum, giving $r^-_1=\tfrac12t_1\sqrt2=t_1/\sqrt2$; hence
$\zeta^-_1=\zeta^-_0r^-_1$ and, by \cref{p1:eq:regular-outer} with $\mathcal C_0'(w)=\mathcal
C_0(w)/ \bigl(w(1-\mathcal C_0(w))\bigr)$, $t^-_1=\mathcal C_0'(\zeta^-_0)\zeta^-_1
=\frac{y_-}{\zeta^-_0(1-y_-)}\zeta^-_0r^-_1$, which is the claim.
\end{proof}

Recomputed at $80$ working digits (see the repair below),
\begin{equation}
\begin{aligned}
 r^-_0&=\phantom{-}0.4686366279248455480872,&
 \zeta^-_0&=-0.9293757352012050888501,\\
 y_-&=-0.5410329414204534926205,&
 t^-_1&=\phantom{-}0.3871501110302429034772,
\end{aligned}
\label{p1:eq:minus-numeric}
\end{equation}
and the first local coefficients are
\begin{equation}
\begin{aligned}
 t^-_2&=-0.03256491851564763414731,&
 t^-_3&=\phantom{-}0.01720042160312209148301,\\
 t^-_4&=\phantom{-}0.10407703951838443852546,&
 t^-_5&=\phantom{-}0.16194736077270236118952,\\
 t^-_6&=-0.06292599818799308689694,&
 t^-_7&=-0.24403103923142135135510 .
\end{aligned}
\label{p1:eq:tminus-numeric}
\end{equation}

\begin{repairbox}
\textbf{\textsf{S2}'s minus-sector constants are printed past their accuracy.} \textsf{S2}
prints $r^-_0$, $z^-_0$, $y_-$ and $b^-_1$ to $40$ significant digits, e.g.\
$r^-_0=0.4686366279248455480872026033545827113234$.  Two independent recomputations here ---
one through the germ recursion at $90$ digits, one by direct summation of
$\tfrac12+\sum_{k\ge3}T(\sigma_-^{k})/k$ at $80$ digits with $1400$ exact tree numbers and no
truncation of the inner sums --- agree with each other and give
$r^-_0=0.468636627924845548087202621931819314\ldots$, which departs from \textsf{S2}'s value
at the \emph{twenty-sixth} significant digit.  The same $25$-digit agreement, and no more, is
found for $\zeta^-_0$, $y_-$, $t^-_1$ and every derived quantity, consistent with a single
truncation of the $k\ge3$ tail in \textsf{S2}'s computation.  All minus-sector constants are
therefore quoted here to at most $22$ digits.
\end{repairbox}

Only the odd $t^-_{2j+1}$ are singular in $z$ at $\sigma_-$, so \cref{p1:thm:general-sector}
with $\beta=j+\tfrac12$ gives
\begin{equation}
 \boxed{\;
 a_n^{(-)}\sim\mathfrak s_-\,\sigma_-^{-n}n^{-3/2}
 \sum_{k\ge0}\frac{B^-_k}{n^{k}},
 \qquad
 B^-_k=\sum_{j=0}^{k}\frac{t^-_{2j+1}}{\Gamma(-j-\tfrac12)}
 g_{k-j}\!\left(j+\tfrac12\right),\;}
 \label{p1:eq:minus-sector}
\end{equation}
where $\mathfrak s_-$ is the multiplier of the local cycle; for the direct negative-axis
continuation it is $1$, and it is displayed only to keep the contour dependence visible in
larger deformations.  Since $\sigma_-^{-n}=(-1)^{n}\rho^{-n/2}$, the sector is alternating,
with
\begin{equation}
\begin{aligned}
 B^-_0&=-0.10921302995631017570549,&
 B^-_1&=-0.03367666220778285336719,\\
 B^-_2&=-0.17900090117304631638903,&
 B^-_3&=-1.64234200495327998300,\\
 B^-_4&=-18.9266458632920839673,&
 B^-_5&=-277.723089846678995649 .
\end{aligned}
\label{p1:eq:Bminus-numeric}
\end{equation}
Relative to the leading dominant term its amplitude is
\begin{equation}
 \frac{B^-_0}{C}=-0.2482543049151649529916\ldots,
 \qquad
 \frac{a_n^{(-)}}{a_n^{(0)}}\sim
 -0.2482543049\ldots\,\bigl(-\sqrt\rho\bigr)^{n},
 \label{p1:eq:minus-relative}
\end{equation}
which is smaller than every power of $1/n$: this is the first explicitly computed genuinely
beyond-all-orders correction to A000081.  Its action is
\begin{equation}
 \Omega_-=\operatorname{Log}\frac{\sigma_-}{\rho}
 =\frac\lambda2+\ii\pi \pmod{2\pi\ii},
 \qquad \frac{\Omega_-}{\lambda}=\frac12+\ii\frac{\pi}{\lambda},
 \qquad \frac\pi\lambda=2.8987964297339787377\ldots .
 \label{p1:eq:minus-action}
\end{equation}

\begin{remark}
\label{p1:rem:other-sectors}
The point $+\sqrt\rho$ lies behind the dominant branch point on real continuation and belongs
to sheet-dependent continuations; complex critical points and overlapping root lifts may occur
farther out.  Their local data are produced by \cref{p1:thm:local-completeness} and their
multipliers depend on the chosen contour.  \Cref{p1:eq:minus-sector} is one explicit sector,
not an assertion that it exhausts modulus $\sqrt\rho$.
\end{remark}

\subsection{Transport of the sectors through the inverse}

Write the completed forward interpolation as
\begin{equation}
 \mathfrak a(x)=\mathcal A(x)\bigl(1+\mathcal E(x)\bigr),
 \qquad
 \mathcal E(x)=\sum_{\Omega\ne0}\e^{-\Omega x}x^{\mu_\Omega}
 \sum_{r\ge0}e_{\Omega,r}x^{-r},
 \label{p1:eq:full-interp}
\end{equation}
with $\mathcal A$ from \cref{p1:eq:interpolant} and the support of $\mathcal E$ contained in
the additive semigroup generated by the actions \cref{p1:eq:action}.  For non-integer $x$ a
factor $\sigma^{-x}$ requires a choice of $\operatorname{Log}\sigma$; different choices agree
at integers and differ between them, so the exponentially small part of a \emph{functional}
inverse is not canonical until a continuation and pairing convention has been fixed.  Fix one,
pairing conjugate branches so that $\mathfrak a$ is real on the positive axis.

\begin{proposition}[Sector transport]
\label{p1:prop:transport}
Let $\nu_0$ be the complete algebraic inverse of $\mathcal A$ (\cref{p1:eq:inverse-series})
and put $\nu_{\mathrm{full}}=\nu_0+\Delta$.  Then $\Delta$ is the unique solution, in the
transseries field ordered by the semigroup of actions, of
\begin{equation}
 \Delta=\Psi_{\nu_0}(\Delta),
 \qquad
 \Psi_{\nu_0}(w):=-\frac{\log\bigl(1+\mathcal E(\nu_0+w)\bigr)}
 {\mathcal Q_{\nu_0}(w)},
 \qquad
 \mathcal Q_{x}(w):=\frac{f(x+w)-f(x)}{w},
 \label{p1:eq:transport-fp}
\end{equation}
with $f=\log\mathcal A$ and $\mathcal Q_x(0)=f'(x)$; its coefficients are given by
\cref{p0:lem:lagrange-fp}, so that in particular
\begin{equation}
 \Delta=-\frac{\mathcal E(\nu_0)}{f'(\nu_0)}
 +\bigO(\mathcal E^{2})
 \sim-\frac{\mathcal E(\nu_0)}{\lambda},
 \label{p1:eq:transport-linear}
\end{equation}
and a single relative forward sector $\mathcal E_\Omega\sim e_{\Omega,0}\e^{-\Omega
x}x^{\mu_\Omega}$ maps to
\begin{equation}
 \Delta_\Omega(y)\sim
 -\frac{e_{\Omega,0}}{\lambda}\,C^{\Omega/\lambda}\,
 y^{-\Omega/\lambda}\,\nu_0^{\,\mu_\Omega-p\Omega/\lambda}.
 \label{p1:eq:transport-y}
\end{equation}
\end{proposition}

\begin{proof}
$f(\nu_0+\Delta)-f(\nu_0)+\log(1+\mathcal E(\nu_0+\Delta))=0$ is the exact equation, and
$f(\nu_0+w)-f(\nu_0)=w\,\mathcal Q_{\nu_0}(w)$ by definition; dividing gives
\cref{p1:eq:transport-fp}.  Since $f'(\nu_0)$ has leading term $\lambda\ne0$ it is invertible
in the coefficient field, and every occurrence of $\Delta$ on the right is either nonlinear or
multiplied by a positive transmonomial, so the support is well ordered and each requested
transmonomial receives finitely many contributions; \cref{p0:lem:lagrange-fp} then gives the
coefficients. \Cref{p1:eq:transport-linear} is the first term.  For \cref{p1:eq:transport-y},
use $y\sim C\e^{\lambda\nu_0}\nu_0^{-p}$ to write
$\e^{-\Omega\nu_0}=(y/C)^{-\Omega/\lambda}\nu_0^{-p\Omega/\lambda}$.
\end{proof}

For the first inherited sector, $\Omega_-=\tfrac\lambda2+\ii\pi$ and $e_{-,0}=B^-_0/C$.
Pairing the two lateral determinations $\pm\ii\pi$ replaces $\e^{-\ii\pi x}$ by $\cos(\pi x)$,
and $\e^{-\lambda\nu_0/2}=C^{1/2}y^{-1/2}\nu_0^{-3/4}$, so
\begin{equation}
 \boxed{\;
 \Delta_-(y)\sim
 0.15193347365953098826\ldots\;
 y^{-1/2}\,\nu_0(y)^{-3/4}\cos\bigl(\pi\nu_0(y)\bigr),\;}
 \label{p1:eq:minus-inverse}
\end{equation}
the amplitude being $\sqrt C\cdot\bigl(-B^-_0/(\lambda C)\bigr)$ with $-B^-_0/(\lambda
C)=0.22906811038403197857\ldots$.  The phase is
\begin{equation}
 \pi\nu_0(y)=\frac{\pi}{\lambda}\log y
 +\frac{3\pi}{2\lambda}\log\log y+\bigO(1),
 \qquad
 \frac{\pi}{\lambda}=2.89879642973397873771\ldots,
 \quad
 \frac{3\pi}{2\lambda}=4.34819464460096810657\ldots,
 \label{p1:eq:minus-phase}
\end{equation}
so the first nonperturbative correction to the inverse has envelope $y^{-1/2}(\log y)^{-3/4}$
and a logarithmic oscillation with a slower $\log\log y$ drift.

\begin{warning}[Two different oscillations]
\label{p1:rem:two-oscillations}
\Cref{p1:eq:minus-inverse} tends to zero and comes from a complex singular action.  The
sawtooth of \cref{p1:eq:sawtooth} is $\bigO(1)$, never tends to zero, and comes from integer
rounding.  They have different origins, different scales, and must not be conflated.
\end{warning}

\section{Other inverse notions: the compositional inverse of $T$}
\label{p1:sec:compositional}

Since $T(z)=z+\bigO(z^{2})$ with integer coefficients and unit linear coefficient, $T$ has a
compositional inverse $Z:=T^{\langle-1\rangle}\in z\Int[[z]]$, unrelated to the growth
inversion of \cref{p1:sec:inverse}.

\begin{proposition}[Coefficients of $Z$]
\label{p1:prop:compositional}
For $n\ge1$,
\begin{equation}
 \Coef{w^{n}}Z(w)=\frac1n\Coef{z^{n-1}}\left(\frac{z}{T(z)}\right)^{n}
 =\frac1n\Coef{z^{n-1}}
 \exp\!\left(-n\sum_{k\ge1}\frac{T(z^{k})}{k}\right),
 \label{p1:eq:Z-lagrange}
\end{equation}
a finite ordinary-Bell expression in $a_1,\ldots,a_{n-1}$.  The first sixteen values,
recomputed by exact rational arithmetic, are
\begin{equation}
 1,\,-1,\,0,\,1,\,-1,\,1,\,-4,\,11,\,-18,\,18,\,0,\,-60,\,189,\,-360,
 \,453,\,-373 .
 \label{p1:eq:Z-values}
\end{equation}
\end{proposition}

\begin{proof}
The first equality is \cref{p0:thm:LB} with $\phi(z)=T(z)/z$, a unit series. The second is
\cref{p1:eq:polya} rearranged, and \cref{p0:def:obell} converts the exponential into a finite
Bell sum. Integrality is automatic for the compositional inverse of a series in
$z+z^{2}\Int[[z]]$.
\end{proof}

\begin{remark}
\textsf{S1} displays $Z(w)=w-w^{2}+w^{4}-w^{5}+w^{6}-4w^{7}+11w^{8}-\cdots$, which agrees with
\cref{p1:eq:Z-values}, and identifies the coefficients with OEIS A050395.  The expansion is
confirmed here; the OEIS identification is reported from \textsf{S1} and has not been
re-checked in this volume.
\end{remark}

Near the dominant value $w=1$ the branch of $Z$ through $Z(1)=\rho$ is \emph{analytic} in
$1-w$, because it undoes the square root: from $T=\zeta\e^{T}$ one has $\zeta=T\e^{-T}$, hence
\begin{equation}
 Z(w)=\zeta^{\langle-1\rangle}\bigl(w\e^{-w}\bigr)
 \label{p1:eq:Z-near-one}
\end{equation}
on that branch, $\zeta$ being injective on $(0,\sqrt\rho)$ by \cref{p1:prop:critical}.  Since
$w\e^{-w}-\e^{-1}$ has a double zero at $w=1$, the local expansion of $Z$ at $w=1$ has no
linear term.  Its coefficients follow from one more Lagrange extraction: writing
$\zeta(\rho+x)-\e^{-1}=x\,\psi(x)$ with $\psi(0)=\zeta'(\rho)\ne0$,
\begin{equation}
 \Coef{\mathsf v^{m}}\bigl(\zeta^{\langle-1\rangle}(\e^{-1}+\mathsf v)
 -\rho\bigr)=\frac1m\Coef{x^{m-1}}\psi(x)^{-m},
 \label{p1:eq:zeta-local-inverse}
\end{equation}
followed by composition with $(1-\varsigma)\e^{-(1-\varsigma)}-\e^{-1}
=-\e^{-1}\sum_{m\ge2}\frac{m-1}{m!}\varsigma^{m}
=-\e^{-1}\frac{\varsigma^{2}}{2}\Phi(\varsigma)$, the same kernel \cref{p1:eq:Phi} that
governed the forward branch.

\section{Algorithms, validation and audit}
\label{p1:sec:algorithms}

\subsection{The pipeline}

\begin{algorithm}[All dominant and inverse coefficients through order $M$]
\label{p1:alg:pipeline}
\leavevmode
\begin{enumerate}
\item Generate $a_1,\ldots,a_N$ and $b_1,\ldots,b_N$ exactly by
      \cref{p1:eq:recurrence} with a divisor sieve, checking the divisibility
      assertion at each step; form $\varrho_r$ by \cref{p1:eq:varrho}.
\item Solve $\log x+1+\sum_{r=2}^{N}\varrho_rx^{r}=0$ for $\rho$ by
      safeguarded Newton or interval bisection.  Choose $N$ so that
      $\rho^{N/2}$ is below the target precision; the omitted tail is
      geometrically small because $R$ is analytic on $|z|<\sqrt\rho$
      (\cref{p1:prop:R-analytic}).
\item Compute $\widehat R_0,\ldots,\widehat R_{2M+1}$ by
      \cref{p1:eq:Rhat}; exponentiate by \cref{p1:eq:E-bell} and form
      $\theta_1,\ldots,\theta_{2M+2}$ by \cref{p1:eq:theta-bell}, then
      $u_j=\theta_{j+1}/\theta_1$.
\item Compute $\omega_1,\ldots,\omega_{2M+1}$ exactly by
      \cref{p1:eq:omega-master} (rationals), set $c_r=-2^{r/2}\omega_r$, and
      form $t_1,t_3,\ldots,t_{2M+1}$ by \cref{p1:eq:t-master}.
\item Compute $g_r(j+\tfrac12)$ exactly for $j+r\le M$ by
      \cref{p1:eq:g-recurrence}, then $\tau_0,\ldots,\tau_M$ by
      \cref{p1:eq:tau-master}, then $C=\tau_0/\sqrt\pi$ and
      $f_m=\tau_m/\tau_0$.
\item Reciprocal: $\bar f_m$ by \cref{p1:eq:recip-coeff}.
\item Inverse: $\varphi_m$ by \cref{p1:eq:phi}, then $d_m$ by
      \cref{p1:eq:d-triangular}, or $P_m(L)$ by the flattening recursion of
      \cref{p0:thm:flattening}.
\item Sectors: apply \cref{p1:thm:local-completeness} outward in modulus,
      transfer each germ by \cref{p1:eq:general-sector}, and invert by
      \cref{p1:prop:transport}.
\end{enumerate}
\end{algorithm}

At every stage the coefficient of order $m$ depends only on data of order at most $m$; no
nonlinear global solve occurs after $\rho$ is known.  The useful power-series primitives are
the three dense recurrences
\begin{equation}
 \Coef{v^{n}}F^{\gamma}=\frac1n\sum_{k=1}^{n}
 \bigl((\gamma+1)k-n\bigr)F_k\Coef{v^{n-k}}F^{\gamma},
 \quad
 E_n=\frac1n\sum_{k=1}^{n}kA_kE_{n-k},
 \quad
 H_n=F_n-\frac1n\sum_{k=1}^{n-1}kH_kF_{n-k},
 \label{p1:eq:dense-recurrences}
\end{equation}
for $F^{\gamma}$ with $F_0=1$, for $E=\e^{A}$ with $A_0=0$, and for $H=\log F$ with $F_0=1$.
They are equivalent to the Bell forms used in the statements and are what an implementation
should actually run.

\begin{warning}[Precision discipline]
\label{p1:rem:precision}
The coefficients grow rapidly and \cref{p1:eq:tau-master} combines alternating contributions
from several Puiseux coefficients, so obtaining $d$ reliable digits of $\tau_M$ can require
substantially more than $d$ working digits.  The tables in this chapter were produced with
$90$--$110$ decimal working digits, with universal objects ($\omega_m$, $g_r(\beta)$, $a_n$)
held in exact rational or integer arithmetic.  The repair recorded at \cref{p1:sub:minus} is
exactly what happens when this discipline lapses.
\end{warning}

\subsection{Numerical validation}
\label{p1:sub:numerics}

Let $A_M(n)$ denote \cref{p1:eq:forward} truncated after $\tau_M n^{-M}$.
\Cref{p1:tab:forward-errors} reports the signed relative error $\bigl(A_M(n)-a_n\bigr)/a_n$
against exact integers from \cref{p1:eq:recurrence}; \cref{p1:tab:inverse-errors} reports the
signed index error $\nu_M(a_n)-n$ for the Lambert-normalised inverse $\nu_M=x_0+\sum_{m\le
M}d_mx_0^{-m}$.  Both tables were regenerated here and agree entry by entry with the
corresponding tables of \textsf{S1} and \textsf{S2} at the displayed precision.

\begin{table}[htbp]
\centering
\caption{Signed relative error $(A_M(n)-a_n)/a_n$ of the forward expansion.}
\label{p1:tab:forward-errors}
\scriptsize
\setlength{\tabcolsep}{4pt}
\begin{tabular}{rcccccc}
\toprule
$n$ & $M=0$ & $M=2$ & $M=4$ & $M=6$ & $M=8$ & $M=10$\\
\midrule
$10$  & $-1.52\cdot10^{-2}$ & $-3.49\cdot10^{-4}$ & $\phantom{-}2.63\cdot10^{-3}$ & $\phantom{-}3.99\cdot10^{-3}$ & $\phantom{-}5.52\cdot10^{-3}$ & $\phantom{-}8.77\cdot10^{-3}$\\
$20$  & $-6.58\cdot10^{-3}$ & $-3.42\cdot10^{-4}$ & $-3.18\cdot10^{-5}$ & $\phantom{-}1.04\cdot10^{-6}$ & $\phantom{-}9.73\cdot10^{-6}$ & $\phantom{-}1.41\cdot10^{-5}$\\
$50$  & $-2.22\cdot10^{-3}$ & $-1.77\cdot10^{-5}$ & $-2.88\cdot10^{-7}$ & $-1.20\cdot10^{-8}$ & $-9.84\cdot10^{-10}$ & $-1.36\cdot10^{-10}$\\
$100$ & $-1.06\cdot10^{-3}$ & $-2.08\cdot10^{-6}$ & $-8.06\cdot10^{-9}$ & $-7.88\cdot10^{-11}$ & $-1.50\cdot10^{-12}$ & $-4.73\cdot10^{-14}$\\
$200$ & $-5.15\cdot10^{-4}$ & $-2.53\cdot10^{-7}$ & $-2.40\cdot10^{-10}$ & $-5.73\cdot10^{-13}$ & $-2.66\cdot10^{-15}$ & $-2.04\cdot10^{-17}$\\
$500$ & $-2.03\cdot10^{-4}$ & $-1.59\cdot10^{-8}$ & $-2.39\cdot10^{-12}$ & $-9.02\cdot10^{-16}$ & $-6.62\cdot10^{-19}$ & $-8.01\cdot10^{-22}$\\
\bottomrule
\end{tabular}
\end{table}

\begin{table}[htbp]
\centering
\caption{Signed index error $\nu_M(a_n)-n$ of the Lambert-normalised
inverse; $M=0$ is the bare carrier $x_0$.}
\label{p1:tab:inverse-errors}
\scriptsize
\setlength{\tabcolsep}{4pt}
\begin{tabular}{rcccccc}
\toprule
$n$ & $M=0$ & $M=1$ & $M=2$ & $M=4$ & $M=6$ & $M=8$\\
\midrule
$10$  & $\phantom{-}1.64\cdot10^{-2}$ & $\phantom{-}7.15\cdot10^{-3}$ & $\phantom{-}1.29\cdot10^{-3}$ & $-2.61\cdot10^{-3}$ & $-4.16\cdot10^{-3}$ & $-5.76\cdot10^{-3}$\\
$20$  & $\phantom{-}6.55\cdot10^{-3}$ & $\phantom{-}1.91\cdot10^{-3}$ & $\phantom{-}4.44\cdot10^{-4}$ & $\phantom{-}3.76\cdot10^{-5}$ & $-2.02\cdot10^{-7}$ & $-9.39\cdot10^{-6}$\\
$50$  & $\phantom{-}2.11\cdot10^{-3}$ & $\phantom{-}2.59\cdot10^{-4}$ & $\phantom{-}2.32\cdot10^{-5}$ & $\phantom{-}3.32\cdot10^{-7}$ & $\phantom{-}1.26\cdot10^{-8}$ & $\phantom{-}9.98\cdot10^{-10}$\\
$100$ & $\phantom{-}9.88\cdot10^{-4}$ & $\phantom{-}6.16\cdot10^{-5}$ & $\phantom{-}2.74\cdot10^{-6}$ & $\phantom{-}9.34\cdot10^{-9}$ & $\phantom{-}8.35\cdot10^{-11}$ & $\phantom{-}1.53\cdot10^{-12}$\\
$200$ & $\phantom{-}4.78\cdot10^{-4}$ & $\phantom{-}1.50\cdot10^{-5}$ & $\phantom{-}3.33\cdot10^{-7}$ & $\phantom{-}2.79\cdot10^{-10}$ & $\phantom{-}6.09\cdot10^{-13}$ & $\phantom{-}2.71\cdot10^{-15}$\\
$500$ & $\phantom{-}1.88\cdot10^{-4}$ & $\phantom{-}2.38\cdot10^{-6}$ & $\phantom{-}2.10\cdot10^{-8}$ & $\phantom{-}2.78\cdot10^{-12}$ & $\phantom{-}9.59\cdot10^{-16}$ & $\phantom{-}6.75\cdot10^{-19}$\\
\bottomrule
\end{tabular}
\end{table}

Two features are worth naming.  First, at $n=10$ and $n=20$ the error stops improving and then
worsens: the expansion is asymptotic, and the least-term truncation of
\cref{p0:thm:optimal-truncation} has been passed.  Second, the inverse errors track the
forward relative errors divided by $f'\approx\lambda$, exactly as
\cref{p1:thm:inverse-remainder} predicts; this is why
\cref{p1:tab:forward-errors,p1:tab:inverse-errors} are so nearly proportional.

\begin{remark}[On comparisons with the literature]
\label{p1:rem:genitrini}
\textsf{S1} states that its first five coefficients ``agree with Genitrini's independently
computed values'' and \textsf{S3} that its Puiseux table ``agrees with the independent table
in [Genitrini 2016]''.  Those comparisons are reported from the drafts and have \emph{not}
been re-performed here; what has been done is a full independent recomputation of every
constant from the exact recurrence, described in \cref{p1:sub:audit}.
\end{remark}

\subsection{Audit}
\label{p1:sub:audit}

Every numerical value printed in this chapter was produced by an independent implementation of
\cref{p1:alg:pipeline} written for this volume: exact integer generation of $a_1,\ldots,a_N$
with $N$ between $700$ and $1400$; exact rational $\varrho_r$, $\omega_m$ and $g_r(\beta)$
(the last from an exact Bernoulli-number table); and $60$--$110$ decimal working digits for
everything else.  The following independent identities were checked and hold to all digits
carried.

\begin{enumerate}
\item $(n-1)\mid\sum_{j<n}b_ja_{n-j}$ at every step, for $n\le1400$;
      $a_{10}=719$, $a_{20}=12\,826\,228$, $a_{50}=425\,976\,989\,835\,141\,038\,353$.
\item $\lambda=1+R(\rho)$, i.e.\ \cref{p1:eq:lambda-check}.
\item $\theta_0=0$, the vanishing of the constant term in
      \cref{p1:eq:theta-def} predicted by $\e\zeta(\rho)=1$.
\item $C^{2}=\theta_1/(2\pi)$, i.e.\ \cref{p1:eq:otter} --- the two sides
      travel through disjoint parts of the pipeline.
\item $c_r=-2^{r/2}\omega_r$ for $1\le r\le18$ in exact arithmetic, and
      $t_1=-\sqrt{2\theta_1}$, $t_2=\tfrac23\theta_1$.
\item \Cref{p1:eq:tau-expanded} re-derived symbolically from
      \cref{p1:eq:tau-master}, and the $g_r(\beta)$ of
      \cref{p1:eq:g-first,p1:eq:g4} re-derived from
      \cref{p1:eq:varkappa}.
\item $\mathrm D_m=\lambda^{m+1}d_m$ and
      $\eta^{\textsf{S2}}_m=\lambda^{m}\varphi_m$ for $1\le m\le8$,
      reconciling \textsf{S1}/\textsf{S3} with \textsf{S2}.
\item $P_1(0)=d_1$, $P_2(0)=d_2$ (a case of \cref{p1:eq:P-shape}).
\item $r^-_1=t_1/\sqrt2$ and the closed form \cref{p1:eq:bminus1}, each
      against the germ recursion; and $r^-_0$ by a second, wholly independent
      summation (see the repair in \cref{p1:sub:minus}).
\item The single-sum \cref{p1:eq:Rhat} against the siblings' double sum, for
      $0\le m\le30$.
\end{enumerate}

\begin{table}[htbp]
\centering
\caption{Principal symbols of this chapter.}
\label{p1:tab:notation}
\small
\begin{tabular}{p{0.17\textwidth}p{0.75\textwidth}}
\toprule
Symbol & Meaning\\
\midrule
$a_n$, $T(z)$ & A000081 and its ordinary generating function.\\
$b_m$, $\varrho_r$ & Divisor transform $\sum_{d\mid m}da_d$; coefficients of $R$, $\varrho_r=(b_r-ra_r)/r$.\\
$R$, $\zeta$ & Disturbance $\sum_{k\ge2}T(z^{k})/k$ and Cayley argument $z\e^{R}$; $T=-\Wlam_0(-\zeta)$.\\
$\rho,\alpha,\lambda,p$ & Otter singularity; $\alpha=\rho^{-1}$, $\lambda=\log\alpha$, $p=\tfrac32$.\\
$\Phi$, $\omega_m$, $c_r$ & Universal Cayley kernel and its two coefficient normalisations, $c_r=-2^{r/2}\omega_r$.\\
$s$, $\theta$, $\theta_m$, $u_j$ & Local coordinate $1-z/\rho$; disturbance $1-\e\zeta$; its coefficients; $u_j=\theta_{j+1}/\theta_1$.\\
$\widehat R_m$, $\ell_j$ & $\Coef{s^m}R(\rho(1-s))$; scaled logarithmic jets $\rho^j(\log\zeta)^{(j)}(\rho)$.\\
$t_n$ & Puiseux coefficients, $T=1+\sum t_ns^{n/2}$.\\
$\varkappa_m(\beta)$, $g_r(\beta)$ & Bernoulli increments and gamma-ratio transfer coefficients.\\
$\tau_m$, $C$, $f_m$, $\bar f_m$ & Forward coefficients; $C=\tau_0/\sqrt\pi$; $f_m=\tau_m/\tau_0$; reciprocal coefficients.\\
$F$, $H$, $\varphi_m$ & $F=1+\sum f_mv^m$; $H=\log F=\sum\varphi_mv^m$; this $H$ is the Part~0 datum $Q$.\\
$x_0$, $d_m$, $\nu$ & Lambert carrier; its corrections; the smooth asymptotic inverse.\\
$\mathcal X$, $L$, $\mathsf q$, $\eta_k$, $P_m$ & Flattening variables and coefficient polynomials.\\
$N(y)$ & Integer staircase $\min\{n\ge1:a_n\ge y\}$.\\
$\sigma_-$, $t^-_m$, $B^-_k$, $\Omega_\sigma$ & The inherited singularity $-\sqrt\rho$, its local and sector coefficients, and the actions.\\
\bottomrule
\end{tabular}
\end{table}

\section{Summary}
\label{p1:sec:summary}

The P\'olya equation becomes transparent after the factorisation $T=\mathcal C\circ\zeta$ with
$\mathcal C(w)=-\Wlam_0(-w)$: the universal layer is one Lagrange coefficient
$\Coef{v^{r-1}}\Phi(v)^{-r/2}$, the problem-specific layer is ordinary Bell composition of the
analytic disturbance $\zeta$, and coefficient transfer is Bernoulli polynomials in the gamma
ratio.  The chain is
\begin{equation}
 a_n\;\to\;\varrho_r\;\to\;\rho\;\to\;\widehat R_m\;\to\;\theta_m
 \;\to\;t_n\;\to\;\tau_m\;\to\;f_m\;\to\;\varphi_m
 \;\to\;d_m\ \text{or}\ P_m(L),
 \label{p1:eq:chain}
\end{equation}
every arrow a finite Bell- or Bernoulli-polynomial formula.  Inversion adds nothing new:
A000081 is the $\alpha=1$ case of \cref{p0:def:model} with $a=\lambda$, $b=-\tfrac32$, $Q=H$,
so \cref{p0:thm:lambert-core} supplies the carrier exactly, \cref{p0:thm:lambert-centered}
every algebraic correction, \cref{p0:thm:flattening} the logarithmic form,
\cref{p0:thm:remainder-transport} the error, and \cref{p0:thm:staircase} the passage to the
integer inverse. Beyond all orders, the recursion supplies a singularity grammar whose first
explicit instance is the alternating sector of relative size $0.2483\ldots\,(-\sqrt\rho)^{n}$
at $\sigma_-=-\sqrt\rho$, transported by \cref{p1:prop:transport} into an inverse correction
of envelope $y^{-1/2}(\log y)^{-3/4}$ with a $\log y$-periodic phase.  What remains open is
not coefficient algebra but analysis: the classification of accessible sheets, of critical
preimages, and of their Stokes connections --- and, prior to that, a proof of $\mathsf
H(\mathcal R)$ for $\mathcal R$ close to $1$.

%% ==================== fragment p2 ====================
\chapter{The double factorial}\label{p2:sec:top}

\begin{sourcebox}
This chapter merges three independent treatments of \(n!!\):
\emph{Double\_Factorial\_Transseries} (henceforth \textsf{S1}),
\emph{Double\_Factorial\_Transseries-3} (\textsf{S2}), and
\emph{double\_factorial\_transseries-2} (\textsf{S3}).
All three reach the same forward and inverse coefficients; they differ in
organisation, in what is proved, and --- decisively --- in notation, where the
letter \(C\) carries three incompatible meanings.  The chapter adopts
\textsf{S3}'s \emph{exact} centred normal form as its spine, \textsf{S2}'s Borel
kernel and sharp remainder, \textsf{S2}'s graded closed coefficient formula, and
\textsf{S1}/\textsf{S2}'s parity-mean centring of the periodic interpolation.
Every displayed rational number below was regenerated with exact rational
arithmetic (\texttt{sympy}) from the defining recurrences, and every numerical
table was recomputed at \(220\)-digit precision (\texttt{mpmath}); see
\cref{p2:sec:numerics}.
\end{sourcebox}

\section{Orientation: which inverse?}
\label{p2:sec:orientation}

For \(n\in\Nat\) the double factorial is
\begin{equation}
 0!!=1!!=1,\qquad n!!=n\,(n-2)!!\quad(n\ge2),
 \label{p2:eq:dfac-recurrence}
\end{equation}
so that \(n!!\) is the product of the positive integers not exceeding \(n\)
of the same parity as \(n\).  The sequence
\(1,1,2,3,8,15,48,105,384,945,\dots\) is OEIS~A006882.  Its elementary closed
forms already contain the whole difficulty:
\begin{equation}
 (2m)!!=2^m\,m!,
 \qquad
 (2m+1)!!=\frac{2^{m+1}\Gamma\!\left(m+\tfrac32\right)}{\sqrt\pi}.
 \label{p2:eq:parity-closed}
\end{equation}
The even subsequence is a gamma function of the halved argument; the odd
subsequence is the \emph{same} gamma function of the halved argument times a
different constant.  The two constants do not agree, and no amount of
reindexing makes them agree.

\begin{warning}
\label{p2:rem:spine}
\(n!!\) is not the restriction to \(\Nat\) of one canonical slowly varying
function.  Its even and odd subsequences are two gamma-type branches whose
quotient is a nonzero constant \(\sqrt{2/\pi}\ne1\).  Consequently
\emph{there is no canonical real inverse of \(n!!\) until an interpolation
is fixed}, and different admissible interpolations give inverses that differ
at an order \emph{larger} than the entire Bernoulli correction grid
(\cref{p2:thm:no-canonical}).  Every statement in this chapter that begins
``the inverse'' names its interpolation first.
\end{warning}

Following \cref{p0:def:three-inverses} we therefore distinguish, for a target
\(y\to+\infty\):

\begin{enumerate}[label=\textup{(\roman*)}]
\item \emph{the two parity-resolved functional inverses}, obtained by inverting
      the canonical gamma continuation of one parity subsequence
      (\cref{p2:sec:branch-inverse});
\item \emph{the functional inverse of a single all-integer interpolation},
      for instance the periodic one recorded by OEIS
      (\cref{p2:sec:oeis-inverse});
\item \emph{the integer staircase inverses}, the least \(n\) with \(n!!\ge y\)
      and the greatest \(n\) with \(n!!\le y\) (\cref{p2:sec:discrete}).
\end{enumerate}

Objects (i) and (iii) are intrinsic to the sequence; object (ii) is not.  This
is the chapter's spine, and \cref{p2:thm:no-canonical} makes it quantitative.

Two asymptotic scales must also be kept apart.  Taking logarithms converts the
growth of \(n!!\) into a phase \(\tfrac12 u(\log u-1)\) whose inversion is
\emph{exactly} a Lambert expression (\cref{p0:prop:factorial-core}); expanding
that Lambert function produces an \emph{outer} layer in powers of
\(1/\log\log y\) with polynomial dependence on \(\log\log\log y\)
(\cref{p0:thm:flattening}).  The Bernoulli corrections generate a far finer
\emph{inner} grid in odd inverse powers of the Lambert core.  The fixed shift
\(-1\) lies strictly between the two levels.  Writing all of this as one
informal series in \(1/\log y\) destroys the ordering; see
\cref{p2:prop:ordering}.

\section{Exact continuations and the parity split}
\label{p2:sec:continuations}

\subsection{The two parity branches}

\begin{definition}[Parity branches]
\label{p2:def:branches}
For the parity label \(\delta\in\{0,1\}\) put
\begin{equation}
 \kappa_\delta=\left(\frac2\pi\right)^{\delta/2},
 \qquad
 \mathcal D_\delta(x)=\kappa_\delta\,2^{x/2}\,
 \Gamma\!\left(\frac x2+1\right)
 \qquad(x>-2),
 \label{p2:eq:branches}
\end{equation}
and write
\begin{equation}
 A_\delta=\log\!\left(\kappa_\delta\sqrt\pi\right)
 =\begin{cases}\tfrac12\log\pi,&\delta=0,\\[0.2em]
               \tfrac12\log2,&\delta=1,\end{cases}
 \qquad
 \e^{A_0}=\sqrt\pi,\quad \e^{A_1}=\sqrt2 .
 \label{p2:eq:Adelta}
\end{equation}
Here \(\delta=0\) labels the even subsequence and \(\delta=1\) the odd one.
\end{definition}

\begin{proposition}[Exact parity interpolation and monotonicity]
\label{p2:prop:branch-interp}
For every integer \(n\ge0\) with \(n\equiv\delta\pmod2\) one has
\(\mathcal D_\delta(n)=n!!\).  Moreover each \(\mathcal D_\delta\) is positive
and strictly increasing on \([0,\infty)\), so it has a single-valued inverse
\(\mathcal D_\delta^{-1}\) on its range; and
\begin{equation}
 \mathcal D_1(x)=\sqrt{\frac2\pi}\;\mathcal D_0(x),
 \qquad
 \mathcal D_1^{-1}(y)=\mathcal D_0^{-1}\!\left(\sqrt{\frac\pi2}\,y\right).
 \label{p2:eq:branch-scaling}
\end{equation}
\end{proposition}

\begin{proof}
If \(n=2m\) then \(\mathcal D_0(2m)=2^m\Gamma(m+1)=2^mm!=(2m)!!\) by
\cref{p2:eq:parity-closed}.  If \(n=2m+1\) then
\[
 \mathcal D_1(2m+1)=\sqrt{\frac2\pi}\,2^{m+1/2}
 \Gamma\!\left(m+\tfrac32\right)
 =\frac{2^{m+1}}{\sqrt\pi}\Gamma\!\left(m+\tfrac32\right)=(2m+1)!!.
\]
Positivity is clear because \(\Gamma>0\) on \((0,\infty)\).  For monotonicity,
\begin{equation}
 \frac{\mathcal D_\delta'(x)}{\mathcal D_\delta(x)}
 =\frac12\left(\log2+\psi\!\left(\frac x2+1\right)\right),
 \qquad \psi=\Gamma'/\Gamma,
 \label{p2:eq:branch-logder}
\end{equation}
and \(\psi\) is strictly increasing on \((0,\infty)\) because
\(\psi'(z)=\sum_{m\ge0}(z+m)^{-2}>0\).  At \(x=0\) the bracket equals
\(\log2-\gamma=0.11593\ldots>0\), hence it is positive for all \(x\ge0\).
\Cref{p2:eq:branch-scaling} is immediate from \cref{p2:eq:branches}, since
\(\kappa_1/\kappa_0=\sqrt{2/\pi}\).
\end{proof}

\begin{remark}
\label{p2:rem:kappa-cancels}
The constant \(\kappa_\delta\) disappears from
\cref{p2:eq:branch-logder}.  Both branches therefore have the same local
conditioning, the same forward correction series, and --- as
\cref{p2:thm:branch-inverse} shows --- literally the same inverse coefficient
polynomials.  Parity survives only as the additive constant \(A_\delta\)
subtracted from \(\log y\).
\end{remark}

\subsection{The periodic OEIS interpolation, and why it is not canonical}

\begin{definition}[Periodic interpolation]
\label{p2:def:oeis}
Put
\begin{equation}
 \mathscr D(x)=2^{x/2}\,\Gamma\!\left(\frac x2+1\right)
 \left(\frac2\pi\right)^{(1-\cos\pi x)/4}.
 \label{p2:eq:oeis}
\end{equation}
\end{definition}

Because \(\tfrac12\sin^2(\pi x/2)=\tfrac14(1-\cos\pi x)\), this is the
interpolation recorded in the OEIS entry for A006882, written there with the
half-argument sine.

\begin{proposition}[Interpolation and eventual monotonicity]
\label{p2:prop:oeis-interp}
\(\mathscr D(n)=n!!\) for every integer \(n\ge0\), and \(\mathscr D\) is
strictly increasing on \([\,1.04,\infty)\).  Its inverse on that ray,
denoted \(\mathscr D^{-1}\), is what ``the inverse of \(\mathscr D\)'' means
below; it is not a globally single-valued inverse near the origin.
\end{proposition}

\begin{proof}
The periodic factor equals \(1\) at even integers and \(\sqrt{2/\pi}\) at odd
integers, so \(\mathscr D\) agrees with \(\mathcal D_0\) at even and with
\(\mathcal D_1\) at odd integers; \cref{p2:prop:branch-interp} finishes the
first claim.  Next,
\begin{equation}
 \frac{\mathscr D'(x)}{\mathscr D(x)}
 =\frac12\left(\log2+\psi\!\left(\frac x2+1\right)\right)
 +\frac\pi4\log\!\left(\frac2\pi\right)\sin\pi x .
 \label{p2:eq:oeis-logder}
\end{equation}
For \(z>\tfrac12\) one has \(\psi(z)>\log\!\left(z-\tfrac12\right)\), hence
\(\psi(x/2+1)>\log\!\left((x+1)/2\right)\) and the first summand exceeds
\(\tfrac12\log(x+1)\).  The second summand has modulus at most
\(\pi a\), where
\begin{equation}
 a:=\frac14\log\frac\pi2=0.112895676322\ldots>0 .
 \label{p2:eq:a-def}
\end{equation}
Thus \(\mathscr D'/\mathscr D>\tfrac12\log(x+1)-\pi a\), which is positive as
soon as \(x+1>\e^{2\pi a}=2.03265\ldots\), i.e.\ for \(x\ge1.04\).
\end{proof}

\begin{remark}
The bound in the proof is not sharp: numerically the last zero of
\cref{p2:eq:oeis-logder} is at \(x=0.6970947\ldots\), so \(\mathscr D\) is in
fact strictly increasing on \([0.70,\infty)\).  Only the proved threshold is
used below.
\end{remark}

The next theorem is the precise form of \cref{p2:rem:spine}.  None of the three
sources states it; all three assert the phenomenon informally.

\begin{theorem}[The inverse transseries is not determined by the sequence]
\label{p2:thm:no-canonical}
For \(\lambda\in\Real\) set
\(\mathscr D_\lambda(x)=\mathscr D(x)\,\e^{\lambda\sin\pi x}\).
Then:
\begin{enumerate}[label=\textup{(\alph*)}]
\item \(\mathscr D_\lambda(n)=n!!\) for every integer \(n\ge0\), and
      \(\mathscr D_\lambda\) is real-analytic on \(\Real\);
\item \(\mathscr D_\lambda\) is strictly increasing on
      \([\,\e^{2\pi(a+|\lambda|)}-1,\infty)\), so it has a large real inverse;
\item with \(T\) and \(\Lambda\) the Lambert data of
      \cref{p2:eq:lambert-core} built from
      \(R_*=\log y-\tfrac14\log(2\pi)\), the leading displacement of that
      inverse beyond the core is
      \begin{equation}
       \mathscr D_\lambda^{-1}(y)-(T-1)
       =\frac{2\sqrt{a^2+\lambda^2}}{\Lambda}
        \cos\!\left(\pi T-\arctan\frac\lambda a\right)
        +\bigO\!\left(\Lambda^{-2}\right).
       \label{p2:eq:no-canonical}
      \end{equation}
\end{enumerate}
In particular the amplitude \(2\sqrt{a^2+\lambda^2}/\Lambda\) of the leading
oscillatory sector is an unbounded function of the interpolation and is
\emph{larger, by the factor \(T\), than the entire Bernoulli grid}
\(\bigO(T^{-1}\Lambda^{-1})\) of \cref{p2:thm:branch-inverse}.
\end{theorem}

\begin{proof}
(a) \(\sin\pi n=0\) for integer \(n\), and \(\e^{\lambda\sin\pi x}\) is entire.

(b) Differentiating logarithmically adds \(\lambda\pi\cos\pi x\) to
\cref{p2:eq:oeis-logder}; repeating the estimate in the proof of
\cref{p2:prop:oeis-interp} gives
\(\mathscr D_\lambda'/\mathscr D_\lambda>\tfrac12\log(x+1)-\pi(a+|\lambda|)\).

(c) By \cref{p2:lem:oeis-centred} below, with \(u=x+1\),
\[
 \log\mathscr D_\lambda(u-1)
 =\frac u2(\log u-1)+\frac14\log(2\pi)-a\cos\pi u
   -\lambda\sin\pi u+\Omega(u),
\]
because \(\cos\pi x=-\cos\pi u\) and \(\sin\pi x=-\sin\pi u\).  Subtracting the
core identity \(\tfrac12T(\log T-1)=R_*\), putting \(u=T+\Delta\) and
multiplying by \(2\) exactly as in \cref{p2:lem:oeis-displacement} gives
\[
 \Lambda\Delta+T\,\varphi(\Delta/T)
 +2\Omega(T+\Delta)
 -2a\cos(\theta+\pi\Delta)-2\lambda\sin(\theta+\pi\Delta)=0,
 \qquad\theta=\pi T .
\]
All terms except the last two are \(\bigO(T^{-1})\) once
\(\Delta=\bigO(\Lambda^{-1})\), so at leading order
\(\Lambda\Delta_0=2a\cos(\theta+\pi\Delta_0)+2\lambda\sin(\theta+\pi\Delta_0)
=2\sqrt{a^2+\lambda^2}\cos(\theta+\pi\Delta_0-\arctan(\lambda/a))\).
Since \(\Delta_0=\bigO(\Lambda^{-1})\), replacing
\(\theta+\pi\Delta_0\) by \(\theta\) costs \(\bigO(\Lambda^{-2})\), which is
\cref{p2:eq:no-canonical}.
\end{proof}

\begin{remark}
The family \(\e^{\lambda\sin\pi x}\) is only the simplest witness.  Any
real-analytic \(g\) vanishing on \(\Nat\) and with \(g'\) bounded produces the
same conclusion, with \(2a\cos\theta\) replaced by the value of
\(-2g\) transported to the core.  What \emph{is} intrinsic is the pair of
parity-branch transseries and the integer staircase; the periodic sector is a
property of the chosen interpolation and nothing else.
\end{remark}

\subsection{The centred coordinate and the exact normal form}

\begin{definition}[Centred coordinate]
\label{p2:def:centred}
Throughout, \(u:=x+1\) and
\begin{equation}
 \Phi(u):=\frac u2(\log u-1),
 \qquad
 \Omega(u):=\log\Gamma\!\left(\frac{u+1}2\right)
 -\frac u2\log\frac u2+\frac u2-\frac12\log(2\pi).
 \label{p2:eq:Omega-def}
\end{equation}
\end{definition}

\begin{theorem}[Exact centred normal form]
\label{p2:thm:normal-form}
For every \(x>-1\) and every \(\delta\in\{0,1\}\), with \(u=x+1\),
\begin{equation}
 \boxed{\;
 \mathcal D_\delta(x)
 =\e^{A_\delta}\left(\frac u\e\right)^{u/2}\e^{\Omega(u)},
 \qquad\text{equivalently}\qquad
 \log\mathcal D_\delta(x)=\Phi(u)+A_\delta+\Omega(u). \;}
 \label{p2:eq:normal-form}
\end{equation}
This is an identity, not an asymptotic statement, and \(\Omega\) does not
depend on \(\delta\).
\end{theorem}

\begin{proof}
Insert \(x=u-1\) into \cref{p2:eq:branches}:
\[
 \log\mathcal D_\delta(x)
 =\log\kappa_\delta+\frac{u-1}2\log2
  +\log\Gamma\!\left(\frac{u+1}2\right).
\]
Subtract \(A_\delta+\Phi(u)\) and use \(A_\delta=\log\kappa_\delta+\tfrac12\log\pi\):
\begin{align*}
 \log\mathcal D_\delta(x)-A_\delta-\Phi(u)
 &=\frac{u-1}2\log2+\log\Gamma\!\left(\frac{u+1}2\right)
   -\frac12\log\pi-\frac u2\log u+\frac u2\\
 &=\log\Gamma\!\left(\frac{u+1}2\right)
   -\frac u2\left(\log u-\log2\right)+\frac u2
   -\frac12\log2-\frac12\log\pi,
\end{align*}
which is exactly \(\Omega(u)\) as defined in \cref{p2:eq:Omega-def}.
\end{proof}

\begin{remark}[Why \(u=x+1\)]
\label{p2:rem:why-centre}
The shift is not cosmetic.  It turns the gamma argument into
\(\tfrac12(u+1)=\tfrac u2+\tfrac12\), and Bernoulli polynomials evaluated at
\(\tfrac12\) annihilate every odd index; consequently \(\Omega\) has only
\emph{odd} inverse powers of \(u\) (\cref{p2:thm:coefficients}) and the
explicit \(\tfrac12\log x\) of the uncentred Stirling form
(\cref{p2:prop:uncentred}) is absorbed into the dominant phase.  This sparse
support is precisely why the inverse lives on the grid
\(T^{-1},T^{-3},T^{-5},\dots\).
\end{remark}

At an integer \(n\) the two constants \(A_0,A_1\) combine into one parity
transmonomial: writing \(\chi_n=(-1)^n\),
\begin{equation}
 A(n)=\frac14\log(2\pi)+\frac{\chi_n}4\log\frac\pi2,
 \qquad
 \e^{A(n)}=(2\pi)^{1/4}\left(\frac\pi2\right)^{\chi_n/4},
 \label{p2:eq:parity-transmonomial}
\end{equation}
so that for every integer \(n\ge0\)
\begin{equation}
 n!!=\left(\frac{n+1}\e\right)^{(n+1)/2}
 \exp\!\bigl(A(n)+\Omega(n+1)\bigr).
 \label{p2:eq:sequence-exact}
\end{equation}
\Cref{p2:eq:sequence-exact} is exact.  The factor \((-1)^n\) is a genuine
transmonomial of the sequence's leading behaviour, not an error term.

\begin{remark}[Notation across the three sources]
\label{p2:rem:notation-clash}
The letter \(C\) means three different things in the siblings, and two Greek
letters collide:
\begin{center}
\small
\begin{tabular}{lccc}
\toprule
this chapter & \textsf{S1} & \textsf{S2} & \textsf{S3}\\
\midrule
\(\kappa_\delta=(2/\pi)^{\delta/2}\) & \(\kappa_\delta\) & \(\kappa_r\) & \(C_\varepsilon\)\\
\(\e^{A_\delta}\in\{\sqrt\pi,\sqrt2\}\) & \(C_\delta\) & \(A_r\) & ---\\
\(A_\delta\in\{\tfrac12\log\pi,\tfrac12\log2\}\) & \(\log C_\delta\) & \(C_r\) & \(A_\varepsilon\)\\
\(c_k\) (centred log.\ coefficient) & \(c_k\) & \(\beta_k\) & \(c_k\)\\
\(a_k=2c_k\) & \(2c_k\) & \(\alpha_k\) & \(a_k\)\\
\(a=\tfrac14\log(\pi/2)\) & \(\rho/2\) & \(a\) & ---\\
\(P_r(L_2)\) (Lambert polynomial) & \(P_r\) & \(u_{r+1}\) & \(\alpha_{r+1}\)\\
\(\tau_m(L_2)\) (reciprocal poly.) & \(h_m\) & \(\tau_m\) & \(\beta_m\)\\
\bottomrule
\end{tabular}
\end{center}
Part~0 adds one further collision: \cref{p0:def:core} writes the factorial
core as \(\Phi_0(X)=X(\kappa\log X+d)\), and that \(\kappa\) is \emph{not} the
branch prefactor \(\kappa_\delta\) of \textsf{S1} and \textsf{S2}.  This chapter
therefore never abbreviates the core parameters, always writing
\(\Phi_0(u)=u\bigl(\tfrac12\log u-\tfrac12\bigr)\) in full; see
\cref{p0:rem:factorial-core}.
In particular \textsf{S2}'s \(\alpha_k,\beta_k\) are Bernoulli data while
\textsf{S3}'s \(\alpha_j,\beta_n\) are outer Lambert polynomials.  Readers
comparing the sources should translate before comparing formulae; the numbers
themselves agree everywhere (\cref{p2:sec:numerics}).
\end{remark}

\section{Parameter dictionary for the Part~0 apparatus}
\label{p2:sec:dictionary}

The chapter uses the Part~0 toolkit rather than re-deriving it.  Its dominant
core is the \emph{factorial} one of \cref{p0:def:core}, not the
exponential--power model \cref{p0:def:model}; see
\cref{p2:rem:not-exponential-power}.  The specialisation is the following.

\begin{center}
\small
\begin{tabular}{p{0.36\textwidth}p{0.27\textwidth}p{0.27\textwidth}}
\toprule
Part~0 datum & even branch \(\delta=0\) & odd branch \(\delta=1\)\\
\midrule
interpolation & \(2^{x/2}\Gamma(\tfrac x2+1)\)
              & \(\sqrt{2/\pi}\;2^{x/2}\Gamma(\tfrac x2+1)\)\\
affine shift \(x_*\) of \cref{p0:def:core} & \(x_*=-1\), taken first & \(x_*=-1\), taken first\\
centred coordinate \(u=x-x_*\) & \(u=x+1\) & \(u=x+1\)\\
admissible core \(\Phi_0\) of \cref{p0:def:core}
  & \(\Phi_0(u)=u\bigl(\tfrac12\log u-\tfrac12\bigr)\) & same\\
additive constant & \(A_0=\tfrac12\log\pi\) & \(A_1=\tfrac12\log2\)\\
perturbation \(Q(s)=\sum_{k\ge1}a_ks^{2k-1}\) & \(a_k=2c_k\), \cref{p2:eq:ck} & same\\
target after constant removal & \(R_0=\log y-\tfrac12\log\pi\)
                              & \(R_1=\log y-\tfrac12\log2\)\\
core (\cref{p0:prop:factorial-core}) & \(T_\delta=2R_\delta/\Wlam_0(2R_\delta/\e)\) & same formula\\
core slope parameter & \(\Lambda_\delta=\log T_\delta\) & same formula\\
inner grid & \(q=T_\delta^{-1}\), \(z=q^2\) & same\\
\bottomrule
\end{tabular}
\end{center}

Concretely: \cref{p0:def:obell,p0:def:ebell} supply the Bell polynomials used
throughout; \cref{p0:lem:power-log} supplies the coefficients of
\((1+w)^{-j}\) and \(\log(1+w)\); \cref{p0:prop:factorial-core} inverts the
dominant phase exactly by \(\Wlam_0\); \cref{p0:thm:core-reversion}, whose
formal engine is \cref{p0:lem:lagrange-fp}, generates the correction blocks and
is specialised once in \cref{p2:lem:normalized};
\cref{p0:thm:flattening} unfolds \(\Wlam_0\) into the outer layer of
\cref{p2:sec:outer}; \cref{p0:def:three-inverses,p0:thm:staircase} govern
\cref{p2:sec:discrete}; \cref{p0:thm:backward-error,p0:thm:remainder-transport}
convert forward residuals into inverse errors; and
\cref{p0:thm:optimal-truncation} fixes the least-term scale of
\cref{p2:sec:beyond}.

\begin{warning}[The double factorial is not of exponential--power type]
\label{p2:rem:not-exponential-power}
It is tempting --- and wrong --- to file \(\mathcal D_\delta\) under
\cref{p0:def:model}, \(F(x)=C\exp(ax^\alpha)x^b\exp Q(x^{-\delta})\) with
\(\alpha\in(0,1]\).  By \cref{p2:thm:normal-form},
\(\log\mathcal D_\delta(x)=\tfrac u2(\log u-1)+A_\delta+\Omega(u)\) with
\(u=x+1\), so the dominant phase is \(\tfrac12u\log u\).  For every
\(\alpha\le1\),
\[
 \frac{u\log u}{u^\alpha}\longrightarrow\infty
 \qquad(u\to+\infty),
\]
so no admissible choice of \((C,a,\alpha,b,\delta,Q)\) reproduces it: the
double factorial lies outside the exponential--power family, in the
\emph{factorial} class of \cref{p0:def:core}, with the parameters that
\cref{p0:rem:factorial-core} records for this chapter.  This is
the structural difference from
\cref{p1:sec:top,p3:sec:top,p4:sec:top}, whose sequences have
\(\log a_n=an^\alpha+b\log n+\cdots\) with \(\alpha\le1\).  What survives from
the exponential--power theory is the \emph{shape} of the reversion, not the
model: \cref{p0:thm:core-reversion} covers both cores, and the resulting
fixed-point map is \cref{p2:eq:master-fp}.  Note also that the shift
\(x_*=-1\) must be applied \emph{before} any further reduction: it is what
removes the \(\tfrac12\log x\) of \cref{p2:eq:uncentred} and produces the
sparse odd-power perturbation (\cref{p2:rem:why-centre}).
None of \textsf{S1}, \textsf{S2}, \textsf{S3} claims membership of the
exponential--power family; all three invert \(\tfrac12u(\log u-1)\) directly.
\end{warning}

\section{Borel structure of the centred correction}
\label{p2:sec:borel}

\begin{sourcebox}
The exact Borel kernel is due to \textsf{S2} and \textsf{S3}, independently and
in the same form; \textsf{S1} has only the formal Bernoulli series.  The sharp
first-omitted-term bound of \cref{p2:thm:remainder} is \textsf{S2}'s;
\textsf{S3} proves the weaker bound with \(\eta(2N+2)\) replaced by \(1\), and
that weaker form is not used here.
\end{sourcebox}

\begin{theorem}[Borel--Laplace representation]
\label{p2:thm:borel}
For \(\Re u>0\),
\begin{equation}
 \Omega(u)=\int_0^\infty\e^{-us}\,\widehat\Omega(s)\,\dd s,
 \qquad
 \widehat\Omega(s)=\frac1{2s^2}\left(\frac s{\sinh s}-1\right),
 \label{p2:eq:borel}
\end{equation}
the singularity of \(\widehat\Omega\) at \(s=0\) being removable.  Moreover
\begin{equation}
 \widehat\Omega(s)=\sum_{m\ge1}\frac{(-1)^m}{s^2+\pi^2m^2},
 \label{p2:eq:mittag}
\end{equation}
so the Borel singularities of \(\Omega\) are exactly the simple poles
\(s=\pm\ii\pi m\), \(m\ge1\), with residues
\begin{equation}
 \Res_{s=\ii\pi m}\widehat\Omega(s)=\frac{(-1)^{m+1}\ii}{2\pi m},
 \qquad
 \Res_{s=-\ii\pi m}\widehat\Omega(s)=\frac{(-1)^{m}\ii}{2\pi m}.
 \label{p2:eq:residues}
\end{equation}
\end{theorem}

\begin{proof}
Put \(t=u/2\) and
\(\widetilde\Omega(t)=\log\Gamma(t+\tfrac12)-t\log t+t-\tfrac12\log(2\pi)\),
so that \(\Omega(u)=\widetilde\Omega(u/2)\) by \cref{p2:eq:Omega-def}.
Binet's first integral and Frullani's formula give, for \(t>0\),
\[
 \psi\!\left(t+\tfrac12\right)-\log t
 =\int_0^\infty\e^{-tr}
 \left(\frac1r-\frac1{2\sinh(r/2)}\right)\dd r,
\]
the integrand being \(\bigO(r)\) at the origin and exponentially damped at
infinity.  On the other hand
\[
 \frac{\dd}{\dd t}\int_0^\infty\e^{-tr}
 \frac{r/(2\sinh(r/2))-1}{r^2}\,\dd r
 =\int_0^\infty\e^{-tr}
 \left(\frac1r-\frac1{2\sinh(r/2)}\right)\dd r ,
\]
differentiation under the integral being justified by the same bounds.  Since
\(\widetilde\Omega'(t)=\psi(t+\tfrac12)-\log t\) and both
\(\widetilde\Omega(t)\) and the displayed integral tend to \(0\) as
\(t\to+\infty\), the two functions coincide.  Substituting \(r=2s\) and
\(t=u/2\) turns the integral into \cref{p2:eq:borel}; analytic continuation in
\(u\) extends it to \(\Re u>0\).  The expansion \(s/\sinh s=1-s^2/6+\bigO(s^4)\)
removes the origin.

For \cref{p2:eq:mittag}, the elementary Mittag--Leffler expansion
\(\dfrac s{\sinh s}=1+2s^2\sum_{m\ge1}\dfrac{(-1)^m}{s^2+\pi^2m^2}\)
gives the claim after division by \(2s^2\).  The residues follow from
\(1/(s^2+\pi^2m^2)=\bigl[(s-\ii\pi m)(s+\ii\pi m)\bigr]^{-1}\).
\end{proof}

\begin{theorem}[The centred logarithmic coefficients]
\label{p2:thm:coefficients}
As \(u\to\infty\) in any closed subsector of \(|\arg u|<\pi/2\),
\begin{equation}
 \Omega(u)\sim\sum_{k\ge1}\frac{c_k}{u^{2k-1}},
 \label{p2:eq:Omega-series}
\end{equation}
where the following three expressions coincide:
\begin{equation}
 \boxed{\;
 c_k=\frac{2^{2k-1}B_{2k}\!\left(\tfrac12\right)}{2k(2k-1)}
 =\frac{\bigl(1-2^{2k-1}\bigr)B_{2k}}{2k(2k-1)}
 =(-1)^k\frac{(2k-2)!\,\eta(2k)}{\pi^{2k}} \;}
 \label{p2:eq:ck}
\end{equation}
with \(B_m(z)\) the Bernoulli polynomials, \(B_m=B_m(0)\), and
\(\eta(s)=(1-2^{1-s})\zeta(s)\) the Dirichlet eta function.  There are no even
inverse powers.
\end{theorem}

\begin{proof}
Expand each summand of \cref{p2:eq:mittag} by
\begin{equation}
 \frac1{s^2+A^2}
 =\sum_{j=0}^{N-1}\frac{(-1)^js^{2j}}{A^{2j+2}}
 +\frac{(-1)^Ns^{2N}}{A^{2N}(s^2+A^2)},
 \qquad A>0,
 \label{p2:eq:partial-fraction-remainder}
\end{equation}
with \(A=\pi m\), sum over \(m\), and integrate term by term against
\(\e^{-us}\), using \(\int_0^\infty\e^{-us}s^{2j}\dd s=(2j)!\,u^{-2j-1}\).
The coefficient of \(u^{-(2j+1)}\) is
\[
 (2j)!\sum_{m\ge1}\frac{(-1)^m(-1)^j}{(\pi m)^{2j+2}}
 =(-1)^{j+1}\frac{(2j)!\,\eta(2j+2)}{\pi^{2j+2}},
\]
which is the third expression in \cref{p2:eq:ck} with \(k=j+1\).  Euler's
evaluation \(\zeta(2k)=(-1)^{k+1}(2\pi)^{2k}B_{2k}/\bigl(2(2k)!\bigr)\)
together with \(\eta(2k)=(1-2^{1-2k})\zeta(2k)\) and
\(B_{2k}(\tfrac12)=(2^{1-2k}-1)B_{2k}\) converts it into the first two.  Only
odd powers \(u^{-(2k-1)}\) occur because \(\widehat\Omega\) is even.
\end{proof}

\begin{theorem}[Sharp alternating remainder]
\label{p2:thm:remainder}
For real \(u>0\) and \(N\ge0\) put
\(R_N(u)=\Omega(u)-\sum_{k=1}^N c_ku^{-(2k-1)}\) (empty sum for \(N=0\)).
Then \(R_N(u)\) has the sign of the first omitted term and is smaller than it
in modulus:
\begin{equation}
 0<(-1)^{N+1}R_N(u)<\frac{|c_{N+1}|}{u^{2N+1}}
 =\frac{(2N)!\,\eta(2N+2)}{\pi^{2N+2}u^{2N+1}} .
 \label{p2:eq:remainder}
\end{equation}
\end{theorem}

\begin{proof}
Insert \cref{p2:eq:partial-fraction-remainder} with \(A=\pi m\) into
\cref{p2:eq:mittag} and Laplace-transform.  The residual kernel is
\((-1)^{N+1}s^{2N}S_N(s)\) with
\[
 S_N(s)=\sum_{m\ge1}\frac{(-1)^{m-1}}
 {(\pi m)^{2N}\bigl((\pi m)^2+s^2\bigr)} .
\]
Writing \(\bigl((\pi m)^2+s^2\bigr)^{-1}=\int_0^\infty
\e^{-s^2v}\e^{-\pi^2m^2v}\dd v\) gives
\[
 S_N(s)=\int_0^\infty\e^{-s^2v}
 \sum_{m\ge1}\frac{(-1)^{m-1}\e^{-\pi^2m^2v}}{(\pi m)^{2N}}\,\dd v .
\]
For each \(v\ge0\) the inner series is alternating with strictly decreasing
positive terms, hence strictly positive.  Therefore \(S_N(s)>0\) and
\(S_N\) is strictly decreasing in \(s^2\), so
\(0<S_N(s)\le S_N(0)=\eta(2N+2)/\pi^{2N+2}\).  Multiplying by
\(s^{2N}\), integrating against \(\e^{-us}\) and using
\(\int_0^\infty\e^{-us}s^{2N}\dd s=(2N)!u^{-(2N+1)}\) gives
\cref{p2:eq:remainder}.
\end{proof}

\begin{corollary}[Certified forward logarithm]
\label{p2:cor:forward-log}
For \(x\to+\infty\), \(u=x+1\) and either parity,
\begin{equation}
 \log\mathcal D_\delta(x)
 =\Phi(u)+A_\delta+\sum_{k=1}^N\frac{c_k}{u^{2k-1}}
  +\theta_N\,\frac{c_{N+1}}{u^{2N+1}}
 \qquad\text{for some }\theta_N=\theta_N(u)\in(0,1).
 \label{p2:eq:forward-certified}
\end{equation}
Thus truncating after any number of terms overshoots or undershoots by
strictly less than the first omitted term, with the sign of that term.  For the
integer sequence replace \(A_\delta\) by \(A(n)\) of
\cref{p2:eq:parity-transmonomial}.
\end{corollary}

\begin{proof}
Combine \cref{p2:thm:normal-form,p2:thm:remainder}.
\end{proof}

\begin{center}
\small
\begin{longtable}{c r r}
\caption{The centred logarithmic coefficients \(c_k\) of \cref{p2:eq:ck} and
the doubled coefficients \(a_k=2c_k\) used in the inverse problem.  Regenerated
in exact rational arithmetic from \(B_{2k}\).}
\label{p2:tab:ck}\\
\toprule
\(k\) & \(c_k\) & \(a_k=2c_k\)\\
\midrule
\endfirsthead
\toprule
\(k\) & \(c_k\) & \(a_k=2c_k\)\\
\midrule
\endhead
1 & \(-1/12\) & \(-1/6\)\\
2 & \(7/360\) & \(7/180\)\\
3 & \(-31/1260\) & \(-31/630\)\\
4 & \(127/1680\) & \(127/840\)\\
5 & \(-511/1188\) & \(-511/594\)\\
6 & \(1414477/360360\) & \(1414477/180180\)\\
7 & \(-8191/156\) & \(-8191/78\)\\
8 & \(118518239/122400\) & \(118518239/61200\)\\
9 & \(-5749691557/244188\) & \(-5749691557/122094\)\\
\bottomrule
\end{longtable}
\end{center}

\begin{remark}[Late coefficients]
\label{p2:rem:late}
From the third form in \cref{p2:eq:ck} and \(1<\eta(2k)<1+2^{-2k}\cdot\tfrac43\),
\begin{equation}
 c_k=(-1)^k\frac{(2k-2)!}{\pi^{2k}}\bigl(1+\bigO(4^{-k})\bigr),
 \label{p2:eq:ck-late}
\end{equation}
so \cref{p2:eq:Omega-series} is Gevrey-\(1\) in the variable \(u^{-2}\).  The
ratio of consecutive term moduli is \(\sim2k(2k-1)/(\pi u)^2\); the least term
occurs near \(2k=\pi u\) and has size \(\asymp u^{-1/2}\e^{-\pi u}\), matching
the distance \(\pi\) from the origin to the nearest Borel poles
\(\pm\ii\pi\).  This is the instance of \cref{p0:thm:optimal-truncation}
relevant here; see \cref{p2:sec:beyond}.
\end{remark}

\section{The forward transseries}
\label{p2:sec:forward}

\subsection{Multiplicative, reciprocal, and arbitrary-power coefficients}

Exponentiation of \cref{p2:eq:Omega-series} is an instance of the exponential
formula, so \cref{p0:def:ebell} does all the work at once for every power.

\begin{definition}
\label{p2:def:gamma-sparse}
Embed the \(c_k\) sparsely on the odd indices,
\begin{equation}
 \gamma_r=\begin{cases}c_{(r+1)/2},&r\ \text{odd},\\ 0,&r\ \text{even},\end{cases}
 \qquad r\ge1,
 \label{p2:eq:gamma-sparse}
\end{equation}
and for a scalar \(\alpha\) define \(s^{(\alpha)}_m\) by
\begin{equation}
 \exp\!\left(\alpha\sum_{r\ge1}\gamma_r\,u^{-r}\right)
 =\sum_{m\ge0}\frac{s^{(\alpha)}_m}{u^m},
 \qquad s_m:=s^{(1)}_m .
 \label{p2:eq:salpha-def}
\end{equation}
\end{definition}

\begin{proposition}[All multiplicative coefficients at once]
\label{p2:prop:mult-coeff}
For every \(\alpha\) and every \(m\ge0\),
\begin{align}
 s^{(\alpha)}_m
 &=\frac1{m!}\,\EB_m\bigl(\alpha\,1!\,\gamma_1,\alpha\,2!\,\gamma_2,\ldots,
     \alpha\,m!\,\gamma_m\bigr)
 \label{p2:eq:salpha-bell}\\
 &=\sum_{\substack{k_1,k_2,\ldots\ge0\\ \sum_{r\ge1}rk_r=m}}
   \prod_{r\ge1}\frac{(\alpha\gamma_r)^{k_r}}{k_r!}
 =\sum_{\substack{k_1,k_2,\ldots\ge0\\ \sum_{j\ge1}(2j-1)k_j=m}}
   \prod_{j\ge1}\frac{(\alpha c_j)^{k_j}}{k_j!},
 \label{p2:eq:salpha-partition}
\end{align}
and equivalently
\begin{equation}
 s^{(\alpha)}_0=1,
 \qquad
 m\,s^{(\alpha)}_m
 =\alpha\!\!\sum_{1\le j\le(m+1)/2}\!\!(2j-1)\,c_j\,s^{(\alpha)}_{m-2j+1}
 \qquad(m\ge1).
 \label{p2:eq:salpha-recurrence}
\end{equation}
Consequently, for \(x\to+\infty\), \(u=x+1\) and every fixed \(M\ge0\),
\begin{equation}
 \mathcal D_\delta(x)^\alpha
 =\e^{\alpha A_\delta}\left(\frac u\e\right)^{\alpha u/2}
 \left(\sum_{m=0}^{M}\frac{s^{(\alpha)}_m}{u^m}
   +\bigO\!\left(u^{-(M+1)}\right)\right).
 \label{p2:eq:mult-transseries}
\end{equation}
The cases \(\alpha=1\) and \(\alpha=-1\) give \(\mathcal D_\delta\) itself and
its reciprocal.
\end{proposition}

\begin{proof}
Put \(t=u^{-1}\) and \(S_\alpha(t)=\exp(\alpha\sum_r\gamma_rt^r)\).
\Cref{p2:eq:salpha-bell} is the complete-Bell generating identity of
\cref{p0:def:ebell}; multinomial expansion of the exponential gives
\cref{p2:eq:salpha-partition}, and the second form there uses
\(\gamma_{2j-1}=c_j\), \(\gamma_{2j}=0\).  Logarithmic differentiation
\(tS_\alpha'=\alpha\bigl(\sum_r r\gamma_rt^r\bigr)S_\alpha\) and comparison of
coefficients give \cref{p2:eq:salpha-recurrence}.  For
\cref{p2:eq:mult-transseries}, take \(N\) with \(2N-1\ge M\); by
\cref{p2:thm:normal-form,p2:thm:remainder},
\(\alpha\Omega(u)=\alpha\sum_{k\le N}c_ku^{1-2k}+\bigO(u^{-(2N+1)})\), and
\(\exp\) of a finite sum of \(\bigO(u^{-1})\) terms plus a
\(\bigO(u^{-(2N+1)})\) error equals \(\sum_{m\le M}s^{(\alpha)}_mu^{-m}
+\bigO(u^{-(M+1)})\), the exponential being Lipschitz on bounded sets.
\end{proof}

\begin{corollary}[Reciprocal sign symmetry]
\label{p2:cor:reciprocal-sign}
\(s^{(-\alpha)}_m=(-1)^m s^{(\alpha)}_m\) for all \(m\ge0\); in particular the
reciprocal coefficients are \(s^{(-1)}_m=(-1)^ms_m\).
\end{corollary}

\begin{proof}
Only odd \(r\) carry a nonzero \(\gamma_r\), so
\(\sum_r\gamma_r(-t)^r=-\sum_r\gamma_rt^r\) and hence
\(S_\alpha(-t)=S_{-\alpha}(t)\).  Reading off \(\Coef{t^m}\) gives the claim.
\end{proof}

\begin{center}
\small
\begin{longtable}{c r r}
\caption{Centred multiplicative coefficients \(s_m=s^{(1)}_m\) of
\cref{p2:eq:mult-transseries} and the uncentred coefficients \(g_m\) of
\cref{p2:prop:uncentred}.  Regenerated from
\cref{p2:eq:salpha-recurrence} in exact rational arithmetic.}
\label{p2:tab:sm}\\
\toprule
\(m\) & \(s_m\) (centred) & \(g_m\) (uncentred)\\
\midrule
\endfirsthead
\toprule
\(m\) & \(s_m\) (centred) & \(g_m\) (uncentred)\\
\midrule
\endhead
0 & \(1\) & \(1\)\\
1 & \(-1/12\) & \(1/6\)\\
2 & \(1/288\) & \(1/72\)\\
3 & \(1003/51840\) & \(-139/6480\)\\
4 & \(-4027/2488320\) & \(-571/155520\)\\
5 & \(-5128423/209018880\) & \(163879/6531840\)\\
6 & \(168359651/75246796800\) & \(5246819/1175731200\)\\
7 & \(68168266699/902961561600\) & \(-534703531/7054387200\)\\
8 & \(-587283555451/86684309913600\) & \(-4483131259/338610585600\)\\
\bottomrule
\end{longtable}
\end{center}

\subsection{The uncentred normal form and the centre-shift identity}

\begin{proposition}[Uncentred expansion]
\label{p2:prop:uncentred}
As \(x\to+\infty\),
\begin{equation}
 \log\mathcal D_\delta(x)
 \sim\frac x2(\log x-1)+\frac12\log x+A_\delta
 +\sum_{k\ge1}\frac{b_k}{x^{2k-1}},
 \qquad
 b_k=\frac{2^{2k-1}B_{2k}}{2k(2k-1)}
   =(-1)^{k-1}\frac{(2k-2)!\,\zeta(2k)}{\pi^{2k}} ,
 \label{p2:eq:uncentred}
\end{equation}
with \(b_1=\tfrac16\), \(b_2=-\tfrac1{45}\), \(b_3=\tfrac8{315}\),
\(b_4=-\tfrac8{105}\), \(b_5=\tfrac{128}{297}\).  Exponentiating,
\(\mathcal D_\delta(x)\sim\e^{A_\delta}\sqrt x\,(x/\e)^{x/2}\sum_m g_mx^{-m}\),
where \(g_m\) is obtained from \cref{p2:eq:salpha-recurrence} on replacing
\(c_j\) by \(b_j\); this is the familiar form, with amplitudes \(\sqrt\pi\)
(even) and \(\sqrt2\) (odd).  The exact Borel kernel of the uncentred
correction is
\(\widehat{\,\mathcal U}(s)=\tfrac{\coth s}{2s}-\tfrac1{2s^2}
=\sum_{m\ge1}\bigl(s^2+\pi^2m^2\bigr)^{-1}\), with the \emph{same} pole
lattice \(\pm\ii\pi m\) but all residues of one sign pattern.
\end{proposition}

\begin{proof}
Apply the standard Stirling expansion to \(\log\Gamma(z+1)\) with
\(z=x/2\) and add \((x/2)\log2+\log\kappa_\delta\); the elementary terms give
the first three summands of \cref{p2:eq:uncentred}, and the corrections scale
as \(z^{-(2k-1)}=2^{2k-1}x^{-(2k-1)}\).  Euler's evaluation of \(\zeta(2k)\)
gives the second form of \(b_k\).  For the kernel, the uncentred correction is
\(\mathcal U(x)=\Omega_\Gamma(x/2)\) with \(\Omega_\Gamma\) Binet's function;
substituting \(t=2s\) in
\(\Omega_\Gamma(z)=\int_0^\infty\e^{-zt}
\bigl(\tfrac{\coth(t/2)}{2t}-\tfrac1{t^2}\bigr)\dd t\) gives the stated kernel,
and \(\tfrac{\coth s}s=\tfrac1{s^2}+2\sum_{m\ge1}(s^2+\pi^2m^2)^{-1}\) gives
the partial fractions.
\end{proof}

\begin{warning}
\label{p2:rem:two-normal-forms}
The two statements
\(n!!\sim \e^{A(n)}\sqrt n\,(n/\e)^{n/2}\) and
\(n!!\sim \e^{A(n)}\bigl((n+1)/\e\bigr)^{(n+1)/2}\)
are both correct and are not in conflict: they use different normal forms.  The
first leaves a factor \(\sqrt n\) and an \emph{uncentred} correction series;
the second absorbs that factor into the phase.  Only the centred form leads to
the sparse inverse grid, and it is the one used from
\cref{p2:sec:branch-inverse} onwards.
\end{warning}

\begin{proposition}[Centre-shift identity]
\label{p2:prop:centre-shift}
For every \(r\ge1\),
\begin{equation}
 \frac{(-1)^{r+1}}{2r(r+1)}
 +(-1)^{r+1}\!\!\sum_{k=1}^{\lfloor(r+1)/2\rfloor}\!\!
 \binom{r-1}{2k-2}c_k
 =\begin{cases}b_{(r+1)/2},& r\ \text{odd},\\[0.2em] 0,& r\ \text{even}.\end{cases}
 \label{p2:eq:centre-shift}
\end{equation}
\end{proposition}

\begin{proof}
Substitute \(u=x+1=x(1+w)\), \(w=x^{-1}\), into
\(\Phi(u)+\sum_kc_ku^{1-2k}\) and compare with \cref{p2:eq:uncentred}.
First, with \(\log u=\log x+\log(1+w)\),
\[
 \Phi(u)=\frac x2(\log x-1)+\frac12(\log x-1)
 +\frac{x+1}2\log(1+w).
\]
Expanding \(\log(1+w)=\sum_{j\ge1}(-1)^{j-1}w^j/j\) and splitting the prefactor
\(\tfrac{x+1}2=\tfrac x2+\tfrac12\) gives
\[
 \frac{x+1}2\log(1+w)
 =\frac12+\frac12\sum_{r\ge1}(-1)^rx^{-r}
   \left(\frac1{r+1}-\frac1r\right)
 =\frac12+\sum_{r\ge1}\frac{(-1)^{r+1}}{2r(r+1)}\,x^{-r},
\]
so that \(\Phi(u)=\tfrac x2(\log x-1)+\tfrac12\log x
+\sum_{r\ge1}(-1)^{r+1}\bigl(2r(r+1)\bigr)^{-1}x^{-r}\); the two halves
\(\pm\tfrac12\) cancel.  Second,
\[
 \sum_{k\ge1}\frac{c_k}{u^{2k-1}}
 =\sum_{k\ge1}c_k\,x^{-(2k-1)}(1+w)^{-(2k-1)}
 =\sum_{r\ge1}x^{-r}\sum_{k\le(r+1)/2}c_k
   \binom{-(2k-1)}{r-2k+1},
\]
and \(\binom{-(2k-1)}{r-2k+1}=(-1)^{r-2k+1}\binom{r-1}{r-2k+1}
=(-1)^{r+1}\binom{r-1}{2k-2}\).  Adding the two coefficient lists and matching
against \cref{p2:eq:uncentred} --- where the coefficient of \(x^{-r}\) is
\(b_{(r+1)/2}\) for odd \(r\) and \(0\) for even \(r\) --- gives
\cref{p2:eq:centre-shift}.
\end{proof}

\begin{remark}
\Cref{p2:eq:centre-shift} was verified symbolically for \(1\le r\le9\): the
even cases return \(0\) and the odd cases return
\(\tfrac16,-\tfrac1{45},\tfrac8{315},-\tfrac8{105},\tfrac{128}{297}\).  It is
the cheapest available certificate that the centring is consistent, because it
fails immediately under any sign or index slip in \cref{p2:eq:ck}.  \textsf{S3}
states the identity without proof; the proof above is supplied here.
\end{remark}

\section{The branchwise functional inverse}
\label{p2:sec:branch-inverse}

\begin{sourcebox}
All three siblings reach the same objects in this section.  The presentation
follows \textsf{S3} for the exact normal form and \textsf{S2} for the graded
closed formula and the sign law; \textsf{S1}'s ``scaling from the half-shifted
inverse Gamma problem'' is subsumed by \cref{p2:thm:multi-scaling}.  The proofs
of \cref{p2:cor:boundary,p2:thm:sign-law} are new: the corresponding source
arguments are sketches (see the chapter report).
\end{sourcebox}

\subsection{The Lambert core}

\begin{lemma}[Exact inversion of the dominant phase]
\label{p2:lem:core}
Let \(R>0\).  The equation \(\Phi(T)=\tfrac12T(\log T-1)=R\) has a unique
solution \(T>\e\), namely
\begin{equation}
 w=\Wlam_0\!\left(\frac{2R}\e\right),
 \qquad
 T=\frac{2R}w=\e^{\,w+1},
 \qquad
 \Lambda:=w+1=\log T,
 \label{p2:eq:lambert-core}
\end{equation}
and then
\begin{equation}
 \Phi'(T)=\frac\Lambda2,
 \qquad
 \frac{\dd T}{\dd R}=\frac2\Lambda,
 \qquad
 T(\Lambda-1)=2R .
 \label{p2:eq:core-identities}
\end{equation}
\end{lemma}

\begin{proof}
This is \cref{p0:prop:factorial-core} for the parameters of
\cref{p0:rem:factorial-core}; for completeness, setting \(w=\log T-1\) turns
\(\Phi(T)=R\) into \(w\e^{\,w+1}=2R\), i.e.\ \(w\e^{w}=2R/\e>0\), whose unique
real solution is \(w=\Wlam_0(2R/\e)>0\), the principal branch being the only
real branch on \((0,\infty)\).  \(\Phi\) is strictly increasing on
\((1,\infty)\) with range \((-\tfrac12,\infty)\), which gives uniqueness of
\(T\).  The identities in \cref{p2:eq:core-identities} follow from
\(\Phi'(T)=\tfrac12\log T\), implicit differentiation of \(\Phi(T)=R\), and
\(T=2R/w=2R/(\Lambda-1)\).
\end{proof}

For the parity branch \(\delta\) one takes
\begin{equation}
 R_\delta=\log y-A_\delta,
 \qquad
 w_\delta,\;T_\delta,\;\Lambda_\delta\ \text{as in \cref{p2:eq:lambert-core}} .
 \label{p2:eq:Rdelta}
\end{equation}
The subscript \(\delta\) is suppressed inside coefficient calculations, where
it never appears.

\subsection{The normalised displacement equation}

\begin{lemma}[The factorial-core reversion kernel]
\label{p2:lem:normalized}
Write \(u=T(1+E)\), \(q=T^{-1}\), \(z=q^2\), and
\begin{equation}
 \varphi(v)=(1+v)\log(1+v)-v=\sum_{m\ge2}\frac{(-1)^m}{m(m-1)}v^m,
 \qquad
 a_k=2c_k .
 \label{p2:eq:varphi}
\end{equation}
Then the equation \(\Phi(u)+\Omega(u)=R\), with \(\Omega\) replaced by its
asymptotic series, is equivalent to the formal normal form
\begin{equation}
 \boxed{\;
 \Lambda E+\varphi(E)+\sum_{k\ge1}a_k\,z^k(1+E)^{-(2k-1)}=0 . \;}
 \label{p2:eq:master}
\end{equation}
Equivalently, in absolute displacement \(U=TE\) with \(t=q\),
\begin{equation}
 U=\Psi(t,U),
 \qquad
 \Psi(t,v)=-\frac1\Lambda\left[
 \frac{\varphi(tv)}t
 +\sum_{k\ge1}a_k\,t^{2k-1}(1+tv)^{-(2k-1)}\right]
 \in t\,\Rat(\Lambda)[[t,v]] ,
 \label{p2:eq:master-fp}
\end{equation}
which is the specialisation of \cref{p0:thm:core-reversion} to the factorial
core, so that \cref{p0:lem:lagrange-fp} applies verbatim.  There is a unique
formal solution, and it has the form
\begin{equation}
 E=\sum_{n\ge1}e_n(\Lambda)\,z^n,
 \qquad\text{hence}\qquad
 u-T=\sum_{n\ge1}\frac{e_n(\Lambda)}{T^{2n-1}} :
 \label{p2:eq:E-series}
\end{equation}
only \emph{odd} inverse powers of the Lambert core occur.
\end{lemma}

\begin{proof}
Because \(\Phi(T)=R\),
\[
 \Phi\bigl(T(1+E)\bigr)-\Phi(T)
 =\frac{T(1+E)}2\bigl(\log T+\log(1+E)-1\bigr)-\frac T2(\log T-1)
 =\frac T2\bigl[\Lambda E+\varphi(E)\bigr],
\]
using \(\Lambda=\log T\) and
\((1+E)\log(1+E)=\varphi(E)+E\).  Also
\(\Omega(u)=\sum_{k}c_kT^{-(2k-1)}(1+E)^{-(2k-1)}\) formally.  Multiplying the
whole equation by \(2/T=2q\) and using \(a_k=2c_k\) and
\(q\cdot q^{2k-1}=z^k\) gives \cref{p2:eq:master}.  Dividing by \(T\) and
solving for \(U=TE\) gives \cref{p2:eq:master-fp}; the bracket lies in
\(t\,\Rat[[t,v]]\) because \(\varphi(tv)=t^2v^2/2+\bigO(t^3)\), so
\(\varphi(tv)/t=\bigO(t)\), and every Bernoulli term carries \(t^{2k-1}\) with
\(k\ge1\).  Existence and uniqueness of the fixed point are
\cref{p0:lem:lagrange-fp}.  Finally, if \(E\) contains only even powers of
\(q\) then so does every term of \cref{p2:eq:master}; an induction on the
\(q\)-valuation therefore forces \(E\in z\,\Rat(\Lambda)[[z]]\).
\end{proof}

\subsection{The all-orders inverse and its remainder}

\begin{theorem}[Complete branchwise inverse transseries]
\label{p2:thm:branch-inverse}
Fix \(\delta\in\{0,1\}\) and let \(R_\delta,T_\delta,\Lambda_\delta\) be as in
\cref{p2:eq:Rdelta}.  Write \(d_{2n-1}:=e_n\) for the coefficients of
\cref{p2:eq:E-series}.  Then for every fixed \(N\ge0\), as \(y\to+\infty\),
\begin{equation}
 \boxed{\;
 \mathcal D_\delta^{-1}(y)
 =T_\delta-1
 +\sum_{n=1}^{N}\frac{d_{2n-1}(\Lambda_\delta)}{T_\delta^{2n-1}}
 +\bigO\!\left(\frac1{T_\delta^{2N+1}\Lambda_\delta}\right) , \;}
 \label{p2:eq:branch-inverse}
\end{equation}
with an implied constant that may be chosen uniformly in \(\delta\in\{0,1\}\).
All parity dependence sits in \(R_\delta=\log y-A_\delta\); the coefficient
functions \(d_{2n-1}\) are universal.
\end{theorem}

\begin{proof}
Let \(u_N=T+\sum_{n=1}^Ne_n(\Lambda)T^{1-2n}\) and \(E_N=(u_N-T)/T\).  By
construction, substituting \(E_N\) into \cref{p2:eq:master} annihilates the
coefficients of \(z^1,\dots,z^N\), so the formal residual is
\(\bigO(z^{N+1})\) with coefficients rational in \(\Lambda\) and bounded as
\(\Lambda\to\infty\).  Undoing the normalisation (multiplying by \(T/2\))
therefore bounds the exact forward residual
\[
 \rho_N:=\Phi(u_N)+A_\delta+\Omega(u_N)-\log y
\]
by \(\bigO\!\left(T\cdot z^{N+1}\right)=\bigO\!\left(T^{-(2N+1)}\right)\),
where replacing \(\Omega\) by its \(N\)-term truncation costs at most
\(|c_{N+1}|u_N^{-(2N+1)}\) by \cref{p2:thm:remainder}, of the same order.
This is the hypothesis of \cref{p0:thm:backward-error}.  For the slope, on the
interval between \(u_N\) and the exact root,
\[
 \frac{\dd}{\dd u}\bigl[\Phi(u)+\Omega(u)\bigr]
 =\frac12\log u+\Omega'(u)
 =\frac\Lambda2\bigl(1+\littleo(1)\bigr),
\]
because \(u=T\bigl(1+\bigO(T^{-2})\bigr)\) and, differentiating
\cref{p2:eq:borel} under the integral sign,
\(\Omega'(u)=-\int_0^\infty s\e^{-us}\widehat\Omega(s)\dd s=\bigO(u^{-2})\).
The mean-value theorem now divides \(\rho_N\) by a quantity asymptotic to
\(\Lambda/2\), which is the instance of \cref{p0:thm:remainder-transport} for
the factorial core and gives \cref{p2:eq:branch-inverse}.  The parity constant
\(A_\delta\) disappears on differentiation, so the estimate is uniform in
\(\delta\).
\end{proof}

\subsection{Two generators for the coefficients}

\Cref{p2:eq:master} is exactly \cref{p0:eq:core-master}, the reversion
problem of \cref{p0:thm:core-reversion}, with the instance data
\begin{equation}
 h(v)=\varphi(v),\quad h_m=\frac{(-1)^m}{m(m-1)},
 \qquad
 f(z,v)=\sum_{k\ge1}\phi_k(z)(1+v)^{-\mu_k},
 \quad \phi_k(z)=a_kz^k,\quad \mu_k=2k-1 .
 \label{p2:eq:instance-data}
\end{equation}
Its triangular recurrence \cref{p0:eq:core-triangular} and its Bell form
\cref{p0:eq:core-bell} therefore apply verbatim; written out, the former is
\begin{equation}
 e_n(\Lambda)=-\frac1\Lambda\,\Coef{z^n}
 \left\{\varphi(E_{n-1})
 +\sum_{k=1}^{n}a_k\,z^k\bigl(1+E_{n-1}\bigr)^{-(2k-1)}\right\},
 \qquad E_{n-1}=\sum_{j<n}e_jz^j .
 \label{p2:eq:recurrence}
\end{equation}
Only the coefficient \emph{structure} is subject-specific.

\begin{theorem}[Structure of the inverse blocks]
\label{p2:thm:recurrence}
For every \(n\ge1\),
\begin{equation}
 e_n(\Lambda)\in\Lambda^{-1}\Rat[\Lambda^{-1}],
 \qquad
 \deg_{\Lambda^{-1}}e_n\le 2n-1,
 \qquad
 e_n(\Lambda)=-\frac{a_n}\Lambda+\bigO(\Lambda^{-2}).
 \label{p2:eq:en-structure}
\end{equation}
\end{theorem}

\begin{proof}
Specialising \cref{p0:eq:core-bell} to \cref{p2:eq:instance-data} gives
\begin{equation}
 e_n=-\frac1\Lambda\left[
  \sum_{m=2}^{n}\frac{(-1)^m}{m(m-1)}\,\OB_{n,m}(e)
  +\sum_{k=1}^{n}a_k\sum_{m=0}^{n-k}(-1)^m\binom{2k+m-2}{m}
     \OB_{n-k,m}(e)\right],
 \label{p2:eq:bell-recurrence}
\end{equation}
in which every occurring \(e_j\) has \(j<n\).  A monomial of \(\OB_{r,m}\) is a
product
\(e_{j_1}\cdots e_{j_m}\) with \(j_1+\dots+j_m=r\); by induction its
\(\Lambda^{-1}\)-degree lies between \(m\) and
\(\sum_\nu(2j_\nu-1)=2r-m\), and the outer factor \(\Lambda^{-1}\) raises this
by one.  In the first sum \(r=n\) and \(m\ge2\), giving degree at most
\(2n-1\); in the second \(r=n-k\), giving at most \(2(n-k)-m+1\le2n-1\).
Rationality is preserved throughout.  Finally the only contribution of
\(\Lambda^{-1}\)-degree exactly \(1\) is the term \(k=n\), \(m=0\), namely
\(-a_n/\Lambda\), because every Bell monomial contains at least one
\(e_j=\bigO(\Lambda^{-1})\).
\end{proof}

\begin{theorem}[Graded closed formula]
\label{p2:thm:closed}
Let \(A(z)=\sum_{k\ge1}a_kz^k\) and put
\begin{equation}
 A_{n,p}:=\Coef{z^n}A(z)^p=\OB_{n,p}(a),
 \qquad
 K_{n;j,p}:=\Coef{v^{j-1}}\varphi(v)^{\,j-p}(1+v)^{-(2n-p)} .
 \label{p2:eq:AK}
\end{equation}
Expand \(d_{2n-1}(\Lambda)=\sum_{j\ge1}c_{n,j}\Lambda^{-j}\).  Then
\begin{equation}
 \boxed{\;
 c_{n,j}=\frac{(-1)^j}{j}
 \sum_{p=\lceil (j+1)/2\rceil}^{\min(j,n)}
 \binom jp\,A_{n,p}\,K_{n;j,p}\;}
 \qquad(1\le j\le 2n-1),
 \label{p2:eq:closed}
\end{equation}
and \(c_{n,j}=0\) for \(j>2n-1\).  The right-hand side is a finite expression
in Bernoulli numbers, binomial coefficients and ordinary partial Bell
polynomials; it computes any single scalar coefficient without generating the
lower blocks.
\end{theorem}

\begin{proof}
Put \(H(z,v)=\varphi(v)+\sum_{k\ge1}a_kz^k(1+v)^{-(2k-1)}\), so that
\cref{p2:eq:master} reads \(E=-\Lambda^{-1}H(z,E)\).  The closed formula
\cref{p0:eq:core-lagrange} of \cref{p0:thm:core-reversion} gives
\[
 e_n=\sum_{j\ge1}\frac1j\,\Coef{z^nv^{j-1}}
 \left(-\frac{H(z,v)}\Lambda\right)^{j}
 =\sum_{j\ge1}\frac{(-1)^j}{j\Lambda^j}\,
   \Coef{z^nv^{j-1}}H(z,v)^j ,
\]
which already exhibits the grading by powers of \(\Lambda^{-1}\).  It remains
to evaluate \(\Coef{z^nv^{j-1}}H^j\).  Expand \(H^j\) binomially, taking \(p\) of the \(j\) factors from the
Bernoulli sector:
\[
 H^j=\sum_{p=0}^{j}\binom jp\,\varphi(v)^{\,j-p}
 \left(\sum_{k\ge1}a_kz^k(1+v)^{-(2k-1)}\right)^{p} .
\]
Extracting \(z^n\) from the \(p\)-th power, the surviving terms have
\(k_1+\dots+k_p=n\); their coefficients sum to \(A_{n,p}\) and they share the
factor \((1+v)^{-\sum_i(2k_i-1)}=(1+v)^{-(2n-p)}\).  Hence
\(\Coef{z^n}H^j=\sum_p\binom jpA_{n,p}\varphi(v)^{j-p}(1+v)^{-(2n-p)}\), and
extracting \(v^{j-1}\) gives \cref{p2:eq:closed}.  (\Cref{p0:eq:core-lagrange}
caps the outer sum at \(j\le2n-1\); the support argument below recovers that
bound independently.)

For the summation range: \(A_{n,p}=0\) unless \(1\le p\le n\), while
\(\varphi(v)^{j-p}\) has \(v\)-valuation \(2(j-p)\), so a nonzero
\(v^{j-1}\)-coefficient forces \(2(j-p)\le j-1\), i.e.\
\(p\ge\lceil(j+1)/2\rceil\).  Combining, \(j\le2p-1\le2n-1\); this re-proves
the degree bound of \cref{p2:eq:en-structure} and shows the expansion
terminates.
\end{proof}

\begin{remark}
\label{p2:rem:closed-checked}
\Cref{p2:eq:closed} was checked against \cref{p2:eq:recurrence} for all
\(1\le n\le5\) and all \(1\le j\le2n-1\) in exact rational arithmetic: all
\(25\) scalars agree.  For explicit evaluation, \(K_{n;j,p}\) unfolds by
\cref{p0:def:obell,p0:lem:power-log} into
\[
 K_{n;j,p}=\sum_{h=2(j-p)}^{j-1}\OB_{h,\,j-p}(\nu)\,(-1)^{j-1-h}
 \binom{2n-p+j-2-h}{j-1-h},
 \qquad
 \nu_1=0,\quad \nu_m=\frac{(-1)^m}{m(m-1)}\ (m\ge2),
\]
with the conventions \(\OB_{0,0}=1\) and \(\OB_{h,0}=0\) for \(h>0\).
\end{remark}

\begin{proposition}[A second, ungraded closed formula]
\label{p2:prop:lagrange-closed}
Let \(S=A^{\langle-1\rangle}\) be the formal compositional inverse of
\(A(z)=\sum_ka_kz^k\), which exists because \(a_1=-\tfrac16\ne0\), with
\begin{equation}
 \Coef{w^m}S(w)=\frac1m\,\Coef{t^{m-1}}\left(\frac t{A(t)}\right)^m .
 \label{p2:eq:S-lagrange}
\end{equation}
Put
\begin{equation}
 G_\Lambda(E)=(1+E)^2\,
 S\!\left(-\frac{\Lambda E+\varphi(E)}{1+E}\right)
 \in E\,\Rat(\Lambda)[[E]],
 \qquad G_\Lambda'(0)=-\frac\Lambda{a_1}=6\Lambda\ne0 .
 \label{p2:eq:G-Lambda}
\end{equation}
Then
\begin{equation}
 e_n(\Lambda)=\frac1n\,\Coef{E^{n-1}}
 \left(\frac E{G_\Lambda(E)}\right)^{n} .
 \label{p2:eq:lagrange-closed}
\end{equation}
\end{proposition}

\begin{proof}
Since \(\sum_ka_kz^k(1+E)^{-(2k-1)}=(1+E)\,A\bigl(z/(1+E)^2\bigr)\),
\cref{p2:eq:master} reads
\(A\bigl(z/(1+E)^2\bigr)=-\bigl(\Lambda E+\varphi(E)\bigr)/(1+E)\).  Applying
\(S\) and multiplying by \((1+E)^2\) gives \(z=G_\Lambda(E)\); differentiating
at \(E=0\), where \(\varphi(0)=\varphi'(0)=0\), gives
\(G_\Lambda'(0)=S'(0)\cdot(-\Lambda)=-\Lambda/a_1\).
\Cref{p2:eq:lagrange-closed,p2:eq:S-lagrange} are then two applications of
\cref{p0:thm:LB}.  Although \(A\) is a divergent Gevrey series, every finite
jet uses only finitely many \(a_k\), so this is a formal identity requiring no
convergence hypothesis.
\end{proof}

\begin{remark}
\Cref{p2:eq:lagrange-closed} was verified against \cref{p2:eq:recurrence} for
\(n\le4\).  It is less useful than \cref{p2:eq:closed} in practice: it does not
separate the powers of \(\Lambda^{-1}\), and it requires reverting \(A\)
first.  Its interest is conceptual --- it exhibits the inverse blocks as a
\emph{double} Lagrange--B\"urmann extraction, the inner one reverting the pure
Bernoulli series.
\end{remark}

\subsection{Boundary coefficients and the global sign law}

\begin{corollary}[The two edges of the coefficient triangle]
\label{p2:cor:boundary}
For every \(n\ge1\),
\begin{equation}
 c_{n,1}=-a_n=-2c_n,
 \qquad
 c_{n,2n-1}=\frac1{3^n}\binom{1/2}{n}
 =-\frac1{2^{n-1}(2n-1)}\binom{2n-1}{n}a_1^n .
 \label{p2:eq:boundary}
\end{equation}
In particular the two closed forms recorded separately by \textsf{S1} and
\textsf{S2} for the deepest coefficient are the same number.
\end{corollary}

\begin{proof}
For \(j=1\) the range in \cref{p2:eq:closed} forces \(p=1\); then
\(A_{n,1}=a_n\), \(K_{n;1,1}=\Coef{v^0}(1+v)^{-(2n-1)}=1\), and the prefactor
is \(-1\).

For \(j=2n-1\), the range forces \(p=n\); then \(A_{n,n}=a_1^n\) and
\[
 K_{n;2n-1,n}=\Coef{v^{2n-2}}\varphi(v)^{\,n-1}(1+v)^{-n}
 =\Coef{v^{2n-2}}\left(\frac{v^2}2\right)^{n-1}=2^{1-n},
\]
because \(\varphi(v)^{n-1}\) has valuation \(2n-2\).  With
\(\tfrac{(-1)^{2n-1}}{2n-1}\binom{2n-1}{n}\) as prefactor this is the second
expression in \cref{p2:eq:boundary}.

Equality of the two expressions is elementary.  Using
\(\binom{1/2}{n}=(-1)^{n-1}\binom{2n}{n}\bigl(4^n(2n-1)\bigr)^{-1}\),
\[
 \frac1{3^n}\binom{1/2}{n}=(-1)^{n-1}\frac{\binom{2n}{n}}{12^n(2n-1)},
\]
while \(a_1=-\tfrac16\) and \(\binom{2n-1}{n}=\tfrac12\binom{2n}{n}\) give
\[
 -\frac{\binom{2n-1}{n}a_1^n}{2^{n-1}(2n-1)}
 =(-1)^{n+1}\frac{\binom{2n}{n}}{2\cdot2^{n-1}6^n(2n-1)}
 =(-1)^{n+1}\frac{\binom{2n}{n}}{12^n(2n-1)} .
\]
\end{proof}

\begin{remark}[The deepest edge as a square root]
\label{p2:rem:deep-edge}
Let \(t_n=c_{n,2n-1}\).  Only the quadratic part of \(\varphi\) and the single
coefficient \(a_1\) can reach \(\Lambda^{-(2n-1)}\): in
\cref{p2:eq:bell-recurrence} the degree bound \(2r-m+1\) attains \(2n-1\) only
for \(r=n\), \(m=2\) in the first sum, since in the second sum
\(2(n-k)-m+1=2n-1\) forces \(k=1\), \(m=0\), where \(\OB_{n-1,0}=0\) for
\(n\ge2\).  Hence \(t_1=-a_1=\tfrac16\) and
\(t_n=-\tfrac12\sum_{j=1}^{n-1}t_jt_{n-j}\) for \(n\ge2\), i.e.\
\(T(z)=\sum_nt_nz^n\) satisfies \(T+\tfrac12T^2=\tfrac16z\), whence
\(T(z)=-1+\sqrt{1+z/3}\) and \(t_n=3^{-n}\binom{1/2}{n}\).  Equivalently, the
maximal-pole part of \cref{p2:eq:master} is
\(\Lambda E+\tfrac12E^2-\tfrac16z=0\) with root
\(E=-\Lambda+\sqrt{\Lambda^2+z/3}\).  \textsf{S1} asserts this without proof;
the degree bookkeeping above supplies it.
\end{remark}

\begin{theorem}[Coefficientwise sign law]
\label{p2:thm:sign-law}
For every \(n\ge1\),
\begin{equation}
 (-1)^{n+1}d_{2n-1}(\Lambda)\in\Lambda^{-1}\Rat_{>0}[\Lambda^{-1}],
 \qquad
 \deg_{\Lambda^{-1}}d_{2n-1}=2n-1 \ \text{exactly}.
 \label{p2:eq:sign-law}
\end{equation}
That is, \(c_{n,j}\ne0\) with sign \((-1)^{n+1}\) for every
\(1\le j\le2n-1\), and \(c_{n,j}=0\) beyond.
\end{theorem}

\begin{proof}
Write \(a_k=(-1)^ka_k^+\) with
\(a_k^+=2(2k-2)!\,\eta(2k)\pi^{-2k}>0\) by \cref{p2:eq:ck}.  Substituting
\(z=-\zeta\) and \(E=-P\) into \cref{p2:eq:master} and using
\(\varphi(-P)=\sum_{m\ge2}P^m/\bigl(m(m-1)\bigr)\) turns it into
\begin{equation}
 \Lambda P=\sum_{m\ge2}\frac{P^m}{m(m-1)}
 +\sum_{k\ge1}a_k^+\,\zeta^k(1-P)^{-(2k-1)} .
 \label{p2:eq:positive-form}
\end{equation}
Every coefficient on the right is \emph{strictly positive}: \(1/(m(m-1))>0\),
\(a_k^+>0\), and \((1-P)^{-(2k-1)}=\sum_{m\ge0}\binom{2k+m-2}{m}P^m\) has
positive coefficients.  Writing \(P=\sum_np_n(\Lambda)\zeta^n\), the
triangular recurrence obtained from \cref{p2:eq:positive-form} expresses
\(p_n\) as \(\Lambda^{-1}\) times a positive-coefficient polynomial in
\(p_1,\dots,p_{n-1}\); by induction each \(p_n\in\Lambda^{-1}
\Rat_{\ge0}[\Lambda^{-1}]\), with no cancellation possible anywhere.  Since
\(E(z)=-P(-z)\), we get \(d_{2n-1}=e_n=(-1)^{n+1}p_n\), which gives the sign
statement once we show that \(p_n\) has a \emph{strictly} positive coefficient
at every \(\Lambda^{-j}\) with \(1\le j\le2n-1\).

We prove that by induction, exhibiting one contributing configuration for each
\(j\); because all contributions are nonnegative, one suffices.  For \(n=1\),
\(p_1=a_1^+/\Lambda\), and \(2n-1=1\).  Let \(n\ge2\) and assume the claim for
all smaller indices.
\begin{itemize}
\item \(j=1\): take the forcing term \(k=n\), \(m=0\) in
      \cref{p2:eq:positive-form}, contributing \(a_n^+\Lambda^{-1}\).
\item \(2\le j\le2n-2\): take the forcing term \(k=1\) together with the
      single factor \(m=1\) of \((1-P)^{-1}\), so that the contribution is
      \(\Lambda^{-1}a_1^+\binom{1}{1}\) times a coefficient of \(p_{n-1}\).
      By induction \(p_{n-1}\) has strictly positive coefficients at
      \(\Lambda^{-j'}\) for every \(1\le j'\le2n-3\), and
      \(j-1\) runs over exactly that range.
\item \(j=2n-1\): take the nonlinear term \(m=2\), i.e.\ a product
      \(p_{n_1}p_{n_2}\) with \(n_1+n_2=n\), \(n_i\ge1\), at their top degrees
      \(2n_i-1\); the resulting \(\Lambda^{-1}\)-degree is
      \(1+(2n_1-1)+(2n_2-1)=2n-1\).
\end{itemize}
Hence every \(c_{n,j}\), \(1\le j\le2n-1\), is nonzero of sign
\((-1)^{n+1}\); \cref{p2:thm:closed} already showed \(c_{n,j}=0\) for
\(j>2n-1\).
\end{proof}

\begin{remark}
The sign law was confirmed symbolically for \(n\le7\): after multiplication by
\((-1)^{n+1}\) all \(1,3,5,7,9,11,13\) coefficients of the respective blocks
are positive rationals, and the degree is exactly \(2n-1\) in each case.
\textsf{S2} states the theorem but its proof gestures at ``a positive
rooted-tree contribution'' without exhibiting one; the attainability induction
above replaces that step.  The law is also the cheapest runtime invariant for a
symbolic implementation (\cref{p2:alg:inverse}).
\end{remark}

\subsection{The explicit blocks}

Repeated use of \cref{p2:eq:recurrence} gives, in exact rational arithmetic,
\begin{align}
 d_1(\Lambda)&=\frac1{6\Lambda},
 \label{p2:eq:d1}\\[0.3em]
 d_3(\Lambda)&=-\frac{14\Lambda^2+10\Lambda+5}{360\Lambda^3},
 \label{p2:eq:d3}\\[0.3em]
 d_5(\Lambda)&=\frac{2232\Lambda^4+1176\Lambda^3+714\Lambda^2
   +350\Lambda+105}{45360\Lambda^5},
 \label{p2:eq:d5}\\[0.3em]
 d_7(\Lambda)&=-\frac1{5443200\Lambda^7}
  \bigl(822960\Lambda^6+292536\Lambda^5+136956\Lambda^4
  \notag\\[-0.2em]
  &\hspace{6.4em}{}+68040\Lambda^3+32970\Lambda^2+12250\Lambda+2625\bigr),
 \label{p2:eq:d7}\\[0.3em]
 d_9(\Lambda)&=\frac1{359251200\Lambda^9}
  \bigl(309052800\Lambda^8+77920128\Lambda^7+26024592\Lambda^6
  \notag\\[-0.2em]
  &\hspace{4.2em}{}+10019592\Lambda^5+4516028\Lambda^4+2023560\Lambda^3
  \notag\\[-0.2em]
  &\hspace{4.2em}{}+812350\Lambda^2+242550\Lambda+40425\bigr),
 \label{p2:eq:d9}\\[0.3em]
 d_{11}(\Lambda)&=-\frac1{5884534656000\Lambda^{11}}
  \bigl(46195687238400\Lambda^{10}+8854370366400\Lambda^9
  \notag\\[-0.2em]
  &\hspace{2.0em}{}+2171643318000\Lambda^8+605515022112\Lambda^7
  +212869408752\Lambda^6
  \notag\\[-0.2em]
  &\hspace{2.0em}{}+84136852800\Lambda^5+35589153600\Lambda^4
  +14234019800\Lambda^3
  \notag\\[-0.2em]
  &\hspace{2.0em}{}+4868113250\Lambda^2+1213962750\Lambda+165540375\bigr).
 \label{p2:eq:d11}
\end{align}
The rapidly growing numerators are an argument for regenerating blocks from
\cref{p2:eq:recurrence} rather than storing a table; \(d_{13}\), which
\textsf{S3} also displays, is omitted here for that reason.  The alternating
block signs and the coefficientwise positivity after multiplication by
\((-1)^{n+1}\) illustrate \cref{p2:thm:sign-law}; the leading and trailing
coefficients illustrate \cref{p2:cor:boundary}.  In separated form,
\begin{align}
 d_1&=\frac1{6\Lambda},
 &d_3&=-\frac7{180\Lambda}-\frac1{36\Lambda^2}-\frac1{72\Lambda^3},
 \notag\\
 d_5&=\frac{31}{630\Lambda}+\frac7{270\Lambda^2}+\frac{17}{1080\Lambda^3}
      +\frac5{648\Lambda^4}+\frac1{432\Lambda^5},
 &&
 \label{p2:eq:d-separated}
\end{align}
\begin{equation}
 d_7=-\frac{127}{840\Lambda}-\frac{4063}{75600\Lambda^2}
  -\frac{11413}{453600\Lambda^3}-\frac1{80\Lambda^4}
  -\frac{157}{25920\Lambda^5}-\frac{35}{15552\Lambda^6}
  -\frac5{10368\Lambda^7}.
 \label{p2:eq:d7-separated}
\end{equation}
The common-denominator forms are less prone to transcription error; the
separated form makes \(c_{n,1}=-a_n\) visible.  Through three blocks,
\begin{align}
 \mathcal D_\delta^{-1}(y)={}&T-1+\frac1{6T\Lambda}
 -\frac{14\Lambda^2+10\Lambda+5}{360\,T^3\Lambda^3}
 \notag\\
 &+\frac{2232\Lambda^4+1176\Lambda^3+714\Lambda^2+350\Lambda+105}
        {45360\,T^5\Lambda^5}
 +\bigO\!\left(\frac1{T^7\Lambda}\right).
 \label{p2:eq:three-blocks}
\end{align}

\begin{remark}[Lambert-only form, and a correction to the usual first guess]
\label{p2:rem:first-guess}
Because \(T=2R_\delta/w\) and \(\Lambda=w+1\), the expansion can be written
without \(T\).  The first three terms are
\begin{align}
 \mathcal D_\delta^{-1}(y)
 \sim{}&\frac{2R_\delta}{w}-1+\frac{w}{12R_\delta(w+1)}
 -\frac{w^3\bigl(14(w+1)^2+10(w+1)+5\bigr)}{2880\,R_\delta^3(w+1)^3}
 \notag\\
 &+\frac{w^5\bigl(2232(w+1)^4+1176(w+1)^3+714(w+1)^2+350(w+1)+105\bigr)}
        {1451520\,R_\delta^5(w+1)^5}+\cdots,
 \label{p2:eq:inverse-in-w}
\end{align}
with \(w=\Wlam_0(2R_\delta/\e)\).  A frequently used rough inversion is
\(x\approx2R_\delta/\Wlam_0(2R_\delta/\e)-1\); its leading correction is
\emph{not} of order \(T^{-1}\) alone but
\[
 \frac1{6\Lambda T}=\frac{w}{12R_\delta(w+1)}
 =\frac1{12R_\delta}\left(1-\frac1{w+1}\right)
 \sim\frac1{12\log y}.
\]
The factor \(\Lambda^{-1}\) is forced by the slope
\(\Phi'(T)=\Lambda/2\) of \cref{p2:eq:core-identities}.
\end{remark}

\section{The outer iterated-logarithmic layer}
\label{p2:sec:outer}

Keeping \(\Wlam_0\) unexpanded is both more accurate and algebraically
cleaner.  To place the result inside a conventional polynomial--logarithmic
transseries one applies \cref{p0:thm:flattening}.  Set
\begin{equation}
 Z=\frac{2R_\delta}\e,\qquad L_1=\log Z,\qquad L_2=\log L_1 .
 \label{p2:eq:L1L2}
\end{equation}

\begin{proposition}[Outer data]
\label{p2:prop:outer-data}
With \(\stirlingfirst{n}{k}\) the unsigned Stirling numbers of the first kind,
\cref{p0:thm:flattening} gives, for \(Z\to+\infty\),
\begin{equation}
 \Wlam_0(Z)\sim L_1-L_2+\sum_{r\ge1}\frac{P_r(L_2)}{L_1^r},
 \qquad
 P_r(t)=\sum_{m=1}^{r}(-1)^{r-m}\stirlingfirst{r}{r-m+1}\frac{t^m}{m!} .
 \label{p2:eq:W-outer}
\end{equation}
Write \(\Wlam_0(Z)/L_1=1+\sum_{j\ge1}\upsilon_j(L_2)L_1^{-j}\) with
\(\upsilon_1=-L_2\) and \(\upsilon_j=P_{j-1}\) for \(j\ge2\), and define the
reciprocal polynomials
\begin{equation}
 \tau_0=1,\qquad \tau_m=-\sum_{j=1}^{m}\upsilon_j\,\tau_{m-j}\quad(m\ge1).
 \label{p2:eq:tau}
\end{equation}
Then
\begin{equation}
 T\sim\frac{2R_\delta}{L_1}\sum_{m\ge0}\frac{\tau_m(L_2)}{L_1^m},
 \qquad
 \Lambda=1+\Wlam_0(Z) .
 \label{p2:eq:T-outer}
\end{equation}
\end{proposition}

The first data, recomputed from \cref{p2:eq:W-outer,p2:eq:tau}, are
\begin{gather}
 P_1=t,\qquad P_2=\tfrac12t^2-t,\qquad P_3=\tfrac13t^3-\tfrac32t^2+t,
 \label{p2:eq:P-first}\\[0.3em]
 \begin{aligned}
 \tau_0&=1, &\tau_1&=L_2, &\tau_2&=L_2^2-L_2,\\
 \tau_3&=L_2^3-\tfrac52L_2^2+L_2,
 &\tau_4&=L_2^4-\tfrac{13}3L_2^3+\tfrac92L_2^2-L_2,\\
 \tau_5&=L_2^5-\tfrac{77}{12}L_2^4+\tfrac{71}6L_2^3-7L_2^2+L_2. & &
 \end{aligned}
 \label{p2:eq:tau-first}
\end{gather}

\begin{theorem}[Every coefficient of the fully expanded transseries]
\label{p2:thm:outer-complete}
Put \(t=L_1^{-1}\),
\begin{equation}
 \mathcal W(t;L_2)=1+\sum_{j\ge1}\upsilon_j(L_2)t^j,
 \qquad
 \mathcal L(t;L_2)=\mathcal W(t;L_2)+t,
 \label{p2:eq:WL-series}
\end{equation}
so that \(\Wlam_0/L_1=\mathcal W\), \(\Lambda/L_1=\mathcal L\) and
\(T=(2R_\delta/L_1)\mathcal W^{-1}\).  With
\(e_n(\Lambda)=\sum_{p=1}^{2n-1}c_{n,p}\Lambda^{-p}\),
\begin{equation}
 \mathcal D_\delta^{-1}(y)
 \sim\frac{2R_\delta}{L_1}\sum_{r\ge0}\frac{\tau_r(L_2)}{L_1^r}-1
 +\sum_{n\ge1}\left(\frac{L_1}{2R_\delta}\right)^{2n-1}
   \sum_{r\ge1}\frac{D_{n,r}(L_2)}{L_1^{r}},
 \label{p2:eq:full-outer}
\end{equation}
where every \(D_{n,r}\in\Rat[L_2]\) is given by the finite extraction
\begin{equation}
 \boxed{\;
 D_{n,r}(L_2)=\Coef{t^r}\left\{
 \mathcal W(t;L_2)^{2n-1}\sum_{p=1}^{2n-1}c_{n,p}\,t^{p}\,
 \mathcal L(t;L_2)^{-p}\right\} . \;}
 \label{p2:eq:Dnr}
\end{equation}
Together with \cref{p2:eq:closed,p2:eq:W-outer,p2:eq:tau} this is a general
formula for every coefficient at every level.  The leading term of the
\(n\)-th inverse-power sector is \(-a_nL_1^{2n-2}(2R_\delta)^{-(2n-1)}\).
\end{theorem}

\begin{proof}
The core term is \cref{p2:eq:T-outer}, and \cref{p2:eq:tau} is reciprocal
series division.  For the \(n\)-th block,
\[
 e_n(\Lambda)\,T^{-(2n-1)}
 =\left(\frac{L_1}{2R_\delta}\right)^{2n-1}\mathcal W^{2n-1}
   \sum_{p=1}^{2n-1}c_{n,p}L_1^{-p}\mathcal L^{-p},
\]
and setting \(t=L_1^{-1}\) and extracting \(\Coef{t^r}\) gives
\cref{p2:eq:Dnr}; at fixed \(r\) only finitely many \(\upsilon_j\) and
\(c_{n,p}\) occur, so the extraction is finite.  The first nonzero power is
\(r=1\), from \(p=1\) and the constant terms \(\mathcal W(0)=\mathcal L(0)=1\),
with \(c_{n,1}=-a_n\) by \cref{p2:cor:boundary}.
\end{proof}

\begin{proposition}[Ordering of the scales]
\label{p2:prop:ordering}
For every fixed \(M\), \(T/L_1^M\to\infty\) as \(y\to\infty\).  Consequently
\begin{equation}
 \frac{T}{L_1^M}\gg 1\gg\frac1{T\Lambda}\gg\frac1{T^3\Lambda}
 \gg\frac1{T^5\Lambda}\gg\cdots,
 \label{p2:eq:ordering}
\end{equation}
so the constant shift \(-1\) is beyond every fixed order of the outer
\(L_1^{-1}\)-expansion, yet dominates the entire inner grid.  In particular the
first Bernoulli displacement satisfies
\(1/(6T\Lambda)\asymp1/(12R_\delta)\asymp1/(12\log y)\).
\end{proposition}

\begin{proof}
\(T=\e^{\Lambda}\) with \(\Lambda=1+\Wlam_0(Z)\sim L_1\), while
\(L_1^M=\e^{ML_2}\) with \(L_2=\log L_1=\littleo(L_1)\);
hence \(\log(T/L_1^M)=\Lambda-ML_2\to\infty\).  The remaining
comparisons are immediate from \(T\to\infty\) and
\cref{p2:eq:core-identities}, which give
\(T\Lambda=2R_\delta\Lambda/(\Lambda-1)\sim2R_\delta\).
\end{proof}

\begin{remark}
This ordering is destroyed if all terms are written as one informal series in
\(1/\log y\): the shift \(-1\), the outer polynomials \(\tau_m(L_2)/L_1^m\)
and the inner blocks \(d_{2n-1}(\Lambda)T^{1-2n}\) then become
indistinguishable.  \textsf{S1}, \textsf{S2} and \textsf{S3} all make this
point; the quantitative version is \cref{p2:eq:ordering}.
\end{remark}

\section{The periodic interpolation and its oscillatory inverse}
\label{p2:sec:oeis-inverse}

\begin{sourcebox}
\textsf{S1} and \textsf{S2} centre the periodic problem at the parity-mean
constant \(\tfrac14\log(2\pi)\); \textsf{S3} centres it at the even-branch
constant \(\tfrac12\log\pi\) and consequently reports a leading displacement
\(2a(1+\cos\theta)/\Lambda\) rather than \(2a\cos\theta/\Lambda\).  The two
are the same function expanded about different cores (they differ by
\(2a/\Lambda+\bigO(\Lambda^{-2})\), which is exactly \(h'\!\cdot a\) for the
core map \(h\)).  This chapter adopts the mean centring, which makes the
carrier a mean-zero oscillation with the cleanest harmonic content.
\end{sourcebox}

\begin{lemma}[Centred exact form of \(\mathscr D\)]
\label{p2:lem:oeis-centred}
With \(u=x+1\), \(a=\tfrac14\log(\pi/2)\) as in \cref{p2:eq:a-def} and
\begin{equation}
 A_*:=\frac14\log(2\pi)=0.4594692666\ldots,
 \qquad \e^{A_*}=(2\pi)^{1/4}=1.5832334871\ldots,
 \label{p2:eq:Astar}
\end{equation}
one has, exactly and for all \(x>-1\),
\begin{equation}
 \log\mathscr D(x)=\Phi(u)+A_*-a\cos(\pi u)+\Omega(u).
 \label{p2:eq:oeis-centred}
\end{equation}
At an integer \(x=n\), \(\cos\pi u=(-1)^{n+1}\) and
\(A_*-a\cos\pi u=A(n)\) of \cref{p2:eq:parity-transmonomial}: the periodic
factor freezes to the correct parity constant.
\end{lemma}

\begin{proof}
By \cref{p2:eq:oeis,p2:thm:normal-form} with \(\delta=0\),
\(\log\mathscr D(x)=\Phi(u)+A_0+\Omega(u)+P(x)\) where
\(P(x)=\tfrac14(1-\cos\pi x)\log(2/\pi)=-a(1-\cos\pi x)\).  Since
\(\cos\pi x=\cos\pi(u-1)=-\cos\pi u\),
\(P=-a(1+\cos\pi u)=-a-a\cos\pi u\), and
\(A_0-a=\tfrac12\log\pi-\tfrac14\log(\pi/2)=\tfrac14\log(2\pi)=A_*\).
\end{proof}

\begin{lemma}[Exact displacement equation]
\label{p2:lem:oeis-displacement}
Let \(R_*=\log y-A_*\) and let \(T,\Lambda\) be the core data of
\cref{p2:eq:lambert-core} built from \(R_*\); put \(q=T^{-1}\),
\(\theta=\pi T\).  Then \(x=\mathscr D^{-1}(y)\) satisfies \(x=T-1+\Delta\)
where \(\Delta\) is the unique small root of
\begin{equation}
 \Lambda\Delta+\frac{\varphi(q\Delta)}{q}
 +2\,\Omega(T+\Delta)-2a\cos(\theta+\pi\Delta)=0 .
 \label{p2:eq:oeis-exact-eq}
\end{equation}
Replacing \(\Omega\) by its asymptotic series turns \cref{p2:eq:oeis-exact-eq}
into the formal equation \(\mathcal F(q,\Delta)=0\) with
\begin{equation}
 \mathcal F(q,v)=\Lambda v+\sum_{m\ge2}\frac{(-1)^m}{m(m-1)}v^mq^{m-1}
 -2a\cos(\theta+\pi v)
 +2\sum_{k\ge1}c_kq^{2k-1}(1+qv)^{-(2k-1)} .
 \label{p2:eq:oeis-F}
\end{equation}
Since the periodic term is \(\bigO(1)\) while every other nonlinear term
carries a positive power of \(q\), one has \(\Delta=\bigO(\Lambda^{-1})\),
\emph{not} \(\bigO(q\Lambda^{-1})\).
\end{lemma}

\begin{proof}
Evaluate \cref{p2:eq:oeis-centred} at \(u=T+\Delta\), set it equal to
\(\log y=R_*+A_*\), subtract \(\Phi(T)=R_*\), and multiply by \(2\).  By the
computation in \cref{p2:lem:normalized} with \(E=q\Delta\),
\(2\bigl[\Phi(T+\Delta)-\Phi(T)\bigr]
=T\bigl[\Lambda q\Delta+\varphi(q\Delta)\bigr]
=\Lambda\Delta+\varphi(q\Delta)/q\), since \(T=q^{-1}\).  This gives
\cref{p2:eq:oeis-exact-eq}.  The last
sum in \cref{p2:eq:oeis-F} is \(2\Omega(T(1+q v))\) expanded.  Finally
\(\varphi(q\Delta)/q=q\Delta^2/2+\bigO(q^2)\), so at \(q=0\) the equation
reduces to \(\Lambda\Delta_0=2a\cos(\theta+\pi\Delta_0)\), forcing
\(\Delta_0\asymp\Lambda^{-1}\).
\end{proof}

\subsection{The periodic carrier}

\begin{theorem}[Existence, uniqueness, and the closed carrier expansion]
\label{p2:thm:carrier}
Fix \(\theta\in\Real\) and put \(\mu=2a/\Lambda\).
\begin{enumerate}[label=\textup{(\alph*)}]
\item If \(\Lambda>2\pi a=0.70934\ldots\), the equation
      \begin{equation}
       \Delta_0=\mu\cos(\theta+\pi\Delta_0)
       \label{p2:eq:carrier-eq}
      \end{equation}
      has exactly one real solution, and \(|\Delta_0|\le\mu\).
\item If moreover \(\mu\le\tfrac15\), equivalently
      \(\Lambda\ge10a=1.12896\ldots\), that solution is the sum of the
      convergent Lagrange--B\"urmann series
      \begin{equation}
       \boxed{\;
       \Delta_0=\sum_{m\ge1}\frac{\mu^m}{m!}
       \left.\frac{\dd^{m-1}}{\dd v^{m-1}}\cos^m(\theta+\pi v)\right|_{v=0}
       =\sum_{m\ge1}\frac{a^m}{m!\,\Lambda^m}
        \sum_{j=0}^{m}\binom mj\bigl(\ii\pi(m-2j)\bigr)^{m-1}
        \e^{\ii(m-2j)\theta} . \;}
       \label{p2:eq:carrier-series}
      \end{equation}
      The second expression is real.  At inverse-logarithmic order \(m\) only
      the Fourier modes \(m,m-2,\dots,-m\) occur.
\item The first three terms are
      \begin{equation}
       \Delta_0=\mu\cos\theta-\frac\pi2\mu^2\sin2\theta
       +\frac{\pi^2}2\mu^3\cos\theta\,(2-3\cos^2\theta)+\bigO(\mu^4).
       \label{p2:eq:carrier-first}
      \end{equation}
\end{enumerate}
\end{theorem}

\begin{proof}
(a) The map \(g(v)=\mu\cos(\theta+\pi v)\) satisfies
\(|g'|\le\pi\mu=2\pi a/\Lambda<1\) and \(|g|\le\mu\), so it is a contraction of
\(\Real\) into \([-\mu,\mu]\); Banach's theorem gives a unique fixed point
there, and any real fixed point lies in \([-\mu,\mu]\).

(b) \Cref{p2:eq:carrier-eq} is \(v=\mu\phi(v)\) with
\(\phi(v)=\cos(\theta+\pi v)\) entire.  By the classical convergence criterion
for Lagrange's series, if there is \(r>0\) with
\(|\mu|\max_{|v|=r}|\phi(v)|<r\), then the series
\(\sum_m\tfrac{\mu^m}{m!}\bigl(\tfrac{\dd}{\dd v}\bigr)^{m-1}\phi(v)^m|_{v=0}\)
converges and its sum is the unique root of \(v=\mu\phi(v)\) in \(|v|<r\).
Here \(|\cos(\theta+\pi v)|\le\cosh(\pi|{\Im}\,v|)\le\cosh(\pi r)\) on
\(|v|=r\); with \(r=0.38\) and \(\mu\le\tfrac15\),
\(\mu\cosh(0.38\pi)\le0.3603<0.38\).  (The optimal choice
\(\max_{r>0}r/\cosh(\pi r)=0.21095\ldots\) at \(r=0.382\) shows
\(\mu<0.2109\) already suffices.)  For the Fourier form, insert
\(\cos^m v=2^{-m}\sum_{j}\binom mj\e^{\ii(m-2j)v}\) into the derivative:
\[
 \left.\frac{\dd^{m-1}}{\dd v^{m-1}}\cos^m(\theta+\pi v)\right|_{v=0}
 =2^{-m}\sum_{j=0}^m\binom mj\bigl(\ii\pi(m-2j)\bigr)^{m-1}
 \e^{\ii(m-2j)\theta},
\]
and \(\mu^m2^{-m}=a^m\Lambda^{-m}\).  Pairing \(j\) with \(m-j\) replaces
\(m-2j\) by its negative, conjugating each summand, so the sum is real.  Only
\(|m-2j|\le m\) occurs, with the same parity as \(m\).

(c) For \(m=1\) the derivative is \(\cos\theta\).  For \(m=2\),
\(\tfrac{\dd}{\dd v}\cos^2=-2\pi\cos\sin=-\pi\sin2(\cdot)\), giving
\(-\tfrac\pi2\mu^2\sin2\theta\).  For \(m=3\),
\(\tfrac{\dd^2}{\dd v^2}\cos^3(\theta+\pi v)|_{0}
=3\pi^2\cos\theta(2-3\cos^2\theta)\), and division by \(3!\) gives the third
term.
\end{proof}

\begin{repairbox}
\textsf{S1} states the carrier expansion as ``convergent for sufficiently
large \(\Lambda\)'' with no threshold and no criterion, and \textsf{S2} proves
only the contraction estimate, from which convergence of the Lagrange series
does not follow.  \Cref{p2:thm:carrier} separates the two claims and supplies
an explicit sufficient condition, \(\mu\le\tfrac15\).  Numerically at
\(\Lambda=3\), \(\theta=2\), the partial sums \(m\le1,\dots,5\) of
\cref{p2:eq:carrier-series} approximate the exact fixed point
\(-0.0256993556\) with errors
\(5.6\cdot10^{-3},1.1\cdot10^{-3},1.8\cdot10^{-4},1.9\cdot10^{-5},
1.9\cdot10^{-6}\), consistent with a convergence ratio \(\approx\mu/0.211\).
\end{repairbox}

\subsection{Algebraic corrections around the carrier}

\begin{theorem}[Triangular periodic-grid recurrence]
\label{p2:thm:periodic-recurrence}
Let
\begin{equation}
 M:=\partial_v\mathcal F(0,\Delta_0)=\Lambda+2\pi a\sin(\theta+\pi\Delta_0),
 \label{p2:eq:M}
\end{equation}
which satisfies \(M\ge\Lambda-2\pi a>0\) for \(\Lambda>2\pi a\).  Seek
\(\Delta=\Delta_0+\sum_{n\ge1}\Delta_nq^n\) with each \(\Delta_n\) a formal
series in \(\Lambda^{-1}\) with trigonometric-polynomial coefficients, and put
\(\Delta_{<n}=\Delta_0+\sum_{j<n}\Delta_jq^j\).  Then the solution of
\(\mathcal F=0\) is unique and generated by
\begin{equation}
 \boxed{\;
 \Delta_n=-\frac1M\,\Coef{q^n}\,
 \mathcal F\bigl(q,\Delta_{<n}(q)\bigr),\qquad n\ge1. \;}
 \label{p2:eq:periodic-recurrence}
\end{equation}
The first two are
\begin{equation}
 \Delta_1=\frac{\tfrac16-\tfrac12\Delta_0^2}{M},
 \qquad
 \Delta_2=-\frac1M\left[
 a\pi^2\cos(\theta+\pi\Delta_0)\,\Delta_1^2
 +\Delta_0\Delta_1+\frac{\Delta_0-\Delta_0^3}{6}\right].
 \label{p2:eq:Delta12}
\end{equation}
\end{theorem}

\begin{proof}
In \(\Coef{q^n}\mathcal F(q,\Delta)\) the unknown \(\Delta_n\) appears only
linearly, through the \(v\)-derivative of the \(q\)-free part
\(\Lambda v-2a\cos(\theta+\pi v)\) evaluated at \(v=\Delta_0\), which is
\(M\); every other contribution involves only \(\Delta_j\) with \(j<n\).
This is \cref{p0:lem:lagrange-fp} applied in the variable \(q\) over the
coefficient ring generated by \(\Lambda^{-1}\) and \(\e^{\pm\ii\theta}\), and
\(M\) is invertible there because \(M=\Lambda(1+\bigO(\Lambda^{-1}))\).

For \(n=1\): the \(m=2\) term of \cref{p2:eq:oeis-F} gives
\(\tfrac12\Delta_0^2\), the cosine gives \(2a\pi\sin(\theta+\pi\Delta_0)
\Delta_1\), and the \(k=1\) Bernoulli term gives \(2c_1=-\tfrac16\); collecting
gives \(M\Delta_1+\tfrac12\Delta_0^2-\tfrac16=0\).  For \(n=2\): the \(m=2\)
term gives \(\Delta_0\Delta_1\), the \(m=3\) term gives
\(-\tfrac16\Delta_0^3\), the second-order Taylor term of the cosine gives
\(a\pi^2\cos(\theta+\pi\Delta_0)\Delta_1^2\), and expanding
\(2c_1q(1+q\Delta)^{-1}\) gives \(-2c_1\Delta_0=\tfrac16\Delta_0\); collecting
gives \cref{p2:eq:Delta12}.
\end{proof}

\begin{corollary}[Asymptotic meaning]
\label{p2:cor:periodic-meaning}
Fix \(N,K\ge0\).  Truncate \(\Delta\) after \(q^N\) and expand each
\(\Delta_n\) in \(\Lambda^{-1}\) through \(\Lambda^{-K}\), obtaining
\(\widetilde\Delta\).  Then, as \(y\to+\infty\),
\begin{equation}
 \mathscr D^{-1}(y)
 =T-1+\widetilde\Delta
 +\bigO\!\left(\frac1{T^{N+1}\Lambda}\right)
 +\bigO\!\left(\frac1{\Lambda^{K+2}}\right),
 \label{p2:eq:oeis-inverse}
\end{equation}
uniformly in the phase \(\theta=\pi T\).
\end{corollary}

\begin{proof}
By construction \cref{p2:eq:periodic-recurrence} cancels the residual of
\(\mathcal F\) through \(q^N\), and \cref{p2:thm:carrier} cancels it through
\(\Lambda^{-K}\) at each fixed \(q\)-order; the replacement of \(\Omega\) by
its truncation is controlled by \cref{p2:thm:remainder}.  The exact
\(v\)-derivative of the left side of \cref{p2:eq:oeis-exact-eq} equals
\(M+\bigO(q)\) and is bounded below by \(\Lambda-2\pi a-\bigO(q)\), a positive
multiple of \(\Lambda\) for large \(T\); all trigonometric coefficients and
their derivatives are bounded, so \cref{p0:thm:backward-error} converts the
residual into a root error after division by a quantity \(\asymp\Lambda\),
uniformly in \(\theta\).
\end{proof}

\begin{proposition}[The parity branches resum the periodic sector]
\label{p2:prop:resummation}
At an integer \(u=n+1\) one has \(A_*-a\cos\pi u=A_\delta\) with
\(\delta\equiv n\pmod2\).  The branchwise normalisation therefore absorbs the
whole \(\bigO(1)\) periodic constant into the target \(R_\delta\) before the
core is formed, so the \(\bigO(\Lambda^{-1})\) carrier disappears and the first
correction is the much smaller \(1/(6T\Lambda)\).  In this precise sense the two
parity transseries are subsequence resummations of the periodic one: they build
the parity phase into the core instead of expanding it perturbatively.
\end{proposition}

\begin{proof}
Immediate from \cref{p2:lem:oeis-centred} and
\cref{p2:thm:branch-inverse,p2:thm:carrier}.
\end{proof}

\begin{proposition}[Phase-coordinate reversion, formal]
\label{p2:prop:phase-reversion}
Let \(h(R)=2R/\Wlam_0(2R/\e)\) be the inverse of \(\Phi\)
(\cref{p2:lem:core}), and set
\(\mathcal A(R)=-a\cos\bigl(\pi h(R)\bigr)+\Omega\bigl(h(R)\bigr)\).  Then the
equation \(\Phi(u)+\mathcal A(\Phi(u))=R_*\) --- which is
\cref{p2:eq:oeis-centred} rewritten in the phase coordinate
\(\sigma=\Phi(u)\) --- has the formal reversion
\begin{equation}
 \sigma=R_*+\sum_{m\ge1}\frac{(-1)^m}{m!}
 \left(\frac{\dd}{\dd R_*}\right)^{m-1}\mathcal A(R_*)^m,
 \qquad
 \mathscr D^{-1}(y)=h(\sigma)-1 .
 \label{p2:eq:phase-reversion}
\end{equation}
\end{proposition}

\begin{proof}
Write \(\sigma=R_*+\varsigma\); then \(\varsigma=-\mathcal A(R_*+\varsigma)\),
and Lagrange's formula for \(\varsigma=\tau\phi(\varsigma)\) with
\(\phi(\varsigma)=-\mathcal A(R_*+\varsigma)\), \(\tau=1\), together with the
translation invariance of \(\dd/\dd\varsigma\), gives
\cref{p2:eq:phase-reversion}.
\end{proof}

\begin{repairbox}
\textsf{S3} states \cref{p2:eq:phase-reversion} as a boxed asymptotic
expansion, writing ``setting \(\tau=1\) as an asymptotic expansion gives''.
That is a formal-to-analytic slide: \(\mathcal A\) is not small in any norm
that makes the reversion series converge, and \(m\)-fold differentiation with
respect to \(R_*\) produces factors \(h'(R_*)=2/\Lambda\) whose accumulated
effect is exactly what needs to be estimated.  The statement is demoted here to
a \emph{formal} proposition.  An asymptotic claim would require, at each order
\(m\), a bound on \(\bigl(\dd/\dd R_*\bigr)^{m-1}\mathcal A^m\) uniform in the
oscillatory phase --- available for the truncated grid of
\cref{p2:cor:periodic-meaning}, which is therefore the statement to use.
Comparing the two routes at leading order: the \(m=1\) term of
\cref{p2:eq:phase-reversion} contributes
\(h'(R_*)\cdot(-\mathcal A)=2a\cos(\pi T)/\Lambda+\bigO(T^{-1}\Lambda^{-1})\),
in agreement with \cref{p2:eq:carrier-first}.
\end{repairbox}

\begin{warning}
\label{p2:rem:not-intrinsic}
The carrier of \cref{p2:thm:carrier} records the choice
\cref{p2:eq:oeis}, not a property of \(n!!\).  By
\cref{p2:thm:no-canonical} another interpolation, agreeing with \(n!!\) at
every nonnegative integer, changes the amplitude of this sector from \(2a\) to
\(2\sqrt{a^2+\lambda^2}\) and shifts its phase, while leaving every integer
value and both parity transseries untouched.
\end{warning}

\section{The discrete inverses}
\label{p2:sec:discrete}

Objects (i) and (iii) of \cref{p2:sec:orientation} are intrinsic to the
sequence.  They are the residue-class staircase of
\cref{p0:def:three-inverses,p0:thm:staircase} at modulus \(r=2\), and nothing
new has to be proved: \cref{p0:rem:parity-instances} already records this
chapter as that instance.  Only the dictionary and the arithmetic caveat are
subject-specific.

\begin{proposition}[Specialisation of the staircase]
\label{p2:prop:discrete}
Take \(r=2\), residue classes \(\rho\in\{0,1\}\), \(A_n=n!!\), and
\(\nu_\rho(y)=\mathcal D_\rho^{-1}(y)\), which is legitimate because
\(\mathcal D_\rho\) interpolates the class-\(\rho\) subsequence and is strictly
increasing (\cref{p2:prop:branch-interp}).  Then
\cref{p0:eq:residue-staircase} reads
\begin{equation}
 N(y):=\min\{n\in\Nat: n!!\ge y\}
 =\min\left\{2\left\lceil\frac{\nu_0(y)}2\right\rceil,\;
   2\left\lceil\frac{\nu_1(y)-1}2\right\rceil+1\right\}
 \qquad(y>1),
 \label{p2:eq:N-exact}
\end{equation}
and its floor counterpart, obtained from \cref{p0:thm:staircase} by replacing
\(\lceil\cdot\rceil\) with \(\lfloor\cdot\rfloor\) and \(\min\) with \(\max\),
is
\begin{equation}
 J(y):=\max\{n\in\Nat: n!!\le y\}
 =\max_{\rho\in\{0,1\}}
 \left\{\rho+2\left\lfloor\frac{\nu_\rho(y)-\rho}2\right\rfloor\right\}.
 \label{p2:eq:J-exact}
\end{equation}
Replacing \(\nu_\rho\) by the truncation of \cref{p2:eq:branch-inverse} with a
certified error bound \(\eta_\rho\) (for instance from
\cref{p2:eq:aposteriori}) determines \(N(y)\) and \(J(y)\) exactly under the
separation condition of \cref{p0:thm:staircase}, and
\cref{p0:prop:certified-rounding} bounds the number of exact evaluations of
\(n!!\) needed by \(\sum_\rho(\lfloor\eta_\rho\rfloor+2)\).  Those evaluations
should be made as \(\log\mathcal D_\rho(n)=\log\kappa_\rho+\tfrac n2\log2
+\log\Gamma(\tfrac n2+1)\), or by \cref{p2:eq:dfac-recurrence} in exact
integers.
\end{proposition}

\begin{proof}
Everything is \cref{p0:thm:staircase} and \cref{p0:prop:certified-rounding}
with \(r=2\); the hypotheses hold by \cref{p2:prop:branch-interp} and, for the
asymptotic input, by \cref{p2:thm:branch-inverse}.
\end{proof}

\begin{warning}[Never form \(y\)]
\label{p2:rem:never-form-y}
\Cref{p2:eq:N-exact}, evaluated with four inverse blocks, was tested at
\(y=n!!\) for every \(6\le n\le199\) and returned the correct \(N(y)=n\) in all
\(194\) cases --- but only when the comparison \(A_n\ge y\) is made on exact
integers or on \(\log y\).  Converting \(n!!\) to a fixed-precision float first
destroys the comparison long before the asymptotics fail: at \(80\) working
digits this already produces spurious off-by-one answers from
\(n\approx120\) onwards, because \(120!!\) has more than \(100\) digits.  The
failure is arithmetic, not asymptotic, and it is easy to misdiagnose as a
defect of the expansion.
\end{warning}

\section{Beyond all algebraic orders}
\label{p2:sec:beyond}

\begin{proposition}[Optimal truncation, forward and inverse]
\label{p2:prop:least-term}
Let \(u>0\).  The terms of \cref{p2:eq:Omega-series} have moduli
\(|c_k|u^{1-2k}\) with consecutive ratio
\(\sim2k(2k-1)/(\pi u)^2\); the least term occurs at \(2k\approx\pi u\) and
\begin{equation}
 \min_k\left|\frac{c_k}{u^{2k-1}}\right|
 =\bigO\!\left(u^{-1/2}\e^{-\pi u}\right).
 \label{p2:eq:least-term}
\end{equation}
Consequently, by \cref{p0:thm:optimal-truncation} and the slope
\(\Phi'(T)=\Lambda/2\), the optimally truncated branchwise inverse has error
\begin{equation}
 \bigO\!\left(\frac{\e^{-\pi T}}{\sqrt T\,\Lambda}\right).
 \label{p2:eq:inverse-least-term}
\end{equation}
\end{proposition}

\begin{proof}
\Cref{p2:eq:ck-late} gives \(|c_k|=(2k-2)!\pi^{-2k}(1+\bigO(4^{-k}))\), whence
the ratio.  At \(2k=\pi u\), Stirling's formula for \((2k-2)!\) gives
\cref{p2:eq:least-term}.  For the inverse, \cref{p2:thm:remainder} bounds the
forward logarithmic residual at optimal truncation by
\cref{p2:eq:least-term} with \(u=T(1+\bigO(T^{-2}))\); dividing by the slope
\(\tfrac12\Lambda+\bigO(T^{-2})\) via
\cref{p0:thm:backward-error,p0:thm:remainder-transport} gives
\cref{p2:eq:inverse-least-term}.
\end{proof}

\begin{proposition}[Transport of a forward exponential sector]
\label{p2:prop:transport}
Suppose an exponentially improved representation of the centred forward phase
retains the algebraic terms \(c_1,\dots,c_K\) and has an additional residual
\(\mathcal R\), and let \(u_{\mathrm{alg}}\) solve the correspondingly
truncated algebraic inverse problem.  Then
\begin{equation}
 u-u_{\mathrm{alg}}
 =-\frac{2\,\mathcal R(u_{\mathrm{alg}})}
   {\log u_{\mathrm{alg}}
    -2\sum_{k=1}^{K}(2k-1)c_k\,u_{\mathrm{alg}}^{-2k}}
 +\bigO\!\left(\frac{\sup|F''|\,\mathcal R^2}{\Lambda^3}
   +\frac{|\mathcal R\,\mathcal R'|}{\Lambda^2}\right),
 \label{p2:eq:transport}
\end{equation}
where \(F\) is the doubled algebraic phase and the supremum is over the
interval between \(u_{\mathrm{alg}}\) and the exact root; the estimate assumes
the displayed denominator is \(\asymp\Lambda\) and
\(|\mathcal R'|=\littleo(\Lambda)\).  At leading order this is
\(-2\mathcal R(T)/\Lambda\).
\end{proposition}

\begin{proof}
Let \(F\) be the doubled algebraic phase minus \(2R\), so
\(F(u_{\mathrm{alg}})=0\) and the exact phase equals
\(F+2\mathcal R\).  Taylor expansion of
\(0=F(u_{\mathrm{alg}}+\Delta)+2\mathcal R(u_{\mathrm{alg}}+\Delta)\) about
\(u_{\mathrm{alg}}\), with \(F'(u_{\mathrm{alg}})\) the stated denominator,
gives the first Newton correction; replacing
\(F'+2\mathcal R'\) by \(F'\) and iterating once produces the two displayed
error terms.
\end{proof}

\begin{warning}[What the exponential sector is, and is not]
\label{p2:rem:terminant}
The Borel poles of \cref{p2:eq:mittag} lie at \(s=\pm\ii\pi m\), so the
exponentials associated with them are \(\e^{\mp\ii\pi mu}\) together with
direction-dependent terminant multipliers, not real exponentials.  On the
positive axis the optimally truncated remainder has modulus
\(\asymp u^{-1/2}\e^{-\pi u}\), but this is the size of a
\emph{terminant-weighted combination}: appending a bare term
\(\mathrm{const}\cdot\e^{-\pi u}\) to the divergent Bernoulli series is not a
valid transseries completion, and conflating the two gives incorrect Stokes
statements.  The unambiguous positive-real completion is the exact integral
\cref{p2:eq:borel} itself, which by \cref{p2:thm:remainder} even provides
two-sided bounds at every order; \cref{p2:eq:transport} then carries any
rigorously normalised forward completion through the inverse map.  All three
sources make this point; \textsf{S3} states it most sharply.
\end{warning}

\section{Complex inverse sheets}
\label{p2:sec:complex}

\begin{proposition}[Local sheets]
\label{p2:prop:sheets}
Fix a branch of \(\log y\) and integers \(\ell\) (logarithmic winding) and
\(k\) (Lambert branch).  Put
\begin{equation}
 R_{\delta,\ell}=\log y+2\pi\ii\ell-A_\delta,
 \quad
 w_{k,\ell}=\Wlam_k\!\left(\frac{2R_{\delta,\ell}}\e\right),
 \quad
 T_{k,\ell}=\frac{2R_{\delta,\ell}}{w_{k,\ell}},
 \quad
 \Lambda_{k,\ell}=w_{k,\ell}+1 .
 \label{p2:eq:complex-core}
\end{equation}
In any sector in which \(T_{k,\ell}\to\infty\), the Stirling expansion for
\(\log\Gamma\) is uniform, and the corresponding solution branch of
\(\mathcal D_\delta(x)=y\) is simple, the \emph{same} universal polynomials
give
\begin{equation}
 x_{\delta;k,\ell}(y)\sim T_{k,\ell}-1
 +\sum_{n\ge1}\frac{d_{2n-1}(\Lambda_{k,\ell})}{T_{k,\ell}^{2n-1}} .
 \label{p2:eq:complex-sheet}
\end{equation}
\end{proposition}

\begin{proof}
Every step of \cref{p2:lem:normalized,p2:thm:recurrence} is an identity of
formal series over \(\Cplx(\Lambda)\) and uses only \(\Phi(T)=R\),
\(\Lambda=\log T\); \cref{p2:eq:complex-core} secures those.  The analytic
statement is the sectorial version of \cref{p2:thm:branch-inverse}, with the
mean-value step replaced by the standard implicit-function estimate in the
sector.
\end{proof}

\begin{remark}[Branch caveats]
\label{p2:rem:branch-caveat}
The pair \((k,\ell)\) is an asymptotic label, not a global classification:
distinct labels can be joined by analytic continuation around the poles of
\(\Gamma\) or the logarithmic cut.  Branch collisions occur exactly where
\(\mathcal D_\delta'=0\), i.e.\ where \(\log2+\psi(x/2+1)=0\); by
\cref{p2:prop:branch-interp} this never happens for real \(x\ge0\), but it does
for complex and for negative real \(x\), and near such points
\cref{p2:eq:complex-sheet} must be replaced by a Puiseux-type local analysis.
\end{remark}

\section{Multifactorials}
\label{p2:sec:multifactorial}

\begin{sourcebox}
Only \textsf{S3} treats multifactorials.  Its normal form is reproduced here
with a proof, and \cref{p2:thm:multi-scaling} is new: it subsumes the
``half-shifted inverse Gamma'' scaling that \textsf{S1} and \textsf{S2} record
separately, by identifying the Gamma case as the member \(p=1\).
\end{sourcebox}

\begin{proposition}[Residue-class branches]
\label{p2:prop:multi-branch}
Let \(p\ge1\) and \(r\in\{1,\dots,p\}\).  For integers \(x\equiv r\pmod p\),
\(x\ge r\), set \(x!_{(p)}=x(x-p)(x-2p)\cdots r\).  Then, with
\begin{equation}
 \mathcal D_{p,r}(x)=C_{p,r}\;p^{x/p}\,\Gamma\!\left(\frac xp+1\right),
 \qquad
 C_{p,r}=\frac{p^{1-r/p}}{\Gamma(r/p)},
 \label{p2:eq:multi-branch}
\end{equation}
one has \(\mathcal D_{p,r}(x)=x!_{(p)}\) for every such \(x\).  For \(p=2\),
\(C_{2,2}=1=\kappa_0\) and \(C_{2,1}=\sqrt{2/\pi}=\kappa_1\); for \(p=1\),
\(C_{1,1}=1\) and \(\mathcal D_{1,1}(x)=\Gamma(x+1)\).
\end{proposition}

\begin{proof}
Write \(x=r+mp\).  Then
\(x!_{(p)}=\prod_{i=0}^{m}(r+ip)=p^{m+1}\prod_{i=0}^m(\tfrac rp+i)
=p^{m+1}\Gamma(\tfrac rp+m+1)/\Gamma(\tfrac rp)\), while
\(\mathcal D_{p,r}(x)=C_{p,r}p^{r/p+m}\Gamma(\tfrac rp+m+1)\).  Equating gives
\(C_{p,r}=p^{1-r/p}/\Gamma(r/p)\).
\end{proof}

\begin{theorem}[Centred multifactorial normal form]
\label{p2:thm:multi-normal}
Put \(u=x+\tfrac p2\) and
\(A_{p,r}=\log C_{p,r}+\tfrac12\log(2\pi/p)\).  Then, as \(x\to+\infty\),
\begin{equation}
 \log\mathcal D_{p,r}(x)=\frac up(\log u-1)+A_{p,r}+\Omega_p(u),
 \qquad
 \Omega_p(u)\sim\sum_{k\ge1}\frac{c^{(p)}_k}{u^{2k-1}},
 \quad
 c^{(p)}_k=\frac{p^{2k-1}B_{2k}(\tfrac12)}{2k(2k-1)},
 \label{p2:eq:multi-normal}
\end{equation}
and \(\Omega_p\) has the exact Borel--Laplace representation
\begin{equation}
 \Omega_p(u)=\int_0^\infty\e^{-us}\,
 \frac{\dfrac{ps}{2\sinh(ps/2)}-1}{ps^2}\,\dd s ,
 \label{p2:eq:multi-borel}
\end{equation}
whose singularities are the simple poles \(s=\pm2\pi\ii m/p\).  Hence the
optimally truncated forward remainder has scale \(\e^{-2\pi u/p}\).  For
\(p=2\) this is \cref{p2:thm:normal-form,p2:thm:borel}; for \(p=1\) it is the
half-shifted Stirling expansion of \(\log\Gamma(x+1)=\log\Gamma(u+\tfrac12)\).
\end{theorem}

\begin{proof}
Put \(t=u/p\), so \(x=pt-\tfrac p2\) and
\(\log\mathcal D_{p,r}=\log C_{p,r}+(t-\tfrac12)\log p
+\log\Gamma(t+\tfrac12)\).  The half-shifted Stirling expansion
\(\log\Gamma(t+\tfrac12)\sim t\log t-t+\tfrac12\log(2\pi)
+\sum_k B_{2k}(\tfrac12)\bigl(2k(2k-1)\bigr)^{-1}t^{1-2k}\) --- the case
\(a=\tfrac12\) of \(\log\Gamma(t+a)\), in which every odd-index term vanishes
because \(B_{2j+1}(\tfrac12)=0\) --- gives
\[
 \log\mathcal D_{p,r}
 \sim\log C_{p,r}-\tfrac12\log p+\tfrac12\log(2\pi)
 +\frac up(\log u-1)
 +\sum_k\frac{p^{2k-1}B_{2k}(\tfrac12)}{2k(2k-1)u^{2k-1}},
\]
since \((t-\tfrac12)\log p+t\log t=(u/p)\log u-\tfrac12\log p\) and
\(t^{1-2k}=p^{2k-1}u^{1-2k}\).  The constant is \(A_{p,r}\).  For
\cref{p2:eq:multi-borel}, \(\Omega_p(u)=\widetilde\Omega(u/p)\) with
\(\widetilde\Omega\) as in the proof of \cref{p2:thm:borel}, where
\(\widetilde\Omega(t)=\int_0^\infty\e^{-tr}
\bigl(r/(2\sinh(r/2))-1\bigr)r^{-2}\dd r\); substituting \(r=ps\) gives
\cref{p2:eq:multi-borel}, and \(\sinh(ps/2)=0\) at \(s=2\pi\ii m/p\).
\end{proof}

\begin{theorem}[Universal multifactorial inverse and its scaling]
\label{p2:thm:multi-scaling}
Let \(R=\log y-A_{p,r}\),
\(v=\Wlam_0(pR/\e)\), \(X=pR/v\), \(\Lambda=v+1=\log X\), so that
\(\tfrac Xp(\log X-1)=R\).  With \(u=X(1+E)\), \(q=X^{-1}\), \(z=q^2\), the
normalised equation is
\begin{equation}
 \Lambda E+\varphi(E)+\sum_{k\ge1}a^{(p)}_kz^k(1+E)^{-(2k-1)}=0,
 \qquad
 a^{(p)}_k:=p\,c^{(p)}_k=p^{2k}\,\frac{B_{2k}(\tfrac12)}{2k(2k-1)} .
 \label{p2:eq:multi-master}
\end{equation}
Every statement of
\cref{p2:thm:recurrence,p2:thm:closed,p2:prop:lagrange-closed,p2:cor:boundary,p2:thm:sign-law}
holds verbatim with \(a_k\) replaced by \(a^{(p)}_k\), and for every fixed
\(N\)
\begin{equation}
 \mathcal D_{p,r}^{-1}(y)=X-\frac p2
 +\sum_{n=1}^{N}\frac{d^{(p)}_{2n-1}(\Lambda)}{X^{2n-1}}
 +\bigO\!\left(\frac1{X^{2N+1}\Lambda}\right).
 \label{p2:eq:multi-inverse}
\end{equation}
Moreover the coefficient polynomials of all multifactorials are one family:
\begin{equation}
 \boxed{\;
 d^{(p)}_{2n-1}(\Lambda)=p^{2n}\,d^{(1)}_{2n-1}(\Lambda)
 \qquad\text{for every }p\ge1,\ n\ge1, \;}
 \label{p2:eq:multi-scaling}
\end{equation}
where \(d^{(1)}\) are the blocks of the half-shifted inverse-Gamma problem.  In
particular \(d_{2n-1}=d^{(2)}_{2n-1}=4^n\,d^{(1)}_{2n-1}\).
\end{theorem}

\begin{proof}
The derivation of \cref{p2:eq:multi-master} repeats
\cref{p2:lem:normalized} with \(\tfrac2T\) replaced by \(\tfrac pX\): the
phase difference is \(\tfrac Xp[\Lambda E+\varphi(E)]\), and
\(\tfrac pX\Omega_p(u)=\sum_kp\,c^{(p)}_kq^{2k}(1+E)^{-(2k-1)}\).  Since
\cref{p2:eq:multi-master} differs from \cref{p2:eq:master} only in the values
of the forcing coefficients, all structural statements transfer.  Now
\(a^{(p)}_k=p^{2k}a^{(1)}_k\) by \cref{p2:eq:multi-normal}, so substituting
\(q\mapsto pq\) (equivalently \(z\mapsto p^2z\)) carries the \(p=1\) equation
into the general one with \(\Lambda\) untouched.  Uniqueness of the formal
solution gives \(E^{(p)}(q)=E^{(1)}(pq)\), i.e.\
\(e^{(p)}_n=p^{2n}e^{(1)}_n\), which is \cref{p2:eq:multi-scaling}.  The
remainder in \cref{p2:eq:multi-inverse} follows from
\cref{p0:thm:backward-error,p0:thm:remainder-transport} exactly as in
\cref{p2:thm:branch-inverse}, the slope now being
\(\tfrac1p\log X=\Lambda/p\).
\end{proof}

\begin{remark}
\Cref{p2:eq:multi-scaling} explains and replaces the ``scaling from the
half-shifted inverse Gamma problem'' of \textsf{S1} (\(p_n=4^nd_{2n-1}\)) and
the identity \(d^{!!}_{2n-1}=4^nd^{\Gamma}_{2n-1}\) of \textsf{S2}: both are
the case \((p,1)=(2,1)\).  It was verified symbolically for \(n\le4\) and
\(p\in\{2,3\}\).  As a coefficient-polynomial identity it says nothing about
the numerical Lambert data, which differ between the two problems; the
identity is in \(\Lambda\), not in \(y\).  It also supplies a strong regression
test: an implementation storing only \(d^{(1)}\) reproduces every
multifactorial.
\end{remark}

\begin{center}
\small
\begin{tabular}{cccccc}
\toprule
\(p\) & object & centre \(u\) & phase & \(a^{(p)}_1\) & Borel poles\\
\midrule
1 & \(\Gamma(x+1)\) & \(x+\tfrac12\) & \(u(\log u-1)\) & \(-\tfrac1{24}\)
  & \(\pm2\pi\ii m\)\\
2 & \(n!!\) & \(x+1\) & \(\tfrac u2(\log u-1)\) & \(-\tfrac16\)
  & \(\pm\pi\ii m\)\\
3 & \(n!!!\) & \(x+\tfrac32\) & \(\tfrac u3(\log u-1)\) & \(-\tfrac38\)
  & \(\pm\tfrac{2\pi}3\ii m\)\\
\(p\) & \(x!_{(p)}\) & \(x+\tfrac p2\) & \(\tfrac up(\log u-1)\)
  & \(-\tfrac{p^2}{24}\) & \(\pm\tfrac{2\pi}p\ii m\)\\
\bottomrule
\end{tabular}
\end{center}

\section{Algorithms and certification}
\label{p2:sec:algorithms}

\begin{algorithm}[Forward coefficient table]
\label{p2:alg:forward}
Given \(M\) and a power \(\alpha\in\Rat\): compute \(c_k\) from
\cref{p2:eq:ck} for \(1\le k\le\lceil M/2\rceil\); set \(\gamma_{2k-1}=c_k\),
\(\gamma_{2k}=0\); run \cref{p2:eq:salpha-recurrence}.  This costs
\(\bigO(M^2)\) rational operations and never forms a symbolic exponential.
\end{algorithm}

\begin{algorithm}[Inverse blocks]
\label{p2:alg:inverse}
On the variable \(z=q^2\), with \(\Lambda\) symbolic:
\begin{lstlisting}
input: N
E := 0
for n := 1,...,N:
    H := (1+E)*log(1+E) - E
    for k := 1,...,n:
        H := H + a[k]*z^k*(1+E)^(-(2*k-1))
    H := truncate(H, z^(n+1))
    e[n] := - coefficient(H, z^n) / Lambda
    E := E + e[n]*z^n
    assert truncate(residual(E), z^(n+1)) = 0      # (2.residual)
    assert sign_invariant(e[n], n)                 # Theorem sign law
return e[1],...,e[N]
\end{lstlisting}
Every product, logarithm and negative power may be truncated above \(z^n\) at
step \(n\); no full symbolic logarithm is ever needed.  Representing each
\(e_n\) as a dense polynomial in \(\Lambda^{-1}\) of degree exactly \(2n-1\)
(\cref{p2:thm:sign-law}) gives an exact allocation size.
\end{algorithm}

The two assertions are worth more than a comparison against a stored table.
The first is the \emph{formal residual certificate}: with
\(E_N=\sum_{n\le N}e_nz^n\),
\begin{equation}
 \Lambda E_N+\varphi(E_N)
 +\sum_{k=1}^{N+1}a_kz^k(1+E_N)^{-(2k-1)}
 =\bigO\!\left(z^{N+1}\right),
 \label{p2:eq:residual-certificate}
\end{equation}
and the coefficient of \(z^{N+1}\), divided by \(-\Lambda\), \emph{is} the next
block.  This single identity catches sign, indexing and centring mistakes at
once.  The second is \cref{p2:thm:sign-law}: after multiplication by
\((-1)^{n+1}\) every one of the \(2n-1\) coefficients must be a positive
rational.  Independent checks available in addition:
\begin{enumerate}[label=\textup{(\arabic*)}]
\item expanding the exact gamma ratio of \cref{p2:thm:normal-form} reproduces
      \(c_k\);
\item expanding the Borel kernel \cref{p2:eq:mittag} reproduces \(c_k\) from
      \(\eta\)-values, and \cref{p2:prop:uncentred} reproduces \(b_k\) from
      \(\zeta\)-values;
\item the centre-shift identity \cref{p2:eq:centre-shift} links the two;
\item \cref{p2:eq:closed} and \cref{p2:eq:lagrange-closed} must agree with
      \cref{p2:eq:recurrence} coefficient by coefficient;
\item \cref{p2:eq:multi-scaling} compares every block with the independently
      generated \(p=1\) family.
\end{enumerate}
Every displayed block \(d_1,\dots,d_{13}\) of this chapter passes all five.

\begin{algorithm}[Stable evaluation from a logarithmic target]
\label{p2:alg:evaluate}
Never form \(y\); accept \(\ell=\log y\).  For branch \(\delta\):
\begin{equation}
 R=\ell-A_\delta,
 \quad
 w=\Wlam_0(2R/\e),
 \quad
 T=2R/w,
 \quad
 \Lambda=w+1,
 \quad
 x\leftarrow T-1+\sum_{n=1}^{N}\frac{d_{2n-1}(\Lambda)}{T^{2n-1}} .
 \label{p2:eq:stable-eval}
\end{equation}
Refine by Newton on the logarithmic residual
\(F_\delta(x)=\log\kappa_\delta+\tfrac x2\log2+\log\Gamma(\tfrac x2+1)-\ell\),
\begin{equation}
 x_{\mathrm{new}}=x-\frac{F_\delta(x)}
 {\tfrac12\bigl(\log2+\psi(\tfrac x2+1)\bigr)},
 \label{p2:eq:newton}
\end{equation}
whose denominator is positive by \cref{p2:eq:branch-logder}; convergence is
quadratic once the asymptotic value is in the basin, which a few blocks always
achieve in the ranges tested.  For \(\mathscr D\), add
\(\tfrac14(1-\cos\pi x)\log(2/\pi)\) to \(F_0\) and its derivative to the
denominator, and start from \cref{p2:eq:oeis-inverse}.
\end{algorithm}

\begin{proposition}[A posteriori certificate]
\label{p2:prop:aposteriori}
Let \(\widetilde x\) be any approximation and \(\rho=F_\delta(\widetilde x)\).
If \(I\) is an interval containing \(\widetilde x\) and
\(\mathcal D_\delta^{-1}(y)\), then
\begin{equation}
 \bigl|\widetilde x-\mathcal D_\delta^{-1}(y)\bigr|
 \le\frac{2|\rho|}{\inf_{\xi\in I}\bigl(\log2+\psi(\xi/2+1)\bigr)} ,
 \label{p2:eq:aposteriori}
\end{equation}
and since \(\psi\) is increasing the infimum is attained at the left endpoint.
At large arguments the denominator is \(\sim\Lambda\), matching
\cref{p2:eq:branch-inverse}.
\end{proposition}

\begin{proof}
This is \cref{p0:thm:backward-error} with the explicit slope
\cref{p2:eq:branch-logder}; the mean-value theorem applied to \(F_\delta\) on
\(I\) gives \cref{p2:eq:aposteriori}.
\end{proof}

\section{Numerical validation}
\label{p2:sec:numerics}

All tables in this section were recomputed at \(220\)-digit precision with
\texttt{mpmath}: the exact target \(\ell=\log\mathcal D_\delta(n)\) is formed
from \(\log\Gamma\), the Lambert data from \cref{p2:eq:stable-eval}, and the
blocks from \cref{p2:eq:recurrence} in exact rational arithmetic.  No rounded
value of \(y\) is ever formed.  Entries are \emph{signed}
(approximation minus exact), which the source tables partly suppress.

\begin{table}[htbp]
\centering\small
\caption{Forward: centred logarithmic expansion, \(\Phi(u)+A_\delta+
\sum_{k\le N}c_ku^{1-2k}-\log\mathcal D_\delta(n)\), \(u=n+1\).  Signs
alternate, confirming \cref{p2:thm:remainder}.}
\label{p2:tab:fwd-log}
\begin{tabular}{c c r r r r r}
\toprule
\(\delta\)&\(n\)&\(N=0\)&\(N=1\)&\(N=2\)&\(N=3\)&\(N=4\)\\
\midrule
0&20 & \(3.96616\cdot10^{-3}\) & \(-2.09362\cdot10^{-6}\)
    & \(5.98269\cdot10^{-9}\) & \(-4.14413\cdot10^{-11}\)
    & \(5.30661\cdot10^{-13}\)\\
1&21 & \(3.78606\cdot10^{-3}\) & \(-1.82137\cdot10^{-6}\)
    & \(4.74399\cdot10^{-9}\) & \(-2.99567\cdot10^{-11}\)
    & \(3.49749\cdot10^{-13}\)\\
0&100& \(8.25064\cdot10^{-4}\) & \(-1.88702\cdot10^{-8}\)
    & \(2.34020\cdot10^{-12}\) & \(-7.04697\cdot10^{-16}\)
    & \(3.92938\cdot10^{-19}\)\\
1&101& \(8.16975\cdot10^{-4}\) & \(-1.83207\cdot10^{-8}\)
    & \(2.22773\cdot10^{-12}\) & \(-6.57742\cdot10^{-16}\)
    & \(3.59602\cdot10^{-19}\)\\
\bottomrule
\end{tabular}
\end{table}

\begin{table}[htbp]
\centering\small
\caption{Forward: relative error \(|F_M/\mathcal D_\delta-1|\) of the centred
multiplicative expansion \cref{p2:eq:mult-transseries} truncated after
\(s_Mu^{-M}\), with \(\delta\equiv x\pmod2\).}
\label{p2:tab:fwd-mult}
\begin{tabular}{r c r r r r r}
\toprule
\(x\)&\(\delta\)&\(M=0\)&\(M=1\)&\(M=3\)&\(M=5\)&\(M=9\)\\
\midrule
10&0& \(7.590\cdot10^{-3}\) & \(4.330\cdot10^{-5}\) & \(2.599\cdot10^{-7}\)
   & \(4.974\cdot10^{-9}\) & \(1.388\cdot10^{-11}\)\\
11&1& \(6.957\cdot10^{-3}\) & \(3.538\cdot10^{-5}\) & \(1.751\cdot10^{-7}\)
   & \(2.782\cdot10^{-9}\) & \(5.454\cdot10^{-12}\)\\
20&0& \(3.974\cdot10^{-3}\) & \(9.988\cdot10^{-6}\) & \(1.432\cdot10^{-8}\)
   & \(6.756\cdot10^{-11}\) & \(1.314\cdot10^{-14}\)\\
21&1& \(3.793\cdot10^{-3}\) & \(9.014\cdot10^{-6}\) & \(1.166\cdot10^{-8}\)
   & \(4.972\cdot10^{-11}\) & \(7.956\cdot10^{-15}\)\\
50&0& \(1.635\cdot10^{-3}\) & \(1.483\cdot10^{-6}\) & \(3.106\cdot10^{-10}\)
   & \(2.113\cdot10^{-13}\) & \(9.594\cdot10^{-19}\)\\
51&1& \(1.604\cdot10^{-3}\) & \(1.424\cdot10^{-6}\) & \(2.861\cdot10^{-10}\)
   & \(1.866\cdot10^{-13}\) & \(7.800\cdot10^{-19}\)\\
\bottomrule
\end{tabular}
\end{table}

\begin{table}[htbp]
\centering\small
\caption{Branchwise inverse: \(x_N-x\) with \(y=\mathcal D_\delta(x)\) exact
and \(x_N\) the truncation of \cref{p2:eq:branch-inverse} after \(N\) blocks.
Signs alternate, as \cref{p2:thm:sign-law} predicts for the dominant term.}
\label{p2:tab:inv}
\begin{tabular}{c r r r r r r}
\toprule
\(\delta\)&\(x\)&\(N=0\)&\(N=1\)&\(N=2\)&\(N=3\)&\(N=5\)\\
\midrule
0& 10 & \(-6.3074\cdot10^{-3}\) & \(1.6444\cdot10^{-5}\)
     & \(-1.6097\cdot10^{-7}\) & \(3.6890\cdot10^{-9}\)
     & \(1.1386\cdot10^{-11}\)\\
1& 11 & \(-5.5808\cdot10^{-3}\) & \(1.2104\cdot10^{-5}\)
     & \(-9.9815\cdot10^{-8}\) & \(1.9338\cdot10^{-9}\)
     & \(4.2632\cdot10^{-12}\)\\
0& 20 & \(-2.6055\cdot10^{-3}\) & \(1.7520\cdot10^{-6}\)
     & \(-4.7751\cdot10^{-9}\) & \(3.1029\cdot10^{-11}\)
     & \(7.6506\cdot10^{-15}\)\\
1& 21 & \(-2.4497\cdot10^{-3}\) & \(1.4956\cdot10^{-6}\)
     & \(-3.7170\cdot10^{-9}\) & \(2.2041\cdot10^{-11}\)
     & \(4.5233\cdot10^{-15}\)\\
0& 50 & \(-8.3109\cdot10^{-4}\) & \(8.9794\cdot10^{-8}\)
     & \(-4.1941\cdot10^{-11}\) & \(4.7156\cdot10^{-14}\)
     & \(3.4429\cdot10^{-19}\)\\
1& 51 & \(-8.1110\cdot10^{-4}\) & \(8.4219\cdot10^{-8}\)
     & \(-3.7845\cdot10^{-11}\) & \(4.0942\cdot10^{-14}\)
     & \(2.7669\cdot10^{-19}\)\\
0&100 & \(-3.5755\cdot10^{-4}\) & \(9.5805\cdot10^{-9}\)
     & \(-1.1469\cdot10^{-12}\) & \(3.3159\cdot10^{-16}\)
     & \(1.5896\cdot10^{-22}\)\\
1&101 & \(-3.5329\cdot10^{-4}\) & \(9.2785\cdot10^{-9}\)
     & \(-1.0892\cdot10^{-12}\) & \(3.0878\cdot10^{-16}\)
     & \(1.4232\cdot10^{-22}\)\\
\bottomrule
\end{tabular}
\end{table}

\begin{table}[htbp]
\centering\small
\caption{Inverse of the periodic interpolation:
\(\bigl(T-1+\widetilde\Delta\bigr)-x\) with \(y=\mathscr D(x)\) exact,
successively adding the carrier \(\Delta_0\) (computed as the exact fixed point
of \cref{p2:eq:carrier-eq}), \(\Delta_1/T\) and \(\Delta_2/T^2\) from
\cref{p2:eq:Delta12}.  The core error matches the predicted carrier amplitude
\(2a/\Lambda\): at \(x=20\), \(2a/\Lambda=0.074080\) against an observed
\(0.071518\); at \(x=100\), \(0.048919\) against \(0.048564\).  No source
validates this sector numerically.}
\label{p2:tab:oeis-inv}
\begin{tabular}{r r r r r}
\toprule
\(x\)&core \(T-1\)&\(+\Delta_0\)&\(+\Delta_1/T\)&\(+\Delta_2/T^2\)\\
\midrule
20     & \(7.1518\cdot10^{-2}\) & \(-2.5603\cdot10^{-3}\)
       & \(-1.2693\cdot10^{-5}\) & \(1.6637\cdot10^{-6}\)\\
21     & \(-7.5539\cdot10^{-2}\) & \(-2.4124\cdot10^{-3}\)
       & \(1.4359\cdot10^{-5}\) & \(1.4221\cdot10^{-6}\)\\
20.5   & \(-2.5254\cdot10^{-3}\) & \(-3.2850\cdot10^{-3}\)
       & \(2.0627\cdot10^{-6}\) & \(2.2808\cdot10^{-6}\)\\
100    & \(4.8564\cdot10^{-2}\) & \(-3.5495\cdot10^{-4}\)
       & \(-2.3071\cdot10^{-7}\) & \(9.4103\cdot10^{-9}\)\\
100.25 & \(3.4218\cdot10^{-2}\) & \(-3.9845\cdot10^{-4}\)
       & \(-1.8912\cdot10^{-7}\) & \(1.0779\cdot10^{-8}\)\\
\bottomrule
\end{tabular}
\end{table}

\begin{table}[htbp]
\centering\small
\caption{Multifactorial inverse \cref{p2:eq:multi-inverse}, using only the
\(p=1\) blocks rescaled by \cref{p2:eq:multi-scaling}.  Entries are
\(x_N-x\) with \(y=\mathcal D_{p,r}(x)\) exact; the \(p=2\) row reproduces
\cref{p2:tab:inv}.}
\label{p2:tab:multi-inv}
\begin{tabular}{c c r r r r r}
\toprule
\(p\)&\(r\)&\(x\)&\(N=0\)&\(N=1\)&\(N=2\)&\(N=3\)\\
\midrule
2&2& 20 & \(-2.60549\cdot10^{-3}\) & \(1.75197\cdot10^{-6}\)
       & \(-4.77514\cdot10^{-9}\) & \(3.10287\cdot10^{-11}\)\\
3&1& 22 & \(-5.05001\cdot10^{-3}\) & \(6.03949\cdot10^{-6}\)
       & \(-2.95265\cdot10^{-8}\) & \(3.44131\cdot10^{-10}\)\\
3&2& 20 & \(-5.67882\cdot10^{-3}\) & \(8.16471\cdot10^{-6}\)
       & \(-4.75923\cdot10^{-8}\) & \(6.60178\cdot10^{-10}\)\\
4&3& 23 & \(-8.27262\cdot10^{-3}\) & \(1.54424\cdot10^{-5}\)
       & \(-1.18248\cdot10^{-7}\) & \(2.15627\cdot10^{-9}\)\\
5&2& 27 & \(-1.04166\cdot10^{-2}\) & \(2.15596\cdot10^{-5}\)
       & \(-1.85311\cdot10^{-7}\) & \(3.79755\cdot10^{-9}\)\\
\bottomrule
\end{tabular}
\end{table}

\begin{remark}
The near-geometric improvement at these moderate arguments is not a claim of
convergence: at fixed \(x\) the Bernoulli-driven expansion diverges, and by
\cref{p2:prop:least-term} the attainable accuracy is bounded below by
\(\asymp\e^{-\pi T}/(\sqrt T\Lambda)\).  In high-precision work, truncate at
the least term or finish with \cref{p2:eq:newton}.  Note also that the
parity dependence is confined to the core: once the correct \(A_\delta\) is
used, the two subsequences show nearly identical error patterns, as
\cref{p2:rem:kappa-cancels} predicts.
\end{remark}

\section{Logical status}
\label{p2:sec:status}

Three kinds of statement have been used, and the chapter keeps them apart.

\paragraph{Exact identities.}
\Cref{p2:eq:branches,p2:eq:oeis,p2:eq:multi-branch} and the interpolation
statements \cref{p2:prop:branch-interp,p2:prop:oeis-interp,p2:prop:multi-branch};
the centred normal forms \cref{p2:eq:normal-form,p2:eq:sequence-exact,p2:eq:oeis-centred};
the Lambert identities \cref{p2:eq:core-identities}; the Borel--Laplace
representations \cref{p2:eq:borel,p2:eq:multi-borel} and the Mittag--Leffler
decomposition \cref{p2:eq:mittag}.  These hold wherever the functions are
defined.

\paragraph{Formal coefficient identities.}
The Bell formulae \cref{p2:eq:salpha-bell,p2:eq:bell-recurrence}, the
recurrences \cref{p2:eq:salpha-recurrence,p2:eq:recurrence,p2:eq:periodic-recurrence},
the closed formulae \cref{p2:eq:closed,p2:eq:lagrange-closed,p2:eq:Dnr}, the
boundary and sign statements \cref{p2:cor:boundary,p2:thm:sign-law}, the
centre-shift identity \cref{p2:eq:centre-shift}, the scaling
\cref{p2:eq:multi-scaling}, and \cref{p2:prop:phase-reversion}.  These are
identities of formal series over \(\Rat(\Lambda)\), or over the ring generated
by \(\Lambda^{-1}\) and \(\e^{\pm\ii\theta}\) in the periodic case.  No
convergence of the Bernoulli series is asserted anywhere.

\paragraph{Poincar\'e asymptotics with remainders.}
\Cref{p2:thm:coefficients,p2:thm:remainder,p2:cor:forward-log} for the forward
direction --- where the remainder is in fact two-sided and sharp;
\cref{p2:thm:branch-inverse,p2:cor:periodic-meaning,p2:eq:multi-inverse} for
the inverse, each obtained from a forward residual by
\cref{p0:thm:backward-error,p0:thm:remainder-transport} together with the
explicit slope \cref{p2:eq:branch-logder}; and
\cref{p2:prop:least-term,p2:prop:transport} for the beyond-all-orders layer.
\Cref{p2:thm:carrier}(b) is a genuine convergence statement, with an explicit
threshold.

\begin{remark}[Practical summary]
\label{p2:rem:summary}
For symbolic work, centre at \(u=x+1\), keep \(\Wlam_0\) exact, generate
\(c_k\) from Bernoulli numbers, exponentiate with complete Bell polynomials,
and invert with the triangular \(z=T^{-2}\) recurrence, checking
\cref{p2:eq:residual-certificate} and the sign law at every step.  Expand
\(\Wlam_0\) into iterated logarithms only as a final presentation layer.  When
the parity of \(n\) is known, use the branch inverse; when it is not, state the
interpolation --- by \cref{p2:thm:no-canonical} its choice changes the leading
sector of the answer.
\end{remark}

%% ==================== fragment p3 ====================
\chapter{The partition numbers (A000041)}\label{p3:sec:top}

\begin{sourcebox}
This chapter merges three independent treatments of the partition numbers filed in the
\texttt{special-function-inversion} collection:
\emph{Partition\_Number\_Transseries\_and\_Asymptotic\_Inverse} (below: \textbf{S1}),
\emph{Partition\_Number\_Transseries\_and\_Inverse} (\textbf{S2}), and
\emph{Partition\_Numbers\_Transseries\_and\_Inverse} (\textbf{S3}). All three take Rademacher's
convergent series as analytic input and build a coefficient calculus on top of it; they differ in
notation, in how the non-integral arithmetic phase is defined, and in which tail estimate is
proved. Everything below that is not attributed is common to all three, restated once. Numerical
constants, coefficient tables and error tables have been recomputed here (exact rationals through
\texttt{fractions}/\texttt{sympy}; $140$-digit \texttt{mpmath} for the tables); the recomputation
is described in \cref{p3:sub:numerics}.
\end{sourcebox}

\section{Orientation: what is specific to \texorpdfstring{$p(n)$}{p(n)}}
\label{p3:sec:orientation}

The partition function $p(n)$ is defined by Euler's product
\begin{equation}\label{p3:eq:euler-product}
 \sum_{n\ge0}p(n)q^{n}=\prod_{m\ge1}\frac1{1-q^{m}},\qquad |q|<1 ,
\end{equation}
and is OEIS A000041. Hardy and Ramanujan proved
\begin{equation}\label{p3:eq:HR}
 p(n)\sim\frac1{4\sqrt3\,n}\exp\!\left(\pi\sqrt{\frac{2n}{3}}\right),
\end{equation}
and Rademacher replaced the asymptotic formula by an exactly convergent series. That single fact
makes this chapter structurally different from \cref{p1:sec:top,p2:sec:top,p4:sec:top}: the
exponentially small sectors are \emph{given}, not reconstructed from the late behaviour of a
divergent perturbative series. Three features are genuinely specific to $p(n)$ and are the subject
of this chapter.

\begin{enumerate}
\item \textbf{An arithmetic sector index.}  The sectors are labelled by a
positive integer $k$ (a root of unity $\e^{2\pi\ii h/k}$) and a sign $\sigma$, and their
amplitudes are the Rademacher sums $A_k(n)$ built from Dedekind sums. These amplitudes are
\emph{periodic} in $n$ with period $k$; they are not constants, and after inversion they become
bounded oscillations whose phase is quadratic in $\log y$.
\item \textbf{Non-well-ordered actions.}  Relative to the dominant sector, the
exponential actions are $1-1/k$ (increasing to $1$) and $1+1/k$ (decreasing to $1$). The second
family has no least element. The support of the completed object is therefore not a locally finite
Hahn grid, and the two signs may not be separated before summing over $k$.
\item \textbf{Convergence where one expects divergence.}  Every fixed sector's
algebraic series in $n^{-1/2}$ has radius of convergence exactly $\sqrt{24}$ in the variable
$n^{-1/2}$, and the reversion series of the dominant sector converges as well. Nothing in this
chapter is a divergent asymptotic series; consequently the least-term machinery of
\cref{p0:thm:optimal-truncation} degenerates here, as recorded in \cref{p3:rem:no-least-term}.
\end{enumerate}

Everything else --- the Lambert core, the reversion around it, the flattening into logarithms, the
passage from a smooth inverse to an integer staircase, and the transport of remainders --- is
\cref{p0:sec:top} material and is \emph{cited}, never re-proved.

\subsection{The exponential--power reduction and the Part~0 dictionary}
\label{p3:sub:dictionary}

Fix throughout
\begin{equation}\label{p3:eq:master-variables}
 N=n-\frac1{24},\qquad
 \lambda=\pi\sqrt{\frac23}=\frac{\pi\sqrt6}{3},\qquad
 \mu=\lambda\sqrt N=\frac\pi6\sqrt{24n-1},\qquad
 \kappa=\frac{\pi^2}{6\sqrt3}.
\end{equation}
The dominant Rademacher sector, derived in \cref{p3:thm:exact-sectors}, is
\begin{equation}\label{p3:eq:dominant-in-N}
 P_{1,+}(n)=\frac{\e^{\lambda\sqrt N}}{4\sqrt3\,N}
 \left(1-\frac1{\lambda\sqrt N}\right).
\end{equation}
Read as a function of $N$, \cref{p3:eq:dominant-in-N} is exactly an instance of the
exponential--power model \cref{p0:def:model}, $F(x)=C\exp(ax^{\alpha})x^{b}\exp Q(x^{-\delta})$,
with
\begin{equation}\label{p3:eq:model-in-N}
 x=N,\quad C=\frac1{4\sqrt3},\quad a=\lambda,\quad \alpha=\frac12,\quad
 b=-1,\quad \delta=\frac12,\quad
 Q(t)=\log\!\left(1-\frac{t}{\lambda}\right)
      =-\sum_{r\ge1}\frac{t^{r}}{r\lambda^{r}} .
\end{equation}
So the partition problem is the case $\alpha=1/2$: a \emph{square-root} phase, not a linear one.
It is reduced to the linear case in two steps: the monomial substitution
$\xi=x^{\alpha}=\sqrt N$ of \cref{p0:lem:alpha-reduction}, followed by the elementary linear
rescaling $\mu=\lambda\xi$ that normalises the exponential rate to $1$.

\begin{lemma}[The $\alpha$-reduction for $p(n)$]\label{p3:lem:alpha-reduction}
\emph{(i)} \Cref{p0:lem:alpha-reduction} applied to \cref{p3:eq:model-in-N} with
$\xi=\sqrt N$ gives an $\alpha=1$ model with data
$\bigl(\tfrac1{4\sqrt3},\ \lambda,\ 1,\ -2,\ 1,\ Q\bigr)$, $Q(t)=\log(1-t/\lambda)$; note
$b/\alpha=-2$ and $\delta/\alpha=1$, so this is the normalisation $\delta=1$ assumed from
\cref{p0:sec:top} onwards.  \emph{(ii)} The further substitution $\mu=\lambda\xi$ rescales any
$\alpha=1$ model with data $(C,a,1,b,1,Q)$ to one with data
$\bigl(Ca^{-b},\,1,\,1,\,b,\,1,\,Q(a\,\cdot)\bigr)$.  Here the two steps give
\begin{equation}\label{p3:eq:model-in-mu}
 P_{1,+}=\kappa\,\frac{\e^{\mu}}{\mu^{2}}\left(1-\frac1\mu\right)
 =C'\e^{a'\mu}\mu^{b'}\exp Q'(\mu^{-\delta'}),
\end{equation}
with the dictionary
\[
C'=\kappa=\frac{\pi^{2}}{6\sqrt3},\qquad a'=1,\qquad \alpha'=1,\qquad b'=-2,\qquad
\delta'=1,\qquad Q'(t)=\log(1-t).
\]
The inverse map back to the index is exact and quadratic:
\begin{equation}\label{p3:eq:n-of-mu}
 n=\frac1{24}+\frac{3\mu^{2}}{2\pi^{2}} .
\end{equation}
\end{lemma}

\begin{proof}
(i) is \cref{p0:lem:alpha-reduction} at $\alpha=\delta=1/2$, $b=-1$, which leaves $C$, $a$ and
the coefficients of $Q$ unchanged and replaces $(b,\delta)$ by $(b/\alpha,\delta/\alpha)$.
For (ii), substituting $\xi=\mu/a$ in $C\e^{a\xi}\xi^{b}\exp Q(\xi^{-1})$ gives
$Ca^{-b}\e^{\mu}\mu^{b}\exp Q(a/\mu)$, and $t\mapsto Q(at)$ again lies in $t\,\Real[[t]]$.
Here $a=\lambda$, $b=-2$, so $Ca^{-b}=\lambda^{2}/(4\sqrt3)=\pi^{2}/(6\sqrt3)=\kappa$ (using
$\lambda^{2}=2\pi^{2}/3$) and $Q'(t)=Q(\lambda t)=\log(1-t)$.  Composing, $\mu=\lambda\sqrt N$
and \cref{p3:eq:dominant-in-N} become \cref{p3:eq:model-in-mu}.  Inverting
$\mu=\lambda\sqrt{n-1/24}$ gives \cref{p3:eq:n-of-mu}, again by $\lambda^{2}=2\pi^{2}/3$.
\end{proof}

\begin{remark}[The model hypotheses hold exactly, and $Q$ converges]
\label{p3:rem:model-exact}
\Cref{p0:def:model} asks for the analytic estimates \cref{p0:eq:model-analytic} and
\cref{p0:eq:model-derivative} in addition to the formal shape.  For the partition envelope both
hold \emph{exactly and trivially}: $F(\mu)=\kappa\e^{\mu}\mu^{-2}(1-1/\mu)$ is an elementary
closed form, so $\log F(\mu)-\log\kappa-\mu+2\log\mu=\log(1-1/\mu)$ is literally the sum of the
series $Q'$, and $(\log F)'=\Phi'$ is computed in closed form in
\cref{p3:lem:Phi-derivatives}; no asymptotic estimate is being differentiated.  In particular
$Q'(t)=\log(1-t)$ is \emph{convergent}, with radius $1$.  \Cref{p0:rem:formal-analytic} states
that $Q$ is divergent of Gevrey type in all four chapters of this volume; that is not so here,
and the discrepancy should be read as a further instance of
\cref{p3:rem:no-least-term} --- the partition chapter is the one place in the volume where the
whole apparatus applies to a convergent object.
\end{remark}

\begin{remark}[Why this is the point of the whole construction]
\label{p3:rem:reduction-is-the-point}
After \cref{p3:lem:alpha-reduction} nothing in the inversion of the dominant sector is
partition-specific. The three sources each rediscover, in their own notation, the Lambert core of
\cref{p0:thm:lambert-core}, the centred reversion of \cref{p0:thm:lambert-centered} and the
logarithmic flattening of \cref{p0:thm:flattening}. We record the dictionary once
(\cref{p3:tab:dictionary}) and cite. What is \emph{not} generic --- and occupies the rest of the
chapter --- is the arithmetic: the sector index, the Dedekind-sum amplitudes, the accumulation of
actions at $1$, and the interaction between the periodic amplitudes and the $\bigO(1)$ shift
produced by the reversion.
\end{remark}

\begin{table}[htbp]
\centering
\caption{Parameter dictionary from \cref{p0:sec:top} to the partition
chapter.  All Part~0 statements are used at these values.}
\label{p3:tab:dictionary}
\small
\begin{tabular}{@{}lll@{}}
\toprule
Part~0 object & Partition value & Reference \\
\midrule
model variable $x$ & $\mu=\pi\sqrt{2(n-1/24)/3}$ & \cref{p3:eq:master-variables}\\
$C$ & $\kappa=\pi^{2}/(6\sqrt3)$ & \cref{p3:lem:alpha-reduction}\\
$a,\ \alpha$ & $1,\ 1$ (after reduction and rescaling; $\lambda,\ 1/2$ before) & \cref{p3:lem:alpha-reduction}\\
$b$ & $-2$ & \cref{p3:lem:alpha-reduction}\\
$Q(t),\ \delta$ & $\log(1-t),\ 1$ & \cref{p3:lem:alpha-reduction}\\
Lambert core $X$ & $-2\Wlam_{-1}\!\left(-\tfrac12\sqrt{\kappa/y}\right)$ & \cref{p3:prop:lambert-core}\\
$\Psi(t,u)$ of \cref{p0:thm:lambert-centered} & $3\log(1+tu)-\log(1-t+tu)$ & \cref{p3:thm:core-correction}\\
reversion coefficients $c_m$ & $1,\tfrac52,\tfrac{13}3,\tfrac{53}{12},-\tfrac{14}5,\dots$ & \cref{p3:tab:cm}\\
index map $x\mapsto n$ & $n=\tfrac1{24}+3\mu^{2}/(2\pi^{2})$ & \cref{p3:eq:n-of-mu}\\
staircase separation & unit spacing of $\Int$; see \cref{p3:thm:rounding} & \cref{p0:thm:staircase}\\
\bottomrule
\end{tabular}
\end{table}

\begin{remark}[Notation across the three sources]\label{p3:rem:notation-clash}
The sources disagree on symbols in a way that has to be stated once, because two of the clashes
are dangerous.
\begin{enumerate}
\item \textbf{The letter $x$.}  S1 writes $N=n-1/24$ for the shifted index and
$x$ for the phase $\pi\sqrt{2N/3}$; S2 writes $x=n-1/24$ for the shifted index and $\mu$ for the
phase; S3 writes $N$ and $\mu$. This volume uses $N$ for the shifted index and $\mu$ for the
phase, i.e.\ the S3 convention, which is the intersection of the other two.
\item \textbf{The shift itself.}  Every displayed \emph{exact} identity in this
chapter is in the shifted variable $N=n-1/24$. Expansions in $n^{-1/2}$
(\cref{p3:sec:half-powers}) re-expand the shift and are the only place where unshifted $n$ appears
inside a coefficient. Confusing the two is the single most common error in this subject: e.g.\
replacing $N$ by $n$ in \cref{p3:eq:dominant-in-N} costs a relative $\bigO(n^{-1/2})$ and
destroys every statement about exponentially small sectors.
\item \textbf{The normalising constant} is $\mathcal{C}$ in S1, $\kappa$ in S2, $D$
in S3; all equal $\pi^{2}/(6\sqrt3)$. We write $\kappa$, reserving $D$ for nothing and $\Phi'$ for
the slope.
\item \textbf{The logarithmic variable.}  S1 and S2 expand in
$R=\log(y/\kappa)$; S3 expands in $H=\log(4y/\kappa)=R+2\log2$. The two expansions have the
\emph{same} coefficient polynomials (see \cref{p3:rem:H-versus-R}) but are different series and
must not be mixed.
\end{enumerate}
\end{remark}

\section{Dedekind sums and the arithmetic amplitudes}\label{p3:sec:arithmetic}

\subsection{Definitions}

\begin{definition}[Sawtooth and Dedekind sum]\label{p3:def:dedekind}
For $u\in\Real$ put
\begin{equation}\label{p3:eq:sawtooth}
 ((u))=\begin{cases} u-\lfloor u\rfloor-\tfrac12,& u\notin\Int,\\ 0,& u\in\Int,\end{cases}
\end{equation}
and for $k\ge1$, $h\in\Int$,
\begin{equation}\label{p3:eq:dedekind}
 s(h,k)=\sum_{r=1}^{k-1}\left(\!\left(\frac rk\right)\!\right)
 \left(\!\left(\frac{hr}{k}\right)\!\right),
\end{equation}
with the empty-sum convention $s(h,1)=0$.
\end{definition}

Throughout, $U_{k}=\{h:0\le h<k,\ \gcd(h,k)=1\}$, so that $U_{1}=\{0\}$.

\begin{definition}[Rademacher amplitude]\label{p3:def:Ak}
For $k\ge1$ and $n\in\Int$,
\begin{equation}\label{p3:eq:Ak}
 A_{k}(n)=\sum_{h\in U_{k}}
 \exp\!\left(\pi\ii\,s(h,k)-\frac{2\pi\ii hn}{k}\right).
\end{equation}
\end{definition}

\begin{lemma}[Elementary properties]\label{p3:lem:Ak-properties}
For every $k\ge1$ and $n\in\Int$:
\begin{enumerate}
\item $A_{k}(n)\in\Real$, and
$A_{k}(n)=\sum_{h\in U_{k}}\cos\!\bigl(\pi s(h,k)-2\pi hn/k\bigr)$;
\item $A_{k}(n+k)=A_{k}(n)$;
\item $|A_{k}(n)|\le\varphi(k)\le k$;
\item $A_{1}(n)=1$ and $A_{2}(n)=(-1)^{n}$.
\end{enumerate}
\end{lemma}

\begin{proof}
(i) For $k\ge2$ the involution $h\mapsto k-h$ is a bijection of $U_{k}$ (note $k-h\ne h$ unless
$k=2h$, impossible for $\gcd(h,k)=1$, $k\ge3$; for $k=2$ the set is $\{1\}$ and $s(1,2)=0$ makes
the single term real). The reciprocity $s(-h,k)=-s(h,k)$, immediate from $((-u))=-((u))$, shows
that the terms at $h$ and $k-h$ are complex conjugates. Hence the sum is real and equals the sum
of the real parts, which is the stated cosine sum. (ii) Replacing $n$ by $n+k$ changes each
exponent by $-2\pi\ii h$, an integer multiple of $2\pi\ii$. (iii) Triangle inequality on
$\varphi(k)$ unimodular terms. (iv) $U_{1}=\{0\}$ and $s(0,1)=0$ give $A_{1}\equiv1$;
$U_{2}=\{1\}$, $s(1,2)=0$ give $A_{2}(n)=\e^{-\pi\ii n}=(-1)^{n}$.
\end{proof}

Only the trivial bound (iii) is used anywhere in this chapter. Weil-type bounds for these
Kloosterman-like sums would sharpen the constants in \cref{p3:thm:tail,p3:thm:inverse-tail} but
change no structure; see \cref{p3:rem:sharper-bounds}.

\begin{proposition}[The first amplitudes]\label{p3:prop:first-Ak}
With $s(1,k)=\dfrac{(k-1)(k-2)}{12k}$,
\begin{align}
 A_{1}(n)&=1, &
 A_{2}(n)&=(-1)^{n}, \nonumber\\
 A_{3}(n)&=2\cos\!\left(\frac\pi{18}-\frac{2\pi n}{3}\right), &
 A_{4}(n)&=2\cos\!\left(\frac\pi8-\frac{\pi n}{2}\right),
 \label{p3:eq:first-Ak}\\
 A_{5}(n)&=2\cos\!\left(\frac\pi5-\frac{2\pi n}{5}\right)+2\cos\frac{4\pi n}{5},&
 A_{6}(n)&=2\cos\!\left(\frac{5\pi}{18}-\frac{\pi n}{3}\right).\nonumber
\end{align}
\end{proposition}

\begin{proof}
$s(1,k)=\sum_{r=1}^{k-1}((r/k))^{2}=\sum_{r=1}^{k-1}(r/k-1/2)^{2}
=\frac{(k-1)(2k-1)}{6k}-\frac{k-1}{2}+\frac{k-1}{4} =\frac{(k-1)(k-2)}{12k}$. Hence
$s(1,3)=\frac1{18}$, $s(1,4)=\frac18$, $s(1,5)=\frac15$, $s(1,6)=\frac5{18}$, and
$s(k-1,k)=-s(1,k)$. For $k=3,4,6$ the set $U_{k}$ is $\{1,k-1\}$, and the two terms are conjugate,
giving twice the real part of the $h=1$ term after reducing $2\pi(k-1)n/k\equiv-2\pi n/k
\pmod{2\pi}$. For $k=5$, $U_{5}=\{1,2,3,4\}$ and a direct evaluation of \cref{p3:eq:dedekind}
gives $s(2,5)=s(3,5)=0$: with
$((1/5),(2/5),(3/5),(4/5))=(-\tfrac3{10},-\tfrac1{10},\tfrac1{10},\tfrac3{10})$, the four products
in $s(2,5)$ are $\tfrac3{100},-\tfrac3{100},-\tfrac3{100},\tfrac3{100}$. Pairing $h=1$ with $h=4$
and $h=2$ with $h=3$ gives the stated two cosines.
\end{proof}

\subsection{The real-analytic interpolation of the amplitudes}
\label{p3:sub:interpolation-amplitudes}

The inverse problem needs the amplitudes at non-integral arguments: the smooth inverse of a smooth
interpolation evaluates $A_{k}$ at a real point. The finite Fourier shape of \cref{p3:eq:Ak}
supplies a canonical choice.

\begin{definition}[Entire amplitude interpolation]\label{p3:def:Atilde}
For $k\ge1$ and $\nu\in\Cplx$ set
\begin{equation}\label{p3:eq:Atilde}
 \widetilde A_{k}(\nu)=\sum_{h\in U_{k}}
 \cos\!\left(\pi s(h,k)-\frac{2\pi h\nu}{k}\right).
\end{equation}
\end{definition}

\begin{lemma}\label{p3:lem:Atilde}
$\widetilde A_{k}$ is entire, $k$-periodic, real on $\Real$, satisfies $\widetilde
A_{k}(n)=A_{k}(n)$ for $n\in\Int$, and $|\widetilde
A_{k}^{(r)}(\nu)|\le\varphi(k)(2\pi/k)^{r}\cosh(2\pi|\Im\nu|) \le k\,(2\pi)^{r}\e^{2\pi|\Im\nu|}$
for every $r\ge0$. Moreover $\widetilde A_{1}\equiv1$ and $\widetilde A_{2}(\nu)=\cos\pi\nu$,
while
\begin{equation}\label{p3:eq:Atilde3}
 \widetilde A_{3}(\nu)=2\cos(\pi\nu)\,
 \cos\!\left(\frac\pi{18}+\frac{\pi\nu}{3}\right).
\end{equation}
\end{lemma}

\begin{proof}
Each summand is entire and $k$-periodic in $\nu$; realness on $\Real$ is manifest; agreement at
integers is \cref{p3:lem:Ak-properties}(i). The derivative bound is termwise, using $|h/k|\le1$
and $|\cos^{(r)}(z)|\le\cosh(\Im z)$. For $k=1$, $U_{1}=\{0\}$ makes the single term $\cos0=1$
identically. For $k=3$, sum-to-product on $A=\frac\pi{18}-\frac{2\pi\nu}3$ and
$B=-\frac\pi{18}-\frac{4\pi\nu}3$ gives $\tfrac{A+B}2=-\pi\nu$ and
$\tfrac{A-B}2=\frac\pi{18}+\frac{\pi\nu}3$, whence \cref{p3:eq:Atilde3}.
\end{proof}

\begin{repairbox}
\textbf{S2, Definition ``Real Rademacher amplitude''.} S2 defines the interpolation as a sum over
$1\le h\le k$ with $\gcd(h,k)=1$. For $k\ge2$ this is the same index set as $U_{k}$, but for $k=1$
it produces the single term $\cos(2\pi\nu)$ instead of the constant $1$. S2 nevertheless uses
$\mathcal A_{1}\equiv1$ throughout (its dominant sector carries no amplitude, and its
relative-correction formula separates the $k=1$ terms with amplitude $1$), so the definition
contradicts its own use: with S2's index set the interpolating function $\mathcal{P}$ would not be
an interpolation of $p$ --- its dominant term would oscillate off-lattice with period $1$ at full
relative strength, and the ``eventual monotonicity'' proposition would be false. \textbf{Fix.}
Take $h$ over $U_{k}=\{0\le h<k:\gcd(h,k)=1\}$ with $U_{1}=\{0\}$, i.e.\ \cref{p3:def:Atilde};
equivalently, decree $\widetilde A_{1}\equiv1$ separately, as S3 does. All of S2's later formulas
are then correct as printed.
\end{repairbox}

\begin{remark}[Non-uniqueness, honestly stated]\label{p3:rem:interp-nonunique}
\cref{p3:def:Atilde} is \emph{a} choice, not \emph{the} choice. Any strictly increasing smooth
$\mathcal{P}$ with $\mathcal{P}(n)=p(n)$ has an inverse, and the algebraic (Lambert--logarithmic)
part of that inverse is the same for all of them, because it is determined by the dominant sector,
which is interpolation-independent once $N$ is treated as a real variable. The exponentially small
sectors are \emph{not} interpolation-independent. \cref{p3:def:Atilde} is singled out only because
it is entire, finite-Fourier, real on $\Real$, and literally the analytic continuation of the
amplitudes already present in Rademacher's theorem. \cref{p3:sec:profinite} gives the alternative
that avoids interpolation altogether.
\end{remark}

\section{Rademacher's formula in exact transseries normal form}
\label{p3:sec:exact}

\subsection{The analytic input}

\begin{theorem}[Hardy--Ramanujan--Rademacher; quoted]\label{p3:thm:rademacher}
For every integer $n\ge1$,
\begin{equation}\label{p3:eq:rademacher}
 p(n)=\frac1{\pi\sqrt2}\sum_{k=1}^{\infty}A_{k}(n)\sqrt k\,
 \frac{\dd}{\dd n}
 \left[\frac{\sinh\!\bigl(\tfrac{\lambda}{k}\sqrt{n-\tfrac1{24}}\bigr)}
 {\sqrt{n-\tfrac1{24}}}\right],
\end{equation}
the series converging when summed in increasing $k$.
\end{theorem}

This is the only unproved input of the chapter. It is Rademacher's theorem (1937, 1938, 1943); a
textbook proof is in Apostol, and a machine-checked proof exists in Isabelle/HOL (Eberl--Paulson,
2026); see \cref{p3:rem:references}. All three sources take it as given and so do we. Note the
factor $\sqrt k$: some normalisations absorb it into the definition of the arithmetic sum, and
mixing conventions is a standard source of a spurious factor $k$.
\Cref{p3:sub:normalization-crosscheck} records the cross-checks that detect such a slip.

\subsection{Exact two-exponential sectors}

\begin{lemma}[Elementary derivative]\label{p3:lem:derivative}
For $a>0$ and $v>0$,
\begin{equation}\label{p3:eq:derivative}
 \frac{\dd}{\dd v}\Bigl(v^{-1/2}\sinh(a\sqrt v)\Bigr)
 =\frac1{4v^{3/2}}
 \Bigl[(a\sqrt v-1)\e^{a\sqrt v}+(a\sqrt v+1)\e^{-a\sqrt v}\Bigr].
\end{equation}
\end{lemma}

\begin{proof}
With $w=a\sqrt v$ one has $\dd w/\dd v=w/(2v)$ and $v^{-1/2}\sinh w=a\,w^{-1}\sinh w$, so
\[
\frac{\dd}{\dd v}\bigl(v^{-1/2}\sinh(a\sqrt v)\bigr) =a\,\frac{\dd}{\dd w}\!\left(\frac{\sinh
w}{w}\right)\frac{w}{2v} =\frac{a}{2v}\left(\cosh w-\frac{\sinh w}{w}\right)
=\frac{1}{2v^{3/2}}\bigl(w\cosh w-\sinh w\bigr),
\]
using $a=w/\sqrt v$ in the last step. Finally $2(w\cosh w-\sinh w)=(w-1)\e^{w}+(w+1)\e^{-w}$.
\end{proof}

\begin{theorem}[Exact sector decomposition]\label{p3:thm:exact-sectors}
Let $N,\lambda,\mu,\kappa$ be as in \cref{p3:eq:master-variables} and set, for $k\ge1$ and
$\sigma\in\{+1,-1\}$,
\begin{equation}\label{p3:eq:sector}
 P_{k,\sigma}(n)
 =\frac{A_{k}(n)}{4\sqrt{3k}\;N}\,
 \e^{\sigma\mu/k}\left(1-\frac{\sigma k}{\mu}\right)
 =\frac{\kappa}{\mu^{2}}\frac{A_{k}(n)}{\sqrt k}\,
 \e^{\sigma\mu/k}\left(1-\frac{\sigma k}{\mu}\right).
\end{equation}
Then, for every integer $n\ge1$,
\begin{equation}\label{p3:eq:sector-sum}
 p(n)=\sum_{k\ge1}\bigl(P_{k,+}(n)+P_{k,-}(n)\bigr),
\end{equation}
the two signs being paired before the sum over $k$ is formed. Equivalently, with the entire paired
kernel
\begin{equation}\label{p3:eq:g-kernel}
 g(z)=\cosh z-\frac{\sinh z}{z}
 =\sum_{m\ge1}\frac{2m}{(2m+1)!}z^{2m},\qquad g(0)=0,
\end{equation}
one has the manifestly convergent form
\begin{equation}\label{p3:eq:paired-sector}
 p(n)=\frac{2\kappa}{\mu^{2}}\sum_{k\ge1}\frac{A_{k}(n)}{\sqrt k}\,
 g\!\left(\frac\mu k\right).
\end{equation}
\end{theorem}

\begin{proof}
Apply \cref{p3:lem:derivative} with $v=N$ and $a=\lambda/k$, so that $a\sqrt v=\mu/k$, inside
\cref{p3:eq:rademacher}. The $k$th summand becomes
\[
\frac{A_{k}(n)\sqrt k}{\pi\sqrt2}\cdot\frac1{4N^{3/2}} \Bigl[\Bigl(\frac\mu k-1\Bigr)\e^{\mu/k}
+\Bigl(\frac\mu k+1\Bigr)\e^{-\mu/k}\Bigr] =\frac{A_{k}(n)\sqrt
k}{\pi\sqrt2}\cdot\frac{\mu}{4kN^{3/2}} \Bigl[\e^{\mu/k}\Bigl(1-\frac k\mu\Bigr)
+\e^{-\mu/k}\Bigl(1+\frac k\mu\Bigr)\Bigr].
\]
Now $\mu/(\pi\sqrt2\,\sqrt N)=\lambda/(\pi\sqrt2)=1/\sqrt3$, so the prefactor is
$A_{k}(n)/(4\sqrt{3k}\,N)$, which is \cref{p3:eq:sector} summed over $\sigma$. Since
$\mu^{2}=\lambda^{2}N=(2\pi^{2}/3)N$, we get $1/(4\sqrt{3k}\,N)=\kappa/(\sqrt k\,\mu^{2})$, the
second form of \cref{p3:eq:sector}. Adding the two signs and using
$\e^{z}(1-1/z)+\e^{-z}(1+1/z)=2g(z)$ with $z=\mu/k$ gives \cref{p3:eq:paired-sector}. All
manipulations are termwise algebra, so the sum is unchanged.
\end{proof}

\begin{remark}[The three sources' normal forms are the same]
\label{p3:rem:three-normal-forms}
S1 writes $p(n)=\frac{\mathcal{C}}{x^{2}}\sum_{k}\frac{A_{k}}{\sqrt k}
\{\e^{x/k}(1-k/x)+\e^{-x/k}(1+k/x)\}$ (its $x$ is our $\mu$); S2 writes the sectors
$P_{k,\sigma}=\frac{A_{k}}{4\sqrt{3k}\,x} \e^{\sigma\mu/k}(1-\sigma k/\mu)$ (its $x$ is our $N$);
S3 writes $p(n)=D\mu^{-3}\sum_{k}A_{k}\sqrt k [(\tfrac\mu k-1)\e^{\mu/k}+(\tfrac\mu
k+1)\e^{-\mu/k}]$. Factoring $\mu/k$ out of S3's bracket and using $\sqrt k\cdot(1/k)=1/\sqrt k$
turns S3 into S1, and $\kappa/\mu^{2}=1/(4\sqrt{3k}\,N)\cdot\sqrt k$ turns S1 into S2. All three
are \cref{p3:eq:sector}. We adopt the second form of \cref{p3:eq:sector} because it displays the
Part~0 model $\kappa\e^{\mu}\mu^{-2}(1-1/\mu)$ at $k=1$, $\sigma=+$ verbatim.
\end{remark}

\subsection{The dominant envelope and the exact relative factorisation}

\begin{definition}[Dominant envelope, actions]\label{p3:def:envelope}
Put
\begin{equation}\label{p3:eq:envelope}
 F(\mu)=P_{1,+}=\kappa\,\frac{\e^{\mu}}{\mu^{2}}\left(1-\frac1\mu\right)
 =\kappa\,\e^{\mu}\,\frac{\mu-1}{\mu^{3}},
\end{equation}
and for $k\ge1$, $\sigma=\pm1$ (excluding $(1,+)$) the \emph{action}
\begin{equation}\label{p3:eq:action}
 \alpha_{k,\sigma}=1-\frac{\sigma}{k}
 \qquad\text{so that}\qquad
 \alpha_{1,-}=2,\quad \alpha_{k,+}=1-\frac1k,\quad \alpha_{k,-}=1+\frac1k .
\end{equation}
\end{definition}

\begin{proposition}[Exact relative factorisation]\label{p3:prop:relative}
For $\mu>1$,
\begin{equation}\label{p3:eq:relative}
 p(n)=F(\mu)\bigl(1+\mathcal{R}(\mu;n)\bigr),
\end{equation}
where, with the amplitude ratios
\begin{equation}\label{p3:eq:r-ratios}
 r_{k,\sigma}(\mu;n)=\frac{A_{k}(n)}{\sqrt k}\,
 \frac{1-\sigma k/\mu}{1-1/\mu},
\end{equation}
one has the finite-level resolved form, valid for each $K\ge1$,
\begin{equation}\label{p3:eq:resolved}
 \mathcal{R}=r_{1,-}\,\e^{-2\mu}
 +\sum_{k=2}^{K}\sum_{\sigma=\pm1}
 r_{k,\sigma}(\mu;n)\,\e^{-\alpha_{k,\sigma}\mu}
 +\mathcal{R}_{>K}(\mu;n),
\end{equation}
and the exact paired form
\begin{equation}\label{p3:eq:R-exact}
 \mathcal{R}(\mu;n)=\frac{1+\mu^{-1}}{1-\mu^{-1}}\e^{-2\mu}
 +\frac{2\e^{-\mu}}{1-\mu^{-1}}
 \sum_{k\ge2}\frac{A_{k}(n)}{\sqrt k}\,g\!\left(\frac\mu k\right),
\qquad
 \mathcal{R}_{>K}=\frac{2\e^{-\mu}}{1-\mu^{-1}}\sum_{k>K}\frac{A_{k}(n)}{\sqrt k}
 g\!\left(\frac\mu k\right).
\end{equation}
\end{proposition}

\begin{proof}
Divide \cref{p3:eq:sector} by \cref{p3:eq:envelope}: $P_{k,\sigma}/F=\frac{A_{k}}{\sqrt
k}\frac{1-\sigma k/\mu}{1-1/\mu} \e^{(\sigma/k-1)\mu}$, which is \cref{p3:eq:r-ratios} times
$\e^{-\alpha_{k,\sigma}\mu}$; the case $k=1,\sigma=-1$ has $A_{1}=1$ and gives the first term of
\cref{p3:eq:R-exact}. For $k\ge2$ the paired sum
$P_{k,+}+P_{k,-}=\frac{2\kappa}{\mu^{2}}\frac{A_{k}}{\sqrt k} g(\mu/k)$ from
\cref{p3:eq:paired-sector}, divided by $F=\kappa\e^{\mu}\mu^{-2}(1-1/\mu)$, gives the second term.
\end{proof}

\begin{remark}[What the Hardy--Ramanujan formula omits]
\Cref{p3:eq:HR} is $F(\mu)$ after replacing $N$ by $n$ and discarding $1-1/\mu$. Each of those two
simplifications costs a relative $\bigO(n^{-1/2})$; together they account for the entire
algebraic half-power series of \cref{p3:sec:half-powers}. \Cref{p3:eq:relative} loses nothing: it
records every Ford-circle arc, in both signs, with its exact algebraic amplitude.
\end{remark}

\section{Ordering of the sectors and tail estimates}\label{p3:sec:tail}

\subsection{The action-one accumulation}

The relative actions of \cref{p3:eq:action} are
\[
\tfrac12,\ \tfrac23,\ \tfrac34,\ \dots\ \uparrow 1 \qquad\text{and}\qquad 2,\ \tfrac32,\
\tfrac43,\ \dots\ \downarrow 1 .
\]
The first family is increasing with supremum $1$; the second is decreasing with infimum $1$ and
\emph{has no least element}. Two consequences must be stated before any transseries language is
used.

\begin{definition}[Rademacher transseries]\label{p3:def:rademacher-transseries}
A \emph{Rademacher transseries} for $p$ is the datum of
\begin{enumerate}
\item for each $K\ge1$, the finite family of resolved sectors
$\{(k,\sigma):k\le K\}$ with the coefficients of \cref{p3:sec:half-powers}, treated as independent
formal exponential variables (so that coefficient extraction takes place in a finitely generated
multivariate power-series ring), together with
\item the exact paired remainder $\mathcal{R}_{>K}$ of \cref{p3:eq:R-exact}, retained
as one analytic block.
\end{enumerate}
The limit $K\to\infty$ is taken analytically, through the convergence in \cref{p3:thm:rademacher},
and never by claiming that the whole support is a locally finite Hahn grid.
\end{definition}

\begin{remark}[Why the pairing is mandatory]\label{p3:rem:pairing}
For $k\gg\mu$ each of $P_{k,+}$ and $P_{k,-}$ separately contains a term proportional to $\sqrt
k\,/\mu$, and $\sum_{k}\sqrt k/\mu$ diverges. Only in the sum do those pieces cancel:
$2g(\mu/k)=\tfrac23(\mu/k)^{2}+O((\mu/k)^{4})$, which makes the $k$th paired term
$\bigO(k^{-3/2})$. Thus \cref{p3:eq:sector-sum} prescribes an order of summation; it is not
an absolutely rearrangeable double series. All three sources say this; S2 and S3 say it most
sharply, and \cref{p3:def:rademacher-transseries} follows S2.
\end{remark}

\begin{remark}[Truncating ``through $\e^{-\mu}$'' is not a prescription]
\label{p3:rem:no-eta-cutoff}
Because the positive actions accumulate at $1$, an instruction such as ``retain all terms down to
relative order $\e^{-\mu}$'' selects infinitely many sectors and is meaningless. A legitimate
truncation names either a Rademacher level $K$, or an action cutoff $\eta<1$ (which then names
finitely many positive sectors, namely $k\le(1-\eta)^{-1}$), or a target accuracy together with
\cref{p3:thm:tail}.
\end{remark}

\subsection{The tail bound}

\begin{lemma}[Kernel bound]\label{p3:lem:g-bound}
For $z\ge0$,
\begin{equation}\label{p3:eq:g-bound}
 0\le g(z)\le\frac{z^{2}}{3}\cosh z\le\frac{z^{2}}{3}\e^{z},
 \qquad\text{hence}\qquad
 g(z)\le\frac{\cosh 1}{3}z^{2}\quad\text{for }0\le z\le1 .
\end{equation}
The factor $\cosh z$ may not be dropped: $g(1)=\e^{-1}=0.3679\ldots>\tfrac13$.
\end{lemma}

\begin{proof}
By \cref{p3:eq:g-kernel}, $g(z)=\sum_{m\ge1}\frac{2m}{(2m+1)!}z^{2m}$ has non-negative
coefficients, whence $g\ge0$. For the upper bound compare with $\frac{z^{2}}{3}\cosh
z=\sum_{m\ge1}\frac{z^{2m}}{3(2m-2)!}$ termwise: the claim is
$\frac{2m}{(2m+1)!}\le\frac1{3(2m-2)!}$, i.e.\ $6m(2m-2)!\le(2m+1)!$, i.e.\ $6m\le(2m+1)2m(2m-1)$,
which holds for all $m\ge1$ with equality at $m=1$. Monotonicity of $\cosh$ gives the last
inequality on $[0,1]$, and $g(1)=\cosh1-\sinh1=\e^{-1}$ shows the factor cannot be dropped.
\end{proof}

\begin{theorem}[Explicit finite-$K$ tail bound]\label{p3:thm:tail}
Let $\mu>1$ and $K\ge1$. Then
\begin{equation}\label{p3:eq:tail}
 \bigl|\mathcal{R}_{>K}(\mu;n)\bigr|
 \le\frac{4\mu^{2}}{3\,(1-\mu^{-1})\sqrt K}\,
 \exp\!\left[-\left(1-\frac1{K+1}\right)\mu\right].
\end{equation}
The same bound holds with $A_{k}(n)$ replaced by $\widetilde A_{k}(\nu)$ for any real $\nu$, since
$|\widetilde A_{k}|\le\varphi(k)\le k$ as well.
\end{theorem}

\begin{proof}
From \cref{p3:eq:R-exact}, $|A_{k}|\le k$ and \cref{p3:lem:g-bound},
\[
|\mathcal{R}_{>K}| \le\frac{2\e^{-\mu}}{1-\mu^{-1}}\sum_{k>K}\sqrt k\,g\!\left(\frac\mu k\right)
\le\frac{2\e^{-\mu}}{1-\mu^{-1}}\cdot\frac{\mu^{2}}{3} \sum_{k>K}k^{-3/2}\e^{\mu/k}
\le\frac{2\mu^{2}\e^{-\mu+\mu/(K+1)}}{3(1-\mu^{-1})}\sum_{k>K}k^{-3/2},
\]
and $\sum_{k>K}k^{-3/2}\le\int_{K}^{\infty}t^{-3/2}\dd t=2K^{-1/2}$.
\end{proof}

\begin{theorem}[Sharper fixed-$K$ tail]\label{p3:thm:tail-sharp}
For each fixed $K\ge1$, as $\mu\to\infty$,
\begin{equation}\label{p3:eq:tail-sharp}
 \sum_{k>K}\bigl(P_{k,+}+P_{k,-}\bigr)
 =\bigO_{K}\!\left(\mu^{-1/2}\e^{\mu/(K+1)}\right),
 \qquad
 \mathcal{R}_{>K}=\bigO_{K}\!\left(\mu^{3/2}\e^{-(1-1/(K+1))\mu}\right).
\end{equation}
\end{theorem}

\begin{proof}
Split at $k=\mu$. For $K<k\le\mu$ use $|A_{k}|\le k$ and
$2g(\mu/k)\le2\cosh(\mu/k)\le2\e^{\mu/(K+1)}$ in \cref{p3:eq:paired-sector}, so those terms
contribute at most $\frac{2\kappa}{\mu^{2}}\cdot2\e^{\mu/(K+1)}\sum_{k\le\mu}\sqrt k
=\bigO(\mu^{-1/2}\e^{\mu/(K+1)})$ since $\sum_{k\le\mu}\sqrt k=\bigO(\mu^{3/2})$. For
$k>\mu$ put $z=\mu/k\in(0,1)$; then $g(z)\le\tfrac13z^{2}\cosh1$ by \cref{p3:lem:g-bound}, and the
$k$th term is $\bigO(\mu^{-2}\sqrt k\cdot\mu^{2}k^{-2})=\bigO(k^{-3/2})$, whose tail is
$\bigO(\mu^{-1/2})$. Dividing by $F\asymp\kappa\mu^{-2}\e^{\mu}$ gives the relative
statement.
\end{proof}

\begin{remark}[Which bound to use]
\Cref{p3:thm:tail} (from S2) is explicit and uniform in $K$; \cref{p3:thm:tail-sharp} (from S1 and
S3) is sharper in $\mu$ by a factor $\mu^{-1/2}$ but only for fixed $K$. Both are needed: the
first to certify a computation, the second to state clean asymptotic orders. S1 and S3 state only
the second, S2 only the first.
\end{remark}

\begin{remark}[Sharper amplitude bounds]\label{p3:rem:sharper-bounds}
All constants above come from $|A_{k}|\le k$. Lehmer's explicit remainder estimates and the
effective error bounds of Banerjee--Paule--Radu--Schneider (\cref{p3:rem:references}) are much
sharper, and Weil-type bounds for the Kloosterman-like sums $A_{k}$ give
$|A_{k}(n)|\ll_{\varepsilon}k^{1/2+\varepsilon}$ for many $k$. Substituting any such bound in the
proofs of \cref{p3:thm:tail,p3:thm:tail-sharp} improves the powers of $\mu$ but changes no
statement of this chapter, all of which are of the form ``a fixed Rademacher level is
exponentially accurate''. For a \emph{specific} residue class some $A_{k}$ vanish, so the first
non-zero omitted sector can be much smaller than the generic estimate.
\end{remark}

\subsection{An exact condensation at the accumulation level}

Right at the accumulation action $1$, an alternative to sector-by-sector bookkeeping is to sum the
kernel expansion first. This is S3's contribution and has no counterpart in S1 or S2.

\begin{definition}[Rademacher--Kloosterman zeta]\label{p3:def:Zeta}
For $\Re s>2$ and $n\in\Int$ put $\displaystyle
\mathcal{Z}_{n}(s)=\sum_{k\ge1}\frac{A_{k}(n)}{k^{s}}$; the series converges absolutely there by
\cref{p3:lem:Ak-properties}(iii). The same definition with $\widetilde A_{k}(\nu)$ defines
$\mathcal{Z}_{\nu}(s)$ for real $\nu$.
\end{definition}

\begin{theorem}[Exact zeta condensation]\label{p3:thm:zeta-condensation}
For every integer $n\ge1$,
\begin{equation}\label{p3:eq:zeta-condensation}
 p(n)=\kappa\sum_{r\ge1}\frac{4r}{(2r+1)!}\,\mu^{2r-2}\,
 \mathcal{Z}_{n}\!\left(2r+\frac12\right),
\end{equation}
and the double series obtained by expanding $\mathcal{Z}_{n}$ converges absolutely.
\end{theorem}

\begin{proof}
Insert $2g(z)=\sum_{r\ge1}\frac{4r}{(2r+1)!}z^{2r+1}\cdot z^{-1}$ --- more precisely
$2g(z)=\sum_{r\ge1}\frac{4r}{(2r+1)!}z^{2r}$ --- into \cref{p3:eq:paired-sector} with $z=\mu/k$:
\[
p(n)=\frac{\kappa}{\mu^{2}}\sum_{k\ge1}\frac{A_{k}(n)}{\sqrt k}
\sum_{r\ge1}\frac{4r}{(2r+1)!}\frac{\mu^{2r}}{k^{2r}}
=\kappa\sum_{k,r}\frac{4r}{(2r+1)!}\mu^{2r-2}\frac{A_{k}(n)}{k^{2r+1/2}} .
\]
Absolute convergence: $\sum_{k}|A_{k}|k^{-2r-1/2}\le\sum_{k}k^{-2r+1/2}<\infty$ for $r\ge1$, and
$\sum_{r\ge1}\frac{4r}{(2r+1)!}\mu^{2r-2}\zeta(2r-\tfrac12)$ converges for every fixed $\mu$
because $\zeta(2r-\tfrac12)\to1$ and $\sum_{r}\frac{4r\mu^{2r-2}}{(2r+1)!}<\infty$. Tonelli's
theorem justifies the interchange, giving \cref{p3:eq:zeta-condensation}.
\end{proof}

\begin{remark}
\Cref{p3:eq:zeta-condensation} orders nothing by asymptotic size --- it is an exact analytic
rearrangement, not a transseries. Its use is precisely where the transseries language fails: it
packages the whole action-one accumulation into Dirichlet values $\mathcal{Z}_{n}(2r+\tfrac12)$,
which can be used as a single block in the inverse formulas of \cref{p3:sec:full-inverse}.
\end{remark}

\section{All-order coefficients in the unshifted variable}\label{p3:sec:half-powers}

The shifted phase $\mu$ makes each sector elementary; the price is that the conventional scale is
$n^{-1/2}$. The entire infinite algebraic array of this subject comes from re-expanding the single
shift $N=n-1/24$. This section gives all of it, in one formula valid for every sector.

\subsection{The sector generating function}

\begin{theorem}[Sectorwise convergent half-power expansion]
\label{p3:thm:sector-gf}
Fix $k\ge1$ and $\sigma\in\{+1,-1\}$, and set $t=n^{-1/2}$. Then
\begin{equation}\label{p3:eq:sector-unshifted}
 P_{k,\sigma}(n)=\frac{A_{k}(n)}{4\sqrt{3k}\;n}\,
 \exp\!\left(\frac{\sigma\lambda\sqrt n}{k}\right)
 \sum_{m\ge0}c^{(k,\sigma)}_{m}\,n^{-m/2},
\end{equation}
where the $c^{(k,\sigma)}_{m}$ are the Taylor coefficients at $t=0$ of
\begin{equation}\label{p3:eq:Omega}
 \Omega_{k,\sigma}(t)
 =\frac1{1-t^{2}/24}\,
 \exp\!\left[\frac{\sigma\lambda}{kt}
 \left(\sqrt{1-\frac{t^{2}}{24}}-1\right)\right]
 \left(1-\frac{\sigma kt}{\lambda\sqrt{1-t^{2}/24}}\right).
\end{equation}
The function $\Omega_{k,\sigma}$ is holomorphic on $|t|<\sqrt{24}$, so the series in
\cref{p3:eq:sector-unshifted} converges for every real $n>1/24$; in particular
$c^{(k,\sigma)}_{0}=1$.
\end{theorem}

\begin{proof}
With $t=n^{-1/2}$ we have $N=n(1-t^{2}/24)$ and $\mu=\lambda\sqrt N=(\lambda/t)\sqrt{1-t^{2}/24}$.
Substituting in \cref{p3:eq:sector},
\[
P_{k,\sigma} =\frac{A_{k}(n)}{4\sqrt{3k}\,n}\cdot\frac1{1-t^{2}/24}\cdot
\e^{\sigma\mu/k}\left(1-\frac{\sigma kt}{\lambda\sqrt{1-t^{2}/24}}\right),
\]
and $\e^{\sigma\mu/k}=\e^{\sigma\lambda/(kt)}
\exp[\tfrac{\sigma\lambda}{kt}(\sqrt{1-t^{2}/24}-1)]$ with
$\e^{\sigma\lambda/(kt)}=\e^{\sigma\lambda\sqrt n/k}$. This is \cref{p3:eq:sector-unshifted} with
\cref{p3:eq:Omega}. Holomorphy: the only candidate singularities of \cref{p3:eq:Omega} in
$|t|<\sqrt{24}$ are $t=0$ and the zeros of $\sqrt{1-t^{2}/24}$, and the latter lie on
$|t|=\sqrt{24}$. At $t=0$ the exponent is
$\frac{\sigma\lambda}{kt}(\sqrt{1-t^{2}/24}-1)=-\frac{\sigma\lambda t}{48k} +\bigO(t^{3})$,
so the singularity is removable, and the last factor is holomorphic there as well. Setting $t=0$
gives $\Omega_{k,\sigma}(0)=1$.
\end{proof}

\begin{repairbox}
\textbf{S1, Theorem ``Universal sector coefficient generator''.} S1 states only that the series
``converges for all sufficiently small $|t|$ and supplies an asymptotic expansion to arbitrary
order''. This understates the truth and, read literally, invites the error the next remark
forbids: it suggests the half-power series is a merely asymptotic object. S2 and S3 both identify
the radius as exactly $\sqrt{24}$, i.e.\ convergence for every real $n>1/24$;
\cref{p3:thm:sector-gf} adopts their statement with the branch-point argument made explicit.
\textbf{Consequence.} Nothing in this chapter is a divergent series in $n^{-1/2}$, and no
least-term truncation is available; see \cref{p3:rem:no-least-term}.
\end{repairbox}

\subsection{A finite closed formula for every sector coefficient}

Write
\begin{equation}\label{p3:eq:b-parameter}
 \beta=\beta_{k,\sigma}=\frac{\sigma\pi}{6k},
 \qquad\text{so that}\qquad
 \frac{\sigma\lambda}{k}=2\sqrt6\,\beta,
 \qquad
 \frac{\sigma k}{\lambda}=\frac{1}{12\beta}\cdot\sqrt{6}\,,
\end{equation}
the two identities following from $\lambda=\pi\sqrt6/3$ and $k=\sigma\pi/(6\beta)$. Thus each
sector is the \emph{same} function of the single parameter $\beta$; the sector label enters
nowhere else.

\begin{theorem}[Finite coefficient formula, every sector]\label{p3:thm:c-finite}
For every $k\ge1$, $\sigma=\pm1$ and $m\ge0$,
\begin{equation}\label{p3:eq:c-finite}
 c^{(k,\sigma)}_{m}
 =\frac1{(-4\sqrt6)^{m}}
 \sum_{j=0}^{\lfloor(m+1)/2\rfloor}
 \binom{m+1}{j}\,\frac{m+1-j}{(m+1-2j)!}\,
 \beta_{k,\sigma}^{\,m-2j},
\end{equation}
terms with a negative factorial index being absent. In particular
$c^{(k,\sigma)}_{m}\in\Rat[\beta,\beta^{-1}]\cdot\sqrt6^{\,m\bmod 2}$, and the dependence on the
sector is only through $\beta_{k,\sigma}=\sigma\pi/(6k)$.
\end{theorem}

\begin{proof}
Put $u=-t/(4\sqrt6)$, so that $t^{2}/24=4u^{2}$ and, by \cref{p3:eq:b-parameter},
$\frac{\sigma\lambda}{kt}=\frac{2\sqrt6\beta}{-4\sqrt6\,u}=-\frac{\beta}{2u}$ and $\frac{\sigma
kt}{\lambda}=-\frac{4\sqrt6\,u\sqrt6}{12\beta}=-\frac{2u}{\beta}$. Writing $s=\sqrt{1-4u^{2}}$,
\cref{p3:eq:Omega} becomes
\begin{equation}\label{p3:eq:Omega-u}
 G(u):=\Omega_{k,\sigma}(-4\sqrt6\,u)
 =\frac{\e^{\beta T}}{s^{2}}\left(1+\frac{2u}{\beta s}\right),
 \qquad
 T=T(u)=\frac{1-s}{2u},
\end{equation}
where we used $-\frac{\beta}{2u}(s-1)=\beta\frac{1-s}{2u}=\beta T$ and $s^{2}=1-4u^{2}$. The
Catalan variable $T$ satisfies
\begin{equation}\label{p3:eq:catalan}
 T=u(1+T^{2}),\qquad u=\frac{T}{1+T^{2}},\qquad
 s=\frac{1-T^{2}}{1+T^{2}},\qquad
 \dd u=\frac{1-T^{2}}{(1+T^{2})^{2}}\,\dd T .
\end{equation}
(The first identity follows from $2uT=1-s$ and $s^{2}=1-4u^{2}$; the rest are algebra.)
Substituting \cref{p3:eq:catalan} into \cref{p3:eq:Omega-u} gives the exact differential identity
\begin{equation}\label{p3:eq:G-differential}
 G(u)\,\dd u
 =\e^{\beta T}\left(\frac{1}{1-T^{2}}+\frac{2T}{\beta(1-T^{2})^{2}}\right)\dd T
 =\frac1\beta\,\dd\!\left(\frac{\e^{\beta T}}{1-T^{2}}\right),
\end{equation}
because $\frac{\dd}{\dd T}\frac{\e^{\beta T}}{1-T^{2}} =\e^{\beta
T}\bigl(\frac{\beta}{1-T^{2}}+\frac{2T}{(1-T^{2})^{2}}\bigr)$. Now let
$D_{m}=\Coef{u^{m}}G(u)=(-4\sqrt6)^{m}c^{(k,\sigma)}_{m}$. Then
\[
D_{m}=\operatorname*{Res}_{u=0}\frac{G(u)\dd u}{u^{m+1}} =\frac1\beta\operatorname*{Res}_{T=0}
\frac{(1+T^{2})^{m+1}}{T^{m+1}} \,\dd\!\left(\frac{\e^{\beta T}}{1-T^{2}}\right),
\]
using $u^{-m-1}=T^{-m-1}(1+T^{2})^{m+1}$. The residue of an exact differential vanishes, so
integrating by parts,
\[
D_{m}=-\frac1\beta\operatorname*{Res}_{T=0} \frac{\e^{\beta T}}{1-T^{2}}
\,\dd\!\left(T^{-m-1}(1+T^{2})^{m+1}\right) =\frac{m+1}{\beta}\operatorname*{Res}_{T=0} \e^{\beta
T}(1+T^{2})^{m}\,\frac{\dd T}{T^{m+2}},
\]
where we used $\dd(T^{-m-1}(1+T^{2})^{m+1})
=(1+T^{2})^{m}T^{-m-2}\bigl(-(m+1)(1+T^{2})+2(m+1)T^{2}\bigr)\dd T
=-(m+1)(1+T^{2})^{m}(1-T^{2})T^{-m-2}\dd T$, and the factor $1-T^{2}$ cancels. Hence
\[
D_{m}=\frac{m+1}{\beta}\Coef{T^{m+1}}\e^{\beta T}(1+T^{2})^{m}
=\frac{m+1}{\beta}\sum_{j\ge0}\binom mj\frac{\beta^{m+1-2j}}{(m+1-2j)!},
\]
the sum running over $0\le j\le\lfloor(m+1)/2\rfloor$. Finally $(m+1)\binom
mj=(m+1-j)\binom{m+1}{j}$ and division by $(-4\sqrt6)^{m}$ gives \cref{p3:eq:c-finite}.
\end{proof}

\begin{remark}[Why a finite sum exists]\label{p3:rem:why-finite}
A square-root shift would normally give an infinite convolution at each order. The Catalan
coordinate $T$ absorbs the square root, and the derivative already present in Rademacher's formula
turns the integrand into the exact differential \cref{p3:eq:G-differential}. After the
change-of-variable Jacobian cancels the denominator $1-T^{2}$, only the polynomial $(1+T^{2})^{m}$
survives.
\end{remark}

\begin{remark}[Provenance and the scope of the generalisation]
\label{p3:rem:osullivan}
For $(k,\sigma)=(1,+1)$, \cref{p3:eq:c-finite} at $\beta=\pi/6$ is O'Sullivan's coefficient
formula, also used by Banerjee, Paule, Radu and Schneider in their error analysis. S1 proves only
that case, by the same Catalan substitution. S2 and S3 both observe that the proof never uses
$k=1$ or $\sigma=+1$ --- the sector enters only through $\beta$ --- and state the general formula;
S3 gives a second, generating-function proof through the identity
$\sum_{j\ge0}\binom{m+2j}{j}q^{j}=C(q)^{m}/\sqrt{1-4q}$ for the Catalan series $C$. The residue
proof above is S2's, written out. We record the generalisation as the principal coefficient
statement of the chapter, since it covers every exponentially small sector at no extra cost.
\end{remark}

The first coefficients, as functions of $\beta$, are
\begin{align}
 c_{0}&=1, &
 c_{1}&=-\frac{\sqrt6\,(\beta^{2}+2)}{24\beta}, &
 c_{2}&=\frac{\beta^{2}+12}{192},\nonumber\\
 c_{3}&=-\frac{\sqrt6\,(\beta^{4}+36\beta^{2}+72)}{13824\,\beta}, &
 c_{4}&=\frac{\beta^{4}+80\beta^{2}+720}{221184}, &
 c_{5}&=-\frac{\sqrt6\,(\beta^{6}+150\beta^{4}+3600\beta^{2}+7200)}
 {26542080\,\beta},\label{p3:eq:c-first}
\end{align}
\begin{equation}\label{p3:eq:c-six}
 c_{6}=\frac{\beta^{6}+252\beta^{4}+12600\beta^{2}+100800}{637009920},
 \qquad
 c_{7}=-\frac{\sqrt6\bigl(\beta^{8}+392\beta^{6}+35280\beta^{4}
 +705600\beta^{2}+1411200\bigr)}{107017666560\,\beta}.
\end{equation}
These were recomputed here by expanding \cref{p3:eq:Omega} symbolically to order $t^{9}$ and
comparing with \cref{p3:eq:c-finite} coefficient by coefficient in exact arithmetic; the two agree
identically. The parity pattern visible in \cref{p3:eq:c-first} --- $c_{m}$ is $\sqrt6$ times an
odd Laurent polynomial in $\beta$ for odd $m$, and a polynomial in $\beta^{2}$ for even $m$ --- is
immediate from \cref{p3:eq:c-finite}, since $m-2j\equiv m$ $(\mathrm{mod}\ 2)$ and the prefactor
contributes $\sqrt6^{\,m}$.

\subsection{The Bell-polynomial route}

\Cref{p3:eq:c-finite} is the fastest generator. For symbolic systems, or when the sector
coefficients are wanted as a byproduct of a general composition engine, the following route uses
only the primitives of \cref{p0:def:obell,p0:def:ebell}.

\begin{proposition}[Logarithmic coefficients of a sector]\label{p3:prop:ell}
Write $\log\Omega_{k,\sigma}(t)=\sum_{j\ge1}\ell_{j}t^{j}$ (suppressing the sector labels) and put
$a=1/24$. Then, for $j\ge1$,
\begin{equation}\label{p3:eq:ell}
 \ell_{j}
 =\mathbf 1_{2\mid j}\,\frac{a^{j/2}}{j/2}
 +\mathbf 1_{2\nmid j}\,\frac{\sigma\lambda}{k}
   \binom{1/2}{(j+1)/2}(-a)^{(j+1)/2}
 -\sum_{\substack{1\le q\le j\\ 2\mid (j-q)}}\frac1q
   \left(\frac{\sigma k}{\lambda}\right)^{\!q}
   \frac{(q/2)_{(j-q)/2}}{\bigl((j-q)/2\bigr)!}\;a^{(j-q)/2},
\end{equation}
where $(x)_{r}=x(x+1)\cdots(x+r-1)$. Consequently the sector coefficients are the complete
exponential Bell polynomials
\begin{equation}\label{p3:eq:c-bell}
 c_{m}=\frac1{m!}\,\EB_{m}\bigl(1!\,\ell_{1},2!\,\ell_{2},\dots,m!\,\ell_{m}\bigr),
\end{equation}
equivalently the triangular recurrence
\begin{equation}\label{p3:eq:c-recurrence}
 c_{0}=1,\qquad c_{m}=\frac1m\sum_{j=1}^{m}j\,\ell_{j}\,c_{m-j}\quad(m\ge1).
\end{equation}
\end{proposition}

\begin{proof}
Taking logarithms in \cref{p3:eq:Omega},
\[
\log\Omega=-\log(1-at^{2}) +\frac{\sigma\lambda}{kt}\bigl(\sqrt{1-at^{2}}-1\bigr)
+\log\!\left(1-\frac{\sigma k}{\lambda}t(1-at^{2})^{-1/2}\right).
\]
The first term contributes $a^{r}/r$ at degree $2r$. For the second,
$\sqrt{1-at^{2}}-1=\sum_{r\ge1}\binom{1/2}{r}(-a)^{r}t^{2r}$, so dividing by $t$ contributes
$\frac{\sigma\lambda}{k}\binom{1/2}{r}(-a)^{r}$ at degree $2r-1$. For the third, expand
$\log(1-w)=-\sum_{q\ge1}w^{q}/q$ with $w=\frac{\sigma k}{\lambda}t(1-at^{2})^{-1/2}$ and use
$(1-at^{2})^{-q/2}=\sum_{r\ge0}\frac{(q/2)_{r}}{r!}a^{r}t^{2r}$; the term of degree $j=q+2r$ is
the displayed sum. \Cref{p3:eq:c-bell} is the exponential formula of \cref{p0:def:ebell}, and
\cref{p3:eq:c-recurrence} follows by comparing coefficients in $\Omega'=(\log\Omega)'\,\Omega$.
\end{proof}

For orientation, at general $(k,\sigma)$ and $a=1/24$,
\begin{align}
 \ell_{1}&=-\frac{\sigma\lambda a}{2k}-\frac{\sigma k}{\lambda}, &
 \ell_{2}&=a-\frac{k^{2}}{2\lambda^{2}},\nonumber\\
 \ell_{3}&=-\frac{\sigma\lambda a^{2}}{8k}-\frac{\sigma k a}{2\lambda}
 -\frac{\sigma k^{3}}{3\lambda^{3}}, &
 \ell_{4}&=\frac{a^{2}}{2}-\frac{k^{2}a}{2\lambda^{2}}
 -\frac{k^{4}}{4\lambda^{4}} .\label{p3:eq:ell-first}
\end{align}

\subsection{The dominant Poincar\'e expansion, and what it does not say}

\begin{definition}\label{p3:def:omega}
$\omega_{m}=c^{(1,+1)}_{m}$, i.e.\ \cref{p3:eq:c-finite} at $\beta=\pi/6$.
\end{definition}

\begin{corollary}\label{p3:cor:poincare}
For each fixed $M\ge1$,
\begin{equation}\label{p3:eq:poincare}
 p(n)=\frac{\e^{\lambda\sqrt n}}{4\sqrt3\,n}
 \left(\sum_{m=0}^{M-1}\omega_{m}n^{-m/2}+\bigO_{M}(n^{-M/2})\right),
\end{equation}
with
\begin{equation}\label{p3:eq:omega-first}
 \omega_{0}=1,\quad
 \omega_{1}=-\frac{\sqrt6\,(\pi^{2}+72)}{144\pi},\quad
 \omega_{2}=\frac{\pi^{2}+432}{6912},\quad
 \omega_{3}=-\frac{\sqrt6\,(\pi^{4}+1296\pi^{2}+93312)}{2985984\,\pi},
\end{equation}
\begin{equation}\label{p3:eq:omega-four}
 \omega_{4}=\frac{\pi^{4}+2880\pi^{2}+933120}{286654464},
\end{equation}
numerically $\omega_{1}=-0.443287976873582391\ldots$, $\omega_{2}=0.063927894155250197\ldots$,
$\omega_{3}=-0.027730933907464250\ldots$, $\omega_{4}=0.003354707463289919\ldots$.
\end{corollary}

\begin{proof}
By \cref{p3:thm:sector-gf} with $(k,\sigma)=(1,+1)$, the displayed sum plus a convergent tail is
exactly $P_{1,+}(n)$ divided by $\e^{\lambda\sqrt n}/(4\sqrt3\,n)$; the tail after $M$ terms is
$\bigO_{M}(n^{-M/2})$ because $\Omega_{1,+}$ is holomorphic at $0$. The difference
$p(n)-P_{1,+}(n)$ is $\bigO(\mu^{3/2}\e^{-\mu/2})$ relative to $P_{1,+}$ by
\cref{p3:thm:tail-sharp} with $K=1$, hence beyond every algebraic order. The closed forms follow
from \cref{p3:eq:c-finite} at $\beta=\pi/6$; they were recomputed here in exact arithmetic and
agree with the tables of S1 and S2 (which print the same numbers in two different but equal
groupings).
\end{proof}

\begin{repairbox}
\textbf{Convergent sector versus asymptotic total.} This is the trap of the subject and all three
sources are, at some point, loose about it. The three statements below are \emph{different}:
\begin{enumerate}
\item $\sum_{m\ge0}\omega_{m}n^{-m/2}$ \emph{converges} for every $n>1/24$
(\cref{p3:thm:sector-gf});
\item its sum is $P_{1,+}(n)\cdot 4\sqrt3\,n\,\e^{-\lambda\sqrt n}$, \emph{not}
$p(n)\cdot4\sqrt3\,n\,\e^{-\lambda\sqrt n}$;
\item \cref{p3:eq:poincare} is an \emph{asymptotic} statement about $p(n)$,
whose error at every fixed order is a sum of two utterly different things: a convergent algebraic
tail of size $\bigO(n^{-M/2})$ and an exponentially small Rademacher remainder of size
$\bigO(\mu^{3/2}\e^{-\mu/2})$ relative to the dominant sector.
\end{enumerate}
S1's ``Dominant Poincar\'e expansion'' subsection writes $p(n)\sim\cdots$ without separating (ii)
from (iii) and, in its Theorem on the sector generating function, weakens (i) to ``sufficiently
small $|t|$''. S3's remark after its dominant Poincar\'e display states the correct trichotomy in
one sentence. S2 states it as a standalone remark and is the presentation adopted here.
\textbf{The practical consequence:} summing the $\omega_{m}$ series to convergence does \emph{not}
compute $p(n)$ to arbitrary accuracy; it computes $P_{1,+}(n)$, and the residual error stalls at
the $\e^{-\mu/2}$ floor.
\end{repairbox}

\begin{remark}[Least-term truncation degenerates here]\label{p3:rem:no-least-term}
\Cref{p0:thm:optimal-truncation} determines the optimal truncation point of a \emph{factorially
divergent} algebraic series and the size of the resulting beyond-all-orders floor. Its hypothesis
fails for $p(n)$: by \cref{p3:thm:sector-gf} the algebraic series of each sector has a finite
positive radius of convergence, so there is no least term and no truncation-induced floor. The
floor that does exist --- the $\e^{-\mu/2}$ Rademacher gap --- is not produced by resummation
ambiguity but is supplied explicitly by the next sector. The partition function is thus the
degenerate, best possible case of \cref{p0:thm:optimal-truncation}: the exponentially small
sectors are known exactly and independently of the perturbative data, and the optimal algebraic
truncation is simply ``as many terms as one wishes''. This is the structural reason the numerical
accuracy in \cref{p3:tab:forward} improves so violently with $K$ but not at all with extra
algebraic orders past the $\e^{-\mu/2}$ level.
\end{remark}

\section{The multiplicative reciprocal \texorpdfstring{$1/p(n)$}{1/p(n)}}
\label{p3:sec:reciprocal}

The phrase ``inverse of the partition numbers'' is ambiguous. The reciprocal is the easy reading
and is disposed of first; it needs no reversion at all.

\begin{proposition}[Exact reciprocal transseries]\label{p3:prop:reciprocal}
Let $\mu>1$ be large enough that $|\mathcal{R}(\mu;n)|<1$, with $\mathcal{R}$ as in
\cref{p3:eq:R-exact}. Then
\begin{equation}\label{p3:eq:reciprocal}
 \frac1{p(n)}=\frac{\mu^{2}}{\kappa}\,\frac{\e^{-\mu}}{1-\mu^{-1}}
 \sum_{r\ge0}(-1)^{r}\mathcal{R}(\mu;n)^{r},
\end{equation}
the series converging. As a formal identity in \cref{p3:def:rademacher-transseries} it holds
without the inequality. If $\mathcal{R}$ is written at finite level as a linear form $\sum_{j\in
J_{K}}r_{j}(\mu;n)\,q_{j}$ in independent exponential variables $q_{j}=\e^{-\alpha_{j}\mu}$, then
for every multi-index $\boldsymbol\alpha=(\alpha_{j})_{j\in J_{K}}\ne\boldsymbol0$,
\begin{equation}\label{p3:eq:reciprocal-multi}
 \Coef{\boldsymbol q^{\boldsymbol\alpha}}\bigl(1+\mathcal{R}\bigr)^{-1}
 =(-1)^{|\boldsymbol\alpha|}\frac{|\boldsymbol\alpha|!}{\boldsymbol\alpha!}
 \prod_{j\in J_{K}}r_{j}(\mu;n)^{\alpha_{j}} .
\end{equation}
\end{proposition}

\begin{proof}
Invert \cref{p3:eq:relative} and expand $(1+\mathcal{R})^{-1}$ geometrically;
$1/F=\mu^{2}\e^{-\mu}/(\kappa(1-\mu^{-1}))$. For \cref{p3:eq:reciprocal-multi}, only the term
$(-1)^{|\boldsymbol\alpha|} \mathcal{R}^{|\boldsymbol\alpha|}$ of the geometric series can produce
$\boldsymbol q^{\boldsymbol\alpha}$, and the multinomial theorem supplies the coefficient.
\end{proof}

\begin{corollary}[Reciprocal Poincar\'e coefficients]\label{p3:cor:reciprocal-poincare}
Define $\eta_{m}$ by $\bigl(\sum_{m\ge0}\omega_{m}t^{m}\bigr)^{-1} =\sum_{m\ge0}\eta_{m}t^{m}$,
i.e.
\begin{equation}\label{p3:eq:eta}
 \eta_{0}=1,\qquad \eta_{m}=-\sum_{j=1}^{m}\omega_{j}\eta_{m-j}
 \quad(m\ge1),
 \qquad\text{equivalently}\qquad
 \eta_{m}=\frac1{m!}\EB_{m}\bigl(-1!\ell_{1},\dots,-m!\ell_{m}\bigr)
\end{equation}
with $\ell_{j}=\ell_{j}^{(1,+1)}$ of \cref{p3:eq:ell}. Then for each fixed $M$,
\begin{equation}\label{p3:eq:reciprocal-poincare}
 \frac1{p(n)}=4\sqrt3\,n\,\e^{-\lambda\sqrt n}
 \left(\sum_{m=0}^{M-1}\eta_{m}n^{-m/2}+\bigO_{M}(n^{-M/2})\right),
\end{equation}
and $\sum_{m\ge0}\eta_{m}t^{m}$ converges on $|t|<\sqrt{24}$ minus the zero set of $\Omega_{1,+}$,
to $1/\Omega_{1,+}(t)$.
\end{corollary}

\begin{proof}
The recurrence is Cauchy inversion of a power series with unit constant term; the Bell form is the
exponential formula applied to $\Omega_{1,+}^{-1}=\exp(-\sum\ell_{j}t^{j})$. Convergence is
inherited from \cref{p3:thm:sector-gf} and $\Omega_{1,+}(0)=1$. The asymptotic statement follows
from \cref{p3:cor:poincare} exactly as there; the same caveats of the repair box after
\cref{p3:cor:poincare} apply verbatim, with ``$P_{1,+}$'' replaced by ``$1/P_{1,+}$''.
\end{proof}

\begin{remark}
The reciprocal therefore requires no new analytic input and no series reversion. It is recorded
because two of the three sources devote a section to it, and because it is the natural place to
notice that the multinomial formula \cref{p3:eq:reciprocal-multi} is the
$|\boldsymbol\alpha|$-linear shadow of the far harder reversion formulas of
\cref{p3:sec:full-inverse}.
\end{remark}

\section{What ``the inverse'' of a discrete sequence means here}
\label{p3:sec:meaning}

\begin{proposition}[Strict monotonicity]\label{p3:prop:strict-increase}
$p(n)<p(n+1)$ for every $n\ge1$.
\end{proposition}

\begin{proof}
Adjoining a part $1$ injects the partitions of $n$ into those of $n+1$, and the one-part partition
$(n+1)$ is not in the image.
\end{proof}

By \cref{p0:def:three-inverses} there are three distinct objects, and this chapter needs all
three.

\begin{enumerate}
\item \textbf{The integer staircase} (the generalised inverse)
      \begin{equation}\label{p3:eq:staircase-def}
       p^{\leftarrow}(y)=\min\{n\ge1:p(n)\ge y\},\qquad y\ge1 .
      \end{equation}
Well defined by \cref{p3:prop:strict-increase}. It is integer valued with unit jumps, so no smooth
expansion can represent it to $\littleo(1)$.
\item \textbf{The functional inverse of a chosen interpolation.}  Any strictly
increasing smooth $\mathcal{P}$ with $\mathcal{P}(n)=p(n)$ has an inverse $\mathcal{N}$, and
$p^{\leftarrow}(y)=\lceil\mathcal{N}(y)\rceil$ (\cref{p3:thm:staircase}). \cref{p3:def:Atilde}
supplies the canonical choice used here.
\item \textbf{The inverse along the range.}  If one knows in advance that
$y=p(n)$ for an integer $n$, the target is that integer. A smooth formula that is accurate to
$\littleo(1/2)$ recovers it by rounding (\cref{p3:thm:rounding}), and the accuracy available is
exponentially small, not merely $\littleo(1)$.
\end{enumerate}

\begin{definition}[The Rademacher interpolation]\label{p3:def:Pcal}
For real $\nu>1/24$ put $\mu(\nu)=\tfrac\pi6\sqrt{24\nu-1}$ and
\begin{equation}\label{p3:eq:Pcal}
 \mathcal{P}(\nu)=\frac{2\kappa}{\mu(\nu)^{2}}\sum_{k\ge1}
 \frac{\widetilde A_{k}(\nu)}{\sqrt k}\,g\!\left(\frac{\mu(\nu)}{k}\right),
\end{equation}
with $g$ as in \cref{p3:eq:g-kernel} and $\widetilde A_{k}$ as in \cref{p3:def:Atilde}.
\end{definition}

\begin{proposition}[Interpolation, regularity, eventual monotonicity]
\label{p3:prop:Pcal}
The series \cref{p3:eq:Pcal} and all its derivatives converge locally uniformly on
$(1/24,\infty)$; $\mathcal{P}$ is real analytic there and $\mathcal{P}(n)=p(n)$ for every integer
$n\ge1$. Moreover there is $\nu_{0}$ with $\mathcal{P}'(\nu)>0$ for $\nu\ge\nu_{0}$, so
$\mathcal{P}$ has a real inverse $\mathcal{N}$ on $[\mathcal{P}(\nu_{0}),\infty)$.
\end{proposition}

\begin{proof}
Fix a compact $I\subset(1/24,\infty)$ and a complex neighbourhood of it on which $\mu$ is
holomorphic and bounded. By \cref{p3:lem:Atilde} every $\nu$-derivative of $\widetilde A_{k}$ is
$\bigO_{r}(k)$ uniformly in $k$ there, while $g(\mu/k)=\bigO(k^{-2})$ uniformly with all
parameter derivatives (from \cref{p3:eq:g-kernel} and $\mu$ bounded). Hence the $k$th summand of
\cref{p3:eq:Pcal} and each of its derivatives is $\bigO(k\cdot k^{-1/2}\cdot
k^{-2})=\bigO(k^{-3/2})$, and the Weierstrass $M$-test gives locally uniform convergence and
legitimises termwise differentiation. Equality at integers is \cref{p3:lem:Atilde} combined with
\cref{p3:thm:exact-sectors}.

Monotonicity: write $\mathcal{P}(\nu)=F(\mu(\nu))(1+\mathcal{R}(\mu(\nu);\nu))$ as in
\cref{p3:prop:relative} with $\widetilde A_{k}$ in place of $A_{k}$. We have
$\dd\mu/\dd\nu=\lambda^{2}/(2\mu)>0$ and $\dd\log F/\dd\mu=1-3/\mu+1/(\mu-1)\to1$. Differentiating
$\mathcal{R}$ term by term and applying \cref{p3:thm:tail} with $K=1$ together with
\cref{p3:lem:Atilde} shows $\mathcal{R}+\partial_{\mu}\mathcal{R}=\bigO(\mu^{3}\e^{-\mu/2})$;
note that differentiating $\widetilde A_{k}(\nu(\mu))$ costs a factor $\bigO(k\mu)$, absorbed
by replacing $k^{-3/2}$ with $k^{-1/2}$ in the tail sum and hence by one more power of $\mu$. Thus
$\mathcal{P}'(\nu)=F'(\mu)\mu'(\nu)\bigl(1+\bigO(\mu^{3}\e^{-\mu/2})\bigr)$, which is
positive for large $\nu$.
\end{proof}

\begin{theorem}[Staircase from the smooth inverse]\label{p3:thm:staircase}
For all sufficiently large real $y$,
\begin{equation}\label{p3:eq:staircase}
 p^{\leftarrow}(y)=\bigl\lceil\mathcal{N}(y)\bigr\rceil,
\end{equation}
and $\mathcal{N}(p(n))=n$ exactly for every integer $n\ge\nu_{0}$.
\end{theorem}

\begin{proof}
This is the specialisation of \cref{p0:thm:staircase} to the unit-spaced lattice $\Int$, whose
separation hypothesis is \cref{p3:prop:strict-increase}. Concretely: let $m=p^{\leftarrow}(y)$, so
$p(m-1)<y\le p(m)$. Since $\mathcal{P}$ is increasing on $[\nu_{0},\infty)$ and
$\mathcal{P}(j)=p(j)$, we get $m-1<\mathcal{N}(y)\le m$, whence $\lceil\mathcal{N}(y)\rceil=m$.
\end{proof}

\begin{remark}[The unavoidable sawtooth]\label{p3:rem:sawtooth}
For general real $y$ the difference $p^{\leftarrow}(y)-\mathcal{N}(y)$ is the ceiling defect, of
size $\bigO(1)$, with Fourier expansion
\begin{equation}\label{p3:eq:ceiling-fourier}
 \lceil v\rceil-v=\frac12+\frac1\pi\sum_{r\ge1}\frac{\sin(2\pi rv)}{r}
 \qquad (v\notin\Int),
\end{equation}
the series converging to the midpoint $1/2$ at integers where the true defect is $0$. Since the
first arithmetic correction to $\mathcal{N}$ is of size $y^{-1/2}$ (\cref{p3:thm:leading-chirp}),
\emph{exponential precision in $\mathcal{N}$ buys nothing for $p^{\leftarrow}$ at generic $y$}:
the order-one sawtooth is an intrinsic final sector of the discrete problem, not an error to be
removed. It matters only along the range $y=p(n)$, where the defect vanishes identically. This
warning is S3's and is the sharpest of the three.
\end{remark}

\section{The Lambert core and the algebraic inverse}\label{p3:sec:lambert}

Throughout this section and the next, $y$ denotes a large ordinate and
\begin{equation}\label{p3:eq:R-ell}
 R=\log\frac y\kappa,\qquad \ell=\log R .
\end{equation}

\subsection{The carrier}

\begin{proposition}[Lambert carrier; specialisation of
\cref{p0:thm:lambert-core}]\label{p3:prop:lambert-core} The large solution $X=X(y)$ of
$\kappa\e^{X}X^{-2}=y$, equivalently of
\begin{equation}\label{p3:eq:carrier-phase}
 X-2\log X=R,
\end{equation}
is
\begin{equation}\label{p3:eq:carrier}
 X(y)=-2\,\Wlam_{-1}\!\left(-\frac12\sqrt{\frac\kappa y}\right)
     =-2\,\Wlam_{-1}\!\left(-\frac{\pi}{\sqrt{24\sqrt3\,y}}\right).
\end{equation}
\end{proposition}

\begin{proof}
This is \cref{p0:thm:lambert-core} at the dictionary values $a=1$, $b=-2$, $C=\kappa$ of
\cref{p3:tab:dictionary}: $X=(b/a)\Wlam_{k}((a/b)\e^{R/b}) =-2\Wlam_{k}(-\tfrac12\e^{-R/2})$ and
$\e^{-R/2}=\sqrt{\kappa/y}$. The branch rule of \cref{p0:thm:lambert-core} selects $k=-1$ because
$X\to+\infty$ forces $-X/2\le-1$. The second form uses
$\tfrac12\sqrt\kappa=\tfrac12\pi/\sqrt{6\sqrt3} =\pi/\sqrt{24\sqrt3}$.
\end{proof}

\subsection{Flattening into logarithms}

\begin{theorem}[Carrier coefficient polynomials]\label{p3:thm:carrier-polys}
As $y\to\infty$,
\begin{equation}\label{p3:eq:carrier-flat}
 X\sim R+2\ell+\sum_{j\ge1}\frac{\Pi_{j}(\ell)}{R^{j}},
\end{equation}
in the sense of \cref{p0:thm:flattening}, where
\begin{equation}\label{p3:eq:Pi-lagrange}
 \Pi_{j}(\ell)=\frac{2^{j+1}}{j+1}\,\Coef{v^{j}}\bigl(\ell+\log(1+v)\bigr)^{j+1}
 =2^{j+1}\sum_{q=1}^{j}\frac{s(j,q)}{(j+1-q)!}\,\ell^{\,j+1-q},
\end{equation}
with $s(j,q)$ the signed Stirling numbers of the first kind,
$v(v-1)\cdots(v-j+1)=\sum_{q}s(j,q)v^{q}$. For each fixed $J$,
\begin{equation}\label{p3:eq:carrier-remainder}
 X=R+2\ell+\sum_{j=1}^{J}\frac{\Pi_{j}(\ell)}{R^{j}}
 +\bigO\!\left(\frac{\ell^{\,J+1}}{R^{J+1}}\right).
\end{equation}
Moreover $\deg\Pi_{j}=j$ with leading coefficient $(-1)^{j-1}2^{j+1}/j$, and $\Pi_{j}(0)=0$.
\end{theorem}

\begin{proof}
Write $X=R(1+v)$ in \cref{p3:eq:carrier-phase}: $R+Rv-2\log R-2\log(1+v)=R$, i.e.
\begin{equation}\label{p3:eq:carrier-fixed-point}
 v=z\bigl(\ell+\log(1+v)\bigr),\qquad z=\frac2R,
\end{equation}
a Lagrange fixed point $v=z\phi(v)$ with $\phi(v)=\ell+\log(1+v)$ and $\phi(0)=\ell\ne0$.
\Cref{p0:thm:LB} gives $v=\sum_{m\ge1}\frac{z^{m}}{m}\Coef{v^{m-1}}\phi(v)^{m}$; multiplying by
$R$ and putting $m=j+1$ (so $z^{m}R=2^{j+1}R^{-j}$) yields the first equality of
\cref{p3:eq:Pi-lagrange}. For the second, expand
$(\ell+\log(1+v))^{j+1}=\sum_{q}\binom{j+1}{q}\ell^{\,j+1-q}\log(1+v)^{q}$ and use
$\log(1+v)^{q}=q!\sum_{m\ge q}s(m,q)v^{m}/m!$; extracting $\Coef{v^{j}}$ keeps only $1\le q\le
j+1$, and the term $q=j+1$ contributes $s(j,j+1)=0$. The surviving coefficient of $\ell^{\,j+1-q}$
is $\frac{2^{j+1}}{j+1}\binom{j+1}{q}q!\,\frac{s(j,q)}{j!} =2^{j+1}\frac{s(j,q)}{(j+1-q)!}$, as
claimed. Since $q\ge1$, every monomial carries a positive power of $\ell$, so $\Pi_{j}(0)=0$; the
top power is $q=1$, with $s(j,1)=(-1)^{j-1}(j-1)!$, giving leading coefficient
$2^{j+1}(-1)^{j-1}(j-1)!/j!=(-1)^{j-1}2^{j+1}/j$. The remainder statement
\cref{p3:eq:carrier-remainder} is the general estimate of \cref{p0:thm:flattening} applied on the
scale $\ell/R\to0$; concretely, \cref{p3:eq:carrier-fixed-point} has an analytic solution
$v=v(z,\ell)$ near $z=0$ for $\ell$ in any fixed compact set, and $\ell/R\to0$ lets Taylor's
theorem with remainder be applied at order $J$.
\end{proof}

The first polynomials, recomputed here from the Lagrange form and independently from the Stirling
form (they agree identically), are
\begin{align}
 \Pi_{1}&=4\ell, &
 \Pi_{2}&=-4\ell(\ell-2), \nonumber\\
 \Pi_{3}&=\tfrac83\ell\,(2\ell^{2}-9\ell+6), &
 \Pi_{4}&=-\tfrac83\ell\,(3\ell^{3}-22\ell^{2}+36\ell-12),
 \label{p3:eq:Pi-first}\\
 \Pi_{5}&=\tfrac{16}{15}\ell\,(12\ell^{4}-125\ell^{3}+350\ell^{2}-300\ell+60),&
 \Pi_{6}&=-\tfrac{16}{15}\ell\,(20\ell^{5}-274\ell^{4}+1125\ell^{3}
 -1700\ell^{2}+900\ell-120),\nonumber
\end{align}
\begin{equation}\label{p3:eq:Pi-seven}
 \Pi_{7}=\tfrac{32}{105}\ell\,\bigl(120\ell^{6}-2058\ell^{5}+11368\ell^{4}
 -25725\ell^{3}+24500\ell^{2}-8820\ell+840\bigr),
\end{equation}
so that, explicitly,
\begin{equation}\label{p3:eq:carrier-explicit}
 X=R+2\ell+\frac{4\ell}{R}+\frac{-4\ell^{2}+8\ell}{R^{2}}
 +\frac{\tfrac{16}3\ell^{3}-24\ell^{2}+16\ell}{R^{3}}
 +\frac{-8\ell^{4}+\tfrac{176}3\ell^{3}-96\ell^{2}+32\ell}{R^{4}}+\cdots.
\end{equation}

\begin{remark}[Three printings of one polynomial family]\label{p3:rem:H-versus-R}
The sources print \cref{p3:eq:Pi-lagrange} in three different shapes, and it is worth recording
once that they are the same family, because a reader comparing them will otherwise suspect an
error.
\begin{itemize}
\item S1: $P_{j}(L)=2^{j+1}\sum_{q=1}^{j}\frac{s(j,q)}{(j+1-q)!}L^{j+1-q}$ ---
identical to \cref{p3:eq:Pi-lagrange}.
\item S2: $P_{m}(\ell)=\frac1{m!}\sum_{q=1}^{m}2^{q}(q-1)!\,s(m,q)
\binom{m}{q-1}(2\ell)^{m-q+1}$. The coefficient of $\ell^{\,m-q+1}$ is
$\frac{2^{q}(q-1)!}{m!}\binom{m}{q-1}2^{m-q+1} =\frac{2^{m+1}s(m,q)}{(m-q+1)!}$ after cancelling
$\binom{m}{q-1}=\frac{m!}{(q-1)!(m-q+1)!}$. Identical.
\item S3: $P_{r}(\ell)=\frac1{r!}\sum_{q=1}^{r}(q-1)!\,s(r,q)\binom{r}{q-1}
\ell^{\,r-q+1}$, appearing in the expansion as $2^{r+1}P_{r}(\ell)$. The same cancellation gives
$2^{r+1}P_{r}=\Pi_{r}$. Identical.
\end{itemize}
\textbf{But the variables differ, and this is a genuine disagreement.} S1 and S2 expand in
$R=\log(y/\kappa)$ with $\ell=\log R$; S3 expands in $H=\log(4y/\kappa)=R+2\log2$ with
$\ell_{\mathrm{S3}}=\log(H/2)$, which is the natural variable for the substitution $X=2w$, $w-\log
w=H/2$. The two series
\[
X=R+2\log R+\sum_{j}\frac{\Pi_{j}(\log R)}{R^{j}} \qquad\text{and}\qquad X=H+2\log\tfrac
H2+\sum_{j}\frac{\Pi_{j}(\log\tfrac H2)}{H^{j}}
\]
are \emph{different} asymptotic expansions of the same function, agreeing order by order only
after re-expanding $H=R+2\log2$. Both are correct; mixing them is not. This volume uses the S1/S2
variables $R,\ell$ throughout, because $R$ makes the phase equation exactly
\cref{p3:eq:carrier-phase} with no additive constant. We verified all of S3's displayed
coefficients (its $u$- and $N_{0}$-expansions in $H$, through $H^{-4}$ and $H^{-3}$ respectively)
by symbolic re-expansion; they are correct as printed in S3's own variables.
\end{remark}

\subsection{Restoring the dominant amplitude: the centred reversion}
\label{p3:sub:core}

The carrier $X$ solves $y=\kappa\e^{X}X^{-2}$; the \emph{dominant core} $\mu_{0}=\mu_{0}(y)$
solves the full dominant sector,
\begin{equation}\label{p3:eq:core-equation}
 y=F(\mu_{0})=\kappa\,\e^{\mu_{0}}\mu_{0}^{-2}\left(1-\frac1{\mu_{0}}\right),
 \qquad\text{i.e.}\qquad
 \Phi(\mu_{0})=R,
\end{equation}
where the \emph{dominant phase} is
\begin{equation}\label{p3:eq:Phi}
 \Phi(\mu)=\mu-2\log\mu+\log\!\left(1-\frac1\mu\right)
          =\mu-3\log\mu+\log(\mu-1),\qquad \mu>1 .
\end{equation}

\begin{lemma}[Derivatives of the dominant phase]\label{p3:lem:Phi-derivatives}
For $\mu>1$,
\begin{equation}\label{p3:eq:Phi-prime}
 \Phi'(\mu)=1-\frac3\mu+\frac1{\mu-1}
          =1-\frac2\mu+\frac1{\mu(\mu-1)}
          =\frac{\mu^{2}-3\mu+3}{\mu(\mu-1)}>0,
\end{equation}
and for $m\ge1$
\begin{equation}\label{p3:eq:Phi-m}
 \Phi^{(m)}(\mu)=\mathbf 1_{m=1}
 +(-1)^{m-1}(m-1)!\left(\frac1{(\mu-1)^{m}}-\frac3{\mu^{m}}\right).
\end{equation}
In particular $\Phi'(\mu)=1+\bigO(\mu^{-1})$, and $\mu^{2}-3\mu+3>0$ for all real $\mu$ since
its discriminant is $-3$. Writing $\varrho_{\pm}=\tfrac12(3\pm\ii\sqrt3)$ for the roots of
$\mu^{2}-3\mu+3$,
\begin{equation}\label{p3:eq:F-prime}
 F'(\mu)=\kappa\,\e^{\mu}\,\frac{\mu^{2}-3\mu+3}{\mu^{4}},
 \qquad
 \frac{\dd^{m}}{\dd\mu^{m}}\log F'(\mu)
 =\mathbf 1_{m=1}+(-1)^{m-1}(m-1)!
 \left(\frac1{(\mu-\varrho_{+})^{m}}+\frac1{(\mu-\varrho_{-})^{m}}
 -\frac4{\mu^{m}}\right).
\end{equation}
\end{lemma}

\begin{proof}
Differentiate \cref{p3:eq:Phi} in the form $\mu-3\log\mu+\log(\mu-1)$ and clear denominators for
the three displayed shapes of $\Phi'$; \cref{p3:eq:Phi-m} is $m-1$ further differentiations of
$-3\mu^{-1}+(\mu-1)^{-1}$. For \cref{p3:eq:F-prime}, $F=\kappa\e^{\mu}(\mu-1)\mu^{-3}$ gives
$F'=F\bigl(1+\tfrac1{\mu-1}-\tfrac3\mu\bigr)=F\,\Phi'(\mu)$, and
$F\Phi'=\kappa\e^{\mu}\frac{\mu-1}{\mu^{3}}\cdot\frac{\mu^{2}-3\mu+3}{\mu(\mu-1)}$. Then $\log
F'=\log\kappa+\mu+\log(\mu-\varrho_{+})+\log(\mu-\varrho_{-}) -4\log\mu$, and differentiate.
\end{proof}

\begin{theorem}[The dominant-core correction; specialisation of
\cref{p0:thm:lambert-centered}]\label{p3:thm:core-correction} Put $z=X^{-1}$. There is a unique
$U\in z\,\Rat[[z]]$ with
\begin{equation}\label{p3:eq:U-implicit}
 U=3\log(1+zU)-\log(1-z+zU),
\end{equation}
and $U$ is analytic in a neighbourhood of $z=0$. The dominant core is
\begin{equation}\label{p3:eq:core}
 \mu_{0}(y)=X(y)+U\bigl(X(y)^{-1}\bigr),\qquad U(z)=\sum_{m\ge1}c_{m}z^{m}.
\end{equation}
Moreover \cref{p3:eq:U-implicit} is exactly the fixed-point equation $U=\Psi(z,U)$ of
\cref{p0:thm:lambert-centered} at the dictionary values of \cref{p3:tab:dictionary}, so
\begin{equation}\label{p3:eq:c-lagrange}
 c_{m}=\sum_{q=1}^{m}\frac1q\,
 \Coef{z^{m}v^{q-1}}\Bigl(3\log(1+zv)-\log(1-z+zv)\Bigr)^{q}
\end{equation}
by \cref{p0:lem:lagrange-fp}, and equivalently
\begin{equation}\label{p3:eq:c-recurrence-core}
 c_{m}=\Coef{z^{m}}\Bigl\{3\log\bigl(1+zU_{<m}\bigr)
 -\log\bigl(1-z+zU_{<m}\bigr)\Bigr\},
 \qquad U_{<m}(z)=\sum_{j=1}^{m-1}c_{j}z^{j},
\end{equation}
a triangular recurrence over $\Rat$.
\end{theorem}

\begin{proof}
Put $\mu=X+U$ and divide \cref{p3:eq:core-equation} by the carrier identity
$y=\kappa\e^{X}X^{-2}$:
\[
\e^{U}(1+zU)^{-2}\left(1-\frac{z}{1+zU}\right)=1, \qquad z=\frac1X ,
\]
since $\mu=X(1+zU)$ and $1/\mu=z/(1+zU)$. As $1-\frac{z}{1+zU}=\frac{1-z+zU}{1+zU}$, taking
logarithms gives $U-3\log(1+zU)+\log(1-z+zU)=0$, i.e.\ \cref{p3:eq:U-implicit}. Write $G(z,U)$ for
the left side: $G(0,0)=0$ and $\partial_{U}G(0,0)=1$, so the analytic implicit function theorem
gives a unique analytic solution near $0$, necessarily with rational Taylor coefficients since $G$
has rational ones.

For the identification with \cref{p0:thm:lambert-centered}: that theorem's fixed-point map at
parameters $a,b,Q$ is $\Psi(t,u)=-\frac1a\bigl[b\log(1+tu)+Q\bigl(\tfrac{t}{1+tu}\bigr)\bigr]$. At
$a=1$, $b=-2$, $Q(s)=\log(1-s)$ this is
\[
\Psi(t,u)=2\log(1+tu)-\log\!\left(1-\frac{t}{1+tu}\right) =2\log(1+tu)-\log(1-t+tu)+\log(1+tu)
=3\log(1+tu)-\log(1-t+tu),
\]
which is the right side of \cref{p3:eq:U-implicit}. \Cref{p3:eq:c-lagrange} is then
\cref{p0:lem:lagrange-fp} verbatim. \Cref{p3:eq:c-recurrence-core} holds because in
\cref{p3:eq:U-implicit} every occurrence of $U$ on the right is multiplied by an explicit $z$, so
$\Coef{z^{m}}$ of the right side involves only $c_{1},\dots,c_{m-1}$.
\end{proof}

\begin{table}[htbp]
\centering
\caption{The centred-reversion coefficients $c_{m}$ of
\cref{p3:eq:core}, recomputed here in exact rational arithmetic from
\cref{p3:eq:c-recurrence-core} (truncated power-series arithmetic over
$\Rat$, $M=20$).  All eighteen agree with the tables printed by the three
sources, which call them $b_{m}$ (S1), $a_{m}$ (S2) and $d_{r}$ (S3).}
\label{p3:tab:cm}
\small
\begin{tabular}{@{}llll@{}}
\toprule
$c_{1}=1$ & $c_{2}=\tfrac52$ & $c_{3}=\tfrac{13}3$ & $c_{4}=\tfrac{53}{12}$\\
$c_{5}=-\tfrac{14}5$ & $c_{6}=-\tfrac{763}{30}$ & $c_{7}=-\tfrac{2179}{35}$
 & $c_{8}=-\tfrac{42289}{630}$\\
$c_{9}=\tfrac{132973}{1260}$ & $c_{10}=\tfrac{3449711}{5040}$
 & $c_{11}=\tfrac{29813143}{18480}$ & $c_{12}=\tfrac{496569589}{356400}$\\
$c_{13}=-\tfrac{4334124833}{926640}$ & $c_{14}=-\tfrac{509723063851}{21621600}$
 & $c_{15}=-\tfrac{298798016297}{5896800}$
 & $c_{16}=-\tfrac{933244850273}{33633600}$\\
$c_{17}=\tfrac{4993013022640963}{23156733600}$
 & $c_{18}=\tfrac{208843210189406957}{231567336000}$ & &\\
\bottomrule
\end{tabular}
\end{table}

Thus
\begin{equation}\label{p3:eq:core-explicit}
 \mu_{0}=X+\frac1X+\frac5{2X^{2}}+\frac{13}{3X^{3}}+\frac{53}{12X^{4}}
 -\frac{14}{5X^{5}}-\frac{763}{30X^{6}}-\frac{2179}{35X^{7}}+\cdots.
\end{equation}

\begin{remark}[Implementation check]\label{p3:rem:U-residual}
The cheapest correctness test for a table of $c_{m}$ is the residual identity: with
$U_{M}=\sum_{m\le M}c_{m}z^{m}$,
\begin{equation}\label{p3:eq:U-residual}
 U_{M}-3\log(1+zU_{M})+\log(1-z+zU_{M})=\bigO(z^{M+1}).
\end{equation}
We verified \cref{p3:eq:U-residual} at $M=12$ symbolically: the residual is $\bigO(z^{13})$
exactly. This test is far more reliable than comparing printed digits, and it is the one we used
to certify \cref{p3:tab:cm}.
\end{remark}

\subsection{The algebraic inverse index}

\begin{definition}\label{p3:def:N0}
$\displaystyle N_{0}(y)=\frac1{24}+\frac{3\,\mu_{0}(y)^{2}}{2\pi^{2}}$, the \emph{algebraic}
(phase-independent) inverse index; and $\displaystyle
N_{X}(y)=\frac1{24}+\frac{3X(y)^{2}}{2\pi^{2}}$, the carrier index.
\end{definition}

\begin{corollary}[All algebraic coefficients of the inverse index]
\label{p3:cor:N0-expansion}
With $c_{m}$ as in \cref{p3:tab:cm}, put
\begin{equation}\label{p3:eq:gamma}
 \gamma_{j}=3c_{j+1}+\frac32\sum_{r=1}^{j-1}c_{r}c_{j-r}\qquad(j\ge1).
\end{equation}
Then
\begin{equation}\label{p3:eq:N0-expansion}
 N_{0}(y)=\frac{3X^{2}}{2\pi^{2}}
 +\left(\frac1{24}+\frac3{\pi^{2}}\right)
 +\frac1{\pi^{2}}\sum_{j\ge1}\frac{\gamma_{j}}{X^{j}},
\end{equation}
with
\begin{equation}\label{p3:eq:gamma-first}
 \gamma_{1}=\tfrac{15}2,\quad \gamma_{2}=\tfrac{29}2,\quad
 \gamma_{3}=\tfrac{83}4,\quad \gamma_{4}=\tfrac{559}{40},\quad
 \gamma_{5}=-\tfrac{611}{20},\quad \gamma_{6}=-\tfrac{112459}{840},
\end{equation}
\begin{equation}\label{p3:eq:gamma-more}
 \gamma_{7}=-\tfrac{101329}{420},\quad
 \gamma_{8}=-\tfrac{228677}{3360},\quad
 \gamma_{9}=\tfrac{1709167}{1680},\quad
 \gamma_{10}=\tfrac{325098619}{92400}.
\end{equation}
\end{corollary}

\begin{proof}
Insert $\mu_{0}=X+U$ into \cref{p3:def:N0}. The cross term $\frac{3}{2\pi^{2}}\cdot
2XU=\frac{3}{\pi^{2}}\sum_{m\ge1}c_{m}X^{1-m}$ contributes $3c_{1}/\pi^{2}=3/\pi^{2}$ to the
constant and $3c_{j+1}/(\pi^{2}X^{j})$ at order $X^{-j}$; the square $\frac{3}{2\pi^{2}}U^{2}$
contributes $\frac{3}{2\pi^{2}}\bigl(\sum_{r=1}^{j-1}c_{r}c_{j-r}\bigr)X^{-j}$ and nothing at
order $X^{0}$ or above. The listed values were recomputed from \cref{p3:eq:gamma};
$\gamma_{1},\dots,\gamma_{8}$ agree with the tables of all three sources, and
$\gamma_{9},\gamma_{10}$ are new here.
\end{proof}

\begin{corollary}[The carrier bias]\label{p3:cor:carrier-bias}
$\displaystyle N_{0}(y)-N_{X}(y)=\frac3{\pi^{2}}+\bigO(X^{-1})$, and
\begin{equation}\label{p3:eq:carrier-bias}
 \frac3{\pi^{2}}=0.30396355092701331433\ldots
\end{equation}
\end{corollary}

\begin{proof}
Subtract $N_{X}=\tfrac1{24}+3X^{2}/(2\pi^{2})$ from \cref{p3:eq:N0-expansion}: the
$3X^{2}/(2\pi^{2})$ and $\tfrac1{24}$ cancel, leaving
$3/\pi^{2}+\pi^{-2}\sum_{j\ge1}\gamma_{j}X^{-j}$. The constant was evaluated here to $20$ digits.
\end{proof}

\begin{remark}[Phase locking: do not expand the arithmetic phase about the
carrier]\label{p3:rem:phase-lock} This is the single most consequential structural point of the
inverse problem, and only S2 states it explicitly.

The amplitudes $\widetilde A_{k}$ are \emph{periodic on the index scale}, with period $k$. The
core correction is small in the phase variable, $\mu_{0}-X=\bigO(X^{-1})$; but $\dd
n/\dd\mu=3\mu/\pi^{2}\asymp X$, so the induced change in the index is $\bigO(1)$ --- exactly
$3/\pi^{2}$ in the limit, by \cref{p3:cor:carrier-bias}. A shift of $0.304$ in the index is
\emph{not} a small perturbation of $\cos(\pi\nu)$. Consequently:
\begin{itemize}
\item Taylor-expanding $\widetilde A_{k}(N(\mu_{0}))$ around $N_{X}$ does not
begin with a small correction, and produces the wrong phase already at the leading exponentially
small order.
\item Every arithmetic amplitude must be evaluated at the \emph{full algebraic
core} $N_{0}(y)$, not at $N_{X}(y)$ and not at any truncation of $N_{0}$ coarser than
$\littleo(1)$. By \cref{p3:eq:N0-expansion}, truncating the $\gamma$-series after $M$ terms
leaves an index error $\bigO(X^{-M})$, so $M\ge1$ already suffices for $\littleo(1)$; $M=0$
(i.e.\ dropping the constant $3/\pi^{2}$) does not.
\item Only the \emph{subsequent} exponentially small displacement may be
Taylor-expanded; that is what \cref{p3:sec:full-inverse} does.
\end{itemize}
The numerical signature of the effect is the first column of \cref{p3:tab:inverse}, which tends to
$-3/\pi^{2}$ and not to $0$.
\end{remark}

\begin{remark}[Elementary-logarithm form]\label{p3:rem:N0-in-R}
Substituting \cref{p3:eq:carrier-flat} into \cref{p3:eq:N0-expansion} gives the pure logarithmic
block
\begin{equation}\label{p3:eq:N0-in-R}
 N_{0}(y)=\frac{3R^{2}}{2\pi^{2}}+\frac{6R\ell}{\pi^{2}}
 +\frac{6\ell^{2}+12\ell+3}{\pi^{2}}+\frac1{24}
 +\frac{3(8\ell^{2}+16\ell+5)}{2\pi^{2}R}
 +\bigO\!\left(\frac{\ell^{3}}{R^{2}}\right),
\end{equation}
with $R=\log(y/\kappa)$ and $\ell=\log R$. We recomputed this block symbolically; it agrees with
S2's display (which stops at the constant) and, after the change of variable $H=R+2\log2$, with
S3's. For \emph{numerical} work \cref{p3:eq:N0-in-R} is markedly inferior to evaluating
\cref{p3:eq:carrier} and \cref{p3:eq:core-explicit} directly, and it is recorded only to exhibit
the transseries shape: $N_{0}$ is a polynomial of degree $2$ in $R$ with coefficients polynomial
in $\ell$, plus a descending series in $R^{-1}$ with $\deg_{\ell}$ growing by one per order.
\end{remark}

\section{The full arithmetic inverse transseries}\label{p3:sec:full-inverse}

Everything so far is phase-independent: $X$, $\mu_{0}$, $N_{0}$ are functions of $y$ alone. The
remaining correction is exponentially small and arithmetically oscillatory, and the oscillation
must be frozen before ``analytic inverse'' has a meaning. The three sources do this in two
different ways, both correct.

\subsection{Two ways to freeze the phase}\label{p3:sub:two-freezings}

\begin{definition}[Residue-frozen finite level; S1]\label{p3:def:frozen}
For $K\ge1$ let $\Lambda_{K}=\operatorname{lcm}(1,2,\dots,K)$ and fix a residue
$r\in\Int/\Lambda_{K}\Int$. Since $A_{k}$ has period $k\mid\Lambda_{K}$ for $k\le K$, the numbers
$A_{k}(r)$ are well defined, and
\begin{equation}\label{p3:eq:PKr}
 P_{K,r}(\mu)=\frac{\kappa}{\mu^{2}}\sum_{k=1}^{K}\frac{A_{k}(r)}{\sqrt k}
 \left[\e^{\mu/k}\Bigl(1-\frac k\mu\Bigr)
      +\e^{-\mu/k}\Bigl(1+\frac k\mu\Bigr)\right]
\end{equation}
is a genuine analytic function of $\mu$ with constant coefficients. For large $\mu$ it is
increasing (the $k=1$ positive term dominates), and we write $\mu_{K,r}(y)$ for its large inverse
branch.
\end{definition}

\begin{definition}[Interpolated level; S2, S3]\label{p3:def:interpolated}
$\mathcal{P}_{K}(\mu)$ is \cref{p3:eq:PKr} with $A_{k}(r)$ replaced by $\widetilde A_{k}(n(\mu))$,
$n(\mu)=\tfrac1{24}+3\mu^{2}/(2\pi^{2})$; its large inverse branch is written $\mu_{K}(y)$.
\end{definition}

\begin{proposition}[The two conventions agree below action one]
\label{p3:prop:two-conventions}
Fix $K$, let $y$ be large, and let $r\equiv\lfloor N_{0}(y)\rceil \pmod{\Lambda_{K}}$, where
$\lfloor\cdot\rceil$ is nearest-integer rounding. Assume $|N_{0}(y)-r|\le\tfrac12$ modulo
$\Lambda_{K}$ has been resolved (which holds along the range $y=p(n)$ for large $n$ by
\cref{p3:thm:rounding}). Then every inverse transmonomial of total action strictly less than $1$
is the same for $\mu_{K,r}$ and $\mu_{K}$; the two first differ at total action $\ge1$.
\end{proposition}

\begin{proof}
Both inverses are constructed as $\mu_{0}+v$ with $v$ exponentially small, and in both the
coefficient of the transmonomial $\e^{-\alpha_{j}\mu_{0}}$ at first order is
$-r_{j}(\mu_{0})/\Phi'(\mu_{0})$ (\cref{p3:thm:first-order} below), where $r_{j}$ differs between
the two conventions only in the amplitude, $A_{k}(r)$ versus $\widetilde
A_{k}(n(\mu_{0}))=\widetilde A_{k}(N_{0}(y))$. Now $\widetilde A_{k}$ is entire and $k$-periodic
with $\widetilde A_{k}(r)=A_{k}(r)$, so $|\widetilde A_{k}(N_{0})-A_{k}(r)|\le \sup|\widetilde
A_{k}'|\cdot|N_{0}-r|$. Along the range, $|N_{0}(p(n))-n|=\bigO(\mu^{5/2}\e^{-\mu/2})$ by
\cref{p3:cor:core-inverse-error}, so the amplitude discrepancy is
$\bigO(k\mu^{5/2}\e^{-\mu/2})$ and its contribution to the sector carries the extra factor
$\e^{-\mu/2}$, i.e.\ additional action $\tfrac12$. Since the smallest primary action is
$\tfrac12$, the first affected total action is $\tfrac12+\tfrac12=1$. Higher-order terms of the
reversion have total action $\ge1$ already.
\end{proof}

\begin{remark}[Which convention to prefer]\label{p3:rem:which-convention}
The residue-frozen convention is intrinsic: it invents no interpolation, and at each finite level
it reduces the problem to an ordinary analytic reversion with \emph{constant} coefficients. Its
cost is that the inverse is defined only after the residue class is known, which
\cref{p3:thm:rounding} supplies. The interpolated convention defines a single function
$\mathcal{N}$ for all real $y$ and connects directly to the staircase (\cref{p3:thm:staircase});
its cost is the interpolation dependence of \cref{p3:rem:interp-nonunique} and the appearance of
the amplitude's own $\mu$-derivatives in every higher coefficient. This chapter develops the
reversion in a form that covers both: all formulas below are stated for a family of \emph{blocks}
$E_{j}(\mu)$, and the two conventions are two choices of block.
\end{remark}

\subsection{Blocks, multi-indices, and the reversion equation}
\label{p3:sub:blocks}

\begin{definition}[Sector blocks]\label{p3:def:blocks}
Fix $K\ge1$ and put
\begin{equation}\label{p3:eq:JK}
 J_{K}=\{(1,-1)\}\cup\{(k,\sigma):2\le k\le K,\ \sigma=\pm1\}.
\end{equation}
For $j=(k,\sigma)\in J_{K}$ set $\alpha_{j}=1-\sigma/k$ as in \cref{p3:eq:action}, and define the
\emph{block}
\begin{equation}\label{p3:eq:block}
 E_{j}(\mu)=\frac{\kappa}{\mu^{2}}\frac{a_{k}(\mu)}{\sqrt k}
 \left(1-\frac{\sigma k}{\mu}\right)\e^{\sigma\mu/k}
 =F(\mu)\,r_{j}(\mu)\,\e^{-\alpha_{j}\mu},
 \qquad
 r_{j}(\mu)=\frac{a_{k}(\mu)}{\sqrt k}\,\frac{1-\sigma k/\mu}{1-1/\mu},
\end{equation}
where $a_{k}(\mu)$ is either the constant $A_{k}(r)$ (residue-frozen convention,
\cref{p3:def:frozen}) or $\widetilde A_{k}(n(\mu))$ (interpolated convention,
\cref{p3:def:interpolated}); $a_{1}\equiv1$ in both. A block may also be the whole paired tail
$\mathcal{R}_{>K}\cdot F$ of \cref{p3:eq:R-exact}, or the zeta-condensation block of
\cref{p3:thm:zeta-condensation}, treated as one indivisible symbol.
\end{definition}

Introduce independent formal variables $\varepsilon_{j}$, $j\in J_{K}$, and write
$\boldsymbol\varepsilon^{\boldsymbol m}=\prod_{j}\varepsilon_{j}^{m_{j}}$, $|\boldsymbol
m|=\sum_{j}m_{j}$, $\boldsymbol m!=\prod_{j}m_{j}!$ for multi-indices $\boldsymbol
m\in\Nat^{J_{K}}$. The forward object at level $K$ is
\begin{equation}\label{p3:eq:forward-blocks}
 \mathcal{P}_{K}(\mu;\boldsymbol\varepsilon)
 =F(\mu)+\sum_{j\in J_{K}}\varepsilon_{j}E_{j}(\mu),
\end{equation}
and the inverse problem is to solve $\mathcal{P}_{K}(\mu;\boldsymbol\varepsilon)=y$ for $\mu$ near
$\mu_{0}=F^{-1}(y)$, as a formal series in $\boldsymbol\varepsilon$ which is finally specialised
at $\varepsilon_{j}=1$. Two equivalent engines are available; both are finite at every prescribed
multi-index.

\subsection{Engine I: additive Lagrange--B\"urmann}

\begin{theorem}[Additive reversion]\label{p3:thm:additive-LB}
Let $F$ and $E$ be analytic near $\mu_{0}$ with $F'(\mu_{0})\ne0$ and $F(\mu_{0})=y$. Write
$\mathcal{D}=\frac1{F'(\mu)}\frac{\dd}{\dd\mu}$ evaluated along $\mu=\mu_{0}(y)$, i.e.\
$\mathcal{D}=\dd/\dd y$ in the coordinate $y$. Then the solution of $F(\mu)+\varepsilon E(\mu)=y$
near $\mu_{0}$ is
\begin{equation}\label{p3:eq:additive-LB}
 \mu(y;\varepsilon)=\mu_{0}(y)
 +\sum_{M\ge1}\frac{(-\varepsilon)^{M}}{M!}\,
 \mathcal{D}^{M-1}\!\left(\frac{E(\mu_{0})^{M}}{F'(\mu_{0})}\right).
\end{equation}
Consequently, for the multivariate problem \cref{p3:eq:forward-blocks}, the coefficient of
$\boldsymbol\varepsilon^{\boldsymbol m}$ with $M=|\boldsymbol m|\ge1$ in $\mu-\mu_{0}$ is
\begin{equation}\label{p3:eq:multi-LB}
 V_{\boldsymbol m}(\mu_{0})=\frac{(-1)^{M}}{\boldsymbol m!}\,
 \mathcal{D}^{M-1}\!\left(\frac{E^{\boldsymbol m}(\mu_{0})}{F'(\mu_{0})}\right)
 =\frac{(-1)^{M}(M-1)!}{\boldsymbol m!}\,
 \operatorname*{Res}_{\mu=\mu_{0}}\frac{E^{\boldsymbol m}(\mu)}{(F(\mu)-F(\mu_{0}))^{M}},
\end{equation}
where $E^{\boldsymbol m}=\prod_{j}E_{j}^{m_{j}}$.
\end{theorem}

\begin{proof}
Let $\gamma$ be a small positively oriented circle about $\mu_{0}$ enclosing no other zero of
$F(\cdot)-y$. For small $|\varepsilon|$, Rouch\'e gives exactly one zero $\mu(y;\varepsilon)$ of
$F-y+\varepsilon E$ inside $\gamma$, and the argument principle gives
\[
\mu(y;\varepsilon)=\frac1{2\pi\ii}\int_{\gamma} \zeta\,\dd_{\zeta}\log\bigl(F(\zeta)-y+\varepsilon
E(\zeta)\bigr).
\]
Subtract the same expression at $\varepsilon=0$, expand $\log(1+\varepsilon
E/(F-y))=\sum_{M\ge1}\frac{(-1)^{M-1}}{M} \varepsilon^{M}E^{M}(F-y)^{-M}$, and integrate by parts
once (which converts $\zeta\,\dd_{\zeta}(\cdot)$ into $-(\cdot)\dd\zeta$). The coefficient of
$\varepsilon^{M}$ is
\[
\frac{(-1)^{M}}{M}\operatorname*{Res}_{\zeta=\mu_{0}}
\frac{E(\zeta)^{M}}{(F(\zeta)-F(\mu_{0}))^{M}} .
\]
Now change coordinate $w=F(\zeta)$, $\zeta=\mu_{0}(w)$, $\dd\zeta=\mu_{0}'(w)\dd w$:
\[
\operatorname*{Res}_{\zeta=\mu_{0}}\frac{E(\zeta)^{M}}{(F(\zeta)-y)^{M}}\dd\zeta
=\operatorname*{Res}_{w=y}\frac{E(\mu_{0}(w))^{M}\mu_{0}'(w)}{(w-y)^{M}}\dd w
=\frac1{(M-1)!}\frac{\dd^{M-1}}{\dd y^{M-1}}
\left(\frac{E(\mu_{0}(y))^{M}}{F'(\mu_{0}(y))}\right),
\]
which is \cref{p3:eq:additive-LB} after dividing by $M$ and using $\dd/\dd y=\mathcal{D}$. The
multivariate statement follows by replacing $E$ with $\sum_{j}\varepsilon_{j}E_{j}$ and expanding
the $M$th power by the multinomial theorem: the coefficient $M!/\boldsymbol m!$ cancels the $M!$
in \cref{p3:eq:additive-LB}.

For the partition blocks this is more than formal. With $E=\sum_{j}E_{j}+\mathcal{R}_{>K}F$ we
have $E/F=\bigO(\mu^{C}\e^{-\mu/2})$ together with the same bound for finitely many
derivatives, by \cref{p3:thm:tail} and \cref{p3:lem:Atilde}; so on a fixed relative neighbourhood
of a large $\mu_{0}$, Rouch\'e applies for all $|\varepsilon|\le1+\eta$ with some $\eta>0$, and
the Taylor series in $\varepsilon$ converges at $\varepsilon=1$.
\end{proof}

\begin{remark}[$\mathcal{D}$ is iterated, not a plain derivative]
\label{p3:rem:D-iterated}
In \cref{p3:eq:additive-LB} the operator $\mathcal{D}^{M-1}$ is the $(M-1)$st power of
$\frac1{F'}\frac{\dd}{\dd\mu}$. Replacing it by $\frac1{(F')^{M-1}}\frac{\dd^{M-1}}{\dd\mu^{M-1}}$
after ``pulling out'' the powers of $F'$ is wrong from $M=3$ onwards, because $F'$ is not
constant. S3 flags this explicitly; it is the most common implementation bug in this formula.
\end{remark}

\begin{corollary}[Bell realisation]\label{p3:cor:multi-Bell}
Let $H_{\boldsymbol m}=E^{\boldsymbol m}/F'$ and $\Lambda^{(\boldsymbol
m)}_{s}=\frac{\dd^{s}}{\dd\mu^{s}}\log H_{\boldsymbol m}$, and let
$\mu_{0}^{(s)}=\dd^{s}\mu_{0}/\dd y^{s}$ be generated by
\begin{equation}\label{p3:eq:inverse-derivatives}
 \mu_{0}^{(1)}=\frac1{F'(\mu_{0})},\qquad
 \mu_{0}^{(s)}=-\frac1{F'(\mu_{0})}\sum_{q=2}^{s}F^{(q)}(\mu_{0})\,
 \EB_{s,q}\bigl(\mu_{0}^{(1)},\dots,\mu_{0}^{(s-q+1)}\bigr),\quad s\ge2 .
\end{equation}
Then, with $M=|\boldsymbol m|\ge1$,
\begin{equation}\label{p3:eq:multi-Bell}
 V_{\boldsymbol m}=\frac{(-1)^{M}H_{\boldsymbol m}(\mu_{0})}{\boldsymbol m!}
 \sum_{s=0}^{M-1}\EB_{M-1,s}\bigl(\mu_{0}^{(1)},\mu_{0}^{(2)},\dots\bigr)\,
 \EB_{s}\bigl(\Lambda^{(\boldsymbol m)}_{1},\dots,\Lambda^{(\boldsymbol m)}_{s}\bigr),
\end{equation}
with the conventions $\EB_{0,0}=\EB_{0}=1$.
\end{corollary}

\begin{proof}
Fa\`a di Bruno (\cref{p0:def:ebell}) gives $\frac{\dd^{M-1}}{\dd y^{M-1}}H_{\boldsymbol
m}(\mu_{0}(y)) =\sum_{s}H_{\boldsymbol m}^{(s)}(\mu_{0})\EB_{M-1,s}(\mu_{0}^{(1)},\dots)$, and
writing $H_{\boldsymbol m}=\exp(\log H_{\boldsymbol m})$ gives $H_{\boldsymbol
m}^{(s)}=H_{\boldsymbol m}\,\EB_{s}(\Lambda^{(\boldsymbol m)}_{1},\dots)$. Substitute in
\cref{p3:eq:multi-LB}. \Cref{p3:eq:inverse-derivatives} is obtained by differentiating
$F(\mu_{0}(y))=y$ $s$ times and solving the Fa\`a di Bruno identity for the term containing
$\mu_{0}^{(s)}$.
\end{proof}

The point of \cref{p3:eq:multi-Bell} is that every ingredient is elementary and closed-form for
the partition blocks. Indeed, splitting each aggregate block into its primitive Fourier
constituents
\begin{equation}\label{p3:eq:primitive-block}
 E_{k,h,\sigma,\epsilon}(\mu)
 =\frac{\kappa}{2\sqrt k}\,\frac{\mu-\sigma k}{\mu^{3}}
 \exp\!\left[\frac{\sigma\mu}{k}
 +\ii\epsilon\Bigl(\pi s(h,k)-\frac{2\pi h}{k}n(\mu)\Bigr)\right],
 \qquad h\in U_{k},\ \epsilon=\pm1,
\end{equation}
and using $n'(\mu)=3\mu/\pi^{2}$, $n''(\mu)=3/\pi^{2}$, $n^{(s)}=0$ for $s\ge3$, one gets the
closed logarithmic jets
\begin{equation}\label{p3:eq:jets}
 \frac{\dd^{s}}{\dd\mu^{s}}\log E_{k,h,\sigma,\epsilon}
 =(-1)^{s-1}(s-1)!\left(\frac1{(\mu-\sigma k)^{s}}-\frac3{\mu^{s}}\right)
 +\mathbf 1_{s=1}\frac\sigma k
 -\ii\epsilon\frac{2\pi h}{k}\,n^{(s)}(\mu),
\end{equation}
while $\frac{\dd^{s}}{\dd\mu^{s}}\log F'$ is \cref{p3:eq:F-prime}. Hence $\Lambda^{(\boldsymbol
m)}_{s}=\sum_{j}m_{j}\,(\text{\cref{p3:eq:jets}}) -(\text{\cref{p3:eq:F-prime}})$ with no hidden
differentiation anywhere. This is S3's presentation and is the most explicit of the three.

\subsection{Engine II: the triangular Newton--Hensel recurrence}

The reversion may equally be organised in the phase coordinate, which is cheaper when many
coefficients are wanted.

\begin{theorem}[Triangular recurrence]\label{p3:thm:triangular}
Put $\mu=\mu_{0}+v$ and
\begin{equation}\label{p3:eq:Q-def}
 Q(v;\boldsymbol\varepsilon)=\log\!\Bigl(1+\sum_{j\in J_{K}}\varepsilon_{j}B_{j}(v)\Bigr),
 \qquad B_{j}(v)=r_{j}(\mu_{0}+v)\,\e^{-\alpha_{j}v},
\end{equation}
the bookkeeping variable $\varepsilon_{j}$ standing for the actual exponential
$\e^{-\alpha_{j}\mu_{0}}$, which is substituted only at the very end; the residual $v$-dependence
of $\e^{-\alpha_{j}(\mu_{0}+v)}$ is the factor $\e^{-\alpha_{j}v}$ inside $B_{j}$. Keeping the
$\varepsilon_{j}$ independent until the last step is what makes resonant exponent sums
($\alpha_{j}+\alpha_{j'}=\alpha_{j''}$, which does occur here) come out right. Then the reversion
equation is
\begin{equation}\label{p3:eq:v-equation}
 \Phi'(\mu_{0})\,v+\sum_{s\ge2}\frac{\Phi^{(s)}(\mu_{0})}{s!}v^{s}
 +Q(v;\boldsymbol\varepsilon)=0,
\end{equation}
and its unique formal solution $v=\sum_{\boldsymbol m\ne\boldsymbol0}V_{\boldsymbol m}
\boldsymbol\varepsilon^{\boldsymbol m}$ obeys
\begin{equation}\label{p3:eq:triangular}
 V_{\boldsymbol m}=-\frac1{\Phi'(\mu_{0})}\,
 \Coef{\boldsymbol\varepsilon^{\boldsymbol m}}\left\{
 \sum_{s\ge2}\frac{\Phi^{(s)}(\mu_{0})}{s!}\,v_{<\boldsymbol m}^{\;s}
 +Q(v_{<\boldsymbol m};\boldsymbol\varepsilon)\right\},
\end{equation}
where $v_{<\boldsymbol m}$ collects the already-computed coefficients. Each right-hand side is a
finite expression in $\mu_{0}$, $(\mu_{0}-1)^{-1}$, the actions $\alpha_{j}$ and the amplitudes
$a_{k}$.
\end{theorem}

\begin{proof}
Taking logarithms in $\mathcal{P}_{K}(\mu;\boldsymbol\varepsilon)=y$ and subtracting
$\Phi(\mu_{0})=R$ gives $\Phi(\mu_{0}+v)-\Phi(\mu_{0})+Q(v;\boldsymbol\varepsilon)=0$; Taylor
expansion of $\Phi$ at $\mu_{0}$ is \cref{p3:eq:v-equation}, legitimate as a formal identity
because $v$ has no constant term in $\boldsymbol\varepsilon$. For triangularity: at a fixed
$\boldsymbol m$, the unknown $V_{\boldsymbol m}$ occurs linearly only in the term
$\Phi'(\mu_{0})v$; each $v^{s}$ with $s\ge2$ needs a decomposition of $\boldsymbol m$ into at
least two non-zero parts and so uses only strictly smaller multi-indices; and every monomial of
$Q$ carries at least one explicit $\varepsilon_{j}$, so an occurrence of $V_{\boldsymbol m}$
inside $Q$ would be multiplied into a strictly larger multi-index. Solving the coefficient
equation gives \cref{p3:eq:triangular}. Finiteness holds because a fixed multi-index has finitely
many additive decompositions.
\end{proof}

\begin{corollary}[Multinomial expansion of $Q$]\label{p3:cor:Q-multi}
For $\boldsymbol m\ne\boldsymbol0$,
\begin{equation}\label{p3:eq:Q-multi}
 \Coef{\boldsymbol\varepsilon^{\boldsymbol m}}Q(v;\boldsymbol\varepsilon)\Big|_{v\ \mathrm{fixed}}
 =(-1)^{|\boldsymbol m|+1}\frac{(|\boldsymbol m|-1)!}{\boldsymbol m!}\prod_{j}B_{j}(v)^{m_{j}} .
\end{equation}
\end{corollary}

\begin{proof}
Only $S^{|\boldsymbol m|}$ in $\log(1+S)=\sum_{d\ge1}(-1)^{d+1}S^{d}/d$ contributes, with
multinomial coefficient $|\boldsymbol m|!/\boldsymbol m!$.
\end{proof}

\begin{corollary}[Closed multi-index Lagrange form in the phase]
\label{p3:cor:phase-LB}
Let $G(v)=\bigl(\Phi(\mu_{0}+v)-\Phi(\mu_{0})\bigr)/v$ for $v\ne0$ and $G(0)=\Phi'(\mu_{0})$. Then
$v=-Q(v;\boldsymbol\varepsilon)/G(v)$, and by \cref{p0:thm:LB},
\begin{equation}\label{p3:eq:phase-LB}
 V_{\boldsymbol m}=\sum_{q=1}^{|\boldsymbol m|}\frac{(-1)^{q}}{q}\,\Coef{v^{q-1}}
 \frac1{G(v)^{q}}
 \sum_{\substack{\boldsymbol m_{1}+\cdots+\boldsymbol m_{q}=\boldsymbol m\\ \boldsymbol m_{i}\ne\boldsymbol 0}}
 \prod_{i=1}^{q}\Coef{\boldsymbol\varepsilon^{\boldsymbol m_{i}}}Q(v;\boldsymbol\varepsilon),
\end{equation}
both sums being finite.
\end{corollary}

\begin{proof}
$G$ is analytic near $0$ with $G(0)=\Phi'(\mu_{0})\ne0$, and \cref{p3:eq:v-equation} reads
$vG(v)+Q(v;\boldsymbol\varepsilon)=0$. Apply \cref{p0:thm:LB} to
$v=z\cdot\bigl(-Q(v;\boldsymbol\varepsilon)/G(v)\bigr)$ and set $z=1$; since every monomial of $Q$
has positive $\boldsymbol\varepsilon$-degree, only $q\le|\boldsymbol m|$ contributes. Expanding
$Q^{q}$ and using \cref{p3:eq:Q-multi} gives \cref{p3:eq:phase-LB}.
\end{proof}

\begin{remark}[Three engines, one answer]
\Cref{p3:thm:additive-LB} works in the ordinate $y$ and is S3's; \cref{p3:cor:phase-LB} works in
the phase $\mu$ and is S2's; \cref{p3:thm:triangular} is S1's. They compute the same coefficients:
all three solve the same analytic equation with the same normalisation $V_{\boldsymbol 0}=0$, and
the solution of a triangular system is unique. In practice \cref{p3:eq:triangular} is preferable
when many multi-indices are wanted (it reuses everything), \cref{p3:eq:multi-LB} when a single
high-order multi-index is wanted (it needs no lower ones), and \cref{p3:eq:multi-Bell} when the
answer is wanted symbolically in $\mu_{0}$.
\end{remark}

\subsection{First-order sectors and the leading parity chirp}
\label{p3:sub:first-order}

\begin{theorem}[The first correction from each sector]\label{p3:thm:first-order}
Let $\boldsymbol e_{j}$ be the unit multi-index at $j=(k,\sigma)\in J_{K}$. Then
\begin{equation}\label{p3:eq:first-V}
 V_{\boldsymbol e_{j}}(\mu_{0})=-\frac{r_{j}(\mu_{0})}{\Phi'(\mu_{0})}
 =-\frac{a_{k}(\mu_{0})}{\sqrt k}\,
 \frac{\mu_{0}(\mu_{0}-\sigma k)}{\mu_{0}^{2}-3\mu_{0}+3},
\end{equation}
so that the corresponding displacement of the inverse phase is $V_{\boldsymbol
e_{j}}(\mu_{0})\,\e^{-\alpha_{j}\mu_{0}}$, and the corresponding displacement of the inverse index
is
\begin{equation}\label{p3:eq:first-N}
 \delta N^{(1)}_{k,\sigma}(y)
 =-\frac{3\,a_{k}(\mu_{0})}{\pi^{2}\sqrt k}\,
 \frac{\mu_{0}^{2}(\mu_{0}-\sigma k)}{\mu_{0}^{2}-3\mu_{0}+3}\,
 \e^{-(1-\sigma/k)\mu_{0}} .
\end{equation}
\end{theorem}

\begin{proof}
Take $\boldsymbol m=\boldsymbol e_{j}$ in \cref{p3:eq:triangular}: the $s\ge2$ terms and all
higher $Q$-monomials vanish, leaving $V_{\boldsymbol
e_{j}}=-\Coef{\varepsilon_{j}}Q(0;\boldsymbol\varepsilon)/\Phi'(\mu_{0})
=-B_{j}(0)/\Phi'(\mu_{0})=-r_{j}(\mu_{0})/\Phi'(\mu_{0})$. Insert $r_{j}=\frac{a_{k}}{\sqrt
k}\frac{1-\sigma k/\mu}{1-1/\mu} =\frac{a_{k}}{\sqrt k}\frac{\mu-\sigma k}{\mu-1}$ and
\cref{p3:eq:Phi-prime}: the factors $\mu_{0}-1$ cancel, giving \cref{p3:eq:first-V}. Finally
$n(\mu)=\frac1{24}+\frac{3\mu^{2}}{2\pi^{2}}$ is exactly quadratic, so
\begin{equation}\label{p3:eq:index-from-phase}
 n(\mu_{0}+v)=N_{0}(y)+\frac{3\mu_{0}}{\pi^{2}}v+\frac3{2\pi^{2}}v^{2},
\end{equation}
and the linear term at first order gives \cref{p3:eq:first-N}.
\end{proof}

\begin{corollary}[Index coefficient of every transmonomial]\label{p3:cor:N-multi}
For every $\boldsymbol m\ne\boldsymbol0$, the coefficient of $\boldsymbol\varepsilon^{\boldsymbol
m}$ in the inverse index is
\begin{equation}\label{p3:eq:N-multi}
 \mathfrak N_{\boldsymbol m}(\mu_{0})=\frac{3\mu_{0}}{\pi^{2}}V_{\boldsymbol m}
 +\frac3{2\pi^{2}}\sum_{\substack{\boldsymbol a+\boldsymbol b=\boldsymbol m\\ \boldsymbol a,\boldsymbol b\ne\boldsymbol0}}
 V_{\boldsymbol a}V_{\boldsymbol b} .
\end{equation}
\end{corollary}

\begin{proof}
Expand \cref{p3:eq:index-from-phase} with $v=\sum V_{\boldsymbol m}
\boldsymbol\varepsilon^{\boldsymbol m}$ and compare coefficients. No higher powers occur because
$n$ is quadratic.
\end{proof}

The first two positive sectors, with $a_{2}=\widetilde A_{2}(N_{0})=\cos(\pi N_{0})$ and
$a_{3}=\widetilde A_{3}(N_{0})=2\cos(\pi N_{0})\cos(\tfrac\pi{18}+\tfrac{\pi N_{0}}3)$
(interpolated convention) or $a_{2}=(-1)^{r}$, $a_{3}=2\cos(\tfrac\pi{18}-\tfrac{2\pi r}3)$
(residue-frozen), are
\begin{align}
 \delta N^{(1)}_{2,+}
 &=-\frac{3\,a_{2}}{\pi^{2}\sqrt2}\,
 \frac{\mu_{0}^{2}(\mu_{0}-2)}{\mu_{0}^{2}-3\mu_{0}+3}\,\e^{-\mu_{0}/2},
 \label{p3:eq:dN2}\\
 \delta N^{(1)}_{3,+}
 &=-\frac{3\,a_{3}}{\pi^{2}\sqrt3}\,
 \frac{\mu_{0}^{2}(\mu_{0}-3)}{\mu_{0}^{2}-3\mu_{0}+3}\,\e^{-2\mu_{0}/3}.
 \label{p3:eq:dN3}
\end{align}
Thus \emph{parity of the index appears at relative scale $\e^{-\mu_{0}/2}$, the residue modulo $3$
at $\e^{-2\mu_{0}/3}$, and the residue modulo $k$ first appears in the growing $k$th sector at
$\e^{-(1-1/k)\mu_{0}}$.}

\subsection{The fractional-power staircase in the ordinate}

The core equation converts exponentials in $\mu_{0}$ into fractional powers of $y$ exactly:
\begin{equation}\label{p3:eq:exp-to-y}
 \e^{-\alpha\mu_{0}}=\left(\frac{\kappa(\mu_{0}-1)}{y\,\mu_{0}^{3}}\right)^{\!\alpha}
 \qquad\text{for every }\alpha,
\end{equation}
directly from $y=\kappa\e^{\mu_{0}}(\mu_{0}-1)\mu_{0}^{-3}$.

\begin{proposition}[Every first-order sector in closed $y$-form]
\label{p3:prop:y-form}
With $\alpha=\alpha_{k,\sigma}$ and $z=1/\mu_{0}$,
\begin{equation}\label{p3:eq:y-form}
 \delta N^{(1)}_{k,\sigma}
 =-\frac{3\kappa^{\alpha}}{\pi^{2}\sqrt k}\,a_{k}(\mu_{0})\;
 y^{-\alpha}\,\mu_{0}^{\,1-2\alpha}\sum_{s\ge0}g^{(k,\sigma)}_{s}\,\mu_{0}^{-s},
\end{equation}
where
\begin{equation}\label{p3:eq:g-gf}
 \sum_{s\ge0}g^{(k,\sigma)}_{s}z^{s}
 =\frac{(1-\sigma kz)(1-z)^{\alpha}}{1-3z+3z^{2}},
 \qquad
 g^{(k,\sigma)}_{s}=\sum_{i=0}^{s}(-1)^{i}\binom{\alpha}{i}
 \bigl(\chi_{s-i}-\sigma k\,\chi_{s-i-1}\bigr),
\end{equation}
with $\chi_{-1}=0$ and $\chi_{s}=\dfrac{\varrho_{+}^{s+1}-\varrho_{-}^{s+1}}
{\varrho_{+}-\varrho_{-}}$ the coefficients of $(1-3z+3z^{2})^{-1}$
($\varrho_{\pm}=\tfrac12(3\pm\ii\sqrt3)$; $\chi_{0}=1,\chi_{1}=3,\chi_{2}=6,
\chi_{3}=9,\chi_{4}=9,\chi_{5}=0,\dots$).
\end{proposition}

\begin{proof}
Insert \cref{p3:eq:exp-to-y} into \cref{p3:eq:first-N} and factor $\mu_{0}^{1-2\alpha}$ out of
$\dfrac{\mu_{0}^{2}(\mu_{0}-\sigma k)(\mu_{0}-1)^{\alpha}}
{(\mu_{0}^{2}-3\mu_{0}+3)\mu_{0}^{3\alpha}}$; with $z=1/\mu_{0}$ the remaining factor is
$\dfrac{(1-\sigma kz)(1-z)^{\alpha}}{1-3z+3z^{2}}$. The coefficient formula is the Cauchy product
of $(1-z)^{\alpha}=\sum_{i}(-1)^{i} \binom\alpha i z^{i}$ with $(1-\sigma kz)/(1-3z+3z^{2})$, and
$\chi_{s}$ is the standard solution of $\chi_{s}=3\chi_{s-1}-3\chi_{s-2}$, $\chi_{0}=1$,
$\chi_{1}=3$, whose closed form is as stated since $\varrho_{\pm}$ are the roots of
$\varrho^{2}=3\varrho-3$.
\end{proof}

\begin{corollary}[The inverse staircase of scales]\label{p3:cor:staircase-scales}
Since $\mu_{0}\sim\log y$, the primary positive sectors contribute, in order,
\begin{equation}\label{p3:eq:scales}
 y^{-1/2},\qquad y^{-2/3}(\log y)^{-1/3},\qquad y^{-3/4}(\log y)^{-1/2},
 \qquad\dots,\qquad y^{-(1-1/k)}(\log y)^{-1+2/k},
\end{equation}
accumulating at the scale $y^{-1}$, where they meet the square of the first sector and the whole
negative family.
\end{corollary}

\begin{proof}
By \cref{p3:eq:y-form} the $(k,+)$ sector has size
$y^{-\alpha}\mu_{0}^{1-2\alpha}(1+\bigO(\mu_{0}^{-1}))$ times a bounded amplitude, with
$\alpha=1-1/k$, so $1-2\alpha=-1+2/k$; and $\mu_{0}\sim\log y$ because
$\Phi(\mu_{0})=\log(y/\kappa)$ with $\Phi(\mu)=\mu(1+\littleo(1))$. Substituting $k=2,3,4$ gives
the first three displayed scales. Since $\alpha_{k,+}\uparrow1$, the scales accumulate at
$y^{-1}$; the square of the $(2,+)$ sector has action $\tfrac12+\tfrac12=1$ and every negative
action is $\ge1$.
\end{proof}

\begin{theorem}[The leading parity chirp]\label{p3:thm:leading-chirp}
In the interpolated convention,
\begin{equation}\label{p3:eq:chirp}
 \mathcal{N}(y)=N_{0}(y)
 -\frac{3\sqrt\kappa}{\pi^{2}\sqrt2}\,\frac{\cos\bigl(\pi N_{0}(y)\bigr)}{\sqrt y}
 \left(1+\bigO\!\left(\frac1{\log y}\right)\right)
 +\bigO\!\left(y^{-2/3}(\log y)^{-1/3}\right),
\end{equation}
with the constant
\begin{equation}\label{p3:eq:chirp-constant}
 \frac{3\sqrt\kappa}{\pi^{2}\sqrt2}=\frac{3^{1/4}}{2\pi}
 =0.2094596846361762626265841\ldots
\end{equation}
Since $N_{0}(y)=\frac{3}{2\pi^{2}}(\log y)^{2}(1+\littleo(1))$, the cosine is a \emph{quadratic
logarithmic chirp}, not a log-periodic oscillation.
\end{theorem}

\begin{proof}
Take $k=2$, $\sigma=+1$, $\alpha=\tfrac12$ in \cref{p3:eq:y-form}: the power
$\mu_{0}^{1-2\alpha}=\mu_{0}^{0}$ is absent, $g^{(2,+)}_{0}=1$, and $a_{2}=\cos(\pi N_{0})$ by
\cref{p3:lem:Atilde}. The constant is $3\kappa^{1/2}/(\pi^{2}\sqrt2)$, and
$\sqrt\kappa=\pi/\sqrt{6\sqrt3}$ gives $\frac{3\pi}{\pi^{2}\sqrt2\sqrt{6\sqrt3}}
=\frac{3}{\pi\sqrt{12\sqrt3}}=\frac{3}{2\pi\cdot3^{3/4}}=\frac{3^{1/4}}{2\pi}$, which we evaluated
to $25$ digits. The relative $\bigO(1/\log y)$ is the $s\ge1$ tail of \cref{p3:eq:y-form},
and the next omitted contribution is the $k=3$ sector at scale $y^{-2/3}(\log y)^{-1/3}$ by
\cref{p3:cor:staircase-scales}; all products of two blocks have action $\ge1$.
\end{proof}

\begin{remark}[Constant check]
S2 prints this constant as $0.2094596846361762626$; the closed form $3^{1/4}/(2\pi)$ is S3's. The
two agree, and both agree with our independent $25$-digit evaluation \cref{p3:eq:chirp-constant}.
S1 does not state the constant, giving instead the equivalent phase-variable form
\cref{p3:eq:dN2}; the three are algebraically identical, as the cancellation of $\mu_{0}-1$ in the
proof of \cref{p3:thm:first-order} shows.
\end{remark}

\section{The profinite arithmetic phase}\label{p3:sec:profinite}

The residue-frozen convention of \cref{p3:def:frozen} organises itself into a projective system,
which is the cleanest way to say what data an ``arithmetically complete'' inverse of $p$ actually
depends on.

\begin{theorem}[Compatibility and the profinite parameter]\label{p3:thm:profinite}
Let $\Lambda_{K}=\operatorname{lcm}(1,\dots,K)$, so $\Lambda_{K}\mid\Lambda_{K+1}$. Suppose
$r_{K}\in\Int/\Lambda_{K}\Int$ satisfy $r_{K+1}\equiv r_{K} \pmod{\Lambda_{K}}$, and let
$\mu_{K,r_{K}}$ be the corresponding inverse branches of \cref{p3:def:frozen}. Then for every $K$,
the transseries of $\mu_{K+1,r_{K+1}}$ obtained by deleting all monomials containing the variables
of level $K+1$ coincides with the transseries of $\mu_{K,r_{K}}$. Consequently the complete
arithmetic inverse is a projective family indexed by the profinite integers $\widehat
r\in\widehat{\Int}=\varprojlim\Int/m\Int$, and an actual partition index $n$ supplies its own
phase through the diagonal embedding $n\mapsto\widehat n$.
\end{theorem}

\begin{proof}
Deleting the level-$(K+1)$ variables from the defining equation \cref{p3:eq:v-equation} at level
$K+1$ produces the defining equation at level $K$ \emph{provided} the retained amplitudes agree,
i.e.\ provided $A_{k}(r_{K+1})=A_{k}(r_{K})$ for $k\le K$. That is exactly the compatibility
hypothesis, since $A_{k}$ has period $k\mid\Lambda_{K}$ and $r_{K+1}\equiv
r_{K}\pmod{\Lambda_{K}}$. Both truncated equations are triangular with the same normalisation
$V_{\boldsymbol0}=0$, and a triangular system has a unique solution, so the solutions agree
coefficient by coefficient. The projective limit statement is the definition of $\widehat{\Int}$:
compatible systems $(r_{K})$ are exactly the elements of
$\varprojlim\Int/\Lambda_{K}\Int=\widehat{\Int}$, because $\{\Lambda_{K}\}$ is cofinal in the
divisibility order on $\Nat$.
\end{proof}

\begin{remark}[Why a profinite parameter and not a real one]
\label{p3:rem:why-profinite}
At every algebraic order all indices share one expansion. Parity enters at relative order
$\e^{-\mu/2}$; the residue mod $3$ at $\e^{-2\mu/3}$; the residue mod $k$ at $\e^{-(1-1/k)\mu}$.
No finite real-valued smooth parameter records all of these congruence classes simultaneously, and
no single real analytic function of $n$ interpolates all $A_{k}$ canonically
(\cref{p3:rem:interp-nonunique}). A profinite integer does so canonically and finitely at every
exponential resolution: each finite action cutoff $\eta<1$ sees only $A_{k}$ for
$k\le(1-\eta)^{-1}$, i.e.\ only $\widehat r$ modulo $\Lambda_{\lfloor(1-\eta)^{-1}\rfloor}$. This
is S1's observation and is absent from S2 and S3.
\end{remark}

\begin{remark}[A structure suggested, not proved]\label{p3:rem:sheaf-conjecture}
S1 proposes an ``arithmetic transseries sheaf'' with base $\widehat{\Int}$ and fibres the
finite-level analytic transseries algebras, and calls \cref{p3:thm:profinite} its elementary case.
We record this as a suggestion. Nothing beyond \cref{p3:thm:profinite} --- namely, projective
compatibility of the coefficient systems --- is established here or in the sources: no topology on
the fibres, no gluing statement, and no continuity of $\widehat r\mapsto\mu_{K,\widehat r}$ in any
topology on $\widehat{\Int}$ has been proved. It is stated as \cref{p3:conj:sheaf} in
\cref{p3:sec:open}, not as a theorem.
\end{remark}

\section{Validity, truncation control, and integer recovery}
\label{p3:sec:validity}

\subsection{Remainder transport}

\begin{theorem}[Inverse error after $K$ Rademacher levels]\label{p3:thm:inverse-tail}
Fix $K\ge1$. Let $y=p(n)$, let $\mu=\lambda\sqrt{n-1/24}$, and let $\mu_{K}$ denote either inverse
branch of \cref{p3:def:frozen} (with $r\equiv n$) or of \cref{p3:def:interpolated}. Then
\begin{equation}\label{p3:eq:inverse-tail}
 \mu_{K}(y)-\mu=\bigO_{K}\!\left(\mu^{3/2}\,
 \e^{-(1-\frac1{K+1})\mu}\right),
 \qquad
 n_{K}(y)-n=\bigO_{K}\!\left(\mu^{5/2}\,\e^{-(1-\frac1{K+1})\mu}\right),
\end{equation}
where $n_{K}=\frac1{24}+3\mu_{K}^{2}/(2\pi^{2})$. Explicitly, with the constant of
\cref{p3:thm:tail}, for $\mu\ge\mu_{*}$ (any threshold at which $\Phi'\ge\tfrac12$ and the omitted
tail is $\le\tfrac12$ in absolute value),
\begin{equation}\label{p3:eq:inverse-tail-explicit}
 |n_{K}(y)-n|\le\frac{3\mu}{\pi^{2}}\cdot\frac{2}{1}\cdot
 \frac{8\mu^{2}}{3(1-\mu^{-1})\sqrt K}
 \exp\!\left[-\Bigl(1-\frac1{K+1}\Bigr)\mu\right]
 \bigl(1+\littleo(1)\bigr).
\end{equation}
\end{theorem}

\begin{proof}
This is \cref{p0:thm:remainder-transport} with the forward remainder supplied by
\cref{p3:thm:tail,p3:thm:tail-sharp}. In detail: let $\mathcal{F}_{K}=\log\mathcal{P}_{K}$ and
$\mathcal{F}=\log p$ regarded as functions of the phase. By \cref{p3:prop:relative},
$\mathcal{F}(\mu)-\mathcal{F}_{K}(\mu) =\log\bigl(1+\mathcal{R}_{>K}/(1+\mathcal{R}_{\le
K})\bigr)$, which by \cref{p3:thm:tail-sharp} and $|\log(1+u)|\le2|u|$ for $|u|\le\tfrac12$ is
$\bigO_{K}(\mu^{3/2}\e^{-(1-1/(K+1))\mu})$. Since
$\mathcal{F}_{K}'(\mu)=\Phi'(\mu)(1+\littleo(1))\to1$, the mean value theorem converts this
logarithmic residual into the phase error, which is the first half of \cref{p3:eq:inverse-tail}.
Finally $\dd n/\dd\mu=3\mu/\pi^{2}$ gives the index error. The explicit form
\cref{p3:eq:inverse-tail-explicit} replaces \cref{p3:thm:tail-sharp} by the explicit
\cref{p3:thm:tail} (losing $\mu^{-1/2}$, gaining a constant) and $\mathcal{F}_{K}'\ge\tfrac12$.
\end{proof}

\begin{corollary}[Beyond-all-orders accuracy of the core inverse]
\label{p3:cor:core-inverse-error}
For $y=p(n)$,
\begin{equation}\label{p3:eq:core-error}
 N_{0}(y)-n=\bigO\!\left(\mu^{5/2}\e^{-\mu/2}\right)
 =\bigO\!\left(n^{5/4}\exp\!\left(-\tfrac12\pi\sqrt{2n/3}\right)\right),
\end{equation}
so $N_{0}(p(n))-n\to0$ faster than every power of $n^{-1}$.
\end{corollary}

\begin{proof}
\Cref{p3:thm:inverse-tail} at $K=1$; note $N_{0}=n_{1}$ because the $K=1$ level of
\cref{p3:def:frozen} contains the sectors $(1,\pm)$, and the $(1,-)$ sector has action $2$, hence
contributes below the $K=1$ error bound.
\end{proof}

\begin{proposition}[A posteriori certificate; specialisation of
\cref{p0:thm:backward-error}]\label{p3:prop:certificate} Let
$\mathcal{F}_{K}(\mu)=\log\mathcal{P}_{K}(\mu)$ and let
$\rho(\widetilde\mu)=\mathcal{F}_{K}(\widetilde\mu)-\log y$ be the residual of a candidate
$\widetilde\mu$. If $\mathcal{F}_{K}'\ge m>0$ on an interval $I$ containing both $\widetilde\mu$
and the exact root $\mu_{K}(y)$, then
\begin{equation}\label{p3:eq:certificate}
 |\widetilde\mu-\mu_{K}(y)|\le\frac{|\rho(\widetilde\mu)|}{m},
 \qquad
 |\widetilde n-n_{K}(y)|\le\frac{3}{2\pi^{2}}
 \bigl(|\widetilde\mu|+|\mu_{K}(y)|\bigr)\frac{|\rho(\widetilde\mu)|}{m}.
\end{equation}
\end{proposition}

\begin{proof}
Mean value theorem applied to $\mathcal{F}_{K}$, then the factorisation
$\widetilde\mu^{2}-\mu_{K}^{2}=(\widetilde\mu-\mu_{K})(\widetilde\mu+\mu_{K})$. Since
$\mathcal{F}_{K}'\to1$, one may take $m=\tfrac12$ beyond an effectively checkable threshold.
\end{proof}

\begin{remark}[Why a residual certificate rather than term counting]
\Cref{p3:prop:certificate} separates the \emph{local inversion} error (how well $\widetilde\mu$
solves the level-$K$ equation) from the \emph{Rademacher tail} (how far the level-$K$ equation is
from the truth). Only the second needs \cref{p3:thm:tail}. This modularity is what makes it
possible to certify a truncated symbolic transseries and a numerically refined Newton root by the
same inequality.
\end{remark}

\subsection{Algebraic truncation}

\begin{proposition}\label{p3:prop:algebraic-truncation}
With $U_{M}(z)=\sum_{m\le M}c_{m}z^{m}$ as in \cref{p3:thm:core-correction},
\begin{equation}\label{p3:eq:algebraic-truncation}
 \mu_{0}-\bigl(X+U_{M}(X^{-1})\bigr)=\bigO(X^{-M-1}),
 \qquad
 N_{0}-\left[\frac1{24}+\frac3{2\pi^{2}}\bigl(X+U_{M}(X^{-1})\bigr)^{2}\right]
 =\bigO(X^{-M}),
\end{equation}
the loss of one order coming from $\dd n/\dd\mu\asymp X$. In particular $M\ge1$ is
\emph{necessary} for the index to be correct to $\littleo(1)$ and hence for the arithmetic
phases of \cref{p3:sec:full-inverse} to be evaluated at the right point
(\cref{p3:rem:phase-lock}).
\end{proposition}

\begin{proof}
$U$ is analytic at $0$ by \cref{p3:thm:core-correction}, so its Taylor remainder after $M$ terms
is $\bigO(z^{M+1})$; substitute $z=X^{-1}$ and use $N_{0}-\widetilde
N=\frac3{2\pi^{2}}(\mu_{0}-\widetilde\mu) (\mu_{0}+\widetilde\mu)$ with
$\mu_{0}+\widetilde\mu=\bigO(X)$.
\end{proof}

\subsection{Integer recovery}

\begin{theorem}[Eventual exact recovery by rounding]\label{p3:thm:rounding}
There is $n_{0}$ such that for every integer $n\ge n_{0}$,
\begin{equation}\label{p3:eq:rounding}
 n=\bigl\lfloor N_{0}(p(n))\bigr\rceil,
\end{equation}
$\lfloor\cdot\rceil$ denoting nearest-integer rounding. Consequently the arithmetic phase
$\widehat n$ of \cref{p3:thm:profinite} is recoverable from $y=p(n)$ \emph{before} any
exponentially small sector is evaluated.
\end{theorem}

\begin{proof}
By \cref{p3:eq:core-error}, $|N_{0}(p(n))-n|\to0$; it is therefore $<\tfrac12$ for all large $n$,
at which point rounding is exact. The last sentence is the observation that $A_{k}(n)$ depends
only on $n\bmod k$.
\end{proof}

\begin{remark}[The threshold is very small in practice]
\label{p3:rem:threshold}
\Cref{p3:thm:rounding} is asymptotic and claims no optimal $n_{0}$. We computed $N_{0}(p(n))-n$
for all $1\le n\le 59$ in $140$-digit arithmetic; the largest absolute value is $0.1723\ldots$ at
$n=1$, and the sequence decreases in magnitude thereafter (through $2.8\times10^{-2}$ at $n=10$,
$4.5\times10^{-4}$ at $n=50$, $1.5\times10^{-5}$ at $n=100$, $4.3\times10^{-17}$ at $n=1000$). So
$n_{0}=1$ is consistent with the data, though we prove only the asymptotic statement; producing a
certified $n_{0}$ requires effective constants in \cref{p3:thm:inverse-tail}, i.e.\ effective
Rademacher remainder bounds, which is listed as \cref{p3:conj:effective-threshold}.
\end{remark}

\begin{remark}[Circularity, avoided]\label{p3:rem:no-circularity}
The arithmetic transseries needs the residue $r$ of $n$, which appears to require knowing $n$. The
order of operations resolves this: \emph{(i)} compute $X$ from \cref{p3:eq:carrier}; \emph{(ii)}
compute $\mu_{0}$ and $N_{0}$, which are phase-independent; \emph{(iii)} round, obtaining $n$ by
\cref{p3:thm:rounding} for $y$ on the range; \emph{(iv)} now every $A_{k}(n)$ is known, and the
exponentially small sectors refine an already identified integer. The arithmetic layer is thus a
\emph{refinement} of a known answer, not a device for guessing the phase. For $y$ off the range,
step (iii) yields $\lceil\mathcal{N}(y)\rceil$ by \cref{p3:thm:staircase} instead, and the sectors
describe $\mathcal{N}$ itself.
\end{remark}

\section{Numerical verification}\label{p3:sub:numerics}

Every number in this chapter was recomputed for this volume. Partition values were generated
exactly by Euler's pentagonal recurrence
\begin{equation}\label{p3:eq:pentagonal}
 p(n)=\sum_{m\ge1}(-1)^{m-1}
 \left[p\!\left(n-\tfrac{m(3m-1)}2\right)+p\!\left(n-\tfrac{m(3m+1)}2\right)\right],
 \qquad p(0)=1,\ p(r)=0\ (r<0),
\end{equation}
in exact integer arithmetic (checked against $p(10)=42$, $p(50)=204226$, $p(100)=190569292$);
Dedekind sums exactly as rationals from \cref{p3:eq:dedekind}; the transcendental factors in
$140$-digit floating point. Coefficient identities were verified in exact rational or exact
symbolic arithmetic, never by floating-point spot checks. The tables check normalisation, signs,
phases and the hierarchy of actions; they are used in no proof.

\begin{table}[htbp]
\centering
\caption{Relative error $\bigl(S_{K}(n)-p(n)\bigr)/p(n)$ of the paired truncation
$S_{K}=\sum_{k\le K}(P_{k,+}+P_{k,-})$, recomputed at $140$ digits.  Accuracy need not
improve monotonically in $K$ at small $n$: the amplitudes $A_{k}(n)$ change sign and can
suppress or enhance an individual sector.}
\label{p3:tab:forward}
\small
\begin{tabular}{@{}rrrrrr@{}}
\toprule
$n$ & $p(n)$ & $K=1$ & $K=2$ & $K=3$ & $K=5$\\
\midrule
$10$ & $42$ & $-8.8621\cdot10^{-3}$ & $1.6606\cdot10^{-3}$ & $3.7391\cdot10^{-4}$ & $3.4722\cdot10^{-4}$\\
$25$ & $1\,958$ & $1.0717\cdot10^{-3}$ & $3.0045\cdot10^{-6}$ & $-6.1015\cdot10^{-5}$ & $-3.4866\cdot10^{-6}$\\
$50$ & $204\,226$ & $-7.3074\cdot10^{-5}$ & $3.9028\cdot10^{-6}$ & $2.1050\cdot10^{-7}$ & $-1.1350\cdot10^{-7}$\\
$100$ & $190\,569\,292$ & $-1.8220\cdot10^{-6}$ & $8.6838\cdot10^{-9}$ & $-4.9491\cdot10^{-9}$ & $3.1454\cdot10^{-10}$\\
$250$ & $2.3079\cdot10^{14}$ & $-1.0760\cdot10^{-9}$ & $7.2267\cdot10^{-13}$ & $4.3021\cdot10^{-14}$ & $2.1256\cdot10^{-15}$\\
$500$ & $2.3002\cdot10^{21}$ & $-2.4389\cdot10^{-13}$ & $1.7537\cdot10^{-17}$ & $-2.0017\cdot10^{-19}$ & $-3.5933\cdot10^{-22}$\\
$1000$ & $2.4061\cdot10^{31}$ & $-1.6997\cdot10^{-18}$ & $1.2574\cdot10^{-24}$ & $-3.4744\cdot10^{-27}$ & $3.9453\cdot10^{-30}$\\
\bottomrule
\end{tabular}
\end{table}

\begin{table}[htbp]
\centering
\caption{Inverse errors at exact ordinates $y=p(n)$; entries are (approximation) $-\,n$.
$N_{X}$ is the bare Lambert carrier index, $N_{0}$ the dominant core, $N_{K}$ the Newton
inverse of the level-$K$ interpolated sum (\cref{p3:def:interpolated}).  The
residue-frozen inverses (\cref{p3:def:frozen}) reproduce the last three columns to all
five digits shown for $n\ge100$; at $n=50$ they differ in the fifth digit of the
$K=3,5$ entries ($-1.2992\cdot10^{-6}$, $7.0057\cdot10^{-7}$), at $n=25$ in the fourth
($2.7926\cdot10^{-4}$, $1.5958\cdot10^{-5}$), at $n=10$ in the third
($-5.3338\cdot10^{-3}$, $-1.2006\cdot10^{-3}$, $-1.1150\cdot10^{-3}$).  That is exactly
the pattern \cref{p3:prop:two-conventions} predicts.}
\label{p3:tab:inverse}
\small
\begin{tabular}{@{}rrrrrr@{}}
\toprule
$n$ & $N_{X}-n$ & $N_{0}-n$ & $N_{2}-n$ & $N_{3}-n$ & $N_{5}-n$\\
\midrule
$10$ & $-3.9910\cdot10^{-1}$ & $2.8447\cdot10^{-2}$ & $-5.3290\cdot10^{-3}$ & $-1.1863\cdot10^{-3}$ & $-1.1015\cdot10^{-3}$\\
$25$ & $-3.7877\cdot10^{-1}$ & $-4.9052\cdot10^{-3}$ & $-1.3751\cdot10^{-5}$ & $2.7902\cdot10^{-4}$ & $1.5941\cdot10^{-5}$\\
$50$ & $-3.5044\cdot10^{-1}$ & $4.5102\cdot10^{-4}$ & $-2.4089\cdot10^{-5}$ & $-1.2993\cdot10^{-6}$ & $7.0059\cdot10^{-7}$\\
$100$ & $-3.3600\cdot10^{-1}$ & $1.5378\cdot10^{-5}$ & $-7.3293\cdot10^{-8}$ & $4.1772\cdot10^{-8}$ & $-2.6548\cdot10^{-9}$\\
$250$ & $-3.2364\cdot10^{-1}$ & $1.3942\cdot10^{-8}$ & $-9.3644\cdot10^{-12}$ & $-5.5747\cdot10^{-13}$ & $-2.7544\cdot10^{-14}$\\
$500$ & $-3.1768\cdot10^{-1}$ & $4.4042\cdot10^{-12}$ & $-3.1668\cdot10^{-16}$ & $3.6146\cdot10^{-18}$ & $6.4888\cdot10^{-21}$\\
$1000$ & $-3.1356\cdot10^{-1}$ & $4.2960\cdot10^{-17}$ & $-3.1780\cdot10^{-23}$ & $8.7815\cdot10^{-26}$ & $-9.9717\cdot10^{-29}$\\
\bottomrule
\end{tabular}
\end{table}

Three features of \cref{p3:tab:inverse} confirm the theory point by point. \emph{(i)} The first
column tends to $-3/\pi^{2}=-0.30396\ldots$, not to $0$: this is \cref{p3:cor:carrier-bias} and
the numerical signature of \cref{p3:rem:phase-lock}. (S1 and S2 tabulate
$N_{X}=\tfrac1{24}+3X^{2}/(2\pi^{2})$ and report $\to-0.304$; S3 tabulates $3X^{2}/(2\pi^{2})$
without the $\tfrac1{24}$ and reports $\to-0.346$. The columns differ by exactly $1/24$ and each
is correct for its own definition; we keep the $\tfrac1{24}$ so that $N_{X}$ and $N_{0}$ are
directly comparable.) \emph{(ii)} Restoring the amplitude factor $1-1/\mu$, i.e.\ passing from $X$
to $\mu_{0}$, removes that entire $\bigO(1)$ bias and leaves an exponentially small error.
\emph{(iii)} The residual after the core alternates with the parity of $n$: at $n=49,50,51$ it is
$-5.2643\cdot10^{-4}$, $+4.5102\cdot10^{-4}$, $-4.0920\cdot10^{-4}$, exactly as \cref{p3:eq:dN2}
predicts through $\widetilde A_{2}=\cos\pi N_{0}$. Adding the single first-order $k=2$ sector
reduces these to $-1.0890\cdot10^{-5}$, $-2.4087\cdot10^{-5}$, $+2.9029\cdot10^{-5}$; adding the
first-order sectors through $k=10$ reduces them to $-2.24\cdot10^{-7}$, $-9.95\cdot10^{-8}$,
$-5.81\cdot10^{-8}$. This checks both the sign and the phase of the leading chirp
\cref{p3:eq:chirp}.

\begin{remark}[What was verified in exact arithmetic]\label{p3:rem:symbolic-checks}
(1) \Cref{p3:eq:c-finite} against the Taylor expansion of \cref{p3:eq:Omega}, symbolically in
$\beta$ through order $t^{9}$: identical; the closed forms \cref{p3:eq:c-first,p3:eq:c-six} were
produced from \cref{p3:eq:c-finite} and match S2 (through $c_{5}$) and S3 (through $c_{8}$). (2)
$\omega_{0},\dots,\omega_{4}$ against both the S1 and the S2 grouping of the same numbers:
identical (they look different only because S1 splits each into partial fractions in $\pi$). (3)
$c_{1},\dots,c_{18}$ of \cref{p3:tab:cm} from \cref{p3:eq:c-recurrence-core} in exact rationals,
plus the residual test \cref{p3:eq:U-residual} at $M=12$; all eighteen agree with the sources. (4)
$\gamma_{1},\dots,\gamma_{8}$ against all three sources; $\gamma_{9},\gamma_{10}$ are new here.
(5) $\Pi_{1},\dots,\Pi_{8}$ from the Lagrange form against the Stirling form and against each of
the three printed shapes, confirming \cref{p3:rem:H-versus-R}; S1's printed $P_{6},P_{7}$ are
correct. (6) S3's $u$-expansion in $H$ through $H^{-4}$ and its $N_{0}$-expansion through
$H^{-3}$, and S2's $N_{0}$-block in $R$ through the constant term: correct as printed. (7)
\Cref{p3:eq:first-Ak} against direct evaluation of \cref{p3:eq:Ak} from exact Dedekind sums. (8)
The chirp constant \cref{p3:eq:chirp-constant} to $25$ digits in both forms. No discrepancy with
any source's arithmetic was found: every repair recorded in this chapter is structural, not
numerical.
\end{remark}

\begin{remark}[Normalisation cross-checks]\label{p3:sub:normalization-crosscheck}
Because \cref{p3:eq:rademacher} carries a factor $\sqrt k$ that some normalisations absorb into
$A_{k}$, a missing or spurious power of $k$ is the easiest error to make and the hardest to see.
Three identifications of one quantity detect it: the $k$th derivative factor of
\cref{p3:eq:rademacher} equals $\frac{\mu}{4kN^{3/2}}[\e^{\mu/k}(1-k/\mu)+\e^{-\mu/k}(1+k/\mu)]$,
and after multiplication by $\sqrt k/(\pi\sqrt2)$ it may be written either as
$\frac1{4\sqrt{3k}\,N}[\cdots]$ or as $\frac{\kappa}{\sqrt k\,\mu^{2}}[\cdots]$; at $k=1$ the
positive term is $\frac{\e^{\mu}}{4\sqrt3\,N}(1-\tfrac1\mu)$, which becomes \cref{p3:eq:HR} after
replacing $N$ by $n$ and dropping $1-1/\mu$. Any of these disagreeing with the others by a power
of $k$ localises the slip.
\end{remark}

\section{Algorithms}\label{p3:sec:algorithms}

All formulas of this chapter are constructive and need only formal power-series arithmetic over
$\Rat(\pi)$ together with exact rational Dedekind sums.

\begin{algorithm}[Forward sector coefficients]\label{p3:alg:forward}
Given $k,\sigma,M$: set $\beta=\sigma\pi/(6k)$ and evaluate the finite sum \cref{p3:eq:c-finite}
for $0\le m\le M$ ($\bigO(m)$ field operations each, no series manipulation; keep $\beta$
symbolic if sector-uniform output is wanted). If the whole array is wanted at once, evaluate
$\ell_{1},\dots,\ell_{M}$ from \cref{p3:eq:ell} and run the triangular recurrence
\cref{p3:eq:c-recurrence} instead. Assemble \cref{p3:eq:sector-unshifted}, pairing the two signs
before summing over large $k$ (\cref{p3:rem:pairing}).
\end{algorithm}

\begin{algorithm}[Phase-independent inverse]\label{p3:alg:core}
Given a large $y$ and an order $M$: (1) compute $X$ from \cref{p3:eq:carrier} with a standard
$\Wlam_{-1}$ routine --- do \emph{not} substitute the logarithmic expansion
\cref{p3:eq:carrier-flat}, which is numerically far inferior; (2) form $\mu_{0}=X+\sum_{m\le
M}c_{m}X^{-m}$ from the universal rationals of \cref{p3:tab:cm}, optionally refined by one Newton
step on $\Phi(\mu)=R$; (3) return $N_{0}=\tfrac1{24}+3\mu_{0}^{2}/(2\pi^{2})$; (4) certify by
evaluating $\rho=\mu_{0}-2\log\mu_{0}+\log(1-\mu_{0}^{-1})-\log(y/\kappa)$ and applying
\cref{p3:prop:certificate}; (5) round (\cref{p3:thm:rounding}) if $y$ is known to lie on the
range, else take $\lceil\cdot\rceil$ for the staircase (\cref{p3:thm:staircase}).
\end{algorithm}

\begin{algorithm}[Arithmetic exponential correction]\label{p3:alg:full}
Given $y$, a level $K$ and a finite set of multi-indices: (1) compute $\mu_{0},N_{0}$ by
\cref{p3:alg:core} and use $N_{0}$, never $N_{X}$, to evaluate amplitudes
(\cref{p3:rem:phase-lock}); (2) compute exact rational $s(h,k)$ for $k\le K$ and hence $A_{k}(r)$
(frozen, with $r$ from \cref{p3:thm:rounding}) or $\widetilde A_{k}(N_{0})$ (interpolated); (3)
build the actions and amplitudes of \cref{p3:eq:block}; (4) traverse the multi-indices in a graded
monomial order using \cref{p3:eq:triangular}, or, for one isolated $\boldsymbol m$, use
\cref{p3:eq:multi-LB} or the Bell form \cref{p3:eq:multi-Bell} with the closed jets
\cref{p3:eq:jets}; (5) substitute $\varepsilon_{j}=\e^{-\alpha_{j}\mu_{0}}$ and convert to the
index by \cref{p3:eq:N-multi}; (6) certify by \cref{p3:prop:certificate} on $\mathcal{P}_{K}$,
adding \cref{p3:thm:tail} if the comparison is with the exact sequence.
\end{algorithm}

\begin{remark}[Modular structure of the computation]\label{p3:rem:reuse}
The layers are independent and should be cached separately: the carrier polynomials $\Pi_{j}$
depend only on the exponent data $(a,b)=(1,-2)$; the rationals $c_{m}$ only on the amplitude
$1-1/\mu$; the inverse-derivative polynomials $\mu_{0}^{(s)}$ only on $\Phi$; and the arithmetic
enters solely through the numbers $a_{k}$ in \cref{p3:eq:block}. Universal tables are computed
once and residue-dependent evaluation is then cheap.
\end{remark}

\section{A template for Rademacher-type sequences}\label{p3:sec:template}

Nothing above used the partition function except through four properties. S3 has the most careful
version of this observation and S1 the most usable one; what follows merges them into a hypothesis
list and a theorem, so that the template is a statement one can check rather than a slogan.

\begin{definition}[Rademacher-type representation]\label{p3:def:rademacher-type}
A sequence $(a(n))_{n\ge n_{1}}$ of positive reals is of \emph{Rademacher type} if:
\begin{enumerate}
\item[(T1)] there is a real analytic, strictly increasing \emph{phase}
$\phi:(\nu_{1},\infty)\to(0,\infty)$ with $\phi\to\infty$ and real analytic inverse, giving the
large variable $\mu=\phi(n)$;
\item[(T2)] the \emph{envelope} $F(\mu)=C\e^{a\mu^{\alpha}}\mu^{b}\exp Q(\mu^{-\delta})$
is an exponential--power model \cref{p0:def:model} with $a>0$, $\alpha\in(0,1]$;
\item[(T3)] there are actions $\alpha_{j}>0$ ($j\in J$) with $\inf_{j}\alpha_{j}>0$,
amplitudes $a_{j}(n)$ taking values in a fixed finite-dimensional space at each truncation level
(e.g.\ periodic in $n$), and $g_{j}$ analytic at $\infty$, with
      \begin{equation}\label{p3:eq:template}
       a(n)=F(\mu)\Bigl(1+\sum_{j\in J'}a_{j}(n)g_{j}(\mu)\e^{-\alpha_{j}\mu}
       +\mathcal{T}_{J'}(\mu;n)\Bigr)\quad\text{for every finite }J'\subset J;
      \end{equation}
\item[(T4)] $\{j:\alpha_{j}\le\eta\}$ is finite for every $\eta<\sup_{j}\alpha_{j}$, and
either $\mathcal{T}_{J'}=\bigO(\mu^{C'}\e^{-\eta'\mu})$ for some $\eta'>\max_{j\in
J'}\alpha_{j}$, or --- if the actions accumulate --- the grouped remainder is retained as one
analytic block.
\end{enumerate}
\end{definition}

\begin{theorem}[Template]\label{p3:thm:template}
Let $(a(n))$ be of Rademacher type. Then:
\begin{enumerate}
\item \textup{(Reduction.)} $\mu=\xi^{1/\alpha}$ of \cref{p0:lem:alpha-reduction} brings
\textup{(T2)} to the linear-phase model with computable $(C',1,b',Q')$.
\item \textup{(Carrier.)} $C'\e^{\xi}\xi^{b'}=y$ is solved exactly by
\cref{p0:thm:lambert-core}; equivalently, for $y=C\e^{a\mu^{\alpha}}\mu^{-\beta}$,
      \begin{equation}\label{p3:eq:general-carrier}
       \mu=\left(-\frac{\beta}{a\alpha}\,\Wlam_{-1}\!\left(-\frac{a\alpha}{\beta}
       \Bigl(\frac Cy\Bigr)^{\alpha/\beta}\right)\right)^{1/\alpha}.
      \end{equation}
\item \textup{(Amplitude.)} The envelope inverse is $\xi_{0}=X+U(X^{-1})$ with $U$ the
unique solution of the fixed-point equation of \cref{p0:thm:lambert-centered} at $(1,b',Q')$;
coefficients by \cref{p0:lem:lagrange-fp} or by isolating the linear occurrence of $U$.
\item \textup{(Phase locking.)} If some $a_{j}$ is periodic on the index scale and
$\dd n/\dd\mu\to\infty$, the amplitudes must be evaluated at $\xi_{0}$, not at $X$: the carrier
bias $n(\xi_{0})-n(X)$ is $\bigO(1)$ in general.
\item \textup{(Sectors.)} With the phase frozen, every inverse transmonomial is given by
\cref{p3:eq:multi-LB} (in $y$) or \cref{p3:eq:phase-LB}/\cref{p3:eq:triangular} (in $\mu$), with
first-order term
      \begin{equation}\label{p3:eq:general-first}
       \mu=\mu_{0}-\frac{a_{j}g_{j}(\mu_{0})}{(\log F)'(\mu_{0})}\e^{-\alpha_{j}\mu_{0}}
       +\bigO\bigl(\e^{-2\alpha_{\min}\mu_{0}}\bigr).
      \end{equation}
\item \textup{(Transport.)} A relative forward remainder $\epsilon$ transports to a phase
error $\epsilon/(\log F)'(\mu_{0})$ and an index error $\epsilon\,n'(\mu_{0})/(\log F)'(\mu_{0})$,
by \cref{p0:thm:remainder-transport}.
\item \textup{(Staircase.)} If $a$ is integer valued with gaps bounded below, the
generalised inverse is the ceiling of the smooth one (\cref{p0:thm:staircase}), with the order-one
sawtooth of \cref{p3:rem:sawtooth} as an irreducible final sector.
\end{enumerate}
\end{theorem}

\begin{proof}
(i)--(iii) are \cref{p0:lem:alpha-reduction,p0:thm:lambert-core,p0:thm:lambert-centered} at the
stated parameters. For the display in (ii): raising $y=C\e^{a\mu^{\alpha}}\mu^{-\beta}$ to the
power $\alpha/\beta$ and rearranging gives $\bigl(\tfrac{a\alpha}{\beta}\mu^{\alpha}\bigr)
\e^{-(a\alpha/\beta)\mu^{\alpha}}=\tfrac{a\alpha}{\beta}(C/y)^{\alpha/\beta}$, which is the
Lambert equation with the stated argument, the branch $-1$ forced by $\mu\to\infty$. (iv) is the
computation of \cref{p3:rem:phase-lock} with $3\mu/\pi^{2}$ replaced by $n'(\mu_{0})$: what
matters is only that $n'\to\infty$ while $\xi_{0}-X\to0$, so their product need not vanish. (v) is
\cref{p3:thm:additive-LB,p3:cor:phase-LB,p3:thm:triangular}, whose proofs used no property of $F$
or of the blocks beyond analyticity and $F'\ne0$; the first-order display is \cref{p3:eq:first-V}
in that generality. (vi) is the mean-value argument of \cref{p3:thm:inverse-tail}, and (vii) is
\cref{p3:thm:staircase}.
\end{proof}

\begin{remark}[The partition case, and the scope of the template]
\label{p3:rem:template-scope}
Taking $\phi(n)=\pi\sqrt{2(n-1/24)/3}$, $(C,a,\alpha,b,Q)$ as in \cref{p3:eq:model-in-N},
$J=\{(k,\sigma)\}$, $\alpha_{k,\sigma}=1-\sigma/k$ and $a_{k}=A_{k}$, hypotheses (T1)--(T3) hold
by \cref{p3:thm:exact-sectors,p3:prop:relative}, and (T4) holds in its second form: the actions
\emph{do} accumulate at $1$, and the grouped block is the paired tail \cref{p3:eq:R-exact}
controlled by \cref{p3:thm:tail}. With $(\alpha,\beta)=(1,2)$ in $\mu$,
\cref{p3:eq:general-carrier} reduces to \cref{p3:eq:carrier}. The template applies verbatim to
coefficient sequences of weakly holomorphic modular forms and eta-quotients with Rademacher-type
expansions; the partition function is unusually clean because the Bessel index $3/2$ makes the
kernel two exponentials with amplitudes linear in $\mu^{-1}$, so step (iii) is the explicit
rational recurrence \cref{p3:eq:c-recurrence-core} rather than a general implicit solve. \emph{No
theorem about any other specific sequence is claimed here}: verifying (T1)--(T4) for a given
family is exactly the work the template does not do.
\end{remark}

\section{Structural consequences and open questions}\label{p3:sec:open}

\subsection{Convergent blocks inside an exponentially complete object}

The exact sector sum \cref{p3:eq:paired-sector} converges; each sector's algebraic series in
$n^{-1/2}$ converges on $|n^{-1/2}|<\sqrt{24}$ (\cref{p3:thm:sector-gf}); and the reversion series
$U$ converges near $0$ (\cref{p3:thm:core-correction}). Divergence is therefore \emph{not} what
creates the non-perturbative sectors here: they are supplied by modularity and the circle method,
independently of the perturbative data. The useful trichotomy is that a \emph{Poincar\'e
expansion} records one dominant sector and forgets all exponentially small information; a
\emph{transseries completion} records the other sectors, whether or not the series in each
diverges; and an \emph{exact transseries representation}, as here, may itself converge. Compare
\cref{p1:sec:top,p2:sec:top,p4:sec:top}, where the algebraic series do diverge and
\cref{p0:thm:optimal-truncation} is the operative tool.

\subsection{Arithmetic amplitudes are not Stokes constants}

The $A_{k}(n)$ weight distinct exponential sectors and govern exponentially small corrections to
the inverse, which is formally the role of Stokes constants. They are not Stokes constants: they
are arithmetic, periodic in $n$ rather than fixed, and no natural analytic continuation in the
index produces them canonically (\cref{p3:rem:interp-nonunique}). \Cref{p3:thm:profinite} says
what replaces the Stokes data: a profinite integer, of which every finite exponential resolution
sees only a finite quotient.

\subsection{Open questions}

\begin{conjecture}[Effective rounding threshold]\label{p3:conj:effective-threshold}
There is an explicit small $n_{0}$ --- the evidence of \cref{p3:rem:threshold} suggests $n_{0}=1$
--- with $n=\lfloor N_{0}(p(n))\rceil$ for all $n\ge n_{0}$, provable by inserting effective
Rademacher remainder bounds (Lehmer; Banerjee--Paule--Radu--Schneider) in place of $|A_{k}|\le k$
in \cref{p3:thm:inverse-tail}.
\end{conjecture}

\begin{conjecture}[Arithmetic transseries sheaf]\label{p3:conj:sheaf}
The projective family of \cref{p3:thm:profinite} extends to a sheaf of transseries algebras over
$\widehat{\Int}$, with $\widehat r\mapsto(\text{level-}K$ inverse transseries$)$ continuous for
the profinite topology and the finite-level algebras gluing over the standard clopen basis. Only
the projective compatibility is proved here (\cref{p3:rem:sheaf-conjecture}).
\end{conjecture}

\begin{conjecture}[Uniform treatment at the accumulation action]
\label{p3:conj:uniform-accumulation}
The zeta condensation \cref{p3:eq:zeta-condensation} can serve as a single block in
\cref{p3:thm:additive-LB}, yielding a description of the inverse uniform across the accumulation
action $1$, with error bounds depending on the analytic behaviour of $\mathcal{Z}_{n}(s)$ rather
than on a sector cutoff. This needs continuation and vertical growth information for the
Kloosterman-type Dirichlet series $\mathcal{Z}_{n}$, which is not established here. S3 raises it
as an optimisation problem rather than a gap, correctly: the exact Rademacher representation is
always available as a fallback.
\end{conjecture}

\begin{remark}[Classifying interpolations]\label{p3:rem:interp-problem}
Different strictly increasing smooth interpolations of $p$ agree at every algebraic order of the
inverse and on the integer range, but differ in the exponentially small off-lattice sectors:
candidates include \cref{p3:def:Atilde}, piecewise-logarithmic interpolation through $(n,\log
p(n))$, Newton or cardinal series with prescribed growth, and continuations built from Poincar\'e
series. Which part of the inverse transseries is interpolation-independent?
\Cref{p3:prop:two-conventions} answers this for one pair of choices --- they agree strictly below
action $1$ --- and a general statement is open.
\end{remark}

\begin{remark}[Formalisation targets]\label{p3:rem:formalisation}
Rademacher's theorem is now machine-checked (Eberl and Paulson, Archive of Formal Proofs, 2026).
The natural next steps, in dependency order, are: (a) the normalisation
\cref{p3:thm:exact-sectors}; (b) the generating function \cref{p3:eq:Omega} and the recurrence
\cref{p3:eq:c-recurrence}; (c) the Catalan/residue proof of \cref{p3:thm:c-finite}; (d) the
implicit construction of $U$ and \cref{p3:eq:c-recurrence-core}; (e) the triangular multivariate
inversion \cref{p3:eq:triangular}; (f) residual-to-root transfer, \cref{p3:prop:certificate},
under an explicit derivative lower bound. Steps (b), (d), (e) are Bell-polynomial infrastructure
shared with \cref{p0:sec:top}.
\end{remark}

\section{Summary}\label{p3:sec:summary}

\begin{enumerate}
\item In the shifted phase $\mu=\pi\sqrt{2(n-1/24)/3}$, Rademacher's series is the exact
sector sum \cref{p3:eq:sector-sum}, each sector having a two-term amplitude linear in $\mu^{-1}$.
At $(k,\sigma)=(1,+1)$ it is the exponential--power model \cref{p0:def:model} with $\alpha=1/2$,
reduced by \cref{p0:lem:alpha-reduction} to $\kappa\e^{\mu}\mu^{-2}(1-\mu^{-1})$
(\cref{p3:lem:alpha-reduction,p3:tab:dictionary}).
\item Re-expansion in $n^{-1/2}$ is governed by one analytic function \cref{p3:eq:Omega};
its coefficients have the finite closed form \cref{p3:eq:c-finite} in \emph{every} sector, which
enters only through $\beta=\sigma\pi/(6k)$. Each such series converges while the total is only
asymptotic --- confusing the two is the characteristic error of the subject.
\item The actions $1\mp1/k$ accumulate at $1$ from both sides, so the completed object is
not a Hahn series and the two signs may not be separated
(\cref{p3:def:rademacher-transseries,p3:rem:pairing}); \cref{p3:thm:tail} bounds the tail
explicitly and \cref{p3:thm:zeta-condensation} packages the accumulation exactly.
\item The inverse carrier is $-2\Wlam_{-1}(-\tfrac12\sqrt{\kappa/y})$ with Stirling
flattening polynomials \cref{p3:eq:Pi-lagrange}; the amplitude is restored by
\cref{p0:thm:lambert-centered} at $\Psi(t,u)=3\log(1+tu)-\log(1-t+tu)$, whose coefficients are
\cref{p3:tab:cm}.
\item Since $\dd n/\dd\mu\asymp\mu$, that reversion produces an $\bigO(1)$ index
shift $3/\pi^{2}$ which must be resummed \emph{before} the periodic amplitudes are evaluated
(\cref{p3:rem:phase-lock}).
\item With the phase frozen --- by a residue class mod $\operatorname{lcm}(1,\dots,K)$ or
by the entire interpolation \cref{p3:def:Atilde}, which agree strictly below action $1$
(\cref{p3:prop:two-conventions}) --- every inverse transmonomial is given by
\cref{p3:eq:multi-LB}, \cref{p3:eq:multi-Bell}, \cref{p3:eq:phase-LB} or \cref{p3:eq:triangular}.
The leading correction is the parity chirp \cref{p3:eq:chirp} with constant $3^{1/4}/(2\pi)$, and
the primary sectors form the staircase \cref{p3:eq:scales}.
\item Along the range the core alone recovers $n$ by rounding
(\cref{p3:thm:rounding}); for general $y$ the staircase is the ceiling of the smooth inverse
(\cref{p3:thm:staircase}), whose order-one sawtooth dominates every arithmetic sector
(\cref{p3:rem:sawtooth}).
\end{enumerate}

\begin{remark}[Works referred to]\label{p3:rem:references}
Analytic input: Rademacher, Proc.\ Nat.\ Acad.\ Sci.\ USA \textbf{23} (1937), 78--84; Proc.\
London Math.\ Soc.\ (2) \textbf{43} (1938), 241--254; Ann.\ of Math.\ \textbf{44} (1943),
416--422; textbook treatment in Apostol, \emph{Modular Functions and Dirichlet Series in Number
Theory}, 2nd ed., Springer GTM 41, 1990; machine-checked in Eberl and Paulson, \emph{Rademacher's
series for the partition function}, Archive of Formal Proofs, 2026. Leading asymptotic: Hardy and
Ramanujan, Proc.\ London Math.\ Soc.\ (2) \textbf{17} (1918), 75--115. Remainder estimates:
Lehmer, J.\ London Math.\ Soc.\ \textbf{12} (1937), 171--176 and Trans.\ Amer.\ Math.\ Soc.\
\textbf{43} (1938), 271--295, \textbf{46} (1939), 362--373; Banerjee, Paule, Radu and Schneider,
arXiv:2209.07887. The dominant-sector case of \cref{p3:eq:c-finite} is O'Sullivan's,
arXiv:2205.13468 and arXiv:2203.02868. Lambert $W$: Corless, Gonnet, Hare, Jeffrey and Knuth,
Adv.\ Comput.\ Math.\ \textbf{5} (1996), 329--359; Jeffrey, Corless, Hare and Knuth,
arXiv:math/9512230. The sequence is OEIS A000041.
\end{remark}

%% ==================== fragment p4 ====================
\chapter{The swing factorial (A056040)}
\label{p4:sec:top}

\begin{sourcebox}
This chapter merges three independent treatments of the swinging factorial:
\emph{All-Orders Asymptotic Transseries for the Swing Factorial and Its Two
Inverse Branches} (cited below as \textsf{S1}), \emph{The Full Asymptotic
Transseries of the Swing Factorial and Its Inverse} (\textsf{S2}), and
\emph{Parity-Resolved Transseries for the Swing Factorial and Its Inverse}
(\textsf{S3}).  The general apparatus of \textsf{S2} --- its universal
coefficient calculus for exponential--power inversion and its
Lambert-normalized reversion theorem --- was promoted to \cref{p0:sec:top};
here it is only \emph{cited} and \emph{specialized}.  \textsf{S1} works in the
half index $m$ with the carrier constant $\log4$; \textsf{S2} and \textsf{S3}
work in the original index with $\log2$ and differ only in the sign convention
for the parity marker.  The chapter adopts the \textsf{S2}/\cref{p0:sec:top}
normalization throughout and gives the translation dictionary in
\cref{p4:rem:dictionary}.  \textsf{S1}'s genuinely distinct contribution --- the
centred (translation-symmetric) normal form --- is kept, but recast so that it
too is a single $\sigma$-uniform statement.

Every rational, polynomial and algebraic coefficient displayed in this chapter
was recomputed from the stated recurrences in exact arithmetic
(\texttt{sympy}, rational Bernoulli input, truncated power series over
$\Rat[\lambda,\lambda^{-1},\sigma]/(\sigma^2-1)$); every decimal in the
validation tables was recomputed at $50$ significant digits.  Where a
recomputation reproduced a source table, that is said; where it did not, the
discrepancy is recorded in a \textsf{repairbox}.
\end{sourcebox}

\section{The sequence, the parity obstruction, and what ``inverse'' can mean}
\label{p4:sec:elementary}

\subsection{Definition and exact parity identities}

\begin{definition}[Swing factorial]
\label{p4:def:swing}
For $n\in\Nat_0$ put
\begin{equation}
  \operatorname{sw}(n)=\frac{n!}{\lfloor n/2\rfloor!^{\,2}}.
  \label{p4:eq:swing-def}
\end{equation}
This is OEIS A056040, the \emph{swinging factorial}.
\end{definition}

Splitting on the parity of $n$ gives at once
\begin{equation}
  \operatorname{sw}(2m)=\binom{2m}{m},
  \qquad
  \operatorname{sw}(2m+1)=(2m+1)\binom{2m}{m},
  \label{p4:eq:bisections}
\end{equation}
whence the one-step recurrence
\begin{equation}
  \operatorname{sw}(0)=1,
  \qquad
  \operatorname{sw}(n)=
  \begin{cases}
    n\,\operatorname{sw}(n-1), & n \text{ odd},\\[2pt]
    \dfrac4n\,\operatorname{sw}(n-1), & n \text{ even},
  \end{cases}
  \label{p4:eq:one-step}
\end{equation}
and the closed product form
\begin{equation}
  \operatorname{sw}(n)=2^{\,n-(n\bmod 2)}\prod_{k=1}^{n}k^{(-1)^{k+1}}.
  \label{p4:eq:swing-product}
\end{equation}
The first values are
\[
  1,\,1,\,2,\,6,\,6,\,30,\,20,\,140,\,70,\,630,\,252,\,2772,\,924,\,12012,
  \,3432,\,51480,\,\ldots
\]

\begin{proposition}[Generating functions]
\label{p4:prop:gf}
\begin{align}
  \sum_{n\ge0}\operatorname{sw}(n)\frac{z^n}{n!}&=(1+z)\,I_0(2z),
  \label{p4:eq:egf}\\
  \sum_{n\ge0}\operatorname{sw}(n)z^{n}
  &=\frac1{\sqrt{1-4z^2}}+\frac{z}{(1-4z^2)^{3/2}},
  \label{p4:eq:ogf}
\end{align}
where $I_0$ is the modified Bessel function of order zero.
\end{proposition}

\begin{proof}
By \cref{p4:eq:bisections} the even exponential part is
$\sum_{m\ge0}z^{2m}/(m!)^2=I_0(2z)$ and the odd exponential part is
$z\,I_0(2z)$, which is \cref{p4:eq:egf}.  For the ordinary series put
$F(v)=\sum_{m\ge0}\binom{2m}{m}v^m=(1-4v)^{-1/2}$; then the even part is
$F(z^2)$ and the odd part is
$z\bigl(2vF'(v)+F(v)\bigr)\big|_{v=z^2}=z(1-4z^2)^{-3/2}$.
\end{proof}

\begin{remark}[The generating function already shows the obstruction]
\label{p4:rem:codominance}
The two singularities $z=\pm\tfrac12$ of \cref{p4:eq:ogf} have equal modulus.
Co-dominant singularities of equal modulus at $\pm\rho$ produce a period-two
oscillation in the coefficient asymptotics, and no single non-oscillatory
singular expansion can represent the whole sequence.  Even/odd projection
separates the two contributions exactly.  This is the generating-function
shadow of everything that follows.
\end{remark}

\subsection{The sequence is not eventually monotone}

\begin{proposition}[Exact adjacent ratios]
\label{p4:prop:ratios}
For $m\ge0$,
\begin{equation}
  \frac{\operatorname{sw}(2m+1)}{\operatorname{sw}(2m)}=2m+1,
  \qquad
  \frac{\operatorname{sw}(2m+2)}{\operatorname{sw}(2m+1)}=\frac{2}{m+1}.
  \label{p4:eq:ratios}
\end{equation}
Consequently $\operatorname{sw}$ rises strictly from every even index to the
next odd index, and falls strictly from every odd index to the next even index
as soon as $m\ge2$.  In particular $\operatorname{sw}$ is not eventually
monotone and has no eventually defined single-valued inverse.
\end{proposition}

\begin{proof}
Both ratios are \cref{p4:eq:one-step}.  The second is $<1$ exactly when
$m+1>2$.
\end{proof}

\begin{warning}[There is no ``the'' inverse swing factorial]
\label{p4:rem:two-sheets}
Any formula presented as \emph{the} inverse of $\operatorname{sw}$ without
naming a parity sheet or a generalized-inverse convention is incomplete: of
the three inverse objects of \cref{p0:def:three-inverses}, none exists
globally for $\operatorname{sw}$ and all three exist branchwise
(\cref{p4:sec:discrete}).
\end{warning}

\subsection{The two analytic branches}

All three sources resolve the obstruction the same way: interpolate the two
bisections separately by gamma quotients.  We use the parity marker
$\sigma\in\{-1,+1\}$ of \textsf{S2}, with
\begin{equation}
  \sigma=-1 \ \text{(even phase)},
  \qquad
  \sigma=+1 \ \text{(odd phase)},
  \qquad
  \sigma_n:=(-1)^{n+1},
  \qquad
  \lambda:=\log2 .
  \label{p4:eq:parity-marker}
\end{equation}
The coefficient ring of the whole chapter is the \emph{parity algebra}
\begin{equation}
  \mathscr R:=\Rat[\lambda,\lambda^{-1},\sigma]\big/(\sigma^2-1),
  \label{p4:eq:parity-algebra}
\end{equation}
in which every formula below is a single element carrying both sheets.  This
is not a notational economy only: $\sigma_n$ is a discrete transmonomial of
modulus one, and it must remain visible at every order of every expansion.

\begin{definition}[Parity-resolved gamma branches]
\label{p4:def:branches}
For $x>0$ and $\sigma\in\{-1,+1\}$ set
\begin{equation}
  \mathcal S_\sigma(x)
  :=\frac{\Gamma(x+1)}{\Gamma\!\bigl(x/2+(3-\sigma)/4\bigr)^{2}} .
  \label{p4:eq:branch-def}
\end{equation}
Explicitly $\mathcal S_-(x)=\Gamma(x+1)/\Gamma(x/2+1)^2$ and
$\mathcal S_+(x)=\Gamma(x+1)/\Gamma\bigl((x+1)/2\bigr)^2$, and
\begin{equation}
  \mathcal S_-(2m)=\operatorname{sw}(2m),
  \qquad
  \mathcal S_+(2m+1)=\operatorname{sw}(2m+1)
  \qquad(m\in\Nat_0).
  \label{p4:eq:branch-interpolates}
\end{equation}
\end{definition}

\begin{proposition}[Duplication normal form and duality]
\label{p4:prop:duplication}
For $x>0$,
\begin{equation}
  \mathcal S_\sigma(x)
  =\frac{2^{x}}{\sqrt\pi}
   \left(\frac{\Gamma(x/2+1)}{\Gamma(x/2+1/2)}\right)^{\!\sigma},
  \label{p4:eq:branch-ratio}
\end{equation}
and consequently
\begin{equation}
  \mathcal S_+(x)\,\mathcal S_-(x)=\frac{4^{x}}{\pi},
  \qquad
  \frac{\mathcal S_\sigma(x+2)}{\mathcal S_\sigma(x)}
  =\frac{16(x+1)(x+2)}{(2x+3-\sigma)^2}
  =\begin{cases}
     \dfrac{4(x+1)}{x+2},&\sigma=-1,\\[8pt]
     \dfrac{4(x+2)}{x+1},&\sigma=+1.
   \end{cases}
  \label{p4:eq:duality-and-step}
\end{equation}
\end{proposition}

\begin{proof}
Legendre duplication in the form
$\Gamma(x+1)=2^{x}\pi^{-1/2}\Gamma\bigl((x+1)/2\bigr)\Gamma(x/2+1)$ turns
\cref{p4:eq:branch-def} into \cref{p4:eq:branch-ratio}; multiplying the two
values of $\sigma$ makes the gamma ratio cancel and leaves $4^x/\pi$.  For the
two-step ratio, $\Gamma(x+3)/\Gamma(x+1)=(x+1)(x+2)$ while the denominator
contributes $\bigl(x/2+(3-\sigma)/4\bigr)^{-2}=16(2x+3-\sigma)^{-2}$.
\end{proof}

\begin{proposition}[Strict monotonicity and existence of the real inverses]
\label{p4:prop:monotone}
For $x>0$,
\begin{equation}
  \frac{\dd}{\dd x}\log\mathcal S_\sigma(x)
  =\lambda+\frac\sigma2
   \Bigl[\psi\bigl(\tfrac x2+1\bigr)-\psi\bigl(\tfrac x2+\tfrac12\bigr)\Bigr],
  \label{p4:eq:logder}
\end{equation}
where $\psi=\Gamma'/\Gamma$.  The bracket decreases strictly from $2\lambda$ at
$x=0$ to $0$ at $x=\infty$.  Hence $\mathcal S_+$ has strictly positive
derivative on $[0,\infty)$, $\mathcal S_-$ has derivative $0$ at $x=0$ and
strictly positive derivative on $(0,\infty)$, and both are strictly increasing
bijections of $[0,\infty)$ onto $[1,\infty)$.  Write
\begin{equation}
  \mathcal I_\sigma:=\mathcal S_\sigma^{-1}:[1,\infty)\to[0,\infty).
  \label{p4:eq:inverse-def}
\end{equation}
\end{proposition}

\begin{proof}
Differentiating \cref{p4:eq:branch-ratio} logarithmically gives
\cref{p4:eq:logder}.  Put
$g(x)=\psi(x/2+1)-\psi(x/2+1/2)$.  Then
$g'(x)=\tfrac12\bigl[\psi'(x/2+1)-\psi'(x/2+1/2)\bigr]<0$ because the
trigamma function $\psi'(t)=\sum_{k\ge0}(t+k)^{-2}$ is strictly decreasing on
$(0,\infty)$; and $g(x)\to0$ as $x\to\infty$ by
$\psi(t+\tfrac12)-\psi(t)\to0$.  Finally $g(0)=\psi(1)-\psi(1/2)=2\log2
=2\lambda$.  Thus $0<g(x)<2\lambda$ for $x>0$.  For $\sigma=+1$,
\cref{p4:eq:logder} equals $\lambda+g/2>0$ on $[0,\infty)$; for $\sigma=-1$ it
equals $\lambda-g/2$, which vanishes only at $x=0$ and is positive afterwards.
Both branches take the value $1$ at $x=0$ and tend to $+\infty$, so each is a
continuous increasing bijection onto $[1,\infty)$.
\end{proof}

\subsection{Why one cannot simply interpolate through both parities}

\begin{proposition}[Every smooth interpolation of the whole sequence oscillates]
\label{p4:prop:oscillatory}
With $R(x)=\Gamma(x/2+1)/\Gamma\bigl((x+1)/2\bigr)$, put
\begin{equation}
  \widehat{\mathcal S}(x)=\frac{2^{x}}{\sqrt\pi}\,R(x)^{-\cos\pi x} .
  \label{p4:eq:oscillatory-interps}
\end{equation}
Then $\widehat{\mathcal S}(n)=\operatorname{sw}(n)$ for every $n\in\Nat_0$,
yet $\widehat{\mathcal S}$ is not eventually monotone: its logarithmic
derivative
\begin{equation}
  \frac{\dd}{\dd x}\log\widehat{\mathcal S}(x)
  =\lambda+\pi\sin(\pi x)\log R(x)-\cos(\pi x)\,\frac{R'(x)}{R(x)}
  \label{p4:eq:oscillatory-derivative}
\end{equation}
has an oscillating term of amplitude $\sim\tfrac\pi2\log x$ against remaining
terms that are $O(1)$, so it changes sign infinitely often.  The same holds
for the convex interpolation
$\tfrac{1+\cos\pi x}2\mathcal S_-+\tfrac{1-\cos\pi x}2\mathcal S_+$.
Consequently the phase-resolved \emph{pair} $(\mathcal S_-,\mathcal S_+)$, and
not any single function, is the asymptotically stable continuation of
$\operatorname{sw}$; the corresponding inverse object is the pair
$(\mathcal I_-,\mathcal I_+)$.
\end{proposition}

\begin{proof}
At $x=n$ one has $-\cos\pi n=(-1)^{n+1}=\sigma_n$, so
$\widehat{\mathcal S}(n)=(2^{n}/\sqrt\pi)R(n)^{\sigma_n}
=\mathcal S_{\sigma_n}(n)=\operatorname{sw}(n)$ by
\cref{p4:eq:branch-ratio} and \cref{p4:eq:branch-interpolates}.
Differentiating
$\log\widehat{\mathcal S}=x\lambda-\tfrac12\log\pi-\cos(\pi x)\log R(x)$ gives
\cref{p4:eq:oscillatory-derivative}; by \cref{p4:cor:R-asymptotic} below,
$\log R(x)=\tfrac12\log(x/2)+O(x^{-1})$ and $R'/R=O(x^{-1})$, so the
oscillating term has amplitude $\tfrac\pi2\log x+O(1)\to\infty$ and dominates.
The convex interpolation is treated the same way after writing it as
$\mathcal S_-(x)$ times a factor whose logarithmic derivative oscillates with
amplitude $\asymp\log x$.
\end{proof}

\section{Exact normal form, Binet reduction, and Borel structure}
\label{p4:sec:normalform}

\subsection{The single scalar correction}

\begin{definition}[The gamma-ratio correction]
\label{p4:def:D}
For $x>0$ set
\begin{equation}
  D(x):=\log\frac{\Gamma(x/2+1)}{\Gamma(x/2+1/2)}-\frac12\log\frac x2 .
  \label{p4:eq:D-def}
\end{equation}
\end{definition}

\begin{theorem}[Exact $\sigma$-uniform normal form]
\label{p4:thm:normal-form}
For $x>0$ and $\sigma\in\{-1,+1\}$,
\begin{equation}
  \boxed{\;
  \mathcal S_\sigma(x)=C_\sigma\,\e^{\lambda x}\,x^{\sigma/2}\,
  \e^{\sigma D(x)},
  \qquad
  C_\sigma=\frac{2^{-\sigma/2}}{\sqrt\pi}.\;}
  \label{p4:eq:normal-form}
\end{equation}
At integer arguments this is the exact parity formula
\begin{equation}
  \log\operatorname{sw}(n)
  =n\lambda-\tfrac12\log\pi
  +\sigma_n\Bigl(\tfrac12\log\tfrac n2+D(n)\Bigr),
  \qquad n\ge1 .
  \label{p4:eq:integer-exact}
\end{equation}
\end{theorem}

\begin{proof}
By \cref{p4:eq:D-def}, $\Gamma(x/2+1)/\Gamma(x/2+1/2)=(x/2)^{1/2}\e^{D(x)}$.
Substituting into \cref{p4:eq:branch-ratio} gives
$\mathcal S_\sigma=(2^x/\sqrt\pi)(x/2)^{\sigma/2}\e^{\sigma D}$, and
$(x/2)^{\sigma/2}=2^{-\sigma/2}x^{\sigma/2}$.  For
\cref{p4:eq:integer-exact}, take logarithms of \cref{p4:eq:normal-form} at
$x=n$, $\sigma=\sigma_n$, and use
$\tfrac12\log n-\tfrac\lambda2=\tfrac12\log\tfrac n2$.
\end{proof}

\Cref{p4:eq:normal-form} is exactly \cref{p0:def:model} with
\begin{equation}
  \boxed{\;
  \alpha=1,\qquad a=\lambda=\log 2,\qquad b=\frac\sigma2,\qquad
  C=C_\sigma=\frac{2^{-\sigma/2}}{\sqrt\pi},\qquad
  q_j=\sigma d_j,\;}
  \label{p4:eq:dictionary}
\end{equation}
where $d_j$ are the coefficients of the asymptotic series of $D$
(\cref{p4:thm:logcoeffs}).  Because $\alpha=1$ no monomial substitution is
needed and \cref{p0:lem:alpha-reduction} is not invoked.  We record the
dictionary once and use it in every specialization below.

\begin{remark}[Dictionary between the three sources]
\label{p4:rem:dictionary}
Let $\varepsilon_{\textsf{S3}}\in\{+1,-1\}$ be the marker of \textsf{S3} and
$\epsilon_{\textsf{S1}}\in\{0,1\}$ that of \textsf{S1}.  Then
\[
  \sigma=-\varepsilon_{\textsf{S3}},
  \qquad
  \sigma=2\epsilon_{\textsf{S1}}-1,
  \qquad
  \textsf{S3}:\ Q=-D,\quad q^{\textsf{S3}}_r=-d_{2r-1},
  \qquad
  \textsf{S1}:\ \alpha_{\textsf{S1}}=\log4=2\lambda .
\]
\textsf{S1} parametrizes each sheet by the half index $m$ with
$S^{\textsf{S1}}_{\epsilon}(m)=\operatorname{sw}(2m+\epsilon)$; as functions
this is the reparametrization $x=2m+\epsilon$ of \cref{p4:def:branches}, so
\[
  S^{\textsf{S1}}_{\epsilon}(m)=\mathcal S_{2\epsilon-1}(2m+\epsilon),
  \qquad
  \mathcal I_{\sigma}(y)=2M^{\textsf{S1}}_{\epsilon}(y)+\epsilon .
\]
The reparametrization is invisible in the forward coefficients up to the
substitution $x\mapsto2m+\epsilon$, but it is \emph{not} invisible in the
inverse coefficients: the two normalizations invert different leading phases
($\lambda x+\tfrac\sigma2\log x$ against
$2\lambda m+\tfrac\sigma2\log m$) and therefore have different Lambert cores.
Only for the even sheet do the two cores agree exactly, through
$\Xcore_-(y)=2w^{\textsf{S1}}_0(y)$; see \cref{p4:rem:halfindex-cores}.
\end{remark}

\subsection{The Binet kernel collapses to a hyperbolic tangent}

\begin{theorem}[Exact Borel--Laplace representation]
\label{p4:thm:laplace}
For $\Re x>0$,
\begin{equation}
  \boxed{\;
  D(x)=\int_0^\infty\e^{-xt}\,\widehat D(t)\,\dd t,
  \qquad
  \widehat D(t)=\frac{\tanh(t/2)}{2t},\;}
  \label{p4:eq:laplace}
\end{equation}
the singularity of $\widehat D$ at $t=0$ being removable with
$\widehat D(0)=\tfrac14$.
\end{theorem}

\begin{proof}
Put $z=x/2$ and
$h(z)=\log\Gamma(z+1)-\log\Gamma(z+\tfrac12)-\tfrac12\log z$, so
$h(x/2)=D(x)$.  From the standard integral for a digamma difference,
\[
  \psi(z+1)-\psi(z+\tfrac12)
  =\int_0^\infty\frac{\e^{-(z+1/2)s}-\e^{-(z+1)s}}{1-\e^{-s}}\,\dd s
  =2\int_0^\infty\frac{\e^{-2zv}}{\e^{v}+1}\,\dd v,
\]
the second equality by $s=2v$ and
$(\e^{-v}-\e^{-2v})/(1-\e^{-2v})=1/(\e^{v}+1)$.  Since
$1/(2z)=\int_0^\infty\e^{-2zv}\dd v$,
\[
  h'(z)=\psi(z+1)-\psi(z+\tfrac12)-\frac1{2z}
  =\int_0^\infty\e^{-2zv}\Bigl(\frac2{\e^{v}+1}-1\Bigr)\dd v
  =-\int_0^\infty\e^{-2zv}\tanh\frac v2\,\dd v .
\]
Both $h(z)$ and the right-hand side of \cref{p4:eq:laplace} tend to $0$ as
$\Re z\to+\infty$, so integrating from $+\infty$ down to $z$,
\[
  h(z)=\int_z^\infty\!\!\int_0^\infty\e^{-2\zeta v}\tanh\tfrac v2
        \,\dd v\,\dd\zeta
      =\int_0^\infty\e^{-2zv}\frac{\tanh(v/2)}{2v}\,\dd v ,
\]
which is \cref{p4:eq:laplace} after $x=2z$.  Fubini applies because the double
integral is absolutely convergent for $\Re z>0$.
\end{proof}

\begin{sourcebox}
\textsf{S3} reaches \cref{p4:eq:laplace} from Binet's first formula, through
the elegant identity
$b(t)-2b(2t)=-\tfrac12\tanh(t/2)$ for
$b(t)=(\e^t-1)^{-1}-t^{-1}+\tfrac12$; \textsf{S2} reaches it from the digamma
integral, as above.  \textsf{S1} keeps Binet's \emph{second} formula
(the $\arctan$ kernel) instead and never simplifies the difference of the two
Binet integrals.  The $\tanh$ collapse is the decisive simplification: it is
what makes the Borel plane elementary, and it is adopted here.
\end{sourcebox}

\begin{corollary}[Sign, monotonicity, and elementary bounds]
\label{p4:cor:D-bounds}
For $x>0$,
\begin{equation}
  0<D(x)<\frac1{4x},
  \qquad
  0<-D'(x)<\frac1{4x^{2}},
  \qquad
  (-1)^{k}D^{(k)}(x)>0\quad(k\ge0).
  \label{p4:eq:D-bounds}
\end{equation}
In particular $D$ is completely monotone and strictly decreasing.
\end{corollary}

\begin{proof}
$\widehat D>0$ on $(0,\infty)$, so
$(-1)^kD^{(k)}(x)=\int_0^\infty\e^{-xt}t^k\widehat D(t)\dd t>0$.  From
$0<\tanh(t/2)<t/2$ we get $0<\widehat D(t)<\tfrac14$, and integrating
$\e^{-xt}t^k$ gives the two displayed bounds.
\end{proof}

\begin{corollary}[Asymptotics of the gamma ratio]
\label{p4:cor:R-asymptotic}
With $R(x)=\Gamma(x/2+1)/\Gamma((x+1)/2)$ as in
\cref{p4:eq:oscillatory-interps},
\[
  \log R(x)=\tfrac12\log(x/2)+D(x)=\tfrac12\log(x/2)+O(x^{-1}),
  \qquad
  \frac{R'(x)}{R(x)}=\frac1{2x}+D'(x)=O(x^{-1}) .
\]
\end{corollary}

\begin{proof}
Immediate from \cref{p4:eq:D-def} and \cref{p4:eq:D-bounds}.
\end{proof}

\subsection{The Borel singularity lattice}

\begin{proposition}[Mittag--Leffler expansion and poles]
\label{p4:prop:borel-lattice}
The Borel transform of \cref{p4:eq:laplace} satisfies
\begin{equation}
  \widehat D(t)=2\sum_{k\ge0}\frac{1}{t^{2}+\pi^{2}(2k+1)^{2}} ,
  \label{p4:eq:partial-fractions}
\end{equation}
is even and meromorphic on $\Cplx$, and has exactly the simple poles
\begin{equation}
  t=\pm\ii\pi(2k+1),\qquad k\ge0,
  \qquad
  \Res_{t=t_k}\widehat D(t)=\frac1{t_k} .
  \label{p4:eq:borel-poles}
\end{equation}
The positive real ray is free of singularities; the nearest singularities lie
at distance $\pi$, so the formal series of $D$ is Gevrey-$1$ with singulant
modulus $\pi$.
\end{proposition}

\begin{proof}
$\tanh(w)=\sum_{k\ge0}2w/\bigl(w^2+\pi^2(k+\tfrac12)^2\bigr)$ is the classical
Mittag--Leffler expansion; putting $w=t/2$ and dividing by $2t$ gives
\cref{p4:eq:partial-fractions}, both sides being even, meromorphic, with the
same poles and residues and the same value $\tfrac14$ at $t=0$ (using
$\sum_{k\ge0}(2k+1)^{-2}=\pi^2/8$).  At $t=t_k$ the function $\tanh(t/2)$ has
residue $2$, so $\widehat D$ has residue $2/(2t_k)=1/t_k$.
\end{proof}

\begin{remark}[What the exponentially small scale is, and is not]
\label{p4:rem:no-real-ambiguity}
The quantity $\e^{-\pi x}$ that governs optimal truncation
(\cref{p4:sec:stokes}) is the envelope set by the distance $\pi$ of the
nearest \emph{off-axis} Borel poles; it is not an ambiguity of the
positive-axis sum, since the ray $\arg t=0$ meets no singularity and
\cref{p4:eq:laplace} is exact.  Genuine Stokes jumps appear only when the
Laplace direction is rotated onto the imaginary rays.  \textsf{S1}, retaining
an unsimplified Binet integral, makes the weaker but correct statement that no
Stokes constant need be assigned on the real axis; the Borel picture explains
why.
\end{remark}

\section{The logarithmic coefficients and a sharp enveloping remainder}
\label{p4:sec:logcoeffs}

\subsection{All coefficients}

We use the Bernoulli convention $z/(\e^z-1)=\sum_{n\ge0}B_nz^n/n!$, so
$B_1=-\tfrac12$, and the Bernoulli polynomials
$t\e^{wt}/(\e^t-1)=\sum_n B_n(w)t^n/n!$.

\begin{theorem}[All logarithmic coefficients]
\label{p4:thm:logcoeffs}
As $x\to\infty$ in any closed subsector of $|\arg x|<\pi/2$,
\begin{equation}
  D(x)\sim\sum_{j\ge1}\frac{d_j}{x^{j}},
  \qquad
  d_{2r}=0,
  \qquad
  d_{2r+1}=\frac{(2^{2r+2}-1)B_{2r+2}}{(2r+1)(2r+2)} .
  \label{p4:eq:d-bernoulli}
\end{equation}
Equivalently
\begin{equation}
  d_{2r+1}=(-1)^{r}\,\frac{2\,(2r)!}{\pi^{2r+2}}\bigl(1-2^{-2r-2}\bigr)
  \zeta(2r+2)
  =(-1)^{r}\,\frac{2\,(2r)!}{\pi^{2r+2}}
   \sum_{k\ge0}\frac{1}{(2k+1)^{2r+2}} ,
  \label{p4:eq:d-zeta}
\end{equation}
and the formal Borel transform of the series is exactly the Taylor expansion
of $\widehat D$:
\begin{equation}
  \widehat D(t)=\sum_{r\ge0}d_{2r+1}\frac{t^{2r}}{(2r)!}
  =\frac{\tanh(t/2)}{2t}.
  \label{p4:eq:borel-taylor}
\end{equation}
\end{theorem}

\begin{proof}
Insert into \cref{p4:eq:partial-fractions} the finite geometric identity
\[
  \frac1{t^2+a^2}=\sum_{r=0}^{N-1}\frac{(-1)^{r}t^{2r}}{a^{2r+2}}
   +\frac{(-1)^{N}t^{2N}}{a^{2N}(t^2+a^2)},
  \qquad a=a_k=\pi(2k+1),
\]
and integrate the finite sum against $\e^{-xt}$ using
$\int_0^\infty\e^{-xt}t^{2r}\dd t=(2r)!x^{-2r-1}$.  This gives the second
expression in \cref{p4:eq:d-zeta}, hence the first by
$\zeta(2r+2)(1-2^{-2r-2})=\sum_{k\ge0}(2k+1)^{-(2r+2)}$, and hence
\cref{p4:eq:d-bernoulli} by Euler's evaluation
$B_{2m}=(-1)^{m-1}2(2m)!\zeta(2m)/(2\pi)^{2m}$.  The unexpanded last term is
the remainder controlled in \cref{p4:thm:envelope}, which in particular proves
the asymptotic statement on $\Re x>0$; the identity \cref{p4:eq:borel-taylor}
is the Taylor expansion of $\widehat D$ at the origin.
\end{proof}

\begin{remark}[Sector of validity: a disagreement between the sources]
\label{p4:rem:sector}
\textsf{S2} and \textsf{S3} state \cref{p4:eq:d-bernoulli} on closed
subsectors of $|\arg x|<\pi/2$, which is what the Laplace representation
\cref{p4:eq:laplace} gives.  \textsf{S1} states the corresponding expansion on
$|\arg m|\le\pi-\delta$.  Both are correct: the wider sector follows not from
the Laplace integral but from the Tricomi--Erd\'elyi expansion of a gamma
ratio,
\begin{equation}
  \log\frac{\Gamma(z+a)}{\Gamma(z+b)}
  \sim(a-b)\log z+\sum_{n\ge1}\frac{(-1)^{n+1}}{n(n+1)z^{n}}
     \bigl(B_{n+1}(a)-B_{n+1}(b)\bigr),
  \qquad |\arg z|\le\pi-\delta,
  \label{p4:eq:tricomi}
\end{equation}
applied with $a=1$, $b=\tfrac12$, $z=x/2$.  The volume states the Laplace
sector where an \emph{exact} representation and a signed remainder are wanted,
and the wider sector where only a Poincar\'e statement is wanted.  The two
routes give the same coefficients: with $B_m(\tfrac12)=(2^{1-m}-1)B_m$ and
$B_m(1)=B_m$ for $m\ne1$, \cref{p4:eq:tricomi} yields
$d_n=2^{n}\,\frac{(-1)^{n+1}}{n(n+1)}\bigl(2-2^{-n}\bigr)B_{n+1}$, which is
\cref{p4:eq:d-bernoulli}.
\end{remark}

\begin{table}[htbp]
\centering
\caption{The logarithmic coefficients $d_{2r+1}$ of \cref{p4:thm:logcoeffs},
recomputed from \cref{p4:eq:d-bernoulli} in exact rational arithmetic and
cross-checked against the Taylor expansion of $\tanh(t/2)/(2t)$.}
\label{p4:tab:d}
\small
\begin{tabular}{@{}rrrrrrrr@{}}
\toprule
$j$ & $1$ & $3$ & $5$ & $7$ & $9$ & $11$ & $13$\\
\midrule
$d_j$ & $\dfrac14$ & $-\dfrac1{24}$ & $\dfrac1{20}$ & $-\dfrac{17}{112}$
      & $\dfrac{31}{36}$ & $-\dfrac{691}{88}$ & $\dfrac{5461}{52}$\\
\bottomrule
\end{tabular}
\end{table}

Thus
\begin{equation}
  D(x)\sim\frac1{4x}-\frac1{24x^{3}}+\frac1{20x^{5}}
  -\frac{17}{112x^{7}}+\frac{31}{36x^{9}}
  -\frac{691}{88x^{11}}+\frac{5461}{52x^{13}}-\cdots
  \label{p4:eq:D-first}
\end{equation}

\subsection{An enveloping remainder, not merely a Poincar\'e one}

\begin{theorem}[Exact remainder and strict next-term envelope]
\label{p4:thm:envelope}
For $x>0$ and $N\ge0$ write
$D(x)=\sum_{r=0}^{N-1}d_{2r+1}x^{-(2r+1)}+R_N(x)$ (empty sum for $N=0$).
Then, with $a_k=\pi(2k+1)$,
\begin{equation}
  R_N(x)=2(-1)^{N}\sum_{k\ge0}a_k^{-2N}
  \int_0^\infty\frac{\e^{-xt}t^{2N}}{t^{2}+a_k^{2}}\,\dd t,
  \label{p4:eq:remainder-integral}
\end{equation}
and consequently
\begin{equation}
  \boxed{\;
  0<(-1)^{N}R_N(x)<\frac{|d_{2N+1}|}{x^{2N+1}}\;}
  \qquad(x>0,\ N\ge0).
  \label{p4:eq:envelope}
\end{equation}
Every truncation therefore brackets $D(x)$ together with its successor.
\end{theorem}

\begin{proof}
Substitute the finite geometric identity of the previous proof into
\cref{p4:eq:partial-fractions} and then into \cref{p4:eq:laplace}; Tonelli's
theorem justifies the interchange since all terms are positive.  This gives
\cref{p4:eq:remainder-integral}, in which every integral is positive.
Replacing $(t^{2}+a_k^{2})^{-1}$ by the strict upper bound $a_k^{-2}$ and
using $\int_0^\infty\e^{-xt}t^{2N}\dd t=(2N)!x^{-2N-1}$ gives
\[
  0<(-1)^{N}R_N(x)<2(2N)!\,x^{-(2N+1)}\sum_{k\ge0}a_k^{-(2N+2)}
  =\frac{|d_{2N+1}|}{x^{2N+1}},
\]
by the last expression in \cref{p4:eq:d-zeta}.
\end{proof}

\begin{remark}
\Cref{p4:thm:envelope} is strictly stronger than the Poincar\'e statement, and
it is what makes every later analytic claim in this chapter effective: all
remainder constants below are explicit rational multiples of the first omitted
$|d_{2N+1}|$.  \textsf{S2} and \textsf{S3} both prove it; \textsf{S1} does not
have it, because it never simplifies the Binet kernel.  It was verified
numerically at $x=5$ and $x=10$ for $0\le N\le4$: the bracketing
\cref{p4:eq:envelope} holds with the two sides agreeing to within a factor
$1.33$ at $x=5$ and $1.09$ at $x=10$.
\end{remark}

\subsection{Large order and least term}

\begin{proposition}[Least term and its envelope]
\label{p4:prop:leastterm}
As $r\to\infty$, $d_{2r+1}\sim(-1)^{r}2(2r)!\pi^{-2r-2}$, so the ratio of
consecutive nonzero term magnitudes in \cref{p4:eq:D-first} is
$\sim(2r+2)(2r+1)/(\pi^{2}x^{2})$.  The least term therefore occurs at
$r\approx\pi x/2$ and has magnitude
\begin{equation}
  \frac{2\sqrt2}{\pi\sqrt x}\,\e^{-\pi x}\bigl(1+O(x^{-1})\bigr).
  \label{p4:eq:least-term}
\end{equation}
By \cref{p4:thm:envelope}, truncation at the least term attains an error of
exactly this scale, with no hidden constant.
\end{proposition}

\begin{proof}
The asymptotic form of $d_{2r+1}$ is \cref{p4:eq:d-zeta} with
$\zeta(2r+2)\to1$.  The ratio statement is immediate.  Setting $2r=\pi x$ and
applying Stirling's formula to $(2r)!$ gives
$2(2r)!\pi^{-2r-2}x^{-2r-1}\sim2\sqrt{2\pi^{2}x}\,\e^{-\pi x}/(\pi^{2}x)
=2\sqrt2\,\e^{-\pi x}/(\pi\sqrt x)$.
\end{proof}

\begin{remark}[Numerical confirmation]
\label{p4:rem:leastterm-numeric}
Recomputed: at $x=10$ the least nonzero term of \cref{p4:eq:D-first} occurs at
$r=15$ (against $\pi x/2=15.708$) with magnitude
$6.5448\times10^{-15}$, while \cref{p4:eq:least-term} predicts
$6.4659\times10^{-15}$; at $x=20$ the least term occurs at $r=31$
(against $31.416$) with magnitude $1.03859\times10^{-28}$ against a predicted
$1.03837\times10^{-28}$.
\end{remark}

\section{The forward transseries}
\label{p4:sec:forward}

\subsection{Exponentiation}

\begin{theorem}[All forward coefficients]
\label{p4:thm:forward}
Define $A_0=1$ and, for $n\ge1$,
\begin{equation}
  A_n=\Coef{t^n}\exp\Bigl(\sum_{j\ge1}d_jt^{j}\Bigr)
     =\frac1{n!}\,\EB_n\bigl(1!\,d_1,2!\,d_2,\ldots,n!\,d_n\bigr)
     =\sum_{\substack{k_1,k_2,\ldots\ge0\\ \sum_j jk_j=n}}
       \prod_{j\ge1}\frac{d_j^{\,k_j}}{k_j!},
  \label{p4:eq:A-formulae}
\end{equation}
with $\EB_n$ the complete exponential Bell polynomial of
\cref{p0:def:ebell}.  Equivalently and most cheaply,
\begin{equation}
  A_0=1,\qquad
  nA_n=\sum_{j=1}^{n}jd_jA_{n-j}
      =\sum_{\substack{1\le j\le n\\ j\ \mathrm{odd}}}jd_jA_{n-j}.
  \label{p4:eq:A-recurrence}
\end{equation}
Then, for every fixed $N\ge1$ and uniformly on closed subsectors of
$\Re x>0$,
\begin{equation}
  \boxed{\;
  \mathcal S_\sigma(x)=C_\sigma\,2^{x}x^{\sigma/2}
  \Bigl(\sum_{n=0}^{N-1}\frac{\sigma^{n}A_n}{x^{n}}+O_N(x^{-N})\Bigr).\;}
  \label{p4:eq:forward}
\end{equation}
\end{theorem}

\begin{proof}
The three expressions in \cref{p4:eq:A-formulae} are the definition of
$\EB_n$, its partition form, and the expansion of the exponential as a
product; \cref{p4:eq:A-recurrence} follows from $A'(t)=A(t)\sum_jjd_jt^{j-1}$
for $A(t)=\exp\sum_jd_jt^j$, together with $d_{2r}=0$.  The parity of the
$\sigma$-dependence is exact: since $\sum_jd_jt^j$ is an odd formal series,
\[
  \exp\Bigl(-\sum_jd_jt^j\Bigr)=\exp\Bigl(\sum_jd_j(-t)^j\Bigr)
  =\sum_{n\ge0}(-1)^{n}A_nt^{n},
\]
so $\exp(\sigma D)$ has coefficients $\sigma^nA_n$ in either sheet.  For the
remainder, \cref{p4:prop:forward-effective} below gives an explicit bound.
\end{proof}

\begin{proposition}[Effective forward remainder]
\label{p4:prop:forward-effective}
Fix $N\ge1$ and let $x\ge x_N:=\max\{1,\,2N\}$.  Then
\begin{equation}
  \Bigl|\e^{\sigma D(x)}-\sum_{n=0}^{N-1}\frac{\sigma^nA_n}{x^{n}}\Bigr|
  \le \frac{K_N}{x^{N}},
  \qquad
  K_N:=\e^{1/4}\Bigl(|d_{2\lceil N/2\rceil+1}|
       +\sum_{n\ge N}\frac{|A_n|}{x_N^{\,n-N}}\Bigr)<\infty,
  \label{p4:eq:forward-effective}
\end{equation}
where the tail sum converges because $A_n=O(C^{n}n!^{1/2})$ for some $C$.
\end{proposition}

\begin{proof}
Write $D_{<M}(x)=\sum_{r<M}d_{2r+1}x^{-(2r+1)}$ with $M=\lceil N/2\rceil$, so
that by \cref{p4:thm:envelope} $|D(x)-D_{<M}(x)|\le|d_{2M+1}|x^{-(2M+1)}\le
|d_{2M+1}|x^{-N}$ for $2M+1\ge N$.  Since $|D(x)|\le\tfrac14$ and
$|D_{<M}(x)|\le\tfrac14$ for $x\ge1$ by \cref{p4:eq:envelope} with $N=0$,
the mean value theorem gives
$|\e^{\sigma D}-\e^{\sigma D_{<M}}|\le\e^{1/4}|D-D_{<M}|$.  Finally
$\e^{\sigma D_{<M}(x)}$ is an absolutely convergent power series in $1/x$ for
$x\ge x_N$ whose coefficients agree with $\sigma^nA_n$ for $n<N$ up to terms
already accounted for, and whose tail is bounded by
$\sum_{n\ge N}|A_n|x^{-n}$.  Collecting the two bounds gives
\cref{p4:eq:forward-effective}.  The growth bound on $A_n$ follows from
$d_{2r+1}=O((2r)!\pi^{-2r})$ and the Bell-polynomial partition form.
\end{proof}

\begin{repairbox}
\textbf{Source claim.}  \textsf{S1} obtains its multiplicative remainder by
writing ``$e^{R_N(m)}=1+O(m^{-N})$, so formal exponentiation also gives the
claimed analytic remainder''.  \textbf{Defect.}  That sentence exponentiates a
\emph{formal} truncation error and silently assumes that the tail of the
exponentiated series is of the same order; no bound on the coefficients $A_n$
enters, although the series $\sum A_nx^{-n}$ diverges.  \textbf{Repair.}
\Cref{p4:prop:forward-effective} supplies the missing ingredients (a bound on
$D-D_{<M}$ from the enveloping theorem, a uniform bound on the exponential's
Lipschitz constant, and an explicit tail bound), so the constant in
\cref{p4:eq:forward} is effective.
\end{repairbox}

\subsection{Coefficient table and the two sectors}

\begin{table}[htbp]
\centering
\caption{The universal forward coefficients $A_n$ of \cref{p4:eq:A-formulae},
recomputed by \cref{p4:eq:A-recurrence} in exact rational arithmetic.  The
table reproduces the corresponding tables of all three sources
(\textsf{S3} tabulates $C_m=(-1)^{n}A_n$; \textsf{S1} tabulates
$A_n/2^{n}$ in its half index).}
\label{p4:tab:A}
\small
\begin{tabular}{@{}rl rl@{}}
\toprule
$n$ & $A_n$ & $n$ & $A_n$\\
\midrule
$0$ & $1$ & $8$ & $-334477/8388608$\\
$1$ & $1/4$ & $9$ & $28717403/33554432$\\
$2$ & $1/32$ & $10$ & $59697183/268435456$\\
$3$ & $-5/128$ & $11$ & $-8400372435/1073741824$\\
$4$ & $-21/2048$ & $12$ & $-34429291905/17179869184$\\
$5$ & $399/8192$ & $13$ & $7199255611995/68719476736$\\
$6$ & $869/65536$ & $14$ & $14631594576045/549755813888$\\
$7$ & $-39325/262144$ & $15$ & $-4251206967062925/2199023255552$\\
\bottomrule
\end{tabular}
\end{table}

Explicitly, with $C_-=\sqrt{2/\pi}$ and $C_+=1/\sqrt{2\pi}$,
\begin{align}
  \mathcal S_-(x)&\sim\sqrt{\frac2\pi}\,\frac{2^{x}}{\sqrt x}
  \Bigl(1-\frac1{4x}+\frac1{32x^{2}}+\frac5{128x^{3}}
        -\frac{21}{2048x^{4}}-\frac{399}{8192x^{5}}
        +\frac{869}{65536x^{6}}+\cdots\Bigr),
  \label{p4:eq:even-forward}\\
  \mathcal S_+(x)&\sim\frac{2^{x}\sqrt x}{\sqrt{2\pi}}
  \Bigl(1+\frac1{4x}+\frac1{32x^{2}}-\frac5{128x^{3}}
        -\frac{21}{2048x^{4}}+\frac{399}{8192x^{5}}
        +\frac{869}{65536x^{6}}+\cdots\Bigr),
  \label{p4:eq:odd-forward}
\end{align}
and for the integer sequence itself,
\begin{equation}
  \boxed{\;
  \operatorname{sw}(n)\sim
  \frac{2^{n}}{\sqrt\pi}\,2^{-\sigma_n/2}\,n^{\sigma_n/2}
  \sum_{k\ge0}\frac{\sigma_n^{\,k}A_k}{n^{k}},
  \qquad \sigma_n=(-1)^{n+1}. \;}
  \label{p4:eq:integer-forward}
\end{equation}

\begin{remark}[Period-two, not oscillatory error]
\label{p4:rem:transmonomial}
\Cref{p4:eq:integer-forward} is a two-phase (period-two) transseries: the
discrete transmonomial $\sigma_n$ has modulus one, so it is never absorbed
into a remainder, and it changes the algebraic scale from $2^nn^{-1/2}$ to
$2^nn^{+1/2}$.  Equivalently, with the idempotents
$\chi_\sigma(n)=(1+\sigma\sigma_n)/2$,
\[
  \operatorname{sw}(n)\sim\sum_{\sigma=\pm1}\chi_\sigma(n)\,
  C_\sigma 2^{n}n^{\sigma/2}\sum_{k\ge0}\sigma^{k}A_kn^{-k}.
\]
\end{remark}

\begin{corollary}[Central binomial coefficient and the reciprocal]
\label{p4:cor:central-and-reciprocal}
Putting $x=2m$ in \cref{p4:eq:even-forward},
\begin{equation}
  \binom{2m}{m}\sim\frac{4^{m}}{\sqrt{\pi m}}
  \Bigl(1-\frac1{8m}+\frac1{128m^{2}}+\frac5{1024m^{3}}
        -\frac{21}{32768m^{4}}-\frac{399}{262144m^{5}}+\cdots\Bigr),
  \label{p4:eq:central-binomial}
\end{equation}
that is, the half-index coefficients of \textsf{S1} are $A_n/2^{n}$ on the
even sheet.  Moreover
\begin{equation}
  \frac1{\mathcal S_\sigma(x)}\sim
  \frac1{C_\sigma}\,2^{-x}x^{-\sigma/2}
  \sum_{n\ge0}\frac{(-\sigma)^{n}A_n}{x^{n}} .
  \label{p4:eq:reciprocal}
\end{equation}
\end{corollary}

\begin{proof}
For \cref{p4:eq:central-binomial} substitute $x=2m$ and
$C_-2^{2m}(2m)^{-1/2}=4^m/\sqrt{\pi m}$.  For \cref{p4:eq:reciprocal} note
that the reciprocal of \cref{p4:eq:normal-form} carries $\e^{-\sigma D}$,
whose coefficients are $(-\sigma)^nA_n$ by the oddness argument in the proof
of \cref{p4:thm:forward}.
\end{proof}

\section{The centred normal form}
\label{p4:sec:centred}

\begin{sourcebox}
This section is \textsf{S1}'s distinctive contribution: a translation of the
index that makes the correction series \emph{sparse in a second way} and
equalizes the two leading constants.  \textsf{S1} states it as two parallel
statements in its half index; it is recast here as a single $\sigma$-uniform
identity in the original index, and the coefficients are recomputed in that
normalization.
\end{sourcebox}

\begin{theorem}[Centred exact normal form]
\label{p4:thm:centred}
Put
\begin{equation}
  u:=\frac x2+\frac14 .
  \label{p4:eq:centred-variable}
\end{equation}
Then, exactly and for both parities,
\begin{equation}
  \boxed{\;
  \mathcal S_\sigma(x)=\frac{4^{u}}{\sqrt{2\pi}}
  \left(\frac{\Gamma(u+3/4)}{\Gamma(u+1/4)}\right)^{\!\sigma},\;}
  \label{p4:eq:centred-exact}
\end{equation}
and, as $u\to\infty$ in $|\arg u|\le\pi-\delta$,
\begin{equation}
  \mathcal S_\sigma(x)\sim\frac{4^{u}}{\sqrt{2\pi}}\;u^{\sigma/2}
  \exp\Bigl(\sigma\sum_{q\ge1}\frac{\nu_{2q}}{u^{2q}}\Bigr),
  \qquad
  \nu_{2q}=\frac{B_{2q+1}(1/4)}{q(2q+1)} .
  \label{p4:eq:centred-expansion}
\end{equation}
In particular \emph{every odd inverse power is absent} from the centred
correction, and the leading constant $1/\sqrt{2\pi}$ is the same on both
sheets.
\end{theorem}

\begin{proof}
From $u=x/2+1/4$ we get $x/2+1=u+3/4$ and $x/2+1/2=u+1/4$, while
$2^{x}=2^{2u-1/2}=4^{u}/\sqrt2$; substituting in \cref{p4:eq:branch-ratio}
gives \cref{p4:eq:centred-exact}.  For \cref{p4:eq:centred-expansion} apply
\cref{p4:eq:tricomi} with $a=\tfrac34$, $b=\tfrac14$, $z=u$: the logarithmic
term is $(a-b)\log u=\tfrac12\log u$, and the $n$-th coefficient is
$\frac{(-1)^{n+1}}{n(n+1)}\bigl(B_{n+1}(\tfrac34)-B_{n+1}(\tfrac14)\bigr)$.
Since $a+b=1$, the reflection formula $B_m(1-w)=(-1)^{m}B_m(w)$ gives
$B_{n+1}(\tfrac14)=(-1)^{n+1}B_{n+1}(\tfrac34)$, so the difference vanishes
for odd $n$ and equals $2B_{n+1}(\tfrac34)=-2B_{n+1}(\tfrac14)$ for even
$n=2q$.  This yields $\nu_{2q}=B_{2q+1}(1/4)/(q(2q+1))$.
\end{proof}

\begin{table}[htbp]
\centering
\caption{Centred logarithmic coefficients $\nu_{2q}$ of
\cref{p4:eq:centred-expansion}, recomputed from Bernoulli polynomials at
$1/4$.  They agree with \textsf{S1}'s $\widehat\lambda^{(0)}_{2q}=-\nu_{2q}$
after the marker translation of \cref{p4:rem:dictionary}.}
\label{p4:tab:nu}
\small
\begin{tabular}{@{}rrrrrr@{}}
\toprule
$2q$ & $2$ & $4$ & $6$ & $8$ & $10$\\
\midrule
$\nu_{2q}$ & $\dfrac1{64}$ & $-\dfrac5{2048}$ & $\dfrac{61}{49152}$
 & $-\dfrac{1385}{1048576}$ & $\dfrac{50521}{20971520}$\\
\bottomrule
\end{tabular}
\end{table}

\begin{remark}[Two different sparsities]
\label{p4:rem:two-sparsities}
The uncentred correction $\sigma D(x)$ contains only \emph{odd} powers of
$1/x$; the centred correction contains only \emph{even} powers of $1/u$.
There is no contradiction: the two statements concern different remainder
functions, because the $\tfrac\sigma2\log$ term is attached to $x$ in one and
to $u$ in the other.  Both are instances of one rule, recorded in
\cref{p4:lem:centring} below, and the swing shifts are precisely the
optimally centred ones.
\end{remark}

\begin{lemma}[Optimal centring of a gamma ratio]
\label{p4:lem:centring}
Let $a,b\in\Real$ with $\beta:=a-b\ne0$, and set $s=(a+b-1)/2$ and $z=w+s$.
Then $\Gamma(w+a)/\Gamma(w+b)=\Gamma(z+\tfrac{1+\beta}2)/
\Gamma(z+\tfrac{1-\beta}2)$, the two shifts sum to $1$, and
\begin{equation}
  \log\frac{\Gamma(z+\frac{1+\beta}2)}{\Gamma(z+\frac{1-\beta}2)}
  \sim\beta\log z-\sum_{q\ge1}
   \frac{B_{2q+1}\bigl(\frac{1-\beta}2\bigr)}{q(2q+1)}\,z^{-2q},
  \qquad |\arg z|\le\pi-\delta,
  \label{p4:eq:centring}
\end{equation}
with no odd inverse powers.
\end{lemma}

\begin{proof}
The shift identity is immediate.  Apply \cref{p4:eq:tricomi} and the argument
of \cref{p4:thm:centred}: complementary shifts make the Bernoulli difference
vanish in odd degree by $B_m(1-w)=(-1)^mB_m(w)$.
\end{proof}

For the swing factorial $a=1$, $b=\tfrac12$, $w=x/2$, so $s=\tfrac14$ and
$z=u$, and $\beta=\tfrac12$: \cref{p4:eq:centred-expansion} is exactly
\cref{p4:eq:centring} raised to the power $\sigma$.  The first centred
expansions are
\begin{align}
  \mathcal S_-(x)&\sim\frac{4^{u}}{\sqrt{2\pi u}}
  \Bigl(1-\frac1{64u^{2}}+\frac{21}{8192u^{4}}
   -\frac{671}{524288u^{6}}+\cdots\Bigr),
  \label{p4:eq:centred-even}\\
  \mathcal S_+(x)&\sim\frac{4^{u}\sqrt u}{\sqrt{2\pi}}
  \Bigl(1+\frac1{64u^{2}}-\frac{19}{8192u^{4}}
   +\frac{631}{524288u^{6}}+\cdots\Bigr),
  \label{p4:eq:centred-odd}
\end{align}
whose coefficients were recomputed from
$\exp\bigl(\sigma\sum_q\nu_{2q}z^{q}\bigr)$ and agree with \textsf{S1}.

\section{The two Lambert cores, and why the branch matters}
\label{p4:sec:lambert}

\subsection{Specialization of the core theorem}

Write
\begin{equation}
  L_\sigma(y):=\log\frac{y}{C_\sigma}
   =\log y+\tfrac12\log\pi+\tfrac\sigma2\lambda .
  \label{p4:eq:L-def}
\end{equation}
Dropping the correction in \cref{p4:eq:normal-form} leaves the \emph{core
equation}
\begin{equation}
  L_\sigma(y)=\lambda \Xcore+\frac\sigma2\log \Xcore,
  \qquad\text{equivalently}\qquad
  C_\sigma\,2^{\Xcore}\Xcore^{\sigma/2}=y .
  \label{p4:eq:core-equation}
\end{equation}

\begin{theorem}[Explicit Lambert cores]
\label{p4:thm:cores}
\Cref{p0:thm:lambert-core} with the dictionary \cref{p4:eq:dictionary} gives
\begin{equation}
  \Xcore_{\sigma,k}(y)=\frac{\sigma}{2\lambda}\,
  \Wlam_k\!\bigl(2\sigma\lambda\,\e^{2\sigma L_\sigma(y)}\bigr),
  \label{p4:eq:core-general}
\end{equation}
and the unique large positive real solutions are
\begin{equation}
  \boxed{\;
  \Xcore_+(y)=\frac{\Wlam_0\!\bigl(4\pi\lambda y^{2}\bigr)}{2\lambda},
  \qquad
  \Xcore_-(y)=-\frac{\Wlam_{-1}\!\bigl(-4\lambda/(\pi y^{2})\bigr)}{2\lambda}.
  \;}
  \label{p4:eq:cores}
\end{equation}
$\Xcore_+$ is defined for all $y>0$; $\Xcore_-$ is defined exactly for
\begin{equation}
  y\ge y_\star:=2\sqrt{\frac{\e\lambda}{\pi}}=1.548870223880\ldots,
  \label{p4:eq:ystar}
\end{equation}
and satisfies $\Xcore_-(y)\ge1/(2\lambda)=0.72134752\ldots$ there.
\end{theorem}

\begin{proof}
With $a=\lambda$, $b=\sigma/2$, \cref{p0:thm:lambert-core} reads
$\Xcore=(b/a)\Wlam_k\bigl((a/b)\e^{L/b}\bigr)$, and $1/\sigma=\sigma$ gives
\cref{p4:eq:core-general}.  For $\sigma=+1$, $C_+=1/\sqrt{2\pi}$ so
$\e^{2L_+}=2\pi y^{2}$ and the argument is $4\pi\lambda y^{2}>0$, on which
$\Wlam_0$ is the only real branch.  For $\sigma=-1$, $C_-=\sqrt{2/\pi}$ so
$\e^{-2L_-}=2/(\pi y^{2})$ and the argument is $-4\lambda/(\pi y^{2})<0$;
it lies in $[-1/\e,0)$ exactly when $y^{2}\ge4\lambda\e/\pi$, i.e.\
$y\ge y_\star$.  On that interval both $\Wlam_0$ and $\Wlam_{-1}$ are real;
$\Wlam_{-1}\le-1$ gives $\Xcore_-\ge1/(2\lambda)$, and
$\Wlam_{-1}(z)\to-\infty$ as $z\to0^-$ gives $\Xcore_-(y)\to\infty$.
\Cref{p4:prop:two-roots} identifies the two roots and shows that $\Wlam_0$
gives the wrong one.
\end{proof}

\subsection{Why $\Wlam_0$ is the wrong root on the even sheet}

\begin{proposition}[The even carrier is not injective]
\label{p4:prop:two-roots}
Let $g_-(\Xcore)=C_-2^{\Xcore}\Xcore^{-1/2}$ be the even carrier of
\cref{p4:eq:core-equation}.  Then $g_-$ is strictly decreasing on
$(0,1/(2\lambda))$, strictly increasing on $(1/(2\lambda),\infty)$, and
\begin{equation}
  \min_{\Xcore>0}g_-(\Xcore)=g_-\Bigl(\frac1{2\lambda}\Bigr)
  =2\sqrt{\frac{\e\lambda}{\pi}}=y_\star .
  \label{p4:eq:carrier-min}
\end{equation}
Hence for $y>y_\star$ the equation $g_-(\Xcore)=y$ has exactly two positive
roots $\Xcore^{\mathrm{small}}(y)<1/(2\lambda)<\Xcore^{\mathrm{large}}(y)$;
they coalesce at $y=y_\star$.  In terms of the Lambert branches,
\begin{equation}
  \Xcore^{\mathrm{large}}(y)
   =-\frac{\Wlam_{-1}\bigl(-4\lambda/(\pi y^{2})\bigr)}{2\lambda}
   =\Xcore_-(y)\xrightarrow[y\to\infty]{}\infty,
  \qquad
  \Xcore^{\mathrm{small}}(y)
   =-\frac{\Wlam_{0}\bigl(-4\lambda/(\pi y^{2})\bigr)}{2\lambda}
   =\frac{2}{\pi y^{2}}\bigl(1+O(y^{-2})\bigr).
  \label{p4:eq:two-roots}
\end{equation}
The principal branch therefore returns a root that \emph{tends to zero}; it is
not a large-index inverse at any order, and no amount of asymptotic correction
can repair the choice.  On the odd sheet the carrier
$g_+(\Xcore)=C_+2^{\Xcore}\Xcore^{1/2}$ is a strictly increasing bijection of
$(0,\infty)$ onto $(0,\infty)$, its Lambert argument is positive, and
$\Wlam_0$ is the only real branch.
\end{proposition}

\begin{proof}
$(\log g_-)'(\Xcore)=\lambda-1/(2\Xcore)$ vanishes only at
$\Xcore=1/(2\lambda)$ and changes sign there from $-$ to $+$; the minimum
value is
$C_-\e^{\lambda/(2\lambda)}(2\lambda)^{1/2}
=\sqrt{2/\pi}\cdot\sqrt\e\cdot\sqrt{2\lambda}=2\sqrt{\e\lambda/\pi}$,
which is \cref{p4:eq:carrier-min}, and equals $y_\star$ of
\cref{p4:eq:ystar}.  Strict monotonicity on each side gives exactly two roots
for $y>y_\star$.  Both are given by \cref{p4:eq:core-general} with $\sigma=-1$
and a real branch of $\Wlam$; since $\Wlam_{-1}\le-1\le\Wlam_0\le0$ on
$[-1/\e,0)$, the branch $\Wlam_{-1}$ gives the root
$\ge1/(2\lambda)$ and $\Wlam_0$ the root $\le1/(2\lambda)$.  Finally
$\Wlam_0(z)=z+O(z^{2})$ as $z\to0$, so
$\Xcore^{\mathrm{small}}=-\bigl(-4\lambda/(\pi y^2)+O(y^{-4})\bigr)/(2\lambda)
=2/(\pi y^{2})+O(y^{-4})$.
\end{proof}

\begin{remark}[Numerical confirmation and the branch point]
\label{p4:rem:small-root}
Recomputed at $50$ digits: for $y=10^{2},10^{3},10^{5}$ the principal-branch
root is $6.366759642\times10^{-5}$, $6.366203342\times10^{-7}$,
$6.366197724\times10^{-11}$, against $2/(\pi y^{2})=6.366197724\times10^{-5}$,
$6.366197724\times10^{-7}$, $6.366197724\times10^{-11}$.  The threshold
$y_\star$ where the two roots coalesce is exactly the value at which the
Lambert argument reaches the branch point $-1/\e$; the agreement of these two
descriptions is a useful consistency check on the whole normalization.  Note
also $\mathcal S_-\bigl(1/(2\lambda)\bigr)=1.15213\ldots<y_\star$, a fact used
in \cref{p4:prop:core-inequalities}.
\end{remark}

\begin{remark}[The branch rule is forced by $\operatorname{sign}b$]
\label{p4:rem:branch-rule}
All three sources agree on the assignment: $\Wlam_{-1}$ on the even sheet,
$\Wlam_0$ on the odd sheet.  In the model \cref{p0:def:model} the rule is
$b<0\Rightarrow\Wlam_{-1}$, $b>0\Rightarrow\Wlam_0$ for the large real root,
and $b=\sigma/2$; so the branch distinction is forced by the sign of the
power of $x$ in the forward asymptotics --- that is, by the fact that the even
bisection carries $n^{-1/2}$ and the odd bisection $n^{+1/2}$.  It is a
consequence of the parity split, not an independent convention.
\end{remark}

\subsection{Two exact inequalities}

The sources deduce the \emph{sign} of the first inverse correction from
$D>0$.  The following sharpening is exact at all $y$, not only to first order,
and costs nothing.

\begin{proposition}[The core brackets the true inverse]
\label{p4:prop:core-inequalities}
For every $y>1$,
\begin{equation}
  \mathcal I_+(y)<\Xcore_+(y),
  \label{p4:eq:odd-bracket}
\end{equation}
and for every $y>y_\star$,
\begin{equation}
  \mathcal I_-(y)>\Xcore_-(y).
  \label{p4:eq:even-bracket}
\end{equation}
\end{proposition}

\begin{proof}
By \cref{p4:eq:normal-form} and \cref{p4:cor:D-bounds},
$\mathcal S_+(x)=g_+(x)\e^{D(x)}>g_+(x)$ for all $x>0$, where $g_+$ is the
odd carrier.  Both $\mathcal S_+$ and $g_+$ are strictly increasing.  Put
$x_1=\mathcal I_+(y)$; then
$g_+(x_1)<\mathcal S_+(x_1)=y=g_+(\Xcore_+(y))$, and $g_+$ increasing gives
$x_1<\Xcore_+(y)$.

For $\sigma=-1$, $\mathcal S_-(x)=g_-(x)\e^{-D(x)}<g_-(x)$.  Let
$x_0=\mathcal I_-(y)$ with $y>y_\star$.  Since
$\mathcal S_-$ is increasing and $\mathcal S_-(1/(2\lambda))
=1.15213\ldots<y_\star<y$, we have $x_0>1/(2\lambda)$; also
$\Xcore_-(y)\ge1/(2\lambda)$ by \cref{p4:thm:cores}.  On
$[1/(2\lambda),\infty)$ the carrier $g_-$ is strictly increasing
(\cref{p4:prop:two-roots}), and
$g_-(x_0)>\mathcal S_-(x_0)=y=g_-(\Xcore_-(y))$, whence
$x_0>\Xcore_-(y)$.
\end{proof}

\begin{remark}
This explains without computation the signs $c_1^{(+1)}=-1/(4\lambda)<0$ and
$c_1^{(-1)}=+1/(4\lambda)>0$ found in \cref{p4:sec:reversion}, and supplies a
rigorous one-sided bracket at every $y$, used as a check in
\cref{p4:alg:invert}.
\end{remark}

\subsection{A scaled form, and the half-index cores}

Putting $w_\sigma=2\lambda\Xcore_\sigma$ turns \cref{p4:eq:core-equation} into
the single scalar equation
\begin{equation}
  w+\sigma\log w=\Lambda_\sigma,
  \qquad
  \Lambda_\sigma:=\log(\pi y^{2})+\sigma\log(4\lambda),
  \label{p4:eq:scaled-core}
\end{equation}
so that $w_+=\Wlam_0(\e^{\Lambda_+})$ and
$w_-=-\Wlam_{-1}(-\e^{-\Lambda_-})$.  This is \textsf{S3}'s normalization,
with its $\varepsilon=-\sigma$; it is convenient in
\cref{p4:sec:flatten}.

\begin{remark}[Half-index cores]
\label{p4:rem:halfindex-cores}
\textsf{S1}'s carriers, in the half index and with $\alpha=\log4$, are
$w^{\textsf{S1}}_0=-\tfrac1{2\alpha}\Wlam_{-1}(-2\alpha/(\pi y^{2}))$ and
$w^{\textsf{S1}}_1=\tfrac1{2\alpha}\Wlam_0(\pi\alpha y^{2}/2)$.  On the even
sheet $n=2m$ exactly, and one checks directly that
$\Xcore_-(y)=2w^{\textsf{S1}}_0(y)$; the two normalizations then agree, and
their inverse coefficients are related by
$c_n^{(-1)}=2^{\,n+1}d_n^{\textsf{S1},(0)}$, which we verified symbolically
for $1\le n\le8$.  On the odd sheet $n=2m+1$, so the two carriers solve
\emph{different} equations ($C_+2^{X}X^{1/2}=y$ against
$K_14^{w}w^{1/2}=y$) and their coefficient lists are \emph{not} termwise
comparable, although both expansions converge asymptotically to the same
function $\mathcal I_+$.  Failing to notice this is the easiest way to
``discover'' a contradiction between \textsf{S1} and \textsf{S2}/\textsf{S3}
where there is none.
\end{remark}

\section{All-orders reversion around the core}
\label{p4:sec:reversion}

\subsection{Specialization of the reversion theorem}

\begin{theorem}[Swing specialization of \cref{p0:thm:lambert-centered}]
\label{p4:thm:reversion}
Let $\Xcore=\Xcore_\sigma(y)$ solve \cref{p4:eq:core-equation}, let
$t=\Xcore^{-1}$, and put
\begin{equation}
  \Psi_\sigma(t,u)=-\frac1\lambda\left[
    \frac\sigma2\log(1+tu)
    +\sigma\sum_{j\ge1}d_j\Bigl(\frac{t}{1+tu}\Bigr)^{j}\right].
  \label{p4:eq:Psi-swing}
\end{equation}
Then the unique $U\in t\,\mathscr R[[t]]$ with $U=\Psi_\sigma(t,U)$ has
coefficients
\begin{equation}
  c_n^{(\sigma)}=\Coef{t^n}U
  =\sum_{m=1}^{n}\frac1m\,\Coef{t^{n}u^{m-1}}\Psi_\sigma(t,u)^{m}
  \label{p4:eq:c-lagrange}
\end{equation}
by \cref{p0:lem:lagrange-fp}, and equivalently, with $\gamma_1=0$ and
$\gamma_r=c^{(\sigma)}_{r-1}$ for $r\ge2$,
\begin{equation}
  c^{(\sigma)}_n=-\frac1\lambda\Biggl[
   \frac\sigma2\sum_{k=1}^{n}\frac{(-1)^{k+1}}{k}\OB_{n,k}(\gamma)
   +\sigma\!\!\sum_{\substack{1\le j\le n\\ j\ \mathrm{odd}}}\!\! d_j
     \sum_{k=0}^{n-j}(-1)^{k}\binom{j+k-1}{k}\OB_{n-j,k}(\gamma)\Biggr],
  \label{p4:eq:c-bell}
\end{equation}
with $\OB$ the ordinary partial Bell polynomials of \cref{p0:def:obell} and
the expansions of \cref{p0:lem:power-log}.  Only
$c^{(\sigma)}_1,\ldots,c^{(\sigma)}_{n-1}$ occur on the right.  The
coefficients lie in the parity algebra $\mathscr R$ of
\cref{p4:eq:parity-algebra}; no logarithm of $L_\sigma$ occurs.
\end{theorem}

\begin{proof}
This is \cref{p0:thm:lambert-centered} with the dictionary
\cref{p4:eq:dictionary}: $a=\lambda$, $b=\sigma/2$, $q_j=\sigma d_j$, which
turns the general $\Psi$ into \cref{p4:eq:Psi-swing}.  The restriction of the
inner sum to odd $j$ is $d_{2r}=0$.
\end{proof}

\subsection{Recomputed coefficients}

The coefficients below were recomputed from \cref{p4:eq:Psi-swing} by
truncated power-series arithmetic over $\Rat[\lambda,\lambda^{-1}]$,
separately for $\sigma=+1$ and $\sigma=-1$, and then recombined into
$\mathscr R$ by
$c=\tfrac12(c^{+}+c^{-})+\sigma\cdot\tfrac12(c^{+}-c^{-})$.  They agree with
\textsf{S2} for $1\le n\le10$ and with \textsf{S3} for $1\le n\le8$ after the
marker translation $\sigma=-\varepsilon_{\textsf{S3}}$.

\begin{align}
  c_1^{(\sigma)}&=-\frac{\sigma}{4\lambda},
  \qquad
  c_2^{(\sigma)}=\frac1{8\lambda^{2}},
  \label{p4:eq:c12}\\
  c_3^{(\sigma)}&=\frac{2\sigma\lambda^{2}-3\lambda-3\sigma}{48\lambda^{3}},
  \qquad
  c_4^{(\sigma)}=\frac{-4\lambda^{2}+15\sigma\lambda+6}{192\lambda^{4}},
  \label{p4:eq:c34}\\
  c_5^{(\sigma)}&=-\frac{96\sigma\lambda^{4}-80\lambda^{3}
    +40\sigma\lambda^{2}+135\lambda+30\sigma}{1920\lambda^{5}},
  \label{p4:eq:c5}\\
  c_6^{(\sigma)}&=\frac{96\lambda^{4}-180\sigma\lambda^{3}+215\lambda^{2}
    +210\sigma\lambda+30}{3840\lambda^{6}},
  \label{p4:eq:c6}\\
  c_7^{(\sigma)}&=\frac{1}{53760\lambda^{7}}\bigl(
    8160\sigma\lambda^{6}-4312\lambda^{5}+1428\sigma\lambda^{4}
    +1050\lambda^{3}\nonumber\\
   &\hspace{31mm}{}-4025\sigma\lambda^{2}-2100\lambda-210\sigma\bigr),
  \label{p4:eq:c7}\\
  c_8^{(\sigma)}&=\frac{1}{645120\lambda^{8}}\bigl(
    -48960\lambda^{6}+56056\sigma\lambda^{5}-40908\lambda^{4}
    +15015\sigma\lambda^{3}\nonumber\\
   &\hspace{31mm}{}+50610\lambda^{2}+17010\sigma\lambda+1260\bigr),
  \label{p4:eq:c8}\\
  c_9^{(\sigma)}&=-\frac{1}{1290240\lambda^{9}}\bigl(
    1111040\sigma\lambda^{8}-413184\lambda^{7}+79616\sigma\lambda^{6}
    +43512\lambda^{5}\nonumber\\
   &\hspace{20mm}{}-83748\sigma\lambda^{4}+83475\lambda^{3}
    +92610\sigma\lambda^{2}+22050\lambda+1260\sigma\bigr),
  \label{p4:eq:c9}\\
  c_{10}^{(\sigma)}&=\frac{1}{5160960\lambda^{10}}\bigl(
    2222080\lambda^{8}-1756032\sigma\lambda^{7}+779152\lambda^{6}
    -203280\sigma\lambda^{5}\nonumber\\
   &\hspace{20mm}{}-162582\lambda^{4}+487305\sigma\lambda^{3}
    +309960\lambda^{2}+55440\sigma\lambda+2520\bigr).
  \label{p4:eq:c10}
\end{align}

Setting $\sigma=\pm1$ and $\lambda=\log2$,
\begin{align}
  \mathcal I_+(y)&\sim \Xcore_+
   -\frac1{4\lambda \Xcore_+}+\frac1{8\lambda^{2}\Xcore_+^{2}}
   +\frac{2\lambda^{2}-3\lambda-3}{48\lambda^{3}\Xcore_+^{3}}
   +\frac{-4\lambda^{2}+15\lambda+6}{192\lambda^{4}\Xcore_+^{4}}+\cdots,
  \label{p4:eq:odd-inverse}\\
  \mathcal I_-(y)&\sim \Xcore_-
   +\frac1{4\lambda \Xcore_-}+\frac1{8\lambda^{2}\Xcore_-^{2}}
   +\frac{-2\lambda^{2}-3\lambda+3}{48\lambda^{3}\Xcore_-^{3}}
   +\frac{-4\lambda^{2}-15\lambda+6}{192\lambda^{4}\Xcore_-^{4}}+\cdots.
  \label{p4:eq:even-inverse}
\end{align}
The signs of the first corrections are those forced by
\cref{p4:prop:core-inequalities}.

\subsection{A differential form of the same coefficients}

All three sources also give an operator form of the reversion, which is the
convenient one for the Stokes analysis of \cref{p4:sec:stokes}.  We state it
as a chapter-local lemma; it is a candidate for promotion, being independent
of the swing factorial.

\begin{lemma}[Perturbed inverse in phase coordinates]
\label{p4:lem:phase-lb}
Let $F_0$ be invertible near the point considered with $F_0'\ne0$, let $g$ be
a formal correction, and let $\Xcore$ satisfy $F_0(\Xcore)=z$.  Introduce
\begin{equation}
  \mathscr D:=\frac1{F_0'(\Xcore)}\frac{\dd}{\dd \Xcore}.
  \label{p4:eq:phase-operator}
\end{equation}
Then the solution $x$ of $F_0(x)+\eta\,g(x)=z$ with $x|_{\eta=0}=\Xcore$ is,
as a formal power series in the marker $\eta$,
\begin{equation}
  x=\Xcore+\sum_{k\ge1}\frac{(-\eta)^{k}}{k!}\,
   \mathscr D^{\,k-1}\!\left(\frac{g(\Xcore)^{k}}{F_0'(\Xcore)}\right).
  \label{p4:eq:phase-lb}
\end{equation}
\end{lemma}

\begin{proof}
Let $J=F_0^{-1}$ and $v=F_0(x)$, so the equation becomes
$v+\eta A(v)=z$ with $A(v)=g(J(v))$.  Lagrange--B\"urmann inversion
(\cref{p0:thm:LB}) applied to the observable $J$ gives
\[
  J(v)=J(z)+\sum_{k\ge1}\frac{(-\eta)^{k}}{k!}
   \Bigl(\frac{\dd}{\dd z}\Bigr)^{k-1}\bigl(J'(z)A(z)^{k}\bigr).
\]
Now $J(z)=\Xcore$, $J'(z)=1/F_0'(\Xcore)$, $A(z)=g(\Xcore)$ and
$\dd/\dd z=\mathscr D$, which is \cref{p4:eq:phase-lb}.
\end{proof}

\begin{proposition}[Differential formula for $c_n^{(\sigma)}$]
\label{p4:prop:c-differential}
Let $\mathcal Q_\sigma(t)=\sigma\sum_{j\ge1}d_jt^{j}$ and
\begin{equation}
  \mathscr L_\sigma:=-\frac{t^{2}}{\lambda+\sigma t/2}\frac{\dd}{\dd t}
  \qquad(t=\Xcore^{-1}),
  \label{p4:eq:L-operator}
\end{equation}
which is \cref{p4:eq:phase-operator} for
$F_{0,\sigma}(\Xcore)=\lambda \Xcore+\tfrac\sigma2\log \Xcore$ expressed in
the variable $t$.  Then, for every $n\ge1$,
\begin{equation}
  \boxed{\;
  c_n^{(\sigma)}=\Coef{t^{n}}
   \sum_{k=1}^{\lfloor(n+1)/2\rfloor}\frac{(-1)^{k}}{k!}\,
   \mathscr L_\sigma^{\,k-1}
   \left(\frac{\mathcal Q_\sigma(t)^{k}}{\lambda+\sigma t/2}\right),\;}
  \label{p4:eq:c-differential}
\end{equation}
a \emph{finite} sum requiring only $d_1,\ldots,d_{n}$.
\end{proposition}

\begin{proof}
Apply \cref{p4:lem:phase-lb} with $F_0=F_{0,\sigma}$, $g=\sigma D$ and
$\eta=1$; the substitution $t=\Xcore^{-1}$ turns
$F_{0,\sigma}'(\Xcore)=\lambda+\sigma/(2\Xcore)$ into
$\lambda+\sigma t/2$ and $\dd/\dd \Xcore$ into $-t^{2}\dd/\dd t$.  Setting
the marker $\eta$ to $1$ is legitimate coefficientwise because
$\mathcal Q_\sigma(t)^{k}=O(t^{k})$ while each application of
$\mathscr L_\sigma$ raises the order in $t$ by one, so the $k$-th summand
begins at $t^{2k-1}$ and can contribute to $\Coef{t^{n}}$ only when
$2k-1\le n$.
\end{proof}

\begin{repairbox}
\textbf{Source claim.}  \textsf{S2} writes its operator expansion as an
\emph{equality},
$\mathcal I_\sigma(y)=\Xcore+\sum_{m\ge1}\frac{(-\sigma)^{m}}{m!}
\mathscr D^{m-1}\bigl[D(\Xcore)^{m}/F_0'(\Xcore)\bigr]$, after setting its
coupling parameter to $1$.  \textbf{Defect.}  With the marker removed there is
no small parameter, the series over $m$ does not converge, and the left-hand
side is a number while the right-hand side is a divergent series; the
statement as printed is false.  \textbf{Repair.}  The correct statement is the
asymptotic one: the $m$-th summand is $O(\Xcore^{-(2m-1)})$, so for each fixed
$M$ the partial sum through $m=M$ differs from $\mathcal I_\sigma(y)$ by
$O(\Xcore^{-(2M+1)})$, and regrouping by powers of $\Xcore^{-1}$ reproduces
\cref{p4:eq:c-differential}.  \textsf{S1} and \textsf{S3} state the same
formula correctly, as a coefficient identity.
\end{repairbox}

\subsection{The analytic statement, with a proof}

\begin{repairbox}
\textbf{Source claim.}  All three sources assert the Poincar\'e statement
``$\mathcal I_\sigma(y)=\Xcore_\sigma(y)+\sum_{n<M}c^{(\sigma)}_n
\Xcore_\sigma(y)^{-n}+O(\Xcore_\sigma(y)^{-M})$''.  \textbf{Defect.}
\textsf{S3} gives no proof at all (``Likewise, for each fixed $M$, \ldots'');
\textsf{S1} gives a two-line sketch that appeals to the mean value theorem but
never bounds the residual of the truncated ansatz, and in particular never
uses any control on the forward remainder at the \emph{perturbed} point;
\textsf{S2} proves the ingredients (a residual-to-index bound and derivative
lower bounds) but does not assemble them.  This is the one place in the three
sources where a formal computation is upgraded to an analytic conclusion
without an argument.  \textbf{Repair.}  \Cref{p4:thm:inverse-analytic} below
supplies the proof, and makes the constant effective by using the enveloping
remainder \cref{p4:thm:envelope}.
\end{repairbox}

\begin{theorem}[Inverse transseries, analytically]
\label{p4:thm:inverse-analytic}
Fix $M\ge1$ and $\sigma\in\{-1,+1\}$.  Then, as $y\to+\infty$,
\begin{equation}
  \boxed{\;
  \mathcal I_\sigma(y)=\Xcore_\sigma(y)
   +\sum_{n=1}^{M-1}\frac{c^{(\sigma)}_n}{\Xcore_\sigma(y)^{n}}
   +O\!\left(\Xcore_\sigma(y)^{-M}\right),\;}
  \label{p4:eq:inverse-analytic}
\end{equation}
with an effective implied constant depending only on $M$ and $\sigma$.  Since
$\Xcore_\sigma(y)=\bigl(\log y\bigr)/\lambda\,(1+o(1))$, the expansion is one
in inverse powers of $\log y$ once the dominant power--logarithmic phase has
been inverted exactly.
\end{theorem}

\begin{proof}
Write $\Xcore=\Xcore_\sigma(y)$ and
$x_M:=\Xcore+\sum_{n=1}^{M-1}c^{(\sigma)}_n\Xcore^{-n}$; for $y$ large,
$x_M=\Xcore(1+O(\Xcore^{-2}))$, so $x_M>1$ and $x_M\asymp\Xcore$.  Consider
the phase residual
\[
  \rho:=\log\mathcal S_\sigma(x_M)-\log y
   =\lambda(x_M-\Xcore)
     +\frac\sigma2\log\frac{x_M}{\Xcore}+\sigma D(x_M),
\]
where we used \cref{p4:eq:normal-form} and
$\log y=\log C_\sigma+\lambda\Xcore+\tfrac\sigma2\log\Xcore$.  Put
$t=\Xcore^{-1}$ and $U_{M-1}(t)=\sum_{n<M}c^{(\sigma)}_nt^{n}$, so that
$x_M=\Xcore(1+tU_{M-1})$ and
\[
  \rho=\lambda\,\Xcore\Bigl[\,\underbrace{tU_{M-1}
    +\frac{\sigma t}{2\lambda}\log(1+tU_{M-1})
    +\frac{\sigma t}{\lambda}D_{<M}(x_M)}_{=:t\,\mathcal E(t)}\,\Bigr]
   +\sigma\bigl(D(x_M)-D_{<M}(x_M)\bigr),
\]
where $D_{<M}$ denotes the truncation of the series
\cref{p4:eq:d-bernoulli} after the terms of order $x^{-(M-1)}$.  Two
estimates finish the proof.

\emph{(i) The algebraic residual.}  Substituting $x_M=\Xcore(1+tU_{M-1})$ and
$D_{<M}(x_M)=\mathcal Q_\sigma\bigl(t/(1+tU_{M-1})\bigr)/\sigma$ truncated,
the bracket $\mathcal E(t)$ is exactly
$U_{M-1}(t)-\Psi_\sigma\bigl(t,U_{M-1}(t)\bigr)$ modulo terms of order
$t^{M}$ coming from the truncation of $\mathcal Q_\sigma$.  By
\cref{p4:thm:reversion} the coefficients of $t^{1},\ldots,t^{M-1}$ in
$U-\Psi_\sigma(t,U)$ vanish when $U$ is replaced by $U_{M-1}$ --- this is the
triangularity of the recurrence --- so $\mathcal E(t)=O(t^{M})$ with an
explicit rational-in-$\lambda$ constant obtained from
$c^{(\sigma)}_1,\ldots,c^{(\sigma)}_{M-1}$ and $d_1,\ldots,d_{M}$.  Hence
$\lambda\Xcore\cdot t\,\mathcal E(t)=O(\Xcore^{-M+1}\cdot\Xcore^{-1})\cdot
\Xcore=O(\Xcore^{-M+1})$; more carefully, $\Xcore\,t\,\mathcal E(t)
=\mathcal E(t)=O(\Xcore^{-M})$.

\emph{(ii) The transcendental residual.}  By \cref{p4:thm:envelope} with
$N=\lceil M/2\rceil$,
$|D(x_M)-D_{<M}(x_M)|\le|d_{2N+1}|\,x_M^{-(2N+1)}\le
c_M\,\Xcore^{-M}$ for an explicit $c_M$, since $2N+1\ge M$ and
$x_M\asymp\Xcore$.

Therefore $|\rho|\le C_M\Xcore^{-M}$ with $C_M$ explicit.  Finally,
$\bigl(\log\mathcal S_\sigma\bigr)'(x)=\lambda+\tfrac\sigma{2x}+\sigma D'(x)
\ge\lambda-\tfrac1{2x}-\tfrac1{4x^{2}}$ by \cref{p4:cor:D-bounds}, which
exceeds $\lambda/2$ once $x\ge2/\lambda$.  \Cref{p0:thm:backward-error}
(residual-to-root conversion) then gives
\[
  |x_M-\mathcal I_\sigma(y)|\le\frac{|\rho|}{\lambda/2}
   \le\frac{2C_M}{\lambda}\,\Xcore^{-M},
\]
which is \cref{p4:eq:inverse-analytic}.  The final sentence follows from
$\lambda\Xcore+\tfrac\sigma2\log\Xcore=L_\sigma(y)$ and
$L_\sigma(y)=\log y+O(1)$.
\end{proof}

\begin{remark}[An exact algebraic certificate]
\label{p4:rem:certificate}
Step (i) of the proof is also the recommended \emph{implementation} check.
Given a candidate truncation $U_{M-1}$, substitute it into
$U-\Psi_\sigma(t,U)$ and expand: the coefficients of
$t^{1},\ldots,t^{M-1}$ must vanish identically in
$\Rat[\lambda,\lambda^{-1},\sigma]/(\sigma^{2}-1)$.  This is an exact
algebraic certificate for every table of $c^{(\sigma)}_n$, and it is how the
tables in \cref{p4:eq:c12}--\cref{p4:eq:c10} were verified here, independently
of the recurrence that produced them.
\end{remark}

\subsection{The centred reversion}

Applying \cref{p0:thm:lambert-centered} to the centred normal form
\cref{p4:eq:centred-expansion} --- that is, with
$a=2\lambda$, $b=\sigma/2$, $C=1/\sqrt{2\pi}$ and $q_j=\sigma\nu_j$
($\nu_j=0$ for odd $j$) --- gives the centred cores
\begin{equation}
  \mathcal U_+(y)=\frac{\Wlam_0\!\bigl(8\pi\lambda y^{2}\bigr)}{4\lambda},
  \qquad
  \mathcal U_-(y)
  =-\frac{\Wlam_{-1}\!\bigl(-2\lambda/(\pi y^{2})\bigr)}{4\lambda},
  \label{p4:eq:centred-cores}
\end{equation}
and the centred inverse expansion
\begin{equation}
  \boxed{\;
  \mathcal I_\sigma(y)\sim 2\,\mathcal U_\sigma(y)-\frac12
   +2\sum_{n\ge2}\frac{\widehat c^{\,(\sigma)}_n}{\mathcal U_\sigma(y)^{n}},
  \qquad \widehat c^{\,(\sigma)}_1=0 .\;}
  \label{p4:eq:centred-inverse}
\end{equation}
The vanishing of $\widehat c^{\,(\sigma)}_1$ is immediate from
$\nu_1=0$.  Recomputed in the present normalization,
\begin{align}
  \widehat c^{\,(\sigma)}_2&=-\frac{\sigma}{128\lambda},
  &\widehat c^{\,(\sigma)}_3&=\frac1{512\lambda^{2}},
  &\widehat c^{\,(\sigma)}_4&=\frac{\sigma(5\lambda^{2}-2)}{4096\lambda^{3}},
  \label{p4:eq:chat234}\\
  \widehat c^{\,(\sigma)}_5&=-\frac{7\lambda^{2}-2}{16384\lambda^{4}},
  &\widehat c^{\,(\sigma)}_6&=-\frac{\sigma\bigl(244\lambda^{4}
     -57\lambda^{2}+12\bigr)}{393216\lambda^{5}},
  &\widehat c^{\,(\sigma)}_7&=\frac{334\lambda^{4}-75\lambda^{2}+12}
     {1572864\lambda^{6}},
  \label{p4:eq:chat567}
\end{align}
and
\begin{equation}
  \widehat c^{\,(\sigma)}_8
  =\frac{\sigma\bigl(4155\lambda^{6}-460\lambda^{4}
    +96\lambda^{2}-12\bigr)}{6291456\lambda^{7}} .
  \label{p4:eq:chat8}
\end{equation}
These agree with \textsf{S1}'s centred table after the substitution
$\alpha=2\lambda$ and the marker translation of \cref{p4:rem:dictionary};
observe that $\widehat c_n$ is $\sigma$-even for odd $n$ and $\sigma$-odd for
even $n$, a reflection symmetry of \cref{p4:eq:centred-exact} that the
uncentred coefficients do not have.

\begin{repairbox}
\textbf{Source claim.}  \textsf{S1} concludes that ``the centred one is often
the preferable computational normal form''.  \textbf{Defect.}  Its own table
already contradicts the unqualified form of this claim: at $m=200$ the centred
core beats the uncentred core by $10^{3}$ ($2.8\times10^{-7}$ against
$4.5\times10^{-4}$), but at $R=4$ the \emph{uncentred} expansion is slightly
better ($7.3\times10^{-16}$ against $1.1\times10^{-15}$).  \textbf{Repair.}
The correct statement, confirmed by the recomputation in
\cref{p4:tab:inverse-errors} and \cref{p4:tab:centred-errors}, is: centring
removes the order-$\Xcore^{-1}$ term entirely, so it improves the \emph{bare
core} by a factor $\asymp\Xcore$ and is the right choice when no corrections
are to be carried; once several corrections are carried, the two
normalizations are comparable, with neither dominating uniformly.
\end{repairbox}

\section{Flattening the core into a polynomial--logarithmic transseries}
\label{p4:sec:flatten}

\subsection{The specialization}

Set
\begin{equation}
  H_\sigma(y):=\log\frac{L_\sigma(y)}{\lambda},
  \label{p4:eq:H-def}
\end{equation}
and recall from \cref{p0:thm:flattening} that the inverse of the model
\cref{p0:def:model} has the unique power--logarithmic form
$x\sim a^{-1}\bigl[L-bH+\sum_{n\ge1}P_n(H)L^{-n}\bigr]$ with $\deg P_n\le n$
and leading term $b^{n+1}H^{n}/n$.  With the dictionary
\cref{p4:eq:dictionary},
\begin{equation}
  \boxed{\;
  \mathcal I_\sigma(y)\sim\frac1\lambda\left[
    L_\sigma-\frac\sigma2H_\sigma
    +\sum_{n\ge1}\frac{P^{(\sigma)}_n(H_\sigma)}{L_\sigma^{\,n}}\right],\;}
  \label{p4:eq:flattened}
\end{equation}
where the $P^{(\sigma)}_n$ are generated by \cref{p0:thm:flattening} from
$b=\sigma/2$, $a=\lambda$, $q_j=\sigma d_j$; the fixed-order meaning of
\cref{p4:eq:flattened} is the analytic clause of \cref{p0:thm:flattening},
namely an error $O\bigl((1+H_\sigma)^{N+1}L_\sigma^{-(N+1)}\bigr)$ after $N$
correction terms, whose hypotheses hold here by
\cref{p4:thm:inverse-analytic}.  The recomputed polynomials are
\begin{align}
  P^{(\sigma)}_1(H)&=\frac{H-\sigma\lambda}{4},
  \label{p4:eq:P1}\\
  P^{(\sigma)}_2(H)&=\frac{\sigma H^{2}-2(\lambda+\sigma)H+2\lambda}{16},
  \label{p4:eq:P2}\\
  P^{(\sigma)}_3(H)&=\frac1{96}\Bigl[2H^{3}-(6\sigma\lambda+9)H^{2}
    +(18\sigma\lambda+6)H
    +4\sigma\lambda^{3}-6\lambda^{2}-6\sigma\lambda\Bigr],
  \label{p4:eq:P3}\\
  P^{(\sigma)}_4(H)&=\frac1{384}\Bigl[3\sigma H^{4}
    -(12\lambda+22\sigma)H^{3}+(66\lambda+36\sigma)H^{2}\nonumber\\
   &\hspace{14mm}{}+(24\lambda^{3}-36\sigma\lambda^{2}-72\lambda-12\sigma)H
    -8\lambda^{3}+30\sigma\lambda^{2}+12\lambda\Bigr],
  \label{p4:eq:P4}\\
  P^{(\sigma)}_5(H)&=\frac1{3840}\Bigl[12H^{5}-(60\sigma\lambda+125)H^{4}
    +(500\sigma\lambda+350)H^{3}\nonumber\\
   &\hspace{14mm}{}
    +(240\sigma\lambda^{3}-360\lambda^{2}-1050\sigma\lambda-300)H^{2}
    \nonumber\\
   &\hspace{14mm}{}
    +(-280\sigma\lambda^{3}+780\lambda^{2}+600\sigma\lambda+60)H
    \nonumber\\
   &\hspace{14mm}{}
    -192\sigma\lambda^{5}+160\lambda^{4}-80\sigma\lambda^{3}
    -270\lambda^{2}-60\sigma\lambda\Bigr].
  \label{p4:eq:P5}
\end{align}
These reproduce \textsf{S2} exactly for $1\le n\le5$; they were recomputed for
$\sigma=\pm1$ separately and recombined in $\mathscr R$, and independently
re-derived through $n=6$.  They also satisfy the cross-normalization identity
$P^{(\sigma)}_n(0)=\lambda^{\,n+1}c^{(\sigma)}_n$ of \cref{p0:thm:flattening}:
for instance $P^{(\sigma)}_1(0)=-\sigma\lambda/4=\lambda^{2}c^{(\sigma)}_1$ and
$P^{(\sigma)}_2(0)=\lambda/8=\lambda^{3}c^{(\sigma)}_2$, which ties
\cref{p4:eq:P1}--\cref{p4:eq:P5} to \cref{p4:eq:c12}--\cref{p4:eq:c10} and
catches any transcription error in either list.

\subsection{The pure Lambert block in closed form}

The pure block --- the case $Q\equiv0$, in which the whole correction series
is suppressed and only the Lambert core survives --- is
\cref{p0:thm:pure-stirling}, proved there in the shortest of the three forms
the sources use.  Specializing its data to the swing branches, $a=\lambda$ and
$b=\sigma/2$, gives
\begin{equation}
  P^{(0)}_n(H)=\Bigl(\frac{\sigma}{2}\Bigr)^{\,n+1}
  \sum_{m=1}^{n}(-1)^{m+1}\,\stirlingfirst{n}{m}\,
  \frac{H^{\,n-m+1}}{(n-m+1)!},
  \label{p4:eq:pure-lambert}
\end{equation}
with $\stirlingfirst{n}{m}$ the unsigned Stirling numbers of the first kind.
Reducing with $\sigma^{2}=1$, the two parities differ only by the overall sign
$\sigma^{\,n+1}$: the even and odd pure blocks coincide for odd $n$ and are
opposite for even $n$.

Writing $b=\sigma/2$, the first few are
\[
  P^{(0)}_1=b^{2}H,\quad
  P^{(0)}_2=\tfrac{b^{3}}2(H^{2}-2H),\quad
  P^{(0)}_3=\tfrac{b^{4}}6(2H^{3}-9H^{2}+6H),\quad
  P^{(0)}_4=\tfrac{b^{5}}{12}(3H^{4}-22H^{3}+36H^{2}-12H),
\]
and \cref{p4:eq:pure-lambert} was recomputed here against the recurrence of
\cref{p0:thm:flattening} for $1\le n\le8$, in exact arithmetic with $b$
symbolic.  In particular the leading
coefficient is $b^{n+1}/n$, which is \cref{p0:thm:flattening}'s degree
statement.  \Cref{p0:thm:pure-stirling} is the precise bridge between Lambert
asymptotics and the first-kind Stirling calculus, and it is stated there
rather than here because all four chapters use it.

\subsection{The core series and its convergence}

The flattening can also be organized as ``expand the core, then compose'',
which is \textsf{S3}'s route and which is constructive at every order.

\begin{proposition}[All coefficients of the scaled core]
\label{p4:prop:core-series}
Let $w$ solve $w+\sigma\log w=\Lambda$ as in \cref{p4:eq:scaled-core}, and
put $\ell=\log\Lambda$.  Define
\begin{equation}
  A_k(\ell):=\frac1{k+1}\Coef{z^{k}}\bigl(\ell+\log(1+z)\bigr)^{k+1},
  \qquad k\ge0 .
  \label{p4:eq:Ak-def}
\end{equation}
Then
\begin{equation}
  w\;=\;\Lambda+\sum_{k\ge0}\frac{(-\sigma)^{k+1}A_k(\ell)}{\Lambda^{k}},
  \label{p4:eq:core-series}
\end{equation}
and, with $s(k,j)$ the signed Stirling numbers of the first kind,
\begin{equation}
  A_0(\ell)=\ell,
  \qquad
  A_k(\ell)=\frac1{(k+1)\,k!}\sum_{j=1}^{k}\binom{k+1}{j}j!\,s(k,j)\,
  \ell^{\,k+1-j}\quad(k\ge1).
  \label{p4:eq:Ak-stirling}
\end{equation}
The first polynomials are
\[
  A_0=A_1=\ell,\quad
  A_2=\tfrac{\ell(2-\ell)}2,\quad
  A_3=\tfrac{\ell(2\ell^{2}-9\ell+6)}6,\quad
  A_4=-\tfrac{\ell(3\ell^{3}-22\ell^{2}+36\ell-12)}{12},
\]
\[
  A_5=\tfrac{\ell(12\ell^{4}-125\ell^{3}+350\ell^{2}-300\ell+60)}{60},
  \quad
  A_6=-\tfrac{\ell(20\ell^{5}-274\ell^{4}+1125\ell^{3}-1700\ell^{2}
    +900\ell-120)}{120} .
\]
\end{proposition}

\begin{proof}
Write $w=\Lambda(1+\upsilon)$ and $z=\Lambda^{-1}$.  The equation becomes
$\Lambda\upsilon=-\sigma\bigl(\ell+\log(1+\upsilon)\bigr)$, i.e.\
$\upsilon=-\sigma z\bigl(\ell+\log(1+\upsilon)\bigr)$.  Lagrange inversion
gives
$\upsilon=\sum_{k\ge0}\frac{(-\sigma z)^{k+1}}{k+1}
\Coef{\upsilon^{k}}(\ell+\log(1+\upsilon))^{k+1}$; multiplying by
$\Lambda=z^{-1}$ yields \cref{p4:eq:core-series} with $A_k$ as in
\cref{p4:eq:Ak-def}.  Expanding the $(k+1)$-st power and using
$\bigl(\log(1+z)\bigr)^{j}=j!\sum_{k\ge j}s(k,j)z^{k}/k!$ gives
\cref{p4:eq:Ak-stirling}.
\end{proof}

The sign in \cref{p4:eq:core-series} matches \textsf{S3} after
$\varepsilon=-\sigma$; explicitly $w_-=\Lambda_-+\ell+\ell/\Lambda_-+\cdots$
on the even sheet ($\Wlam_{-1}$) and
$w_+=\Lambda_+-\ell+\ell/\Lambda_+-\cdots$ on the odd sheet ($\Wlam_0$).

\begin{proposition}[An explicit convergence criterion]
\label{p4:prop:core-convergence}
Fix $\rho\in(0,1)$ and put $M_\rho(\ell)=|\ell|+\log\frac1{1-\rho}$.  Then
\begin{equation}
  |A_k(\ell)|\le\frac{M_\rho(\ell)^{\,k+1}}{(k+1)\rho^{k}},
  \label{p4:eq:Ak-bound}
\end{equation}
so the series in \cref{p4:eq:core-series} converges absolutely whenever
$M_\rho(\ell)<\rho|\Lambda|$ for some $\rho\in(0,1)$, and its sum is then the
solution branch $w\sim\Lambda$.  In particular, for real $\Lambda\ge10$ the
choice $\rho=\tfrac12$ already gives convergence.
\end{proposition}

\begin{proof}
On $|z|=\rho$ one has $|\log(1+z)|\le\sum_{k\ge1}\rho^{k}/k
=\log\frac1{1-\rho}$, so $|\ell+\log(1+z)|\le M_\rho(\ell)$ there; Cauchy's
estimate applied to \cref{p4:eq:Ak-def} gives \cref{p4:eq:Ak-bound}.  Summing
the geometric-type majorant gives absolute convergence for
$M_\rho(\ell)<\rho|\Lambda|$.  The sum solves
$w+\sigma\log w=\Lambda$ because Lagrange's theorem identifies the series with
the unique analytic solution $\upsilon(z)$ of
$\upsilon=-\sigma z(\ell+\log(1+\upsilon))$ vanishing at $z=0$; the branch
with $w\sim\Lambda$ is the one selected.  For $\Lambda=10$,
$\ell=\log10=2.3026$ and $M_{1/2}=2.3026+0.6931=2.9957<5$.
\end{proof}

\begin{repairbox}
\textbf{Source claim.}  \textsf{S3} asserts that ``the pure Lambert
logarithmic expansion has a nonzero convergence domain for sufficiently large
argument, as in the classical analysis of Lambert $W$'', citing Corless et
al.\ but giving no criterion.  \textbf{Defect.}  The cited classical result
concerns the doubly-indexed series in $(\log z)^{-1}$ and $\log\log z$; the
single series \cref{p4:eq:core-series}, whose coefficients are polynomials in
the slowly varying $\ell$, needs its own statement, since convergence depends
on the ratio $\ell/\Lambda$ and fails for $\ell$ comparable to $\Lambda$.
\textbf{Repair.}  \Cref{p4:prop:core-convergence} supplies an explicit,
proved criterion, which also delimits the regime where
\cref{p4:eq:core-series} may be summed rather than truncated.
\end{repairbox}

\subsection{Composition, and the equivalence of the two normalizations}

\begin{theorem}[Constructive equivalence of the two inverse normalizations]
\label{p4:thm:equivalence}
With $v_j=(-\sigma)^{j}A_{j-1}(\ell)$ and
\begin{equation}
  R^{(m)}_r(\ell)=\Coef{z^{r}}\Bigl(1+\sum_{j\ge1}v_jz^{j}\Bigr)^{-m}
  =\sum_{k=0}^{r}(-1)^{k}\binom{m+k-1}{k}\OB_{r,k}(v),
  \qquad R^{(m)}_0=1,
  \label{p4:eq:R-composition}
\end{equation}
the Lambert-normalized expansion \cref{p4:eq:inverse-analytic} and the
flattened expansion \cref{p4:eq:flattened} are term-by-term equivalent: in the
scaled variables of \cref{p4:eq:scaled-core},
\begin{equation}
  \mathcal I_\sigma(y)\sim\frac{\Lambda_\sigma}{2\lambda}
   -\frac{\sigma\,\ell_\sigma}{2\lambda}
   +\sum_{N\ge1}\frac{\Pi^{(\sigma)}_N(\ell_\sigma)}{\Lambda_\sigma^{\,N}},
  \qquad
  \Pi^{(\sigma)}_N(\ell)=\frac{(-\sigma)^{N+1}A_N(\ell)}{2\lambda}
   +\sum_{m=1}^{N}c^{(\sigma)}_m(2\lambda)^{m}R^{(m)}_{N-m}(\ell),
  \label{p4:eq:Pi-composition}
\end{equation}
with $\deg\Pi^{(\sigma)}_N=N$ and leading coefficient
$(-1)^{N-1}(-\sigma)^{N+1}/(2\lambda N)$.
\end{theorem}

\begin{proof}
Insert $\Xcore_\sigma=w_\sigma/(2\lambda)$ into
\cref{p4:eq:inverse-analytic}, write $w=\Lambda(1+V)$ with
$V=\sum_jv_jz^{j}$ from \cref{p4:eq:core-series}, and collect the coefficient
of $\Lambda^{-N}$: the core contributes the first term of
\cref{p4:eq:Pi-composition} and the $m$-th correction contributes
$c^{(\sigma)}_m(2\lambda)^{m}\Coef{z^{N-m}}(1+V)^{-m}$.  The Bell form of
$R^{(m)}_r$ is \cref{p0:lem:power-log}.  Degrees: $A_N$ has degree $N$ with
leading coefficient $(-1)^{N-1}/N$ by \cref{p4:eq:Ak-stirling}, while every
correction with $m\ge1$ contributes degree at most $N-m$.
\end{proof}

\begin{repairbox}
\textbf{Source claim.}  \textsf{S1} asserts the equivalence of its two inverse
normalizations in prose: ``expanding the Lambert carrier in the scale
$X^{-1}\Real[\ell]$ and then substituting that expansion \ldots\ reproduces
the recurrence coefficient by coefficient.''  \textbf{Defect.}  No proof and
no formula; the assertion is exactly the sort of ``successive matching''
that the coefficient calculus is supposed to eliminate.  \textbf{Repair.}
\Cref{p4:thm:equivalence}, which follows \textsf{S3}, makes the equivalence a
finite explicit composition with a proof, and identifies the degrees and
leading coefficients.
\end{repairbox}

\begin{remark}[Three flattening normalizations, and which one the volume uses]
\label{p4:rem:three-flattenings}
The three sources flatten in three different variables: \textsf{S2} in
$L_\sigma=\log(y/C_\sigma)$ with $H_\sigma=\log(L_\sigma/\lambda)$;
\textsf{S3} in $\Lambda_\sigma=\log(\pi y^{2})+\sigma\log(4\lambda)$ with
$\ell_\sigma=\log\Lambda_\sigma$; \textsf{S1} in
$X_\epsilon=\log(y\alpha^{\beta_\epsilon}/K_\epsilon)$ in the half index.
Since $\Lambda_\sigma=2L_\sigma+\sigma(\log(4\lambda)-\lambda)\ne2L_\sigma$,
these are genuinely different variables, the three coefficient lists are not
comparable termwise, and none is in error --- all were recomputed and are
correct in their own normalization.  The volume adopts \textsf{S2}'s variables
in \cref{p4:eq:flattened} (they are those of \cref{p0:thm:flattening}) and
records \textsf{S3}'s in \cref{p4:eq:Pi-composition}, where the composition
proof is cleanest.
\end{remark}

\begin{corollary}[Leading elementary forms]
\label{p4:cor:leading-inverse}
As $y\to\infty$,
\begin{equation}
  \mathcal I_\sigma(y)=\frac1\lambda\Bigl(L_\sigma(y)
   -\frac\sigma2\log\frac{L_\sigma(y)}{\lambda}
   +O\Bigl(\frac{\log L_\sigma(y)}{L_\sigma(y)}\Bigr)\Bigr).
  \label{p4:eq:leading-inverse}
\end{equation}
The opposite signs of the $\tfrac12\log L$ term on the two sheets are the
inverse-side signature of the forward powers $n^{-1/2}$ and $n^{+1/2}$.
\end{corollary}

\section{Asymptotic meaning, optimal truncation, and Stokes structure}
\label{p4:sec:stokes}

\subsection{What each expansion means at fixed order}

For each fixed $M$ the forward statement is \cref{p4:eq:forward} with the
effective constant of \cref{p4:prop:forward-effective}, and the inverse
statement is \cref{p4:eq:inverse-analytic}.  Both are Poincar\'e statements
about divergent series: the coefficients $d_j$ grow factorially
(\cref{p4:prop:leastterm}), hence so do the $A_n$ and, through the reversion,
the $c^{(\sigma)}_n$.  The divergence is not a defect but a datum: the Borel
transform and the exact sum of the logarithmic series are known in closed form
(\cref{p4:thm:laplace}, \cref{p4:prop:borel-lattice}).  By contrast the pure
core series \cref{p4:eq:core-series} \emph{converges} in the regime delimited
by \cref{p4:prop:core-convergence}; the full inverse is asymptotic only
because of the gamma-ratio correction.

\subsection{Optimal truncation and its image in the inverse variable}

\begin{proposition}[Least-term truncation, forward and inverse]
\label{p4:prop:optimal}
\Cref{p0:thm:optimal-truncation} applied to \cref{p4:eq:D-first} with the
singulant modulus $\pi$ of \cref{p4:prop:borel-lattice} gives the least-term
index $r\approx\pi x/2$ and the beyond-all-orders floor
\begin{equation}
  \bigl|R_{N}(x)\bigr|_{N\approx\pi x/2}
  =\frac{2\sqrt2}{\pi\sqrt x}\,\e^{-\pi x}\bigl(1+O(x^{-1})\bigr),
  \label{p4:eq:floor-forward}
\end{equation}
which by \cref{p4:thm:envelope} is an equality up to the stated factor and not
merely an upper bound.  Transporting it through the inverse by
\cref{p0:thm:remainder-transport} and
$\bigl(\log\mathcal S_\sigma\bigr)'=\lambda+O(x^{-1})$ gives the inverse floor
\begin{equation}
  O\!\left(\Xcore_\sigma^{-1/2}\e^{-\pi \Xcore_\sigma}\right)
  =O\!\left(y^{-\pi/\lambda}\,
    (\log y)^{-1/2+\sigma\pi/(2\lambda)}\right),
  \qquad
  \frac\pi\lambda=4.5323601418\ldots
  \label{p4:eq:floor-inverse}
\end{equation}
\end{proposition}

\begin{proof}
\Cref{p4:eq:floor-forward} is \cref{p4:prop:leastterm}.  For the image, use the
core equation in the form $\e^{\lambda \Xcore}=(y/C_\sigma)\Xcore^{-\sigma/2}$,
so that
$\e^{-\pi \Xcore}=(y/C_\sigma)^{-\pi/\lambda}\Xcore^{\sigma\pi/(2\lambda)}$;
multiply by $\Xcore^{-1/2}$ and use $\Xcore\asymp\log y$.
\end{proof}

\begin{remark}
Exponential smallness in the index becomes \emph{algebraic} smallness in the
level: an $\e^{-\pi x}$ forward sector maps to a $y^{-\pi/\lambda}$ inverse
sector.  This is why hyperasymptotic corrections to the inverse live on a
strictly finer grid than the algebraic powers of $1/\log y$ produced by
\cref{p4:eq:flattened} --- they are invisible to that scale at every order.
\end{remark}

\subsection{Lateral sums and the aggregate Stokes multipliers}

For a direction $\theta$ let
$\mathcal S_{\theta^{\pm}}D(x)
 =\int_0^{\e^{\ii(\theta\pm0)}\infty}\e^{-xt}\widehat D(t)\,\dd t$
whenever the rotated Laplace integral converges, and let
$\operatorname{Disc}_\theta$ denote the lower-minus-upper lateral difference.

\begin{proposition}[Bridge equations]
\label{p4:prop:stokes}
With the poles and residues of \cref{p4:eq:borel-poles},
\begin{align}
  \operatorname{Disc}_{\pi/2}D(x)
   &=2\sum_{k\ge0}\frac{\e^{-\ii\pi(2k+1)x}}{2k+1}
   =2\operatorname{artanh}\bigl(\e^{-\ii\pi x}\bigr)
   \qquad(\Im x<0),
  \label{p4:eq:upper-stokes}\\
  \operatorname{Disc}_{-\pi/2}D(x)
   &=-2\sum_{k\ge0}\frac{\e^{\ii\pi(2k+1)x}}{2k+1}
   =-2\operatorname{artanh}\bigl(\e^{\ii\pi x}\bigr)
   \qquad(\Im x>0),
  \label{p4:eq:lower-stokes}
\end{align}
the closed forms giving the analytic continuations beyond those half-planes.
Consequently
$\operatorname{Disc}_\theta\log\mathcal S_\sigma
 =\sigma\operatorname{Disc}_\theta D$, and the multiplicative Stokes factor of
the forward branch across the ray $\theta=\pm\pi/2$ is
$\exp\bigl(\sigma\operatorname{Disc}_{\pm\pi/2}D(x)\bigr)$, that is
\begin{equation}
  \mathfrak M^{(\sigma)}_{+}(x)
   =\left(\frac{1+\e^{-\ii\pi x}}{1-\e^{-\ii\pi x}}\right)^{\!\sigma},
  \qquad
  \mathfrak M^{(\sigma)}_{-}(x)
   =\left(\frac{1-\e^{\ii\pi x}}{1+\e^{\ii\pi x}}\right)^{\!\sigma}.
  \label{p4:eq:stokes-multipliers}
\end{equation}
On the positive real axis no lateral choice arises, because $\arg t=0$ is a
regular direction.
\end{proposition}

\begin{proof}
Rotating the Laplace contour past the ray through $t_k=\ii\pi(2k+1)$ picks up
$2\pi\ii$ times the residue of $\e^{-xt}\widehat D(t)$ there, namely
$2\pi\ii\,\e^{-t_kx}/t_k=2\e^{-\ii\pi(2k+1)x}/(2k+1)$; summing over $k$ and
using $\operatorname{artanh}z=\sum_{k\ge0}z^{2k+1}/(2k+1)$ gives
\cref{p4:eq:upper-stokes}.  The series converges for
$|\e^{-\ii\pi x}|=\e^{\pi\Im x}<1$, that is $\Im x<0$.  The lower ray is the
conjugate computation.  The last two displays follow from
$\log\mathcal S_\sigma=\log C_\sigma+\lambda x+\tfrac\sigma2\log x+\sigma D$,
in which only $D$ is discontinuous, together with
$2\operatorname{artanh}z=\log\frac{1+z}{1-z}$.
\end{proof}

\begin{remark}[Convention dependence]
\label{p4:rem:stokes-convention}
\textsf{S2} defines the discontinuity as lower minus upper lateral sum;
\textsf{S3} instead adopts the residue convention in which crossing a singular
ray contributes $+2\pi\ii$ times the enclosed residues, and works with
$Q=-D$.  The two sets of formulae differ by an overall sign, hence by a
reciprocal in the multipliers, and both are internally consistent.  A Stokes
generator is only defined relative to a stated convention, so the volume
states the convention explicitly in \cref{p4:prop:stokes} and does not declare
either source's signs canonical.
\end{remark}

\subsection{Transport of an exponential sector through inversion}

\begin{proposition}[Sector transport]
\label{p4:prop:sector-transport}
Let $F$ be one lateral branch of $\log\mathcal S_\sigma$ and let the other be
$F+\Delta$.  As a formal expansion in the sector $\Delta$,
\begin{equation}
  x_{\mathrm{new}}=x+\sum_{k\ge1}\frac{(-1)^{k}}{k!}
   \left(\frac1{F'(x)}\frac{\dd}{\dd x}\right)^{\!k-1}
   \left(\frac{\Delta(x)^{k}}{F'(x)}\right),
  \label{p4:eq:sector-transport}
\end{equation}
with $\Delta=\sigma\operatorname{Disc}_\theta D$ and
$F'(x)=\lambda+\tfrac{\sigma}{2x}+\sigma D'(x)$.  Quantitatively, and with no
formal series at all: if $|\Delta|\le\delta$ on the interval joining $x$ and
$x_{\mathrm{new}}$ and $F'\ge\mu>0$ there, then
\begin{equation}
  |x_{\mathrm{new}}-x|\le\frac{\delta}{\mu},
  \qquad
  x_{\mathrm{new}}-x=-\frac{\Delta(x)}{F'(x)}+O(\delta^{2}),
  \label{p4:eq:sector-transport-quant}
\end{equation}
the implied constant depending only on $\sup|(F^{-1})''|$ on that interval.
\end{proposition}

\begin{proof}
\Cref{p4:eq:sector-transport} is \cref{p4:lem:phase-lb} with $F_0=F$,
$g=\Delta$ and $\eta=1$, read as an identity of formal series in the sector
marker.  The first inequality in \cref{p4:eq:sector-transport-quant} is
\cref{p0:thm:backward-error}; the second is Taylor's theorem for $F^{-1}$ at
the point $F(x)$.
\end{proof}

\begin{repairbox}
\textbf{Source claim.}  \textsf{S1} states its transport rule as a
\emph{theorem} (``Transseries-sector transport through inversion'') whose
hypothesis is that ``$E$ is smaller than the algebraic scale retained in
$\Phi_{\mathrm{alg}}$'', whose conclusion is prefaced by the word
``formally'', and whose displayed first terms carry an unjustified error
symbol $O(E^{3},E^{2}E',E(E')^{2})$.  \textbf{Defect.}  The hypothesis is
never used; the conclusion is a formal-series identity, not an analytic one;
and the error term is asserted rather than proved.  A statement that holds
only formally must not stand in a \texttt{theorem} environment.
\textbf{Repair.}  \Cref{p4:prop:sector-transport} separates the two claims:
the all-orders identity is formal (and is exactly \cref{p4:lem:phase-lb}),
while the quantitative first-order statement is proved by the mean value
theorem under stated hypotheses.  \textsf{S2} and \textsf{S3} already state
their versions correctly, as formal expansions.
\end{repairbox}

\begin{repairbox}
\textbf{Source claim.}  \textsf{S1}'s section on complex inverse sheets
asserts that for every Lambert branch with $|w_{\epsilon,k}|\to\infty$ in a
suitable sector, ``the same coefficients $d^{(\epsilon)}_r$ give
$M_{\epsilon,k}(y)\sim w_{\epsilon,k}(y)+\sum_r d^{(\epsilon)}_r
w_{\epsilon,k}(y)^{-r}$''.  \textbf{Defect.}  For $k\ne0,-1$ nothing has been
said, let alone proved, about $M_{\epsilon,k}$ being an inverse of anything:
the existence of an analytic branch of $\mathcal S_\sigma^{-1}$ along the
corresponding path is a hypothesis, not a consequence, and it fails when the
path meets a pole of the gamma quotient or a Stokes direction.
\textbf{Repair.}  The statement is demoted to
\cref{p4:rem:complex-branches}, with the two missing hypotheses displayed.
\textsf{S2} states the same thing correctly, as a \emph{formal} complex branch
whose analytic validity is a separate requirement.
\end{repairbox}

\begin{remark}[Complex branches, formally]
\label{p4:rem:complex-branches}
For $k\in\Int$ let $\Xcore_{\sigma,k}$ be as in \cref{p4:eq:core-general}
after fixing branches of the powers and logarithms in $L_\sigma$.  Suppose
that along a path with $|\Xcore_{\sigma,k}(y)|\to\infty$ \emph{(i)} the
forward expansion \cref{p4:eq:forward} is valid, which requires the path to
stay in a sector free of the poles of $\Gamma(x/2+(3-\sigma)/4)^{-2}$ and away
from the Stokes directions of \cref{p4:prop:stokes}, and \emph{(ii)} an
analytic branch $\mathcal I_{\sigma,k}$ of the inverse exists along it.  Then
that branch satisfies
$\mathcal I_{\sigma,k}(y)\sim \Xcore_{\sigma,k}(y)
 +\sum_{n\ge1}c^{(\sigma)}_n\Xcore_{\sigma,k}(y)^{-n}$
with the \emph{same} coefficients, by the proof of
\cref{p4:thm:inverse-analytic}.  The coefficient sequence does not depend on
$k$: all branch dependence sits in $\Wlam_k$ and in the chosen logarithm.
Without hypotheses (i) and (ii) the display is a formal statement only.
\end{remark}

\section{Discrete inversion and index recovery}
\label{p4:sec:discrete}

\subsection{The three inverse objects, parity by parity}

\Cref{p0:def:three-inverses} attaches three objects to a discrete sequence:
the functional inverse of a chosen interpolation, the integer staircase, and
the round-to-nearest generalized inverse.

\begin{warning}[None of the three exists for the unsplit sequence]
\label{p4:rem:no-global}
By \cref{p4:prop:ratios} the set $\{n:\operatorname{sw}(n)\le y\}$ is not an
initial segment of $\Nat_0$ --- for instance
$\operatorname{sw}(5)=30>20=\operatorname{sw}(6)$ --- so the staircase of
$\operatorname{sw}$ is not monotone and does not invert it; the
round-to-nearest inverse is ambiguous whenever two indices of opposite parity
are nearly equidistant in value; and by \cref{p4:prop:oscillatory} no
eventually monotone interpolation exists at all.  All three objects exist, are
unique, and obey \cref{p0:thm:staircase} \emph{on each parity class
separately}.
\end{warning}

\begin{proposition}[Exact branchwise staircases]
\label{p4:prop:staircase}
For $y\ge1$ put
\begin{equation}
  N^{\le}_-(y)=\max\{n\ \text{even}:\operatorname{sw}(n)\le y\},
  \qquad
  N^{\le}_+(y)=\max\{n\ \text{odd}:\operatorname{sw}(n)\le y\},
  \label{p4:eq:Nle-def}
\end{equation}
both sets being non-empty because $\operatorname{sw}(0)=\operatorname{sw}(1)=1$.
Then
\begin{equation}
  \boxed{\;
  N^{\le}_-(y)=2\left\lfloor\frac{\mathcal I_-(y)}{2}\right\rfloor,
  \qquad
  N^{\le}_+(y)=2\left\lfloor\frac{\mathcal I_+(y)-1}{2}\right\rfloor+1,\;}
  \label{p4:eq:staircase}
\end{equation}
and the least index of each parity that reaches the level $y$ is
\begin{equation}
  N^{\ge}_-(y)=2\left\lceil\frac{\mathcal I_-(y)}{2}\right\rceil,
  \qquad
  N^{\ge}_+(y)=2\left\lceil\frac{\mathcal I_+(y)-1}{2}\right\rceil+1 .
  \label{p4:eq:staircase-up}
\end{equation}
These identities are \emph{exact}, not asymptotic.  The quantity
\cref{p4:eq:staircase-up} is the staircase $N_\ast$ of
\cref{p0:def:three-inverses}, and the displayed formula for it is
\cref{p0:thm:staircase}'s exact-rounding identity
$N_\ast=\lceil\nu\rceil$, applied to the bisection
$(\operatorname{sw}(2m+\epsilon))_{m\ge0}$ with its admissible interpolation
$m\mapsto\mathcal S_\sigma(2m+\epsilon)$.
\end{proposition}

\begin{proof}
By \cref{p4:eq:branch-interpolates} and \cref{p4:prop:monotone},
$\mathcal S_-$ is a strictly increasing bijection agreeing with
$\operatorname{sw}$ at the even integers, so for even $n$ we have
$\operatorname{sw}(n)\le y\iff\mathcal S_-(n)\le y\iff n\le\mathcal I_-(y)$.
The largest such even $n$ is $2\lfloor\mathcal I_-(y)/2\rfloor$ and the least
even $n\ge\mathcal I_-(y)$ is $2\lceil\mathcal I_-(y)/2\rceil$.  The odd case
is identical after the shift $n\mapsto n-1$.  The boundary case in which
$\mathcal I_\sigma(y)$ is an integer of the relevant parity is covered,
because then $y$ is exactly a sequence value and floor and ceiling return it.
\end{proof}

The round-to-nearest inverses on the two parity classes are
\begin{equation}
  n^{\mathrm{near}}_-(y)
   =2\left\lfloor\frac{\mathcal I_-(y)}2+\frac12\right\rfloor,
  \qquad
  n^{\mathrm{near}}_+(y)
   =2\left\lfloor\frac{\mathcal I_+(y)-1}2+\frac12\right\rfloor+1 .
  \label{p4:eq:nearest}
\end{equation}

\begin{remark}[The rounding must be certified, not trusted]
\label{p4:rem:certify-rounding}
\Cref{p4:eq:staircase}--\cref{p4:eq:nearest} are exact statements about the
real numbers $\mathcal I_\sigma(y)$, whereas a truncated transseries supplies
only an approximation to them.  When $y$ lies exponentially close to a
sequence value the floor can flip.  The remedy is cheap: use the transseries
to produce a candidate integer, then decide the inequality
$\operatorname{sw}(n)\le y$ exactly, either in integer arithmetic or by
comparing $\log\mathcal S_\sigma(n)$ with $\log y$ at the candidate and its
same-parity neighbour, working through $\log\Gamma$ to stay in range.
\Cref{p0:thm:staircase} is the general separation condition and
\cref{p0:thm:backward-error} converts the logarithmic residual into a
certified bracket for $\mathcal I_\sigma(y)$.
\end{remark}

\subsection{Which parity wins, and by how much}

\begin{theorem}[The sheet gap]
\label{p4:thm:sheet-gap}
Put $Y=\log(y\sqrt\pi)$ and $h=\log(Y/\lambda)$, so that
$L_\sigma=Y+\sigma\lambda/2$.  Then, as $y\to\infty$,
\begin{equation}
  \boxed{\;
  \mathcal I_-(y)-\mathcal I_+(y)
   =\frac h\lambda-1+\frac1{2Y}+O\!\left(\frac{h^{2}}{Y^{2}}\right)
   =\log_2\!\left(\frac{\log(y\sqrt\pi)}{\log2}\right)-1+o(1).\;}
  \label{p4:eq:sheet-gap}
\end{equation}
In particular the gap tends to $+\infty$: every sufficiently large level is
reached first on the \emph{odd} sheet, about $\log_2\log y$ indices before the
even sheet catches up.
\end{theorem}

\begin{proof}
Subtract the two instances of \cref{p4:eq:flattened} truncated after $n=1$.
The leading terms contribute $(L_--L_+)/\lambda=-1$.  Next,
$H_\sigma=\log\bigl((Y+\sigma\lambda/2)/\lambda\bigr)
 =h+\sigma\lambda/(2Y)+O(Y^{-2})$, and the two contributions
$-(\sigma/2\lambda)H_\sigma$ combine to
$(H_-+H_+)/(2\lambda)=h/\lambda+O(Y^{-2})$, the $\sigma$-odd parts
cancelling.  Finally, with $P^{(\sigma)}_1=(H-\sigma\lambda)/4$ from
\cref{p4:eq:P1},
\[
  \frac1\lambda\left[\frac{P^{(-1)}_1(H_-)}{L_-}
   -\frac{P^{(+1)}_1(H_+)}{L_+}\right]
  =\frac1{4\lambda}\left[\frac{H_-+\lambda}{L_-}-\frac{H_+-\lambda}{L_+}
   \right]
  =\frac1{4\lambda}\cdot\frac{2\lambda}{Y}+O\!\left(\frac{h}{Y^{2}}\right)
  =\frac1{2Y}+O\!\left(\frac{h}{Y^{2}}\right).
\]
Collecting the three displays gives \cref{p4:eq:sheet-gap}; the second form
follows from $h/\lambda=\log_2(Y/\lambda)$ and $1/(2Y)\to0$.  The error term
$O(h^{2}/Y^{2})$ absorbs the $n=2$ terms of \cref{p4:eq:flattened}, whose
polynomials have degree $2$ in $H_\sigma$.
\end{proof}

\begin{remark}[Numerical confirmation, and agreement with \textsf{S1}]
\label{p4:rem:gap-numeric}
Recomputed with eight Lambert-normalized corrections on each sheet: for
$\log y=10,30,100,400$ the gap $\mathcal I_-(y)-\mathcal I_+(y)$ equals
$2.97806$, $4.47863$, $6.18565$, $8.17591$, against the prediction
$h/\lambda-1+1/(2Y)$ equal to $2.97829$, $4.47928$, $6.18583$, $8.17593$.
\textsf{S1} states the same result in its half index as
$N_0-N_1=\tfrac2\alpha\log\log y+\tfrac2\alpha\log\tfrac2\alpha-1+o(1)$ with
$\alpha=2\lambda$; since $\tfrac2\alpha=\tfrac1\lambda$ and
$\log\tfrac2\alpha=-\log\lambda$, this reads
$\tfrac1\lambda(\log\log y-\log\lambda)-1+o(1)$, which is
\cref{p4:eq:sheet-gap}.  Directly from the sequence, the least index reaching
$y=10,10^{3},10^{6},10^{12}$ is $5,11,21,39$ on the odd sheet against
$6,14,24,44$ on the even sheet: gaps $1,3,3,5$, consistent with the very slow
$\log_2\log y$ growth.
\end{remark}

\begin{corollary}[Counting function]
\label{p4:cor:counting}
Let $N(y):=\#\{n\in\Nat_0:\operatorname{sw}(n)\le y\}$.  Then
\begin{equation}
  N(y)=\frac{\mathcal I_-(y)+\mathcal I_+(y)}{2}+O(1)
   =\log_2\bigl(y\sqrt\pi\bigr)+O(1).
  \label{p4:eq:counting}
\end{equation}
\end{corollary}

\begin{proof}
By \cref{p4:prop:staircase} the even indices counted are
$0,2,\ldots,2\lfloor\mathcal I_-(y)/2\rfloor$, that is
$\mathcal I_-(y)/2+O(1)$ of them, and similarly there are
$\mathcal I_+(y)/2+O(1)$ odd ones; this is the first equality.  For the
second, average the two instances of \cref{p4:eq:flattened}: as in the proof
of \cref{p4:thm:sheet-gap} the $\sigma$-odd term $-(\sigma/2\lambda)H_\sigma$
cancels in the average, leaving
$\tfrac1{2\lambda}(L_-+L_+)+O(\log\log y/\log y)=Y/\lambda+O(1)$, and
$Y/\lambda=\log_2(y\sqrt\pi)$.
\end{proof}

\section{Algorithms, certification, and validation}
\label{p4:sec:algorithms}

\subsection{Generating the coefficients}

\begin{algorithm}[Exact coefficient generation]
\label{p4:alg:coefficients}
Every coefficient family of this chapter is produced by a finite exact
recurrence over $\Rat[\lambda,\lambda^{-1},\sigma]/(\sigma^{2}-1)$:
\begin{enumerate}
\item $d_{2r+1}$ from \cref{p4:eq:d-bernoulli} (rational Bernoulli input;
      $d_{2r}=0$);
\item $A_n$ from the triangular recurrence \cref{p4:eq:A-recurrence};
\item $c^{(\sigma)}_n$ from \cref{p4:eq:c-bell}, or by solving
      $U=\Psi_\sigma(t,U)$ coefficientwise, or from the differential formula
      \cref{p4:eq:c-differential};
\item $P^{(\sigma)}_n(H)$ from \cref{p0:thm:flattening} with the dictionary
      \cref{p4:eq:dictionary}, or from \cref{p4:eq:Pi-composition} after
      \cref{p4:prop:core-series};
\item the exact certificate of \cref{p4:rem:certificate} after every step.
\end{enumerate}
Truncating every intermediate series at order $n$ is exact for the $n$-th
coefficient, because all three recurrences are triangular; at order $N$ the
cost is $\widetilde{O}(N^{2})$ coefficient operations.  No numerical
fitting and no order-specific ansatz occurs anywhere.
\end{algorithm}

\begin{remark}[How this chapter's tables were produced]
\label{p4:rem:how-computed}
The three sources supply three implementations of the above.  For this chapter
every coefficient was regenerated by a fourth, independent one: truncated
power-series arithmetic over exact rationals with $\lambda$ symbolic, run
separately at $\sigma=\pm1$ and recombined in $\mathscr R$.  Every displayed
coefficient agreed with the source that printed it; the corrections recorded
in this chapter concern statements, never numbers.
\end{remark}

\subsection{Certified numerical inversion}

\begin{proposition}[Derivative bounds for certification]
\label{p4:prop:derivative-bounds}
For $x\ge1$,
\begin{equation}
  \bigl(\log\mathcal S_+\bigr)'(x)>\lambda+\frac1{2x}-\frac1{4x^{2}}>\lambda,
  \qquad
  \bigl(\log\mathcal S_-\bigr)'(x)>\lambda-\frac1{2x},
  \label{p4:eq:derivative-bounds}
\end{equation}
the exact value being \cref{p4:eq:logder}.  Hence if
$r_*=\log\mathcal S_\sigma(x_*)-\log y$ and
$\bigl(\log\mathcal S_\sigma\bigr)'\ge\mu>0$ on the interval joining $x_*$ and
$\mathcal I_\sigma(y)$, then
$|x_*-\mathcal I_\sigma(y)|\le|r_*|/\mu$ by \cref{p0:thm:backward-error}.
\end{proposition}

\begin{proof}
$\bigl(\log\mathcal S_\sigma\bigr)'(x)=\lambda+\tfrac\sigma{2x}+\sigma D'(x)$
by \cref{p4:eq:normal-form}; apply $0<-D'(x)<1/(4x^{2})$ from
\cref{p4:cor:D-bounds}, and $1/(2x)-1/(4x^{2})>0$ for $x>1/2$.
\end{proof}

\begin{algorithm}[Branchwise inversion of a level]
\label{p4:alg:invert}
Given a large $y$ and a parity $\sigma$:
\begin{enumerate}
\item Evaluate the Lambert core:
      $\Xcore_+=\Wlam_0(4\pi\lambda y^{2})/(2\lambda)$ or
      $\Xcore_-=-\Wlam_{-1}(-4\lambda/(\pi y^{2}))/(2\lambda)$.
      \emph{The branch choice is not optional} (\cref{p4:prop:two-roots}).
\item Add corrections $\sum_n c^{(\sigma)}_n\Xcore^{-n}$ by Horner nesting in
      $1/\Xcore$, stopping at or before the least term rather than summing
      indefinitely; alternatively use the centred form
      \cref{p4:eq:centred-inverse}, whose bare core is already accurate to
      $O(\mathcal U^{-2})$.
\item Certify: compute $r_*=\log\mathcal S_\sigma(x_*)-\log y$ through
      $\log\Gamma$ and bound the error by
      \cref{p4:prop:derivative-bounds}.  The inequalities
      \cref{p4:eq:odd-bracket} and \cref{p4:eq:even-bracket} give an
      independent one-sided check that costs nothing.
\item For higher accuracy run Newton's method on the logarithmic equation,
      \[
        x_{j+1}=x_j-\frac{\log\mathcal S_\sigma(x_j)-\log y}
        {\lambda+\tfrac\sigma2\bigl[\psi(x_j/2+1)-\psi(x_j/2+1/2)\bigr]},
      \]
      which is stable even when $\mathcal S_\sigma(x)$ is far outside
      floating-point range.  \Cref{p4:prop:monotone} guarantees a positive
      denominator, and the transseries seed lies in the basin for large $y$.
\item For a discrete answer, round by \cref{p4:eq:staircase} and verify the
      final inequality exactly (\cref{p4:rem:certify-rounding}).
\end{enumerate}
\end{algorithm}

\subsection{Validation}

The tables below were recomputed at $50$ significant digits against exact
gamma-quotient reference values, with $y$ generated as
$\mathcal S_\sigma(x_0)$ so that the exact inverse is $x_0$.  They reproduce
the corresponding tables of \textsf{S2} and \textsf{S3} to every printed
digit, and, after the reparametrization of \cref{p4:rem:dictionary}, those of
\textsf{S1} as well.

\begin{table}[htbp]
\centering
\caption{Relative error of the forward expansion \cref{p4:eq:forward} after
the term $\sigma^{M}A_Mx^{-M}$.  Recomputed.}
\label{p4:tab:forward-errors}
\small
\begin{tabular}{@{}rrrrrrr@{}}
\toprule
$x$&$\sigma$&$M=0$&$M=2$&$M=4$&$M=6$&$M=8$\\
\midrule
$10$&$-1$&$2.5273\!\times\!10^{-2}$&$3.8527\!\times\!10^{-5}$&$4.7159\!\times\!10^{-7}$&$1.4186\!\times\!10^{-8}$&$7.8542\!\times\!10^{-10}$\\
$10$&$+1$&$2.4650\!\times\!10^{-2}$&$3.8626\!\times\!10^{-5}$&$4.7375\!\times\!10^{-7}$&$1.4233\!\times\!10^{-8}$&$7.8709\!\times\!10^{-10}$\\
$20$&$-1$&$1.2573\!\times\!10^{-2}$&$4.8642\!\times\!10^{-6}$&$1.5087\!\times\!10^{-8}$&$1.1546\!\times\!10^{-10}$&$1.6336\!\times\!10^{-12}$\\
$20$&$+1$&$1.2417\!\times\!10^{-2}$&$4.8704\!\times\!10^{-6}$&$1.5121\!\times\!10^{-8}$&$1.1565\!\times\!10^{-10}$&$1.6353\!\times\!10^{-12}$\\
$50$&$-1$&$5.0122\!\times\!10^{-3}$&$3.1226\!\times\!10^{-7}$&$1.5560\!\times\!10^{-10}$&$1.9152\!\times\!10^{-13}$&$4.3651\!\times\!10^{-16}$\\
$50$&$+1$&$4.9872\!\times\!10^{-3}$&$3.1242\!\times\!10^{-7}$&$1.5573\!\times\!10^{-10}$&$1.9164\!\times\!10^{-13}$&$4.3668\!\times\!10^{-16}$\\
\bottomrule
\end{tabular}
\end{table}

\begin{table}[htbp]
\centering
\caption{Absolute error
$\bigl|\mathcal I^{[K]}_\sigma(\mathcal S_\sigma(x_0))-x_0\bigr|$ of the
Lambert-normalized inverse \cref{p4:eq:inverse-analytic} with $K$ corrections
retained; $K=0$ is the bare core.  Recomputed.}
\label{p4:tab:inverse-errors}
\small
\begin{tabular}{@{}rrrrrrr@{}}
\toprule
$x_0$&$\sigma$&$K=0$&$K=2$&$K=4$&$K=6$&$K=8$\\
\midrule
$10$&$-1$&$3.8813\!\times\!10^{-2}$&$1.6744\!\times\!10^{-5}$&$2.9057\!\times\!10^{-7}$&$2.4840\!\times\!10^{-8}$&$1.7249\!\times\!10^{-9}$\\
$10$&$+1$&$3.3589\!\times\!10^{-2}$&$2.2677\!\times\!10^{-4}$&$3.9282\!\times\!10^{-6}$&$5.5615\!\times\!10^{-8}$&$2.0051\!\times\!10^{-9}$\\
$20$&$-1$&$1.8701\!\times\!10^{-2}$&$1.2039\!\times\!10^{-6}$&$6.9374\!\times\!10^{-9}$&$1.8267\!\times\!10^{-10}$&$3.2855\!\times\!10^{-12}$\\
$20$&$+1$&$1.7399\!\times\!10^{-2}$&$3.0220\!\times\!10^{-5}$&$1.3176\!\times\!10^{-7}$&$4.7142\!\times\!10^{-10}$&$4.3075\!\times\!10^{-12}$\\
$50$&$-1$&$7.3186\!\times\!10^{-3}$&$4.2970\!\times\!10^{-8}$&$5.6559\!\times\!10^{-11}$&$2.8529\!\times\!10^{-13}$&$8.3851\!\times\!10^{-16}$\\
$50$&$+1$&$7.1104\!\times\!10^{-3}$&$2.0095\!\times\!10^{-6}$&$1.4043\!\times\!10^{-9}$&$8.1104\!\times\!10^{-13}$&$1.1816\!\times\!10^{-15}$\\
$100$&$-1$&$3.6329\!\times\!10^{-3}$&$3.9487\!\times\!10^{-9}$&$1.6091\!\times\!10^{-12}$&$2.1882\!\times\!10^{-15}$&$1.6191\!\times\!10^{-18}$\\
$100$&$+1$&$3.5808\!\times\!10^{-3}$&$2.5440\!\times\!10^{-7}$&$4.4455\!\times\!10^{-11}$&$6.4401\!\times\!10^{-15}$&$2.3379\!\times\!10^{-18}$\\
\bottomrule
\end{tabular}
\end{table}

\begin{table}[htbp]
\centering
\caption{The same errors for the centred inverse
\cref{p4:eq:centred-inverse}.  The bare core is better by a factor
$\asymp\Xcore$, because $\widehat c^{\,(\sigma)}_1=0$; once corrections are
retained the two normalizations are comparable.  Recomputed.}
\label{p4:tab:centred-errors}
\small
\begin{tabular}{@{}rrrrrrr@{}}
\toprule
$x_0$&$\sigma$&$K=0$&$K=2$&$K=4$&$K=6$&$K=8$\\
\midrule
$10$&$-1$&$8.7330\!\times\!10^{-4}$&$5.5308\!\times\!10^{-5}$&$1.1593\!\times\!10^{-7}$&$2.7056\!\times\!10^{-9}$&$4.4842\!\times\!10^{-11}$\\
$10$&$+1$&$7.6102\!\times\!10^{-4}$&$5.6715\!\times\!10^{-5}$&$2.3489\!\times\!10^{-7}$&$7.9121\!\times\!10^{-9}$&$4.0347\!\times\!10^{-10}$\\
$20$&$-1$&$2.2206\!\times\!10^{-4}$&$7.4914\!\times\!10^{-6}$&$5.2119\!\times\!10^{-9}$&$3.7661\!\times\!10^{-11}$&$3.3070\!\times\!10^{-13}$\\
$20$&$+1$&$2.0696\!\times\!10^{-4}$&$7.5956\!\times\!10^{-6}$&$7.4297\!\times\!10^{-9}$&$6.3564\!\times\!10^{-11}$&$8.0880\!\times\!10^{-13}$\\
$50$&$-1$&$3.5860\!\times\!10^{-5}$&$5.0352\!\times\!10^{-7}$&$6.5130\!\times\!10^{-11}$&$8.3102\!\times\!10^{-14}$&$1.4333\!\times\!10^{-16}$\\
$50$&$+1$&$3.4850\!\times\!10^{-5}$&$5.0641\!\times\!10^{-7}$&$7.5153\!\times\!10^{-11}$&$1.0250\!\times\!10^{-13}$&$2.0275\!\times\!10^{-16}$\\
$100$&$-1$&$8.9913\!\times\!10^{-6}$&$6.3982\!\times\!10^{-8}$&$2.1684\!\times\!10^{-12}$&$7.1217\!\times\!10^{-16}$&$3.2358\!\times\!10^{-19}$\\
$100$&$+1$&$8.8632\!\times\!10^{-6}$&$6.4167\!\times\!10^{-8}$&$2.3300\!\times\!10^{-12}$&$7.9120\!\times\!10^{-16}$&$3.8477\!\times\!10^{-19}$\\
\bottomrule
\end{tabular}
\end{table}

\begin{warning}[These tables are not convergence]
\label{p4:rem:not-convergence}
The decreasing entries at fixed $x_0$ do not indicate convergence as
$K\to\infty$.  Both the forward and the inverse series are divergent, with the
beyond-all-orders floor \cref{p4:eq:floor-inverse}; the tables stop far short
of the least term (at $x=20$ the least term of \cref{p4:eq:D-first} sits at
$r=31$, of size $10^{-28}$).  The non-monotone entries in the odd column at
$x_0=10$ are ordinary behaviour of an asymptotic series at small argument, not
an error.  Beyond the least term one should use the exact Laplace
representation \cref{p4:eq:laplace}, or a Newton step against $\log\Gamma$, or
an exponentially improved representation transported by
\cref{p4:prop:sector-transport}.
\end{warning}

\section{Summary of the coefficient hierarchy}
\label{p4:sec:summary}

The whole calculation is a chain of constructive coefficient transformations,
each with a closed extraction formula and a triangular recurrence:
\begin{equation}
  \boxed{\;
  B_{2r+2}
  \;\xrightarrow{\ \text{gamma ratio}\ }\;
  d_{2r+1}
  \;\xrightarrow{\ \exp,\ \EB\ }\;
  A_n
  \;\xrightarrow{\ \text{Lagrange},\ \OB\ }\;
  c^{(\sigma)}_n
  \;\xrightarrow{\ \text{unfold }\Wlam\ }\;
  P^{(\sigma)}_n(H),\;}
  \label{p4:eq:hierarchy}
\end{equation}
while the Borel kernel $\tanh(t/2)/(2t)$ supplies, in one stroke, the exact
sum of the first stage, the sign and size of every remainder
(\cref{p4:thm:envelope}), the singulant $\pi$, the least-term scale, and the
complete Stokes data.  Two facts are specific to the swing factorial and
cannot be relegated to the general theory of \cref{p0:sec:top}:

\begin{enumerate}
\item \emph{Parity is leading-order data.}  The forward object is a
period-two transseries \cref{p4:eq:integer-forward} whose discrete
transmonomial $\sigma_n$ switches the algebraic scale between $2^{n}n^{-1/2}$
and $2^{n}n^{+1/2}$; the inverse object is a \emph{pair} of monotone branches,
and \cref{p4:prop:oscillatory} shows that this is forced, not chosen.
\item \emph{The two sheets need different Lambert branches.}  $\Wlam_{-1}$ on
the even sheet, $\Wlam_0$ on the odd sheet, the assignment being forced by
$\operatorname{sign}b=\sigma$ in \cref{p0:thm:lambert-core}.  The principal
branch on the even sheet does not merely lose accuracy: by
\cref{p4:prop:two-roots} it returns the root
$2/(\pi y^{2})\bigl(1+O(y^{-2})\bigr)\to0$, so that choice is fatal before any
correction is computed.
\end{enumerate}

%% ==================== fragment p5 ====================
\chapter{Synthesis: what the four subjects share, and where they part}
\label{p5:sec:top}

\begin{sourcebox}
This chapter belongs to the consolidation rather than to any source.  It
compares the four inversions carried out in
\cref{p1:sec:top,p2:sec:top,p3:sec:top,p4:sec:top}, states precisely how much
of each is an instance of \cref{p0:sec:top}, and records what the twelve
articles left open.  Nothing here is asserted beyond what the earlier chapters
prove.
\end{sourcebox}

\section{One inversion, four dictionaries}\label{p5:sec:one-inversion}

The four subjects of this volume were written about independently, and their
\emph{forward} analyses have almost nothing in common: a singularity analysis
of a Pólya functional equation, a pair of gamma-quotient branches, a
convergent Rademacher series over root-of-unity sectors, and a second, quite
different pair of gamma quotients.  Their \emph{inverse} analyses are, to a
degree that the individual articles could not see, the same computation.

The reason is \cref{p0:def:core}.  In each subject one first isolates a
\emph{dominant core} $\Phi_0$ --- an increasing function whose inverse is
available in closed form --- and writes the exact equation for the index $x$
in terms of the observed value $y$ as
\[
  \Phi_0(x)+\text{(a correction that is small relative to $\Phi_0'$)}=L(y).
\]
The core is then inverted exactly by the Lambert $\Wlam$ function, and the
correction is removed to all orders by one triangular recurrence,
\cref{p0:thm:core-reversion} (in the exponential--power case,
\cref{p0:thm:lambert-centered}).  The recurrence does not know which sequence
it came from: it sees only the core data and the coefficients $q_j$ of the
correction.

\Cref{p0:tab:dictionary} records the resulting dictionary.  Three points in it
deserve emphasis, because each was invisible from inside a single subject.

\paragraph{The partition numbers are the same problem after one substitution.}
Of the four, the partition numbers look least like the others: the
Hardy--Ramanujan phase is $\pi\sqrt{2N/3}$, a \emph{square root} of the
variable, where the other three carry a phase linear in the index or in
$x\log x$.  The three partition articles accordingly build their inversion by
hand around the equation they write as $\rho-2\log\rho=L$.  By
\cref{p0:lem:alpha-reduction} that equation \emph{is} the linear-phase core
$a\xi+b\log\xi$ in the variable $\xi=\sqrt{N}$, with $b/\alpha=-2$; the
substitution costs nothing and removes the apparent exception.  What the
substitution does \emph{not} do is commute with the arithmetic shift, and
\cref{p0:rem:affine-shift} shows why the order matters: re-expanding
$N=n-\tfrac1{24}$ in powers of $n$ after the reduction manufactures an entire
second tower of exponents $n^{\alpha-k}$, which is exactly the infinite
algebraic array that one source reports and attributes, correctly, to the
shift alone.

\paragraph{The two ``factorials'' are not the same kind of object.}
It is tempting --- and this volume's own brief made the mistake --- to expect
the double factorial and the swing factorial to sit side by side.  They do
not.  The swing factorial $n!/\lfloor n/2\rfloor!^{2}$ has its factorials
cancel to leave growth $2^{n}n^{\pm1/2}$, so it is an exponential--power model
with $a=\log2$ (\cref{p4:sec:top}).  The double factorial does not: its
logarithm is $\tfrac{s}{2}(\log s-1)+\bigO(\log s)$ with $s=n+1$, a phase of
factorial type, and since $s\log s/s^{\alpha}\to\infty$ for every
$\alpha\le1$, no choice of model data reproduces it
(\cref{p2:sec:top}).  What survives the difference is exactly the axiomatized
core: both $aX+b\log X$ and $X(\kappa\log X+d)$ are admissible in the sense of
\cref{p0:def:core}, both are inverted exactly by $\Wlam$
(\cref{p0:thm:lambert-core,p0:prop:factorial-core}), and the reversion above
them is one theorem, \cref{p0:thm:core-reversion}.  The right statement is
therefore not ``four instances of one model'' but ``four instances of one
\emph{core calculus}, of which three share a model''.

\paragraph{The Lambert branch is determined by a sign.}
Each of the twelve articles verifies by hand, for its own case, which branch
of $\Wlam$ produces the large root; several remark that the appearance of
$\Wlam_{-1}$ is essential and surprising.  With the signed convention of
\cref{p0:rem:conventions} it is neither: by \cref{p0:cor:branch-rule} the
branch is $\Wlam_0$ when $b>0$ and $\Wlam_{-1}$ when $b<0$, with the fold at
$X_\ast=\lvert b\rvert/a$ and the threshold
$L_c=\lvert b\rvert(1-\log(\lvert b\rvert/a))$.  Reading
\cref{p0:tab:dictionary} down the $b$ column now predicts the branch in every
chapter, including the two branches of \cref{p4:sec:top}, which differ only in
the sign of $\sigma$.

\section{Where the common frame stops}\label{p5:sec:where-it-stops}

The shared apparatus covers the passage from an exact forward description to
an all-orders inverse.  It covers neither end of that passage, and it is worth
being explicit about which parts of the four chapters are irreducibly their
own.

\begin{enumerate}[label=(\alph*),leftmargin=*]
\item \emph{Getting the forward description at all.}  The four forward
      analyses share no method.  \Cref{p1:sec:top} needs a singularity
      analysis of $T(z)=z\exp\sum_{k\ge1}T(z^k)/k$, the location and nature of
      the Otter singularity, and a transfer theorem;
      \cref{p2:sec:top,p4:sec:top} need Binet's representation and Stirling's
      series for gamma quotients; \cref{p3:sec:top} needs the circle method,
      or rather its exact endpoint, Rademacher's convergent series.  Nothing
      in \cref{p0:sec:top} produces any of these, and nothing in one of them
      helps with another.
\item \emph{The arithmetic layer.}  The sectors of \cref{p3:sec:top} are
      indexed by a denominator $k$ and carry Kloosterman-type amplitudes built
      from Dedekind sums.  This structure has no analogue in the other three
      subjects, whose exponentially small corrections come from a competing
      singularity (\cref{p1:sec:top}) or from Borel geometry
      (\cref{p2:sec:top,p4:sec:top}).  The frame transports each sector
      through the inversion once the sector is given; it does not explain
      where the sectors come from.
\item \emph{Parity.}  Two of the four sequences are not the restriction of any
      single smooth interpolation, and in both cases this is forced, not
      accidental: the even and odd subsequences have different exponents
      ($n^{\pm1/2}$ for the swing factorial) or different gamma
      representations (the double factorial).  \Cref{p0:def:three-inverses}
      and \cref{p0:thm:staircase}(4) handle the consequence uniformly --- the
      integer inverse is a minimum over residue classes --- but the
      \emph{fact} of the split is a property of the sequence.
\item \emph{Which interpolation.}  For a discrete sequence the asymptotic
      inverse is well defined only after an interpolation is chosen, and the
      choice is not innocent.  \Cref{p0:rem:three-differ} separates the three
      inverse objects; \cref{p2:sec:top} shows what changes when the OEIS
      periodic interpolation is used instead of the branchwise gamma
      continuation, namely a log-periodic oscillation in the inverse that is
      absent from either branch.  One partition source records the sharper
      version of the same phenomenon: the algebraic part of the inverse is
      interpolation-independent, the nonperturbative part is not.
\end{enumerate}

\section{The three inverse objects, subject by subject}
\label{p5:sec:three-objects}

\Cref{p0:def:three-inverses} distinguishes the staircase inverse $N_\ast$, the
inverse of a chosen interpolation, and the asymptotic inverse.  The four
chapters do not stand in the same relation to them, and the differences are
not cosmetic.

\begin{center}
\small
\begin{tabular}{@{}lccc@{}}
\toprule
Chapter & interpolation available & parity split & exact staircase recovery \\
\midrule
\ref{p1:sec:top} rooted trees & from the singular expansion & no & along the range \\
\ref{p2:sec:top} double factorial & two branches, or OEIS & yes & per residue class \\
\ref{p3:sec:top} partitions & from Rademacher & no & along the range \\
\ref{p4:sec:top} swing factorial & two branches & yes & per residue class \\
\bottomrule
\end{tabular}
\end{center}

Two cautions apply to the whole table.  First, by \cref{p0:thm:staircase}(5)
the staircase $N_\ast$ is \emph{not} uniformly approximable by any continuous
function: $\sup_{y\ge Y}\lvert N_\ast(y)-g(y)\rvert\ge\tfrac12$ for every
continuous $g$ and every $Y$.  Statements of the form ``$N_\ast(y)=\nu(y)+
\littleo(1)$'' are therefore false as uniform statements, however natural they
look, and the volume asserts only the exact ceiling identity, recovery along
the range $y=A_n$, and the residue-class refinement.  Second, the error floor
of the asymptotic inverse is not a defect of the reversion but is inherited:
by \cref{p0:thm:optimal-truncation} the forward series has a least-term
envelope, and \cref{p0:thm:backward-error} converts it into a floor for the
inverse.  No amount of further reversion crosses it.

\section{What this consolidation establishes}\label{p5:sec:new}

Several statements in this volume are not in any of its twelve sources.  They
are listed here because a reader who knows the sources should be able to see
quickly what has changed; the corrections themselves are boxed where they
occur and collected in \cref{app:repairs}.

\begin{enumerate}[label=(\arabic*),leftmargin=*]
\item \emph{A correct proof of the Lagrange fixed-point formula.}  Five of the
      twelve articles derive
      $\Coef{t^n}U=\sum_{m\le n}\tfrac1m\Coef{t^nu^{m-1}}\Psi^m$ by
      introducing a marker $z$ and applying Lagrange--Bürmann in $z$.  That
      application is not available: Lagrange--Bürmann needs $\varphi(0)$
      invertible, and here $\varphi(0)=\Psi(t,0)\in tR[[t]]$ is a non-unit.
      The identity is nonetheless true, and \cref{p0:lem:LB-universal} proves
      it by universality --- it is a polynomial identity over $\Rat$ in the
      coefficients of $\varphi$, valid in a localization into which the
      relevant domain injects.  Every reversion in this volume rests on this
      lemma, so the repair is not cosmetic.
\item \emph{The branch rule} \cref{p0:cor:branch-rule}, with the fold point
      and threshold, replacing four case analyses.
\item \emph{The normalization constant of the flattened blocks.}  The sources
      state $P_n(0)=D_n$ in a case where $a=1$; the correct general statement
      is $P_n(0)=a^{n+1}c_n$ (\cref{p0:thm:flattening}(3)), and since $a\ne1$
      in every chapter of this volume the factor is needed throughout.
\item \emph{A quantitative obstruction for the staircase},
      \cref{p0:thm:staircase}(5), in place of an $\littleo(1)$ claim that is
      false uniformly.
\item \emph{Backward error as a certificate.}  \Cref{p0:thm:backward-error}
      upgrades the sources' conditional bound --- which presupposes that a
      root exists --- to an existence statement, by a sign change and the
      intermediate value theorem.
\item \emph{One closed form where the sources had three.}
      \Cref{p0:thm:pure-stirling} is proved equal to all three published forms
      of the pure Lambert block, and the equality was verified in exact
      rational arithmetic for $1\le n\le8$.
\item \emph{A non-recursive formula for the factorial-core corrections.}
      \Cref{p0:lem:lagrange-fp-quadratic} extends the fixed-point formula to
      $\Psi=h+f$ with $h\in u^2R[[u]]$, which is what \cref{p2:sec:top}
      needs; the source reaches the same coefficients only through an ad hoc
      compositional inverse.
\end{enumerate}

\section{Open questions}\label{p5:sec:open}

\begin{enumerate}[label=(\roman*),leftmargin=*]
\item \emph{A Stokes-transport theorem.}  The mechanism by which an
      exponentially small forward sector becomes an exponentially small
      inverse sector is uniform --- it is the perturbation-operator form
      \cref{p0:prop:operator-form} --- but a theorem needs each subject's
      Borel geometry as a hypothesis.  The four chapters supply four separate
      instances; the general statement is not proved here.
\item \emph{The counting function.}  One source states that the number of
      indices with $A_n\le y$ equals the average of the branch inverses plus
      $\bigO(1)$.  The statement is plausible for every subject with a parity
      split, but it needs a per-subject monotonicity hypothesis and is not
      proved in this volume.
\item \emph{Rademacher-type sequences in general.}  \Cref{p3:sec:top} gives a
      template for inverting a sequence presented as a convergent sum over
      arithmetic sectors.  Whether the template applies to the other classical
      Rademacher-type expansions is untested here.
\item \emph{Multifactorials.}  One double-factorial source sketches the
      $m$-fold generalization $n!^{(m)}$.  The core is again of factorial
      type, so \cref{p0:thm:core-reversion} applies verbatim; what is not done
      is the residue-class analysis of the staircase for general $m$, which
      \cref{p0:thm:staircase}(4) reduces to a finite computation but does not
      carry out.
\item \emph{Relation to the formalized Lambert layer of this repository.}  The
      surrounding development contains a Lean-verified treatment of both real
      branches of $\Wlam$, a lower-branch bracket with an explicit remainder
      rate, and the saddle expansions of the Fabius function.  The Lambert
      core these four chapters are built on is therefore already formalized to
      a degree none of the sources assumes.  Which of the reversion machinery
      above is within reach of that layer --- \cref{p0:thm:lambert-core} and
      \cref{p0:cor:branch-rule} look closest --- is an open and, in this
      repository, a natural question.
\end{enumerate}

\appendix

\chapter{Provenance}\label{app:provenance}

This volume consolidates twelve independently written articles, three on each
of its four subjects.  All twelve were filed in this repository on
3~September 2026 as a single quick-intake batch, at
\texttt{drafts/series-and-transseries/special-function-inversion/}; the
directories below are relative to that path.  Nothing in the twelve was
edited in place: this volume is a new document, and the sources remain
available in git history.

\begin{center}
\small
\begin{longtable}{@{}p{0.30\textwidth}p{0.10\textwidth}p{0.12\textwidth}p{0.38\textwidth}@{}}
\toprule
\textbf{Source directory} & \textbf{Source} & \textbf{PDF} & \textbf{Absorbed into} \\
\midrule
\endfirsthead
\toprule
\textbf{Source directory} & \textbf{Source} & \textbf{PDF} & \textbf{Absorbed into} \\
\midrule
\endhead
\multicolumn{4}{@{}l@{}}{\emph{Chapter~\ref{p1:sec:top}: Butcher--Pólya rooted trees}}\\
\texttt{Butcher\_Tree\_Transseries} & 1\,674 lines & 21 pp, A4 & Chapter~\ref{p1:sec:top} \\
\texttt{Butcher\_Tree\_Transseries-2} & 2\,853 lines & 62 pp, A4 & Chapter~\ref{p1:sec:top} \\
\texttt{Butcher\_Tree\_Transseries-3} & 2\,382 lines & 35 pp, A4 & Chapter~\ref{p1:sec:top} \\
\addlinespace
\multicolumn{4}{@{}l@{}}{\emph{Chapter~\ref{p2:sec:top}: the double factorial}}\\
\texttt{Double\_Factorial\_Transseries} & 1\,552 lines & 18 pp, A4 & Chapter~\ref{p2:sec:top} \\
\texttt{Double\_Factorial\_Transseries-3} & 2\,287 lines & 28 pp, A4 & Chapter~\ref{p2:sec:top} \\
\texttt{double\_factorial\_transseries-2} & 2\,138 lines & 44 pp, A4 & Chapter~\ref{p2:sec:top} \\
\addlinespace
\multicolumn{4}{@{}l@{}}{\emph{Chapter~\ref{p3:sec:top}: the partition numbers}}\\
\texttt{Partition\_Number\_Transseries\_and\_Asymptotic\_Inverse} & 1\,901 lines & 32 pp, A4 & Chapter~\ref{p3:sec:top} \\
\texttt{Partition\_Number\_Transseries\_and\_Inverse} & 2\,020 lines & 27 pp, A4 & Chapter~\ref{p3:sec:top} \\
\texttt{Partition\_Numbers\_Transseries\_and\_Inverse} & 1\,905 lines & 33 pp, Letter & Chapter~\ref{p3:sec:top} \\
\addlinespace
\multicolumn{4}{@{}l@{}}{\emph{Chapter~\ref{p4:sec:top}: the swing factorial}}\\
\texttt{Swing\_Factorial\_Transseries} & 1\,885 lines & 28 pp, A4 & Chapter~\ref{p4:sec:top} \\
\texttt{Swing\_Factorial\_Transseries\_Article} & 1\,824 lines & 25 pp, A4 & Chapters~\ref{p0:sec:top} and~\ref{p4:sec:top} \\
\texttt{Swing\_Factorial\_Transseries\_and\_Inverse} & 1\,389 lines & 22 pp, A4 & Chapter~\ref{p4:sec:top} \\
\bottomrule
\end{longtable}
\end{center}

\noindent
The shared apparatus of Chapter~\ref{p0:sec:top} is drawn from all twelve.  Its
organizing normal form, and the Lagrange fixed-point route to the reversion
coefficients, follow the presentation in
\texttt{Swing\_Factorial\_Transseries\_Article} most closely; the
\(\alpha\)-reduction that brings the partition numbers into the same frame is
this volume's, since no source needed it.

\chapter{Repairs to the sources}\label{app:repairs}

Every correction this volume makes to one of its twelve sources is boxed in
the text where it occurs and listed here.  A repair is recorded when a source
states something false, states an asymptotic conclusion that its argument does
not support, or asserts a named result without proof.  Differences of
convention are not repairs and are recorded in the text as remarks instead.

\section{Chapter~\ref{p0:sec:top}: the common apparatus}

\paragraph{A proof that does not apply (five sources).}
\textsf{Swing\_Factorial\_Transseries\_Article},
\textsf{Partition\_Number\_Transseries\_and\_Inverse} (twice),
\textsf{Partition\_Numbers\_Transseries\_and\_Inverse},
\textsf{Partition\_Number\_Transseries\_and\_Asymptotic\_Inverse} and
\textsf{double\_factorial\_transseries-2} all prove the Lagrange fixed-point
formula
\(\Coef{t^n}U=\sum_{m\le n}\tfrac1m\Coef{t^nu^{m-1}}\Psi^m\)
by introducing an auxiliary marker \(z\), solving \(U=z\Psi(t,U)\), and
applying Lagrange--Bürmann in \(z\).  Lagrange--Bürmann requires the relevant
\(\varphi(0)\) to be invertible; here \(\varphi(0)=\Psi(t,0)\in tR[[t]]\) is a
non-unit, and the residue proof genuinely breaks, since \(\psi=w/\varphi(w)\)
is then not an element of \(w\,R[[t]][[w]]\).  The identity is true;
\cref{p0:lem:LB-universal} proves it by universality instead.  Every reversion
in this volume depends on the lemma, so the repair is structural rather than
cosmetic.

\paragraph{An asymptotic claim that is false uniformly.}
Several sources state, or leave the reader to infer, that the integer inverse
satisfies \(N_\ast(y)=\nu(y)+\littleo(1)\).  No continuous function
approximates a staircase uniformly: \cref{p0:thm:staircase}(5) shows
\(\sup_{y\ge Y}\lvert N_\ast(y)-g(y)\rvert\ge\tfrac12\) for every continuous
\(g\) and every \(Y\).  The volume asserts only the exact ceiling identity,
nearest-integer recovery \emph{along the range}, and the residue-class
refinement.

\paragraph{A normalization constant.}
The sources record \(P_n(0)=D_n\) for the flattened blocks, in a case where
\(a=1\).  The general statement is \(P_n(0)=a^{\,n+1}c_n\)
(\cref{p0:thm:flattening}(3)); since \(a\neq1\) in every chapter of this
volume the factor is needed throughout.

\paragraph{Two truncation conventions in one article.}
\textsf{Swing\_Factorial\_Transseries\_Article} uses the exclusive convention
in its forward error table and the inclusive one in its inverse table, so
those two tables are not on a common scale.  The volume is inclusive
throughout (\cref{p0:rem:conventions}); tables inherited from that article's
forward section are re-indexed.

\paragraph{A claim of this volume's own, corrected.}
The brief under which \cref{p0:sec:top} was written asserted that the case
\(\alpha=1\) of \cref{p0:def:model} covers the double factorial.  It does not:
\(\log\bigl((x+1)!!\bigr)=\tfrac{s}{2}(\log s-1)+\bigO(\log s)\) with
\(s=x+1\), and \(s\log s/s^{\alpha}\to\infty\) for every \(\alpha\le1\).  The
sources never claimed membership.  \Cref{p0:def:core} axiomatizes the
dominant core instead, and \cref{p0:prop:factorial-core} inverts the factorial
phase exactly.

\paragraph{Corrupt source text, not reproduced.}
\textsf{Swing\_Factorial\_Transseries\_Article} equation
\texttt{eq:inverse-Stokes-leading} and
\textsf{Butcher\_Tree\_Transseries-2} lines 2482--2483 contain corrupt control
sequences.  The affected statements are re-derived here rather than copied.

\section{Chapter~\ref{p2:sec:top}: the double factorial}

\paragraph{Convergence asserted without a threshold.}
The first source states that the periodic carrier ``has the convergent local
Lagrange--Bürmann expansion'' for sufficiently large \(\Lambda\), with no
threshold and no proof of convergence as opposed to formal solvability; the
second proves only a contraction estimate, from which convergence of the
Lagrange series does not follow.  \Cref{p2:thm:carrier} separates the two:
existence and uniqueness of the real root by contraction for
\(\Lambda>2\pi a\), and convergence of the Lagrange series under the explicit
sufficient condition \(\mu=2a/\Lambda\le1/5\).

\paragraph{Two assertions in place of proofs.}
The deepest-pole coefficient (``only the quadratic kernel and the first
Bernoulli coefficient can contribute'') is exactly the step needing the degree
bookkeeping, and is supplied here by induction on the
\(\Lambda^{-1}\)-degree.  The coefficientwise sign law is argued in one source
with ``use a positive rooted-tree contribution \dots; such a tree exists'',
exhibiting no such object; \cref{p2:thm:sign-law} keeps the correct positivity
reduction and replaces the strictness step with an explicit attainability
induction, verified symbolically for \(n\le7\).

\paragraph{A formal-to-analytic slide.}
The third source presents an all-orders phase reversion as a theorem with an
asymptotic conclusion, introduced by ``setting \(\tau=1\) as an asymptotic
expansion''.  The identity is formal; making it asymptotic needs estimates,
uniform in the oscillatory phase, that are neither stated nor proved.  It is
demoted here to the formal \cref{p2:prop:phase-reversion}, with a proved
remainder supplied separately.

\paragraph{A silent three-way inconsistency.}
One source centres the periodic inverse at the even-branch constant, the other
two at the parity mean; the leading displacements consequently differ by
\(2a/\Lambda\).  Both are correct for their own centring, but no source says
so, and a reader comparing them would conclude that one is wrong.  The volume
adopts the mean centring and states the relation.

\paragraph{A notation collision across three sources.}
The letter \(C\) denotes the multiplicative amplitude, the logarithmic
constant, and the branch prefactor in the three siblings respectively, and
\(\alpha,\beta\) carry two incompatible meanings.  A crosswalk table fixes the
volume's usage; the core parameters are never abbreviated, to avoid a fourth
collision with \(\kappa\) in \cref{p0:def:core}.

\paragraph{A qualitative warning made quantitative.}
All three sources warn that a different smooth interpolation ``can change this
sector'' without exhibiting one.  \Cref{p2:thm:no-canonical} exhibits the
family \(\mathcal D(x)\e^{\lambda\sin\pi x}\), which interpolates \(n!!\) at
every non-negative integer and changes the leading oscillatory amplitude from
\(2a/\Lambda\) to \(2\sqrt{a^2+\lambda^2}/\Lambda\) --- a change larger than
the entire Bernoulli grid.  This is what makes the chapter's organizing claim
a theorem rather than a caution.

\section{Chapter~\ref{p1:sec:top}: the rooted-tree numbers}

\paragraph{Digits printed beyond the accuracy of the computation.}
One source prints the constants of the inherited \(-\sqrt\rho\) sector to
forty decimal digits, but only about twenty-five of them are correct.  Two
independent recomputations here --- including a from-scratch direct summation
at eighty digits --- give
\[
  r_0^{-}=0.468636627924845548087202621931819314\ldots,
\]
and the volume prints only digits it has verified.  This is the one place in
the twelve articles where a printed numeral is wrong; it is a precision
failure rather than a mistake in a formula.

\paragraph{An audit appendix that contradicts its own main text.}
The same source's audit appendix reports \(D_3=-3.582177326832123\ldots\),
disagreeing with the value its main text derives.  The main text is correct;
the appendix entry is not, and the volume uses the recomputed value.

\paragraph{A symbol used for two different constants.}
One source writes \(\alpha\) for the growth constant where the other two write
\(\lambda\), and the volume's Chapter~\ref{p0:sec:top} uses \(\alpha\) for the
phase exponent of \cref{p0:def:model}.  The chapter fixes \(\lambda\) for the
growth constant throughout and never uses \(\alpha\) for it.

\paragraph{An asymptotic statement with a formal proof.}
A source states its inverse-coefficient theorem asymptotically but proves it
only formally, and separately displays a global transseries with no hypothesis
under which it could hold.  The first is proved here in its asymptotic form;
the second is replaced by a conditional theorem with a finite radius and an
explicit hypothesis.

\paragraph{A circular derivation of the singularity bounds.}
The elementary bounds \(\tfrac14\le\rho\le\e^{-1}\) are used in the sources in
a way that presupposes what they are used to prove.  They are established here
independently, from \(x\le T\e^{-T}\le\e^{-1}\) and the Catalan bound, and the
critical characterization is then obtained non-circularly.

\section{Chapter~\ref{p4:sec:top}: the swing factorial}

\paragraph{A divergent series written as an equality.}
One source sets its bookkeeping coupling to \(1\) and prints the resulting
operator expansion as an equality.  There is then no small parameter, the sum
diverges, and a number is equated to a divergent series; as printed the
statement is false.  The \(m\)-th summand is \(\bigO(X^{-(2m-1)})\), so the
correct statement is asymptotic, and regrouping gives the finite differential
form used here.  The other two sources state the same expansion correctly.

\paragraph{A theorem whose hypothesis is never used.}
A sector-transport result is labelled a theorem, its hypothesis is not used in
the argument, and its conclusion is prefaced ``formally''.  It is split here
into the formal all-orders identity, which is what the proof gives, and a
quantitative first-order statement proved by the mean value theorem under
stated hypotheses.

\paragraph{An overclaim about complex branches.}
One source asserts that every Lambert branch with \(\lvert w\rvert\to\infty\)
yields an inverse branch with the same coefficients.  For \(k\ne0,-1\) the
existence of an analytic branch of the inverse along the path is a hypothesis,
not a consequence, and it fails at gamma poles and Stokes directions.  The
claim is demoted to a remark with the two hypotheses displayed.

\paragraph{The analytic inverse statement, unproved in all three sources.}
One source writes ``likewise, for each fixed \(M\), \dots'' with no proof; the
second sketches an argument that never bounds the residual of the truncated
ansatz nor uses forward control at the perturbed point; the third proves the
ingredients but never assembles them.  This is the one formal-to-analytic
slide common to all three.  \Cref{p4:thm:inverse-analytic} proves the
statement in full, with an effective constant.

\paragraph{A recommendation contradicted by its own table.}
A remark asserts that the centred normalization is ``often the preferable
computational normal form''.  The same article's table shows the centred core
better by a factor \(\approx10^{3}\) at one end and the uncentred form
slightly better at the other.  The corrected statement is that centring kills
the \(X^{-1}\) term and so improves the \emph{bare core} by a factor
\(\asymp X\), while with several corrections retained the two are comparable
and neither dominates uniformly.

\paragraph{Two convergence claims without criteria.}
An equivalence of the Lambert and logarithmic normalizations is asserted in
prose with no formula and no proof; it is proved here as a finite explicit
composition.  A convergence claim appeals to ``the classical analysis of
Lambert \(\Wlam\)'', but the classical result concerns a doubly indexed
series, whereas the series at hand has coefficients polynomial in
\(\ell\) and converges only when \(\ell/L\) is small; an explicit criterion is
proved instead.

\section{Chapter~\ref{p3:sec:top}: the partition numbers}

\paragraph{A definition that contradicts its own use.}
The source abbreviated \textsf{S2} in \cref{p3:sec:top} defines the real
Rademacher amplitude by a sum over \(1\le h\le k\) with \(\gcd(h,k)=1\).  For
\(k\ge2\) that is the intended index set, but at \(k=1\) it returns
\(\cos(2\pi\nu)\) rather than the constant \(1\), while the same article uses
\(A_1\equiv1\) throughout.  With the printed definition the interpolant would
not interpolate \(p\): its dominant term would oscillate off-lattice at full
relative strength, and the article's own eventual-monotonicity statement would
be false.  The volume takes \(h\in U_k\) with \(U_1=\{0\}\); every later
formula of that source is then correct as printed.

\paragraph{An understated radius.}
\textsf{S1} states only that its universal sector-coefficient series
``converges for all sufficiently small \(\lvert t\rvert\)''.  The radius is
exactly \(\sqrt{24}\) in \(t=n^{-1/2}\), the nearest singularities being the
branch points of \(\sqrt{1-t^2/24}\); convergence therefore holds for every
real \(n>1/24\).

\paragraph{A convergent series and an asymptotic total, conflated.}
All three partition sources are loose about the distinction.  The algebraic
series \(\sum_m\omega_mn^{-m/2}\) converges, but its sum is the leading
Rademacher block, \emph{not} \(p(n)\); the Poincaré statement about \(p(n)\)
carries an error that is the sum of a convergent algebraic tail and an
exponentially small Rademacher remainder.  Summing the algebraic series to
convergence does not compute \(p(n)\): the residual stalls at the
\(\e^{-\mu/2}\) floor.

\paragraph{A false intermediate bound.}
\textsf{S1}'s kernel estimate asserts \(g(z)\le z^2/3\) on \([0,1]\), which
fails at \(z=1\) where \(g(1)=\e^{-1}=0.3679\ldots>1/3\).  The chapter proves
\(g(z)\le(z^2/3)\cosh z\) instead, which is what the tail bound actually
needs.

\paragraph{An overclaim in this volume's own Chapter~\ref{p0:sec:top}.}
\Cref{p0:rem:formal-analytic} originally said that the correction series \(Q\)
is divergent of Gevrey type ``in all four chapters''.  It is not: for the
partition numbers \(Q(t)=\log(1-t)\) has radius of convergence \(1\).  The
remark now distinguishes the three divergent chapters from
\cref{p3:sec:top}, where \cref{p0:thm:optimal-truncation} has no content and
the error floor is set by the Rademacher remainder instead.


\end{document}
